\begin{bt}%[1D2V2-7]
% nguon: Archive/2024/TOAN_11_XE/TOAN_11_XE/data4/11C2/KNTT-2.OTC.tex #493
%2.27. 
Từ $0$ giờ đến $12$ giờ trưa, chuông của một chiếc đồng hồ quả lắc sẽ đánh bao nhiêu tiếng, biết rằng nó chỉ đánh chuông báo giờ và số tiếng chuông bằng số giờ?
\loigiai{
Lúc $1$ giờ đồng hồ đánh $1$ tiếng chuông.\\
Lúc $2$ giờ đồng hồ đánh $2$ tiếng chuông.\\
......\\
Lúc $12$ giờ trưa đồng hồ đánh $12$ tiếng chuông.\\
Do đó, từ $0$ giờ đến $12$ giờ trưa, đồng hồ đánh số tiếng chuông là
$1+ 2+ 3+ .... + 11+ 12$.\\
Đây là tổng $12$ số hạng đầu tiên của cấp số cộng có số hạng đầu $u_1= 1$, công sai $d = 1$.\\
Vậy tổng số tiếng chuông đồng hồ trong khoảng thời gian từ $0$ đến $12$ giờ trưa là $\dfrac{12\cdot (1+12)}{2}=78$.
}
\end{bt}

\begin{bt}%[1D2V3-8]
% nguon: Archive/2024/TOAN_11_XE/TOAN_11_XE/data4/11C2/KNTT-2.OTC.tex #494
%2.28. 
Tế bào E.Coli trong điều kiện nuôi cấy thích hợp cứ $20$ phút lại phân đôi một lần. Hỏi sau $24$ giờ, tế bào ban đầu sẽ phân chia thành bao nhiếu tế bào?
\loigiai{
Lần phân chia thứ nhất, $1$ tế bào thành $2$ tế bào, số tế bào lần $1$ phân chia là $u_1 = 2$.\\
Lần phân chia thứ hai $2$, số tế bào lần $2$ phân chia là  $u_2=2\cdot 2 = u_1 \cdot 2$.\\
Lần phân chia thứ $3$ có  $4$ tế bào phân chia, số tế bào lần $3$ phân chia là $u_3=2\cdot u_2$.\\
Như vậy một tế bào phân đôi sẽ tạo thành cấp số nhân có công bội là $2$, số hạng đầu là $u_1=2$.\\
 Sau $n$ lần phân chia từ một tế bào phân được thành $u_n=2^{n-1}u_1$.\\
Đổi $24$ giờ $=24 \cdot 60 =  72 \cdot 20$ (phút)  $\Rightarrow 24$ giờ gấp $72$ lần $20$ phút. \\
Do đó, sau $24$ giờ số tế bào nhận được là $u_{72}=2^{71}\cdot 2 = 2^{72}$ (tế bào).
}
\end{bt}

\begin{bt}%[1D2V3-7]
% nguon: Archive/2024/TOAN_11_XE/TOAN_11_XE/data4/11C2/KNTT-2.OTC.tex #495
%2.29. 
Chứng minh rằng
	\begin{enumerate}
	\item Trong một cấp số cộng $\left(u_n\right)$, mỗi số hạng (trừ số hạng đầu và số hạng cuối, nếu có) đều là trung bình cộng của hai số hạng đứng kề với nó, nghĩa là
	$$u_k=\frac{u_{k-1}+u_{k+1}}{2}, \text{ với } k \geq 2.	$$
	\item Trong một cấp số nhân, bình phương của mỗi số hạng (trừ số hạng đầu và số hạng cuối, nếu có) đều là tích của hai số hạng đứng kề với nó, nghĩa là
	$$	u_k^2=u_{k-1}\cdot u_{k+1}, \text{ với } k \geq 2 .$$
\end{enumerate}
\loigiai{
		\begin{enumerate}
	\item Giả sử $(u_n)$ là một cấp số cộng, có công sai $d$.
	 Khi đó $u_n=u_1+(n-1)d$. \\
	 Với $k \geq 2$, ta có $u_{k-1}+u_{k+1} = u_1+(k-2)d+u_1+kd = 2u_1+2(k-1)d=2u_k$.\\
	 Vậy $u_k=\dfrac{u_{k-1}+u_{k+1}}{2}, \text{ với } k \geq 2.	$
\item Giả sử $(u_n)$ là một cấp số nhân, có công bội $q$.
Khi đó $u_n=q^{n-1}u_1$. \\
Với $k \geq 2$, ta có $u_{k-1}u_{k+1} = q^{k-2}u_1q^{k}u_1 = q^{2k-2}u_1^2=(u_k)^2$.\\
Vậy $	u_k^2=u_{k-1}\cdot u_{k+1}, \text{ với } k \geq 2 .$
\end{enumerate}
}
\end{bt}

\begin{bt}%[1D2V3-7]
% nguon: Archive/2024/TOAN_11_XE/TOAN_11_XE/data4/11C2/KNTT-2.OTC.tex #496
%2.30.
	 Tìm ba số, biết theo thứ tự đó chúng lập thành cấp số cộng và có tổng bằng $21$, và nếu lần lượt cộng thêm các số $2$; $3$; $9$ vào ba số đó thì được ba số lập thành một cấp số nhân.
	 \loigiai{
Giả sử cấp số cộng cần tìm là $x,\,y, \, z$ . Theo tính chất của cấp số cộng ta có  $x+z=2y$.\\
Kết hợp giả thiết ta có $x+y+z=21 \Rightarrow 3y=21 \Leftrightarrow y=7$.\\
Gọi $d$ là công sai của cấp số cộng $x,\,y, \, z$ thì $x=7-d$, $z=7+d$.\\
Sau khi cộng thêm các số $2$, $3$, $9$ vào ba số $x$, $y$, $z$ ta được ba số $ x + 2$, $y + 3$, $z + 9$ hay $9 - d$, $10$, $16 - d$.\\
Theo tính chất của cấp số nhân ta có\\ $(9-d)(16+d)=100 \Leftrightarrow 144-7d-d^2=100 \Leftrightarrow d^2+7d-44=0  \Leftrightarrow \hoac{& d=4 \\&
d=-11.}$\\
Với $d = 4$ ta được cấp số cộng $3$, $7$, $11$. \\
Với $d = -11$ ta được cấp số cộng $18$, $7$, $-4$.  
}
\end{bt}

\begin{bt}%[1D2V2-7]
% nguon: Archive/2024/TOAN_11_XE/TOAN_11_XE/data4/11C2/KNTT-2.OTC.tex #497
 %2.31.	
 Mặt sàn tầng một (tầng trệt) của một ngôi nhà cao hơn mặt sân $0,5 \mathrm{~m}$. Cầu thang đi từ tầng một lên tầng hai gồm $25$ bậc, mỗi bậc cao $16 \mathrm{~cm}$.
 	\begin{enumerate}
\item Viết công thức để tìm độ cao của bậc cầu thang thứ $n$ so với mặt sân.
\item  Tính độ cao của sân tầng hai so với mặt sân.
\end{enumerate}
\loigiai{ 	\begin{enumerate} 
\item Mỗi bậc thang cao $16$ cm = $0{,}16$ m $\Rightarrow n$ bậc thang cao $0{,}16\cdot n$ (m).\\
Vì mặt bằng sàn cao hơn mặt sân $0{,}5$ m nên công thức tính độ cao của bậc $n$ so với mặt sân sẽ là $h_n = (0{,} 5 + 0{,}16n)$ (m).
\item Độ cao của sàn tầng hai so với mặt sân ứng với $n = 25$ là $h_{25} = 0{,}5 + 0{,}16.25 = 4,5$ (m)
\end{enumerate}
}
\end{bt}

\begin{bt}%[1D2C3-6]
% nguon: Archive/2024/TOAN_11_XE/TOAN_11_XE/data4/11C2/KNTT-2.OTC.tex #498
 %2.32. 
	Một hình vuông màu vàng có cạnh $1$ đơn vị dài được chia thành chín hình vuông nhỏ hơn và hình vuông ở chính giữa được tô màu xanh như Hình $2.1$ Mỗi hình vuông màu vàng nhỏ hơn lại được chia thành chín hình vuông con, và mỗi hình vuông con ở chính giữa lại được tô màu xanh. Nếu quá trình này được tiếp tục lặp lại năm lần, thì tổng diện tích các hình vuông được tô màu xanh là bao nhiêu?
	\begin{center}
			\begin{tikzpicture}[>=stealth,thick,scale=0.7]
			\def\n{1}
			\def\a{3}
			\pgfmathsetmacro{\m}{int(3^(\n))}
			\def\hv#1{
				\ifnum#1>0
				\fill[blue!50] (-\a/3,-\a/3) rectangle (\a/3,\a/3);
				\pgfmathtruncatemacro{\k}{#1-1}
				\foreach \i in {0,...,3}{\begin{scope}[shift={(90*\i:2)},scale=1/3]\hv{\k}\end{scope}}
				\foreach \i in {0,...,3}{\begin{scope}[shift={(45+90*\i:{4/sqrt(2)})},scale=1/3]\hv{\k}\end{scope}}
				\fi
			}
			\fill[orange!50] (-\a,-\a) rectangle (\a,\a);
			\hv{\n}
			\foreach \i in {0,1,...,\m}{
				\draw[blue!50] 
				({-\a+2*\i *\a/\m},\a)--++(270:2*\a)
				(\a,{-\a+2*\i *\a/\m})--++(180:2*\a)
				;
			}
		\end{tikzpicture}
		\begin{tikzpicture}[>=stealth,thick,scale=0.7]
			\def\n{2}
			\def\a{3}
			\pgfmathsetmacro{\m}{int(3^(\n))}
			\def\hv#1{
				\ifnum#1>0
				\fill[blue!50] (-\a/3,-\a/3) rectangle (\a/3,\a/3);
				\pgfmathtruncatemacro{\k}{#1-1}
				\foreach \i in {0,...,3}{\begin{scope}[shift={(90*\i:2)},scale=1/3]\hv{\k}\end{scope}}
				\foreach \i in {0,...,3}{\begin{scope}[shift={(45+90*\i:{4/sqrt(2)})},scale=1/3]\hv{\k}\end{scope}}
				\fi
			}
			\fill[orange!50] (-\a,-\a) rectangle (\a,\a);
			\hv{\n}
			\foreach \i in {0,1,...,\m}{
				\draw[blue!50] 
				({-\a+2*\i *\a/\m},\a)--++(270:2*\a)
				(\a,{-\a+2*\i *\a/\m})--++(180:2*\a)
				;
			}
		\end{tikzpicture}
		\begin{tikzpicture}[>=stealth,thick,scale=0.7]
			\def\n{4}
			\def\a{3}
			\def\hv#1{
				\ifnum#1>0
				\fill[blue!50] (-\a/3,-\a/3) rectangle (\a/3,\a/3);
				\pgfmathtruncatemacro{\k}{#1-1}
				\foreach \i in {0,...,3}{\begin{scope}[shift={(90*\i:2)},scale=1/3]\hv{\k}\end{scope}}
				\foreach \i in {0,...,3}{\begin{scope}[shift={(45+90*\i:{4/sqrt(2)})},scale=1/3]\hv{\k}\end{scope}}
				\fi
			}
			\fill[orange!50] (-\a,-\a) rectangle (\a,\a);
			\hv{\n}
		\end{tikzpicture}
	\end{center}
	\loigiai{
Lần phân chia thứ nhất, $1$ hình vuông thành $9$ hình vuông con, diện tích hình vuông tô màu xanh là $u_1=\dfrac{1}{9}$.\\
Lần phân chia thứ hai, $8$ hình vuông thành $9$ hình vuông con, diện tích hình vuông tô màu xanh tăng thêm là $u_2=\dfrac{1}{9}\left(\dfrac{8}{9}\right)$.\\
Lần phân chia thứ ba, $8^2$ hình vuông thành $9$ hình vuông con, diện tích hình vuông tô màu xanh tăng thêm là $u_3=\dfrac{1}{9}\left(\dfrac{8}{9}\right)^2$.\\
Lần phân chia thứ tư, $8^3$ hình vuông thành $9$ hình vuông con, diện tích hình vuông tô màu xanh tăng thêm là $u_4=\dfrac{1}{9}\left(\dfrac{8}{9}\right)^3$.\\
Lần phân chia thứ năm, $8^4$ hình vuông thành $9$ hình vuông con, diện tích hình vuông tô màu xanh tăng thêm là $u_5=\dfrac{1}{9}\left(\dfrac{8}{9}\right)^4$.\\
Như vậy diện tích các hình vuông tăng thêm sau mỗi lần chia  tạo thành cấp số nhân có công bội là $q=\dfrac{8}{9}$, số hạng đầu là $u_1=\dfrac{1}{9}$.\\
Do đó, tổng diện tích hình vuông tô màu xanh sau $5$ lần chia là\\
\[u_1+u_2+u_3+u_4+u_5=\dfrac{1-q^5}{1-q}\cdot u_1=\dfrac{1-\left(\dfrac{8}{9}\right)^5}{1-\dfrac{8}{9}}\cdot \dfrac{1}{9}=\dfrac{26281}{39366}.\]	

}
\end{bt}

\begin{bt}%[1D2H2-4]
% nguon: Archive/2024/TOAN_11_XE/TOAN_11_XE/data4/11C2/KNTT-2.OTC.tex #502
 
	Tứ giác $ABCD$ có số đo bốn góc $A,B,C,D$ theo thứ tự lập thành cấp số cộng. Biết số đo góc $C$ gấp $5$ lần số đo góc $A$. Tính số đo các góc của tứ giác $ABCD$ theo đơn vị độ.
	\loigiai{
		Do bốn góc $A,B,C,D$ theo thứ tự lập thành cấp số cộng nên số đo $4$ góc của tứ giác $ABCD$ lần lượt là $u_1,u_1+d,u_1+2d,u_1+3d$. \\
		Theo giả thiết, ta có: $\heva{u_1+u_1+d+u_1+2d+u_1+3d=360 \\ u_1+2d=5u_1}  \Leftrightarrow  \heva{4u_1+6d=360 \\ -4u_1+2d=0} \Leftrightarrow \heva{u_1=22,5 \\ d=45.}$\\
		Vậy số đo các góc của tứ giác lần lượt là $22,5^\circ; 67,5^\circ; 112,5^\circ; 157,5^\circ$.
	}
\end{bt}

\begin{bt}%[1D2V2-7]
% nguon: Archive/2024/TOAN_11_XE/TOAN_11_XE/data4/11C2/KNTT-2.OTC.tex #503
 
	Người ta trồng cây theo các hàng ngang với quy luật: ở hàng thứ nhất có $1$ cây, ở hàng thứ hai có $2$ cây, ở hàng thứ ba có $3$ cây, $\ldots$ ở hàng thứ $n$ có $n$ cây. Biết rằng người ta trồng hết $4950$ cây. Hỏi số hàng cây được trồng theo cách trên là bao nhiêu?
	\loigiai{
		Số cây ở mỗi hàng lập thành một cấp số cộng với $u_1=1$ và công sai $d=1$.\\
		Theo giả thiết ta có: $S_n=4950 \Leftrightarrow n+\dfrac{(n-1)n}{2} =4950  \Leftrightarrow n^2 +n -4950=0 \Leftrightarrow \hoac{n&=-99 (l)\\n&=100.}$\\
		Vậy có $100$ hàng cây được trồng.
	}
\end{bt}

\begin{bt}%[1D2V3-8]
% nguon: Archive/2024/TOAN_11_XE/TOAN_11_XE/data4/11C2/KNTT-2.OTC.tex #504
 
	Một cái tháp có $11$ tầng. Diện tích của mặt sàn tầng $2$ bằng nửa diện tích của mặt đáy tháp và diện tích của mặt sàn mỗi tầng bằng nửa diện tích của mặt sàn mỗi tầng ngay bên dưới. Biết mặt đáy tháp có diện tích là $12 288m^2$. Tính diện tích của mặt sàn tầng trên cùng của tháp theo đơn vị mét vuông.
	\loigiai{
		Do diện tích của mặt sàn tính từ tầng một lập thành một cấp số nhân với $u_2=\dfrac{1}{2}.12288=6144$ và $q=\dfrac{1}{2}$.\\
		Ta có $\heva{u_2&=6144 \\ q&=\dfrac{1}{2}}  \Leftrightarrow \heva{u_1&=12288 \\ q&=\dfrac{1}{2}.}$\\
		Do đó  $u_{11}=u_1.q^{10}=12288.\dfrac{1}{12^{10}}=12m^2$.
		Vậy diện tích của mặt sàn tầng trên cùng là	$12m^2$.
	}
\end{bt}

\begin{bt}%[1D2V3-8]
% nguon: Archive/2024/TOAN_11_XE/TOAN_11_XE/data4/11C2/KNTT-2.OTC.tex #505
 
	Một khay nước có nhiệt độ $23^\circ$ được đặt vào ngăn đá của tủ lạnh. Biết sau mỗi giờ, nhiệt độ của nước giảm $20\%$. Tính nhiệt độ của khay nước đó sau $6$ giờ theo đơn vị độ $C$.
	\loigiai{
		Nhiệt độ sau mỗi giờ của khay nước theo thứ tự lập thành cấp số nhân với $u_1=23$ và $q=(1-20\%)$.\\
		Ta có $u_6=u_1.q^5=23.(1-20\%)^5 \approx 7,5$.\\
		Nhiệt độ của khay nước sau $6$ giờ là $ \approx 7,5^\circ $.
	}
	
\end{bt}

\begin{bt}%[1D2V3-8]
% nguon: Archive/2024/TOAN_11_XE/TOAN_11_XE/data4/11C2/KNTT-2.OTC.tex #507
 
	Ông An vay ngân hàng $1$ tỉ đồng với lãi suất $12\%/$năm. Ông đã trả nợ theo cách: Bắt đầu từ tháng thứ nhất sau khi vay, cuối tháng ông trả ngân hàng số tiền là $a$ (đồng) và đã trả hết nợ sau đúng $2$ năm kể từ ngày vay. Hỏi số tiền mỗi tháng mà ông An phải trả là bao nhiêu đồng (làm tròn kết quả đến hàng nghìn)?
	\loigiai{
		Do lãi suất là $12\%$/năm tương đương với lãi là $1\%$/tháng.\\
		Sau $1$ tháng, ông An còn nợ là: $10^9.(1+1\%)-a=10^9.(1,01)-S_1$.\\
		Sau $2$ tháng, ông An còn nợ là: $10^9.(1.01)^2-a.(1.01)-a=10^9(1,01)^2-S_2$.\\
		Sau $3$ tháng, ông An còn nợ là: $10^9.(1.01)^3-a(1.01)^2-a(1.01)-a=10^9.(1.01)^3-S_3$.\\
		Sau $24$ tháng, ông An còn nợ là: $10^9.(1.01)^{24}-S_{24}=0$.\\
		Do đó $S_{24}=10^9.(1.01)^{24}  \Leftrightarrow a.\dfrac{1-(1.01)^{24}}{1-(1.01)}=10^9.(1.01)^{24} \Leftrightarrow a =\dfrac{10^9.(1.01)^{24}.0.01}{(1.01)^{24}-1}\approx 47073472,22$.\\
		Vậy mỗi tháng ông An phải trả $47073500$.
	}
\end{bt}

\begin{bt}%[1D2N3-2]
% nguon: Archive/2024/TOAN_11_XE/TOAN_11_XE/data4/11C2/KNTT-2.3.tex #511
%[DCHT Toán 11 - KNTT -Tên Huỳnh Thanh Chí]
	Dãy số $ 1;1;1;1;\ldots $ có phải là một cấp số nhân hay không?
	\dapso{Dãy số $ 1;1;1;1;\ldots $ là một cấp số nhân.}
	\loigiai{
		Dễ thấy $ \dfrac{u_2}{u_1}=\dfrac{u_3}{u_2}=\ldots=1 $ là một số không đổi.\\
		Do đó dãy số $ 1;1;1;1;\ldots $ là một cấp số nhân.
	}
\end{bt}

\begin{bt}%[1D2H3-2]
% nguon: Archive/2024/TOAN_11_XE/TOAN_11_XE/data4/11C2/KNTT-2.3.tex #512
%[DCHT Toán 11 - KNTT -Tên Huỳnh Thanh Chí]
	Dãy số $ \heva{& u_1=3\\ & u_{n+1}=\dfrac{9}{u_n}} $ có phải là một cấp số nhân không? Nếu có, hãy tìm công bội của cấp số nhân đó.
	\dapso{$ (u_n) $ là một cấp số nhân với công bội $ q=1 $.}
	\loigiai{
		Xét dãy số $ \heva{& u_1=3\\ & u_{n+1}=\dfrac{9}{u_n}} $ có
		$ \dfrac{u_{n+1}}{u_n}=\dfrac{9}{u_n}:\dfrac{9}{u_{n-1}}=\dfrac{u_{n-1}}{u_n}\Rightarrow u_{n+1}=u_{n-1},\forall n\ge 2 $.\\
		Do đó ta có $ \heva{& u_1=u_3=u_5=\ldots=u_{2n+1}=\ldots \quad (1)\\ & u_2=u_4=u_6=\ldots=u_{2n}=\ldots \quad (2).} $\\
		Theo đề bài ta có $ u_1=3 \Rightarrow u_2=\dfrac{9}{u_1}=3 $ (3).\\
		Từ $ (1), (2) $ và $ (3) $ suy ra $ u_1=u_2=u_3=u_4=\ldots=u_{2n}=u_{2n+1}=\ldots $.\\
		Do đó $ (u_n) $ là một cấp số nhân với công bội $ q=1 $.
	}
\end{bt}

\begin{bt}%[1D2H3-2]
% nguon: Archive/2024/TOAN_11_XE/TOAN_11_XE/data4/11C2/KNTT-2.3.tex #513
%[DCHT Toán 11 - KNTT -Tên Huỳnh Thanh Chí]
	Cho $ (u_n) $ là cấp số nhân có công bội $ q\ne 0,u_1\ne 0 $. Chứng minh rằng dãy số $ (v_n) $ với $ v_n=u_nu_{2n} $ cũng là một cấp số nhân.
	\dapso{$ (v_n) $ là một cấp số nhân với công bội là $ q^3 $.}
	\loigiai{
		Ta có $ \dfrac{v_n}{v_{n-1}}=\dfrac{u_nu_{2n}}{u_{n-1}u_{2(n-1)}}=\dfrac{u_1q^{n-1}\cdot u_1q^{2n-1}}{u_1q^{n-2}\cdot u_1q^{2n-3}}=q^3 $. Do đó $ (v_n) $ là một cấp số nhân với công bội là $ q^3 $.
	}
\end{bt}

\begin{bt}%[1D2V3-2]
% nguon: Archive/2024/TOAN_11_XE/TOAN_11_XE/data4/11C2/KNTT-2.3.tex #514
%[DCHT Toán 11 - KNTT -Tên Huỳnh Thanh Chí]
	Cho dãy số $ (u_n) $ được xác định bởi $ \heva{& u_1=2\\ & u_{n+1}=4u_n+9},\forall n\in\mathbb{N}^* $. Chứng minh rằng dãy số $ (v_n) $ xác định bởi $ v_n=u_n+3,\forall n\in\mathbb{N}^* $ là một cấp số nhân. Hãy xác định số hạng đầu và công bội của cấp số nhân đó.
	\dapso{$ (v_n) $ là cấp số nhân với công bội $ q=4 $ và số hạng đầu $ v_1=5 $.}
	\loigiai{
		Ta có $ v_n=u_n+3 $ (1) và $ v_{n+1}=u_{n+1}+3 $ (2).\\
		Theo đề ta có $ u_{n+1}=4u_n+9 \Rightarrow u_{n+1}+3=4(u_n+3) $ (3).\\
		Thay (1) và (2)	vào (3) ta được $ v_{n+1}=4v_n \Rightarrow \dfrac{v_{n+1}}{v_n}=4,\forall n\in\mathbb{N}^* $.\\
		Suy ra $ (v_n) $ là cấp số nhân với công bội $ q=4 $ và số hàng đầu $ v_1=u_1+3=2+3=5 $.
	}
\end{bt}

\begin{bt}%[1D2N3-3]
% nguon: Archive/2024/TOAN_11_XE/TOAN_11_XE/data4/11C2/KNTT-2.3.tex #515
%[DCHT Toán 11 - KNTT -Tên Huỳnh Thanh Chí]
	Tìm số hạng tổng quát của dãy số $ 2;4;8;16;32;\ldots $, biết dãy $ (u_n) $ là một cấp số nhân.  
	\dapso{$ u_n=2\cdot 2^{n-1} $.}
	\loigiai{
		Vì dãy số $ (u_n) $ là một cấp số nhân nên $ q=\dfrac{u_2}{u_1}=\dfrac{u_3}{u_2}=\ldots=2 $ và số hạng đầu $ u_1=2 $.\\
		Do đó dãy số $ 2;4;8;16;32;\ldots $ là một cấp số nhân có số hạng tổng quát là $ u_n=u_1q^{n-1}=2\cdot 2^{n-1} $.
	}
\end{bt}

\begin{bt}%[1D2H3-3]
% nguon: Archive/2024/TOAN_11_XE/TOAN_11_XE/data4/11C2/KNTT-2.3.tex #516
%[DCHT Toán 11 - KNTT -Tên Huỳnh Thanh Chí]
	Tìm số hạng đầu, công bội và số hạng tổng quát của cấp số nhân, biết $ \heva{& u_1+u_5=51\\ & u_2+u_6=102.} $
	\dapso{$ u_1=3 $, $ q=2 $ và $ u_n=3\cdot 2^{n-1} $.}
	\loigiai{
		Vì $ (u_n) $ là một cấp số nhân nên $ u_n=u_1\cdot q^{n-1} $.\\
		Ta có $ \heva{& u_1+u_5=51\\ & u_2+u_6=102}\Leftrightarrow\heva{& u_1+u_1q^4=51\\ & u_1q+u_1q^5=102}\Leftrightarrow\heva{& u_1(1+q^4)=51 \qquad (1)\\ & u_1q(1+q^4)=102 \qquad (2).} $\\
		Chia từng vế của $ (2) $ cho $ (1) $ ta được $ \dfrac{u_1q(1+q^4)}{u_1(1+q^4)}=\dfrac{102}{51} \Leftrightarrow q=2 $.\\
		Suy ra $ u_1=\dfrac{51}{1+q^4}=\dfrac{51}{17}=3 $, $ u_n=u_1\cdot q^{n-1}=3\cdot 2^{n-1} $.\\
		Vậy $ u_1=3 $, $ q=2 $ và $ u_n=3\cdot 2^{n-1} $.
	}
\end{bt}

\begin{bt}%[1D2H3-3]
% nguon: Archive/2024/TOAN_11_XE/TOAN_11_XE/data4/11C2/KNTT-2.3.tex #517
%[DCHT Toán 11 - KNTT -Tên Huỳnh Thanh Chí]
	Tìm số hạng đầu, công bội và số hạng tổng quát của cấp số nhân, biết $ \heva{& u_1+u_6=30\\ & u_2+u_7=120.} $
	\dapso{$ u_1=\dfrac{6}{205} $, $ q=4 $ và $ u_n=\dfrac{6}{205}\cdot 4^{n-1} $.}
	\loigiai{
		Vì $ (u_n) $ là một cấp số nhân nên $ u_n=u_1\cdot q^{n-1} $.\\
		Ta có $ \heva{& u_1+u_6=30\\ & u_2+u_7=120}\Leftrightarrow\heva{& u_1+u_1q^5=30\\ & u_1q+u_1q^6=102}\Leftrightarrow\heva{& u_1(1+q^5)=30 \qquad (1)\\ & u_1q(1+q^5)=120 \qquad (2).} $\\
		Chia từng vế của $ (2) $ cho $ (1) $ ta được $ \dfrac{u_1q(1+q^5)}{u_1(1+q^5)}=\dfrac{120}{30} \Leftrightarrow q=4 $.\\
		Suy ra $ u_1=\dfrac{30}{1+q^5}=\dfrac{30}{1+4^5}=\dfrac{6}{205} $, $ u_n=u_1\cdot q^{n-1}=\dfrac{6}{205}\cdot 4^{n-1} $.\\
		Vậy $ u_1=\dfrac{6}{205} $, $ q=4 $ và $ u_n=\dfrac{6}{205}\cdot 4^{n-1} $.
	}
\end{bt}

\begin{bt}%[1D2H3-3]
% nguon: Archive/2024/TOAN_11_XE/TOAN_11_XE/data4/11C2/KNTT-2.3.tex #518
%[DCHT Toán 11 - KNTT -Tên Huỳnh Thanh Chí]
	Tìm số hạng đầu, công bội và số hạng tổng quát của cấp số nhân, biết $ \heva{& u_3=40\\ & u_6=160.} $
	\dapso{$ u_1=\dfrac{40}{9} $, $ q=3 $ và $ u_n=40\cdot 3^{n-3} $.}
	\loigiai{
		Vì $ (u_n) $ là một cấp số nhân nên $ u_n=u_1\cdot q^{n-1} $.\\	
		Ta có $ \heva{& u_3=40\\ & u_6=1080}\Leftrightarrow \heva{& u_1q^2=40 \qquad (1)\\ & u_1q^5=1080\qquad (2).} $\\
		Chia từng vế của $ (2) $ cho $ (1) $ ta được $ \dfrac{u_1q^5}{u_1q^2}=\dfrac{1080}{40} \Leftrightarrow q^3=27 \Leftrightarrow q=3 $.\\
		Suy ra $ u_1=\dfrac{40}{q^2}=\dfrac{40}{3^2}=\dfrac{40}{9} $, $ u_n=u_1\cdot q^{n-1}=\dfrac{40}{9}\cdot 3^{n-1}=40\cdot 3^{n-3} $.\\
		Vậy $ u_1=\dfrac{40}{9} $, $ q=3 $ và $ u_n=40\cdot 3^{n-3} $.
	}
\end{bt}

\begin{bt}%[1D2V3-3]
% nguon: Archive/2024/TOAN_11_XE/TOAN_11_XE/data4/11C2/KNTT-2.3.tex #519
%[DCHT Toán 11 - KNTT -Tên Huỳnh Thanh Chí]
	Tìm số hạng đầu, công bội và số hạng tổng quát của cấp số nhân có công bội $ q\in \mathbb{Z},q\ne 0 $, biết $ \heva{& u_2+u_4=10\\ & u_1+u_3+u_5=-21.} $
	\dapso{$ u_1=-1 $, $ q=-2 $ và $ u_n=\dfrac{(-2)^n}{2} $.}
	\loigiai{
		Vì $ (u_n) $ là một cấp số nhân nên $ u_n=u_1\cdot q^{n-1} $ với $ q\in \mathbb{Z},q\ne 0 $.\\
		Ta có $ \heva{& u_2+u_4=10\\ & u_1+u_3+u_5=-21}\Leftrightarrow\heva{& u_1q+u_1q^3=10\\ & u_1q+u_1q^2+u_1q^4=-21}\Leftrightarrow\heva{& u_1(q+q^3)=10 \qquad (1)\\ & u_1(1+q^2+q^4)=-21 \qquad (2).} $\\
		Chia từng vế của $ (2) $ cho $ (1) $ ta được 
		\allowdisplaybreaks
		\begin{eqnarray*}
			\dfrac{u_1(1+q^2+q^4)}{u_1(q+q^3)}=\dfrac{-21}{10} 
			&\Leftrightarrow& 10q^4+21q^3+10q^2+21q+10=0 \\
			&\Leftrightarrow& (q+2)(2q+1)(5q^2-2q+5)=0 \\
			&\Leftrightarrow& \hoac{& q=-2 \text{ (thỏa mãn)}\\ & q=-\dfrac{1}{2} \text{ (loại)}.}
		\end{eqnarray*}
		Suy ra $ u_1=\dfrac{10}{q+q^3}=-1 $, $ u_n=u_1\cdot q^{n-1}=(-1)\cdot (-2)^{n-1}=-(-2)^{n-1}=-\dfrac{(-2)^n}{-2}=\dfrac{(-2)^n}{2} $.\\
		Vậy $ u_1=-1 $, $ q=-2 $ và $ u_n=\dfrac{(-2)^n}{2} $.
	}
\end{bt}

\begin{bt}%[1D2N3-4]
% nguon: Archive/2024/TOAN_11_XE/TOAN_11_XE/data4/11C2/KNTT-2.3.tex #520

	Cho $u_n$ là CSN thỏa $u_1=2$; $u_4=16$. Tìm số hạng thứ $5$ của CSN.
	\loigiai{
		Do $u_n$ là CSN nên ta có $u_4=u_1\cdot q^3 \Rightarrow q^3=\dfrac{u_4}{u_1}=8 \Rightarrow q=2$. \\
		Vậy $u_5=u_1\cdot q^4=2\cdot 2^4=32$.
	}
\end{bt}

\begin{bt}%[1D2H3-2]
% nguon: Archive/2024/TOAN_11_XE/TOAN_11_XE/data4/11C2/KNTT-2.3.tex #521

	Cho cấp số nhân $(u_n)$ có $\heva{&u_4+u_6=-540 \\ &u_3+u_5=180}$. Tính số hạng đầu $u_1$ và công bội $q$ của cấp Số nhân.
	\loigiai{
		Ta có $\heva{&u_4 + u_6=-540 \\ &u_3+u_5=180}
		\Leftrightarrow \heva{&u_1q^3(1+q^2)=-540 \\ &u_1q^2(1+q^2)=180}
		\Leftrightarrow \heva{&u_1=2 \\ &q=-3.}$ \\
		Vậy $\heva{&u_1=2 \\ &q=-3}$ là số hạng cần tìm.
	}
\end{bt}

\begin{bt}%[1D2H3-4]
% nguon: Archive/2024/TOAN_11_XE/TOAN_11_XE/data4/11C2/KNTT-2.3.tex #522

	Cho cấp số nhân có $u_1=-3$, $q=\dfrac{2}{3}$. Số $\dfrac{-96}{243}$ là số hạng thứ mấy của cấp số nhân?
	\loigiai{
		Giả sử số $\dfrac{-96}{243}$ là số hạng thứ $n$ của cấp số nhân.\\
		Ta có: $u_1\cdot q^{n-1}=\dfrac{-96}{243}\Leftrightarrow(-3)\left(\dfrac{2}{3}\right)^{n-1}=\dfrac{-96}{243}\Leftrightarrow n=6$.\\
		Vậy số $\dfrac{-96}{243}$ là số hạng thứ $6$ của cấp số nhân.}
\end{bt}

\begin{bt}%[1D2H3-3]
% nguon: Archive/2024/TOAN_11_XE/TOAN_11_XE/data4/11C2/KNTT-2.3.tex #523

	Cấp số nhân $\left(u_{n}\right)$ có số hạng tổng quát là $u_n=\dfrac{3}{5} \cdot 2^{n-1}, n \in \mathbb{N}^*$. Số hạng đầu tiên và công bội của cấp số nhân đó là
	\loigiai{
		Ta có $u_{1}=\dfrac{3}{5} \cdot 2^{1-1}=\dfrac{3}{5}$ và $u_{2}=\dfrac{3}{5} \cdot 2^{2-1}=\dfrac{6}{5} \Rightarrow q=\dfrac{u_{2}}{u_{1}}=2$.\\
		Vậy $u_{1}=\dfrac{3}{5}$ và $q=2$.
	}
\end{bt}

\begin{bt}%[1D2V3-8]
% nguon: Archive/2024/TOAN_11_XE/TOAN_11_XE/data4/11C2/KNTT-2.3.tex #524

	Người ta thiết kế một cái tháp gồm $10$ tầng theo cách: Diện tích bề mặt trên của mỗi tầng bằng nửa diện tích bề mặt trên của tầng ngay bên dưới và diện tích bề mặt của tầng 1 bằng nửa diện tích bề mặt đế tháp. Biết diện tích bề mặt đế tháp là $12\, 288$ m$^2$, tính diện tích bề mặt trên cùng của tháp.
	\loigiai{
		Gọi $S$ là diện tích mặt đế và $T_1, T_2, \ldots, T_{10}$ là diện tích bề mặt của tầng 1, tầng 2, \ldots, tầng 10.\\
		Khi đó, ta có
		\allowdisplaybreaks
		\begin{eqnarray*}
			&T_1&=\dfrac{1}{2}\cdot S;\\
			&T_2&=\dfrac{1}{2}\cdot T_1=\dfrac{1}{2^2}\cdot S;\\
			&T_3&=\dfrac{1}{2}\cdot T_2=\dfrac{1}{2^3}\cdot S;\\
			&\vdots &\\
			&T_{10}&=\dfrac{1}{2}\cdot T_9=\dfrac{1}{2^{10}}\cdot S.
		\end{eqnarray*}
		Vậy diện tích bề mặt trên cùng của tháp là $T_{10}=\dfrac{1}{2^{10}}\cdot 12288=12$ m$^2$.
	}
\end{bt}

\begin{bt}%[1D2H3-4]
% nguon: Archive/2024/TOAN_11_XE/TOAN_11_XE/data4/11C2/KNTT-2.3.tex #525

	Cho dãy $3,x,12,y$. Tìm $x,y$ để dãy là CSN.
	\loigiai{
		Dãy là CSN khi $\heva{&x^2=3\cdot 12 \\ &12^2=x\cdot y}\Leftrightarrow 
		\heva{&x=6 \\ &y=24}$ hoặc $\heva{&x=-6 \\ &y=-24.}$
	}
\end{bt}

\begin{bt}%[1D2H3-4]
% nguon: Archive/2024/TOAN_11_XE/TOAN_11_XE/data4/11C2/KNTT-2.3.tex #526

	Cho dãy  $x-1, 2x, 4x+3$. Tìm $x$ để dãy là CSN. 
	\loigiai{
		Dãy là CSN khi $(2x)^2=(x-1)(4x+3) \Leftrightarrow x=-3$.
	} 
\end{bt}

\begin{bt}%[1D2V3-7]
% nguon: Archive/2024/TOAN_11_XE/TOAN_11_XE/data4/11C2/KNTT-2.3.tex #527

	Các số $x+6y$, $5x+2y$, $8x+y$ theo thứ tự đó lập thành một cấp số cộng, đồng thời, các số $x+\dfrac{5}{3}$, $y-1$, $2x-3y$ theo thứ tự đó lập thành một cấp số nhân. Hãy tìm $x$ và $y$.
	\loigiai{
		\begin{itemize}
			\item Ba số $x+6y$, $5x+2y$, $8x+y$ lập thành cấp số cộng nên $(x+6y)+(8x+y)=2(5x+2y)\Leftrightarrow x=3y$.
			\item Ba số $x+\dfrac{5}{3}$, $y-1$, $2x-3y$ lập thành cấp số nhân nên $\left(x+\dfrac{5}{3}\right)(2x-3y)=\left(y-1\right)^2$.
		\end{itemize}
		Thay $x=3y$ vào ta được $8y^2+7y-1=0\Leftrightarrow y=-1$ hoặc $y=\dfrac{1}{8}$.\\
		Với $y=-1$ thì $x=-3$; với $y=\dfrac{1}{8}$ thì $x=\dfrac{3}{8}$.}
\end{bt}

\begin{bt}%[1D2V3-5]
% nguon: Archive/2024/TOAN_11_XE/TOAN_11_XE/data4/11C2/KNTT-2.3.tex #528

	Tìm tất cả các giá trị của tham số $m$ để phương trình sau có ba nghiệm phân biệt lập thành một cấp số nhân $x^3-7x^2+2\left(m^2+6m\right)x-8=0$.
	\loigiai{
		+ \textbf{Điều kiện cần:} \\
		Giả sử phương trình đã cho có ba nghiệm phân biệt $x_1$,$x_2$,$x_3$ lập thành một cấp số nhân.\\
		Theo định lý Vi-ét, ta có $x_1x_2x_3=8$.\\
		Theo tính chất của cấp số nhân, ta có $x_1x_3=x_2^2$. Suy ra  $x_2^3=8\Leftrightarrow x_2=2$.\\
		Với nghiệm $x=2$, ta có $m^2+6m-7=0\Leftrightarrow\hoac{&m=1\\&m=-7.}$ \\
		+ \textbf{Điều kiện đủ:}\\
		Với $m=1$ hoặc $m=-7$ thì $m^2+6m=7$.\\
		Khi đó phương trình ban đầu trở thành $x^3-7x^2+14x-8=0$.\\
		Giải phương trình này, ta được các nghiệm là $1$,$2$,$4$. Hiển nhiên ba nghiệm này lập thành một cấp số nhân với công bội $q=2$.\\
		Vậy $m=1$ và $m=-7$ là các giá trị cần tìm.}
\end{bt}

\begin{bt}%[1D2V3-7]
% nguon: Archive/2024/TOAN_11_XE/TOAN_11_XE/data4/11C2/KNTT-2.3.tex #529

	Các số $x+6y$, $5x+2y$, $8x+y$ theo thứ tự đó lập thành một cấp số cộng; đồng thời các số $x-1$, $y+2$, $x-3y$ theo thứ tự đó lập thành một cấp số nhân. Tính $x^2+y^2$.
	\loigiai{
		Theo giả thiết ta có
		\[
		\heva{&(x+6y)+(8x+y)=2(5x+2y)\\&(x-1)(x-3y)=\left(y+2\right)^2}\Leftrightarrow\heva{&x=3y\\&(3y-1)(3y-3y)=\left(y+2\right)^2}\Leftrightarrow \heva{&x=3y\\&0=\left(y+2\right)^2}\Leftrightarrow \heva{&x=-6\\&y=-2.}
		\]
		Vậy $x^2+y^2=40$.}
\end{bt}

\begin{bt}%[1D2N3-6]
% nguon: Archive/2024/TOAN_11_XE/TOAN_11_XE/data4/11C2/KNTT-2.3.tex #530
%[DCHT Toán 11 - KNTT -Đỗ Chí Tâm] 
	Tính tổng 10 số hạng đầu tiên của cấp số nhân $(u_n)$, biết $u_1=-3$ và công bội $q=-2$. \dapso{$S_{10}=1023$}
	\loigiai{
		Ta có: $S_{10}=\dfrac{u_1\left(1-q^{10}\right)}{1-q}=1023$.}
\end{bt}

\begin{bt}%[1D2H3-6]
% nguon: Archive/2024/TOAN_11_XE/TOAN_11_XE/data4/11C2/KNTT-2.3.tex #531
%[DCHT Toán 11 - KNTT -Đỗ Chí Tâm] 
	Tính tổng $8$ số hạng đầu tiên của cấp số nhân $(u_n),$ biết $u_1=3$ và $u_2=6$. \dapso{$S_8=765$}
	\loigiai{
		Ta có: $u_2=u_1.q \Leftrightarrow 6=3.q \Leftrightarrow q=2$\\
		$\hspace*{1.2cm}  S_8=u_1\dfrac{1-q^8}{1-q}=3.\dfrac{1-2^8}{1-2}=765. $
	}
\end{bt}

\begin{bt}%[1D3H1-5]
% nguon: Archive/2024/TOAN_11_XE/TOAN_11_XE/data4/11C2/KNTT-2.3.tex #532
%[DCHT Toán 11 - KNTT -Đỗ Chí Tâm] 
	Tính tổng vô hạn $S=1+\dfrac{1}{2}+\dfrac{1}{2^2}+...+\dfrac{1}{2^n}+...$ \dapso{$S=2$}
	\loigiai{
		Đây là tổng của cấp số nhân lùi vô hạn, với $u_1=1, q=\dfrac{1}{2}$. Khi đó
		$$S=\dfrac{u_1}{1-q}=\dfrac{1}{1-\dfrac{1}{2}}=2.$$
	}
\end{bt}

\begin{bt}%[1D2V3-6]
% nguon: Archive/2024/TOAN_11_XE/TOAN_11_XE/data4/11C2/KNTT-2.3.tex #533
%[DCHT Toán 11 - KNTT -Đỗ Chí Tâm] 
	Tính tổng $200$ số hạng đầu tiên của dãy số $(u_n)$ biết $\heva{&u_1=1\\&u_{n+1}=3u_n}$.
	\loigiai{
		Dễ thấy dãy đã cho là một cấp số nhân với công bội $q=3; u_1=1$.\\
		Từ đó $S_{200}=u_1\dfrac{q^{200}-1}{q-1}$ $=\dfrac{3^{200}-1}{2}$.
	}
\end{bt}

\begin{bt}%[1D2V3-6]
% nguon: Archive/2024/TOAN_11_XE/TOAN_11_XE/data4/11C2/KNTT-2.3.tex #534
%[DCHT Toán 11 - KNTT -Đỗ Chí Tâm]
	Một cấp số nhân có số hạng đầu $u_1=3$, công bội $q=2$. Biết $S_n=765$, tìm $n$. \dapso{$n=8$}
	\loigiai{ 
		Áp dụng công thức tính tổng của cấp số nhân ta có
		$S_n=\dfrac{u_1( 1-q^n )}{1-q}=\dfrac{3.( 1-2^n )}{1-2}=765 \Leftrightarrow n=8$.} 
\end{bt}

\begin{bt}%[1D2V3-7]
% nguon: Archive/2024/TOAN_11_XE/TOAN_11_XE/data4/11C2/KNTT-2.3.tex #535
%[DCHT Toán 11 - KNTT -Đỗ Chí Tâm] 
	Ba số $x, y, z$ theo thứ tự đó lập thành một CSN với công bội $q (q\ne 1)$, đồng thời các số $x, 2y, 3z$ theo thứ tự đó lập thành một CSC với công sai $d$ . Hãy tìm $q$?
	\loigiai{
		Ta có $x+3z=2.2y \Leftrightarrow x+3xq^2=2.2xq\Leftrightarrow 1+3q^2=4q \Leftrightarrow \hoac{q&=\dfrac{1}{3}\\q&=1 (L)}$
	}
\end{bt}

\begin{bt}%[1D2V3-7]
% nguon: Archive/2024/TOAN_11_XE/TOAN_11_XE/data4/11C2/KNTT-2.3.tex #536
%[DCHT Toán 11 - KNTT -Đỗ Chí Tâm] 
	Biết rằng $a, b, c$ là ba số hạng liên tiếp của một CSC và $a, c, b$ là ba số hạng liên tiếp của một CSN, đồng thời $a+b+c=30$. Tìm $a,b,c$.
	\loigiai{
		Theo đề ta có  $\heva{&a+c=2b\,\,\,\hspace*{0.8cm}(1)\\&ab=c^2\,\,\,\hspace*{1.2cm}(2)\\&a+b+c=30\,\,\, (3)}$\\
		Từ $(1)$ và $(3)$ ta được $3b=30\Leftrightarrow b=10$\\
		Thay $b=10$ vào $(1), (2)$ ta được $\heva{a+c=20\\10a=c^2}\Leftrightarrow \hoac{&c=10, a=10\,\, (L)\\ &c=-20, a=40 (N)}$\\
		Vậy $a=40, b=10, c=-20$
	}
\end{bt}

\begin{bt}%[1D2V3-7]
% nguon: Archive/2024/TOAN_11_XE/TOAN_11_XE/data4/11C2/KNTT-2.3.tex #537
%[DCHT Toán 11 - KNTT -Đỗ Chí Tâm] 
	Ba số $x, y, z$ theo thứ tự đó lập thành một CSN. Ba số $x, y-4 , z$ theo thứ tự đó lập thành CSN. Đồng thời các số $x, y-4 , z-9$ theo thứ tự đó lập thành CSC. Tìm $x,y,z$.
	
	\loigiai{
		Theo đề ta có  $\heva{&xz=y^2\,\,\,\hspace*{2.2cm}(1)\\&xz=(y-4)^2\,\,\,\hspace*{1.4cm}(2)\\&x+(z-9)=2(y-4)\,\,\, (3)}$\\
		Từ $(1)$ và $(2)$ ta có $y^2=(y-4)^2\Leftrightarrow y=2$\\
		Thế $y=2$ vào $(1)$ và $(3)$ ta được $\heva{xz=4\\x+z=5}\Rightarrow x=4, z=1$ hoặc $x=1, z=4$\\
		Vậy có 2 bộ $(x,y,z)$ thỏa yêu cầu bài toán là $(1,2,4)$ và $(4,2,1)$
	}
\end{bt}

\begin{bt}%[1D2V3-7]
% nguon: Archive/2024/TOAN_11_XE/TOAN_11_XE/data4/11C2/KNTT-2.3.tex #538
%[DCHT Toán 11 - KNTT -Đỗ Chí Tâm] 
	Cho $a,b,c$ là ba số hạng liên tiếp của một CSN và $a,b,c-4$ là ba số hạng liên tiếp của một CSC, đồng thời $a,b-1,c-5$ là ba số hạng liên tiếp của một CSN. Tìm $a,b,c$ biết $a,b,c$ là các số nguyên. 
	
	\loigiai{
		Theo đề ta có  $\heva{&ac=b^2\,\,\,\hspace*{1.9cm}(1)\\&a+c-4=2b\,\,\,\hspace*{0.8cm}(2)\\&a(c-5)=(b-1)^2\,\,\,\,(3)}$\\
		Thay $(1)$ vào $(3)$: $b^2-5a=b^2-2b+1\Leftrightarrow b=\dfrac{5a+1}{2}$\\
		Thay vào $(2)$ ta được $a+c-4=5a+1\Leftrightarrow c=4a+5$\\
		Thế $b, c$ theo $a$ vào $(1)$ ta được $9a^2-10a+1=0\Leftrightarrow a=1 \vee a=\dfrac{1}{9} (L)$
		Vậy $a=1, b=3, c=9$
	}
\end{bt}

\begin{bt}%[1D2C3-7]
% nguon: Archive/2024/TOAN_11_XE/TOAN_11_XE/data4/11C2/KNTT-2.3.tex #539
%[DCHT Toán 11 - KNTT -Đỗ Chí Tâm]
	Cho $4$ số nguyên dương, trong đó $3$ số đầu lập thành một CSC, $3$ số hạng sau thành lập CSN.
	Biết rằng tổng của số hạng đầu và số hạng cuối là $37$, tổng của hai số hạng giữa là $36$. Tìm tổng $4$ số đó
	\loigiai{
		Gọi 4 số cần tìm lần lượt là $a, b, c, d$\\
		$a, b, c$ là $3$ số hạng liên tiếp của CSC. Ta có $a+c=2b\,\,\, (1)$\\
		$b,c,d$ là $3$ số hạng liên tiếp của CSN. Ta có $bd=c^2\,\,\, (2)$\\
		Theo giả thuyết ta có $\heva{a+d&=37\,\,\, (3)\\b+c&=36\,\,\, (4)}$\\
		Từ $(4)\Rightarrow b=36-c$ thay vào $(1)$ ta được $a=72-3c$, thay $a$ vào $(3)$ ta được $d=-35+3c$\\
		Thế $b,d$ vào $(2)$ ta được $(36-c)(-35+3c)=c^2\Rightarrow c=20 \vee c=\dfrac{63}{4} (L)$\\
		Vậy $c=20, a=12, b=16, d=95\Rightarrow S=a+b+c+d=143$
	}
\end{bt}

\begin{bt}%[1D2H3-8]
% nguon: Archive/2024/TOAN_11_XE/TOAN_11_XE/data4/11C2/KNTT-2.3.tex #540
%[DCHT Toán 11 - KNTT - Dung Phuong] 
	Dân số trung bình của Việt Nam năm $2020$ là $97{,}6$ triệu người, tỉ lệ tăng dân số là $1{,}14 \% /$năm.
	\begin{flushright}
		\textit{(Nguồn: Niên giám thống kê của Việt Nam năm 2020, NXB Thống kê, 2021)}
	\end{flushright}
	Giả sử tỉ lệ tăng dân số không đổi qua các năm.
	\begin{enumerate}
		\item Sau 1 năm, dân số của Việt Nam sẽ là bao nhiêu triệu người (làm tròn kết quả đến hàng phần mười)?
		\item Viết công thức tính dân số Việt Nam sau $n$ năm kể từ năm $2020$.
	\end{enumerate}
	\dapso{$\approx 98{,}7$ triệu người}
	\loigiai{
		\begin{enumerate}
			\item Sau 1 năm, dân số của Việt Nam sẽ là
			\allowdisplaybreaks
			\begin{eqnarray*}
				u_1&=&97{,}6+97{,}6 \cdot 0{,}0114=97{,}6 \cdot(1+0{,}0114)\\
				&=&97{,}6 \cdot 1{,}0114 \approx 98{,}7 (\text{triệu người}).
			\end{eqnarray*} 
			\item Gọi $u_n$ là dân số của Việt Nam sau $n$ năm.\\
			Do tỉ lệ tăng dân số hàng năm là $1{,}14 \%$ nên ta có
			\allowdisplaybreaks
			\begin{eqnarray*}
				u_n &=&u_{n-1}+u_{n-1} \cdot 0{,}0114=u_{n-1} \cdot(1+0{,}0114) \\
				&=&u_{n-1} \cdot 1{,}0114\ \text{với}\ n \geq 2.
			\end{eqnarray*}
			Do đó, $\left(u_n\right)$ là cấp số nhân có số hạng đầu $u_1=97{,}6\cdot 1{,}0114$, công bội $q=1{,}0114$.\\
			Vậy dân số của Việt Nam sau $n$ năm kể từ năm $2020$ là
			\[u_n=97{,}6 \cdot 1{,}0114 \cdot 1{,}0114^{n-1}=97{,}6 \cdot 1{,}0114^n\ (\text{triệu người}). \]
		\end{enumerate}
	}
\end{bt}

\begin{bt}%[1D2H3-8]
% nguon: Archive/2024/TOAN_11_XE/TOAN_11_XE/data4/11C2/KNTT-2.3.tex #541
%[DCHT Toán 11 - KNTT - Dung Phuong] 
	Bác Linh gửi vào ngân hàng $100$ triệu đồng tiền tiết kiệm với hình thức lãi kép, kì hạn 1 năm với lãi suất $6 \% /$năm. Viết công thức tính số tiền (cả gốc và lãi) mà bác Linh có được sau $n$ năm (giả sử lãi suất không thay đổi qua các năm).
	\dapso{$100 \cdot 1{,}06^{n-1}$ triệu đồng}
	\loigiai{
		Gọi $u_n$ là số tiền (cả gốc lẫn lãi) mà bác Linh có được sau $n$ năm.\\
		Do lãi suất 1 năm là $6\%$ nên ta có
		\allowdisplaybreaks
		\begin{eqnarray*}
			u_n &=&u_{n-1}+u_{n-1} \cdot 0{,}06=u_{n-1} \cdot(1+0{,}06) \\
			&=&u_{n-1} \cdot 1{,}06\ \text{với}\ n \geq 2.
		\end{eqnarray*}
		Do đó, $\left(u_n\right)$ là cấp số nhân có số hạng đầu $u_1=100$ (triệu đồng), công bội $q=1{,}06$.\\
		Vậy số tiền mà bác Linh có được sau $n$ năm là
		\[u_n=100 \cdot 1{,}06^{n-1}\ (\text{triệu đồng}). \]
	}
\end{bt}

\begin{bt}%[1D2V3-6]
% nguon: Archive/2024/TOAN_11_XE/TOAN_11_XE/data4/11C2/KNTT-2.3.tex #542
%[DCHT Toán 11 - KNTT - Dung Phuong] 
	Một hình vuông màu vàng có cạnh $1$ đơn vị dài được chia thành chín hình vuông nhỏ hơn và hình vuông ở chính giữa được tô màu xanh như Hình $2.1$ Mỗi hình vuông màu vàng nhỏ hơn lại được chia thành chín hình vuông con, và mỗi hình vuông con ở chính giữa lại được tô màu xanh. Nếu quá trình này được tiếp tục lặp lại năm lần, thì tổng diện tích các hình vuông được tô màu xanh là bao nhiêu?
	\dapso{$\dfrac{26281}{39366}$}
	\begin{center}
		\begin{tikzpicture}[>=stealth,thick,scale=0.7]
			\def\n{1}
			\def\a{3}
			\pgfmathsetmacro{\m}{int(3^(\n))}
			\def\hv#1{
				\ifnum#1>0
				\fill[blue!50] (-\a/3,-\a/3) rectangle (\a/3,\a/3);
				\pgfmathtruncatemacro{\k}{#1-1}
				\foreach \i in {0,...,3}{\begin{scope}[shift={(90*\i:2)},scale=1/3]\hv{\k}\end{scope}}
				\foreach \i in {0,...,3}{\begin{scope}[shift={(45+90*\i:{4/sqrt(2)})},scale=1/3]\hv{\k}\end{scope}}
				\fi
			}
			\fill[orange!50] (-\a,-\a) rectangle (\a,\a);
			\hv{\n}
			\foreach \i in {0,1,...,\m}{
				\draw[blue!50] 
				({-\a+2*\i *\a/\m},\a)--++(270:2*\a)
				(\a,{-\a+2*\i *\a/\m})--++(180:2*\a)
				;
			}
		\end{tikzpicture}
		\begin{tikzpicture}[>=stealth,thick,scale=0.7]
			\def\n{2}
			\def\a{3}
			\pgfmathsetmacro{\m}{int(3^(\n))}
			\def\hv#1{
				\ifnum#1>0
				\fill[blue!50] (-\a/3,-\a/3) rectangle (\a/3,\a/3);
				\pgfmathtruncatemacro{\k}{#1-1}
				\foreach \i in {0,...,3}{\begin{scope}[shift={(90*\i:2)},scale=1/3]\hv{\k}\end{scope}}
				\foreach \i in {0,...,3}{\begin{scope}[shift={(45+90*\i:{4/sqrt(2)})},scale=1/3]\hv{\k}\end{scope}}
				\fi
			}
			\fill[orange!50] (-\a,-\a) rectangle (\a,\a);
			\hv{\n}
			\foreach \i in {0,1,...,\m}{
				\draw[blue!50] 
				({-\a+2*\i *\a/\m},\a)--++(270:2*\a)
				(\a,{-\a+2*\i *\a/\m})--++(180:2*\a)
				;
			}
		\end{tikzpicture}
		\begin{tikzpicture}[>=stealth,thick,scale=0.7]
			\def\n{4}
			\def\a{3}
			\def\hv#1{
				\ifnum#1>0
				\fill[blue!50] (-\a/3,-\a/3) rectangle (\a/3,\a/3);
				\pgfmathtruncatemacro{\k}{#1-1}
				\foreach \i in {0,...,3}{\begin{scope}[shift={(90*\i:2)},scale=1/3]\hv{\k}\end{scope}}
				\foreach \i in {0,...,3}{\begin{scope}[shift={(45+90*\i:{4/sqrt(2)})},scale=1/3]\hv{\k}\end{scope}}
				\fi
			}
			\fill[orange!50] (-\a,-\a) rectangle (\a,\a);
			\hv{\n}
		\end{tikzpicture}
	\end{center}
	\loigiai{
		Lần phân chia thứ nhất, $1$ hình vuông thành $9$ hình vuông con, diện tích hình vuông tô màu xanh là $u_1=\dfrac{1}{9}$.\\
		Lần phân chia thứ hai, $8$ hình vuông thành $9$ hình vuông con, diện tích hình vuông tô màu xanh tăng thêm là $u_2=\dfrac{1}{9}\left(\dfrac{8}{9}\right)$.\\
		Lần phân chia thứ ba, $8^2$ hình vuông thành $9$ hình vuông con, diện tích hình vuông tô màu xanh tăng thêm là $u_3=\dfrac{1}{9}\left(\dfrac{8}{9}\right)^2$.\\
		Lần phân chia thứ tư, $8^3$ hình vuông thành $9$ hình vuông con, diện tích hình vuông tô màu xanh tăng thêm là $u_4=\dfrac{1}{9}\left(\dfrac{8}{9}\right)^3$.\\
		Lần phân chia thứ năm, $8^4$ hình vuông thành $9$ hình vuông con, diện tích hình vuông tô màu xanh tăng thêm là $u_5=\dfrac{1}{9}\left(\dfrac{8}{9}\right)^4$.\\
		Như vậy diện tích các hình vuông tăng thêm sau mỗi lần chia  tạo thành cấp số nhân có công bội là $q=\dfrac{8}{9}$, số hạng đầu là $u_1=\dfrac{1}{9}$.\\
		Do đó, tổng diện tích hình vuông tô màu xanh sau $5$ lần chia là\\
		\[u_1+u_2+u_3+u_4+u_5=\dfrac{1-q^5}{1-q}\cdot u_1=\dfrac{1-\left(\dfrac{8}{9}\right)^5}{1-\dfrac{8}{9}}\cdot \dfrac{1}{9}=\dfrac{26281}{39366}.\]	
		
	}
\end{bt}

\begin{bt}%[1D2H3-8]
% nguon: Archive/2024/TOAN_11_XE/TOAN_11_XE/data4/11C2/KNTT-2.3.tex #543
 %[DCHT Toán 11 - KNTT - Dung Phuong] 
	Một khay nước có nhiệt độ $23^\circ$ được đặt vào ngăn đá của tủ lạnh. Biết sau mỗi giờ, nhiệt độ của nước giảm $20\%$. Tính nhiệt độ của khay nước đó sau $6$ giờ theo đơn vị độ $C$.
	\dapso{$ \approx 7,5^\circ $ gam}
	\loigiai{
		Nhiệt độ sau mỗi giờ của khay nước theo thứ tự lập thành cấp số nhân với $u_1=23$ và $q=(1-20\%)$.\\
		Ta có $u_6=u_1.q^5=23.(1-20\%)^5 \approx 7,5$.\\
		Nhiệt độ của khay nước sau $6$ giờ là $ \approx 7,5^\circ $.
	}
	
\end{bt}

\begin{bt}%[1D2H3-8]
% nguon: Archive/2024/TOAN_11_XE/TOAN_11_XE/data4/11C2/KNTT-2.3.tex #544
%[DCHT Toán 11 - KNTT - Dung Phuong] 
	Chu kì bán rã của nguyên tố phóng xạ poloni 210 là 138 ngày, nghĩa là sau 138 ngày, khối lượng của nguyên tố đó chi còn một nửa (theo: https://vi.wikipedia.org/wiki/ Poloni-210). Tính khối lượng còn lại của 20 gam poloni 210 sau:
	\begin{listEX}[2]
		\item [a)]  690 ngày;
		\item [b)] 7314 ngày (khoảng 20 năm).
	\end{listEX}
	\dapso{$\dfrac{20}{2^{53}}$ gam}
	\loigiai{
		\begin{listEX}[1]
			\item [a)] Ta có $\dfrac{690}{138}=5$ suy ra khối lượng còn lại sau 690 này là $\dfrac{20}{2^5}=0{,}625$ gam;
			\item [b)] Ta có $\dfrac{7314}{138}=53$ suy ra khối lượng còn lại sau 7314 này là $\dfrac{20}{2^{53}}$ gam.
		\end{listEX}
	}
\end{bt}

\begin{bt}%[1D2V3-8]
% nguon: Archive/2024/TOAN_11_XE/TOAN_11_XE/data4/11C2/KNTT-2.3.tex #545
%[DCHT Toán 11 - KNTT - Dung Phuong] 
	Tế bào E.Coli trong điều kiện nuôi cấy thích hợp cứ $20$ phút lại phân đôi một lần. Hỏi sau $24$ giờ, tế bào ban đầu sẽ phân chia thành bao nhiếu tế bào?
	\dapso{$2^{72}$}
	\loigiai{
		Lần phân chia thứ nhất, $1$ tế bào thành $2$ tế bào, số tế bào lần $1$ phân chia là $u_1 = 2$.\\
		Lần phân chia thứ hai $2$, số tế bào lần $2$ phân chia là  $u_2=2\cdot 2 = u_1 \cdot 2$.\\
		Lần phân chia thứ $3$ có  $4$ tế bào phân chia, số tế bào lần $3$ phân chia là $u_3=2\cdot u_2$.\\
		Như vậy một tế bào phân đôi sẽ tạo thành cấp số nhân có công bội là $2$, số hạng đầu là $u_1=2$.\\
		Sau $n$ lần phân chia từ một tế bào phân được thành $u_n=2^{n-1}u_1$.\\
		Đổi $24$ giờ $=24 \cdot 60 =  72 \cdot 20$ (phút)  $\Rightarrow 24$ giờ gấp $72$ lần $20$ phút. \\
		Do đó, sau $24$ giờ số tế bào nhận được là $u_{72}=2^{71}\cdot 2 = 2^{72}$ (tế bào).
	}
\end{bt}

\begin{bt}%[1D2H3-3]
% nguon: Archive/2024/TOAN_11_XE/TOAN_11_XE/data4/11C2/KNTT-2.3.tex #546

	Xác định công bội, số hạng thứ $5$, số hạng tổng quát và số hạng thứ $100$ của mỗi cấp số nhân sau:
	\begin{enumEX}{2}
	\item $1$, $4$, $16$, $\ldots$;
	\item $2$, $-\dfrac{1}{2}$, $\dfrac{1}{8}$, $\ldots$
	\end{enumEX}
	\loigiai{
	\begin{enumerate}
	\item Ta có $u_1=1$, $u_2=4$, $u_3=16$, $u_4=16\cdot 4=64$, $u_5=64\cdot 4=256$.\\
	Ta có $u_1=1$, $q=4$. \\
	Số hạng tổng quát của cấp số nhân là $u_n=u_1\cdot q^{n-1}=1\cdot 4^{n-1}$, với $n\ge 2$.\\
	Số hạng thứ $100$ là $u_{100}=1\cdot 4^{100-1} =4^{99}$.
	\item Ta có $u_1=2$, $u_2=-\dfrac{1}{2}$, $u_3=\dfrac{1}{8}$, $u_4=\dfrac{1}{8}\cdot \left(-\dfrac{1}{4} \right) =-\dfrac{1}{32}$, $u_5=-\dfrac{1}{32}\cdot \left(-\dfrac{1}{4} \right)=\dfrac{1}{128}$.\\
	Ta có $u_1=2$, $q=-\dfrac{1}{4}$. \\
	Số hạng tổng quát của cấp số nhân là $u_n=u_1\cdot q^{n-1}=2\cdot\left( -\dfrac{1}{4} \right) ^{n-1}$, với $n\ge 2$.\\
	Số hạng thứ $100$ là $u_{100}=2\cdot \left( -\dfrac{1}{4} \right) ^{100-1} =-\left(\dfrac{1}{2} \right) ^{197}$.
	\end{enumerate}
	}
\end{bt}

\begin{bt}%[1D2H2-2]
% nguon: Archive/2024/TOAN_11_XE/TOAN_11_XE/data4/11C2/KNTT-2.3.tex #547

	Viết năm số hạng đầu của mỗi dãy số $\left(u_{n}\right)$ sau và xem nó có phải là cấp số nhân không. Nếu nó là cấp số nhân, hãy tìm công bội $q$ và viết công thức số hạng tổng quát của nó dưới dạng $u_{n}=u_{1} \cdot q^{n-1}$.
	\begin{enumEX}{2}
	\item $u_{n}=5 n$;
	\item $u_{n}=5^{n}$;
	\item $u_{1}=1$, $u_{n}=n u_{n-1}$;
	\item $u_{1}=1$, $u_{n}=5 u_{n-1}$.
	\end{enumEX}
	\loigiai{
	\begin{enumerate}
	\item Ta có $u_1=5$, $u_2=10$, $u_3=15$, $u_4=20$, $u_5=25$.\\
	Đây không là một cấp số nhân.
	\item Ta có $u_1=5$, $u_2=25$, $u_3=125$, $u_4=625$, $u_5=3125$.\\
	Đây là cấp số nhân có $q=5$, số hạng tổng quát $u_n=u_1\cdot q^{n-1}=5\cdot 5^{n-1}$, với $n\ge 2$.
	\item Ta có $u_1=1$, $u_2=2$, $u_3=6$, $u_4=12$, $u_5=60$.\\
	Đây không là một cấp số nhân.
	\item Ta có $u_1=1$, $u_2=5$, $u_3=25$, $u_4=125$, $u_5=625$.\\
	Đây là cấp số nhân có $q=5$, số hạng tổng quát $u_n=u_1\cdot q^{n-1}=1\cdot 5^{n-1}=5^{n-1}$, với $n\ge 2$.
	\end{enumerate}
	}
\end{bt}

\begin{bt}%[1D2H3-4]
% nguon: Archive/2024/TOAN_11_XE/TOAN_11_XE/data4/11C2/KNTT-2.3.tex #548

	Một cấp số nhân có số hạng thứ $6$ bằng $96$ và số hạng thứ $3$ bằng $12$. Tìm số hạng thứ $50$ của cấp số nhân này.
	\loigiai{
	Số hạng tổng quát của cấp nhân là $u_n=u_1\cdot q^{n-1}$, với $n\ge 2$.\\
	Theo bài toán, ta có 
	$$\heva{&u_6=u_1\cdot q^5=96\\&u_3=u_1\cdot q^2=12}\Leftrightarrow \heva{&q^3=8\\&u_1\cdot q^2=12}\Leftrightarrow \heva{&q=2\\&u_1\cdot 2^2=12}\Leftrightarrow \heva{&q=2\\&u_1=3.}$$
	Vậy số hạng tổng quát của cấp số nhân là $u_n=3\cdot 2^{n-1}$, với $n\ge 2$.\\
	Số hạng thứ $50$ là $u_{50}=3\cdot 2^{49}$.
	}
\end{bt}

\begin{bt}%[1D2H3-6]
% nguon: Archive/2024/TOAN_11_XE/TOAN_11_XE/data4/11C2/KNTT-2.3.tex #549

	Một cấp số nhân có số hạng đầu bằng $5$ và công bội bằng $2$. Hỏi phải lấy tổng của bao nhiêu số hạng đầu của cấp số nhân này để có tổng bằng $5115$?
	\loigiai{
	Số hạng tổng quát của cấp số nhân là $u_n=5\cdot 2^{n-1}$, với $n\ge 2$.\\
	Tổng của $n$ số hạng đầu tiên của cấp số nhân là $S_n=u_1\dfrac{q^n-1}{q-1}=5\cdot\dfrac{2^n-1}{2-1}=5\cdot\left( 2^n-1\right) $.\\
	Theo bài toán, ta có $S_n=5\cdot\left( 2^n-1\right)=5115\Leftrightarrow 2^n=1024\Leftrightarrow n=10$.\\
	Vậy phải lấy tổng $10$ số hạng đầu mới thỏa yêu cầu bài toán.
	}
\end{bt}

\begin{bt}%[1D2V3-8]
% nguon: Archive/2024/TOAN_11_XE/TOAN_11_XE/data4/11C2/KNTT-2.3.tex #550
 
	Một công ty xây dựng mua một chiếc máy ủi với giá $3$ tỉ đồng. Cứ sau mỗi năm sử dụng, giá trị của chiếc máy ủi này lại giảm $20 \%$ so với giá trị của nó trong năm liền trước đó. Tìm giá trị còn lại của chiếc máy ủi đó sau $5$ năm sử dụng.
	\loigiai{
	Gọi $(u_n)$ là dãy số biểu diễn giá trị của chiếc máy ủi theo từng năm. \\
	Dãy số này là một cấp số nhân có $u_1=3$, $q=0{,}2$.\\
	Số hạng tổng quát của cấp số nhân này là $u_n=3\cdot 0{,}2^{n-1}$.\\
	Ta có $u_5=3\cdot 0{,}2^4=\dfrac{3}{325}=4{,}8\cdot10^{-3}$.\\
	Tương ứng giá trị của chiếc máy ủi sau $5$ năm là $480$ triệu đồng.
	}
\end{bt}

\begin{bt}%[1D2V3-8]
% nguon: Archive/2024/TOAN_11_XE/TOAN_11_XE/data4/11C2/KNTT-2.3.tex #551
 
	Vào năm 2020, dân số của một quốc gia là khoảng $97$ triệu người và tốc độ tăng trưởng dân số là $0{,}91 \%$. Nếu tốc độ tăng trưởng dân số này được giữ nguyên hằng năm, hãy ước tính dân số của quốc gia đó vào năm 2030.
	\loigiai{
	Dân số năm 2021 tăng lên so với năm 2020 là $97 \cdot 0{,}91 \% $ triệu người.\\
	Dân số năm 2021 là 
	\begin{center}
	$97 + 97 \cdot 0{,}91 \% = 97\cdot (1+0{,}91 \%)$ triệu người.
	\end{center}
	Dân số năm 2022 tăng lên so với năm 2021 là $97\cdot (1+0{,}91 \%)\cdot 0{,}91 \% $ triệu người.\\
	Dân số năm 2022 là 
	\begin{center}
	$97\cdot (1+0{,}91 \%) + 97\cdot (1+0{,}91 \%) \cdot 0{,}91 \% = 97\cdot (1+0{,}91 \%)^2$ triệu người.
	\end{center}
	Dân số năm 2023 tăng lên so với năm 2021 là $97\cdot (1+0{,}91 \%)^2\cdot 0{,}91 \% $ triệu người.\\
	Dân số năm 2023 là
	\begin{center}
	$97\cdot (1+0{,}91 \%)^2 + 97\cdot (1+0{,}91 \%)^2\cdot 0{,}91 \% = 97\cdot (1+0{,}91 \%)^3$ triệu người.
	\end{center}
	Tương tự vậy ta có dân số năm 2030 là $97\cdot (1+0{,}91 \%)^{10} = 106{,}1973784$ triệu người.
	}
\end{bt}

\begin{bt}%[1D2V3-8]
% nguon: Archive/2024/TOAN_11_XE/TOAN_11_XE/data4/11C2/KNTT-2.3.tex #552
 
	Một loại thuốc được dùng mỗi ngày một lần. Lúc đầu nồng độ thuốc trong máu của bệnh nhân tăng nhanh, nhưng mỗi liều kế tiếp có tác dụng ít hơn liều trước đó. Lượng thuốc trong máu ở ngày thứ nhất là $50 \,\mathrm{mg}$, và mỗi ngày sau đó giảm chỉ còn một nửa so với ngày kề trước đó. Tính tổng lượng thuốc (tính bằng $\mathrm{mg}$) trong máu của bệnh nhân sau khi dùng thuốc $10$ ngày liên tiếp.
	\loigiai{
	Gọi $(u_n)$ là dãy số biểu diễn giá trị của lượng thuốc trong máu của bệnh nhân theo từng ngày. \\
	Dãy số này là một cấp số nhân có $u_1=50$, $q=\dfrac{1}{2}$.\\
	Tổng của $n$ số hạng đầu tiên của cấp số nhân là $S_n=u_1\dfrac{1-q^n}{1-q}$.\\
	Theo bài toán, ta có $S_{10}=50 \cdot\dfrac{1-\left(\dfrac{1}{2}\right)^{10}}{1-\dfrac{1}{2}} \approx 99{,}902$.\\
	Vậy tổng lượng thuốc trong máu của bệnh nhân sau khi dùng thuốc $10$ ngày liên tiếp là $99{,}902$ mg.	
	}
\end{bt}

\begin{bt}%[1D2H3-3]
% nguon: Archive/2024/TOAN_11_XE/TOAN_11_XE/data4/11C2/KNTT-2.3.tex #554
%[Dự án Tbt hóa SGK 11, Cánh Diều, 2023]%[Tiêu Khánh Văn]
	Chứng minh mỗi dãy số $\left(u_n\right)$ với số hạng tổng quát như sau là cấp số nhân:
	\begin{enumEX}[a)]{3} 
		\item $u_n=\dfrac{-3}{4}\cdot 2^n$
		\item $u_n=\dfrac{5}{3^n}$;
		\item $u_n=(-0{,}75)^n$.
	\end{enumEX} 
	\loigiai{
		\begin{enumEX}[a)]{1} 
			\item Với mọi $n\in\mathbb{N}^*$, ta có 
			$$\begin{aligned}
				u_{n+1}&=\dfrac{-3}{4}\cdot 2^{n+1}\\
				&=\dfrac{-3}{4}\cdot 2^n\cdot2\\
				&=2u_n.
			\end{aligned}$$
			Vậy dãy số $\left(u_n\right)$ đã cho là cấp số nhân.
			\item  Với mọi $n\in\mathbb{N}^*$, ta có 
			$$\begin{aligned}
				u_{n+1}&=\dfrac{5}{3^{n+1}}\\
				&=\dfrac{5}{3^n\cdot 3}\\
				&=\dfrac{1}{3}u_n.
			\end{aligned}$$
			Vậy dãy số $\left(u_n\right)$ đã cho là cấp số nhân.
			\item Với mọi $n\in\mathbb{N}^*$, ta có 
			$$\begin{aligned}
				u_{n+1}&=(-0{,}75)^{n+1}\\
				&=(-0{,}75)^n\cdot (-0{,}75)\\
				&=(-0{,}75)u_n.
			\end{aligned}$$
			Vậy dãy số $\left(u_n\right)$ đã cho là cấp số nhân.
		\end{enumEX} 
	}
\end{bt}

\begin{bt}%[1D2H3-2]
% nguon: Archive/2024/TOAN_11_XE/TOAN_11_XE/data4/11C2/KNTT-2.3.tex #555
%[Dự án Tex hóa SGK 11, Cánh Diều, 2023]%[Tiêu Khánh Văn]
	Cho cấp số nhân $\left(u_n\right)$ với số hạng đầu $u_1=-5$, công bội $q=2$.
	\begin{enumEX}[a)]{1} 
		\item Tìm $u_9$.
		\item Số $-320$ là số hạng thứ bao nhiêu của cấp số nhân trên?
		\item Số $160$ có phải là một số hạng của cấp số nhân trên không? 
	\end{enumEX} 
	\loigiai{
		\begin{enumEX}[a)]{1} 
			\item Ta có $u_9=u_1\cdot q^8 = -5\cdot 2^8=-1280$.
			\item Giả sử số $-320$ là số hạng thứ $n$ của cấp số nhân trên. Khi đó ta có 
			\begin{eqnarray*}   
				u_n=-320& \Leftrightarrow & u_1\cdot q^{n-1} =-320\\
				& \Leftrightarrow & -5 \cdot 2^{n-1} =-320\\
				& \Leftrightarrow & 2^{n-1} = 64\\
				& \Leftrightarrow & 2^{n-1} = 2^6\\
				& \Leftrightarrow & n-1=6\\
				& \Leftrightarrow & n=7.
			\end{eqnarray*}
			Vậy số $-320$ là số hạng thứ $7$ của cấp số nhân trên.
			\item Giả sử số $160$ là số hạng thứ $n$ của cấp số nhân trên. Khi đó ta có 
			\begin{eqnarray*}   
				u_n=160& \Leftrightarrow & u_1\cdot q^{n-1} =160\\
				& \Leftrightarrow & -5 \cdot 2^{n-1} =160\\
				& \Leftrightarrow & 2^{n-1} = -32.
			\end{eqnarray*} 
			Vì $2^{n-1}>0$, $\forall n\in\mathbb{N}^*$ và $-32<0$ nên không có giá trị nào của $n$ thỏa mãn $u_n=160$.\\
			Vậy số $160$ không phải số hạng nào của cấp số nhân đã cho.
		\end{enumEX} 
	}
\end{bt}

\begin{bt}%[1D2H3-3]
% nguon: Archive/2024/TOAN_11_XE/TOAN_11_XE/data4/11C2/KNTT-2.3.tex #556
%[Dự án Tex hóa SGK 11, Cánh Diều, 2023]%[Tiêu Khánh Văn]
	Cho cấp số nhân $\left(u_n\right)$ với $u_1=3$, $u_3=\dfrac{27}{4}$.
	\begin{enumEX}[a)]{1} 
		\item Tìm công bội $q$ và viết năm số hạng đầu của cấp số nhân trên.
		\item Tính tổng $10$ số hạng đầu của cấp số nhân trên. 
	\end{enumEX} 
	\loigiai{
		\begin{enumEX}[a)]{1} 
			\item Ta có $u_3=\dfrac{27}{4} \Leftrightarrow u_1q^2=\dfrac{27}{4} \Leftrightarrow 3q^2=\dfrac{27}{4} \Leftrightarrow q=\pm\dfrac{3}{2}$.
			\begin{itemize}
				\item Với $q=\dfrac{3}{2}$. Năm số hạng đầu của cấp số nhân đã cho là 
				\begin{center}
					$u_1=3$, $u_2=u_1q=\dfrac{9}{2}$, $u_3=u_2q=\dfrac{27}{4}$, $u_4=u_3q=\dfrac{81}{8}$, $u_5=u_4q=\dfrac{243}{16}$.
				\end{center}
				\item Với $q=-\dfrac{3}{2}$. Năm số hạng đầu của cấp số nhân đã cho là 
				\begin{center}
					$u_1=3$, $u_2=u_1q=-\dfrac{9}{2}$, $u_3=u_2q=\dfrac{27}{4}$, $u_4=u_3q=-\dfrac{81}{8}$, $u_5=u_4q=\dfrac{243}{16}$.
				\end{center}
			\end{itemize}
			\item 
			\begin{itemize}
				\item Với $q=\dfrac{3}{2}$. Tổng $10$ số hạng đầu của cấp số nhân đã cho là
				$$S_{10}=\dfrac{u_1(1-q^{10})}{1-q}=\dfrac{3\left[1-\left(\dfrac{3}{2}\right)^{10}\right]}{1-\dfrac{3}{2}}=\dfrac{174075}{512}\approx 340.$$
				\item Với $q=-\dfrac{3}{2}$. Tổng $10$ số hạng đầu của cấp số nhân đã cho là
				$$S_{10}=\dfrac{u_1(1-q^{10})}{1-q}=\dfrac{3\left[1-\left(-\dfrac{3}{2}\right)^{10}\right]}{1+\dfrac{3}{2}}=-\dfrac{34815}{512}\approx -68.$$
			\end{itemize}
		\end{enumEX}
	}
\end{bt}

\begin{bt}%[1D2V3-8]
% nguon: Archive/2024/TOAN_11_XE/TOAN_11_XE/data4/11C2/KNTT-2.3.tex #557
%[Dự án Tex hóa SGK 11, Cánh Diều, 2023]%[Tiêu Khánh Văn]
	Một tỉnh có $2$ triệu dân vào năm 2020 với tỉ lệ tăng dân số là $1$ \%/năm. Gọi $u_n$ là số dân của tỉnh đó sau $n$ năm. Giả sử tỉ lệ tăng dân số là không đổi.
	\begin{enumEX}[a)]{1} 
		\item Viết công thức tính số dân của tỉnh đó sau $n$ năm kể từ năm 2020.
		\item Tính số dân của tỉnh đó sau $10$ năm kể từ năm 2020.  
	\end{enumEX}
	\loigiai{
		\begin{enumEX}[a)]{1} 
			\item Với $u_n$ là số dân của tỉnh đó sau $n$ năm. \\
			Ta có $u_1=2$ (triệu dân).\\
			$u_{n+1}=u_n+u_n\cdot 0{,}01 = 1{,}01u_n$. \\
			Do đó, $(u_n)$ là cấp số nhân với số hạng đầu $u_1=2$ và công bội $q=1{,}01$. \\
			Vậy công thức tính số dân của tỉnh đó sau $n$ năm là $u_n=u_1q^{n-1}\Rightarrow u_n=2\cdot 1{,}01^{n-1}$.
			\item Số dân của tỉnh đó sau $10$ năm kể từ năm 2020 là $u_{10}=2\cdot 1{,}01^9 = 2{,}187$ (triệu dân).
		\end{enumEX}
	}
\end{bt}

\begin{bt}%[1D2V3-8]
% nguon: Archive/2024/TOAN_11_XE/TOAN_11_XE/data4/11C2/KNTT-2.3.tex #558
%[Dự án Tex hóa SGK 11, Cánh Diều, 2023]%[Tiêu Khánh Văn]
	Một gia đình mua một chiếc ô tô giá $800$ triệu đồng. Trung bình sau mỗi năm sử dụng, giá trị còn lại của ô tô giảm đi $4 \%$ (so với năm trước đó).
	\begin{enumEX}[a)]{1} 
		\item Viết công thức tính giá trị của ô tô sau $1$ năm, $2$ năm sử dụng.
		\item Viết công thức tính giá trị của ô tô sau $n$ năm sử dụng.
		\item Sau $10$ năm, giá trị của ô tô ước tính còn bao nhiêu triệu đồng? 
	\end{enumEX} 
	\loigiai{
		Gọi $u_n$ là giá trị còn lại của ô tô sau $n$ năm sử dụng. 
		\begin{enumEX}[a)]{1} 
			\item Giá trị của ô tô sau $1$ năm sử dụng là $u_1=800-800\cdot0{,}04=800\cdot0{,}96=768$ triệu đồng.\\
			Giá trị của ô tô sau $2$ năm sử dụng là $u_2=u_1-u_1\cdot0{,}04=u_1\cdot0{,}96=737{,}28$ triệu đồng.
			\item Ta có $u_n=u_{n-1}-u_{n-1}\cdot0{,}04=u_{n-1}\cdot0{,}96$. \\
			Do đó, $(u_n)$ là cấp số nhân với số hạng đầu $u_1=768$ và công bội $q=0{,}96$. \\
			Vậy sau $n$ năm sử dụng, giá trị còn lại của chiếc ô tô là $u_n=u_1q^{n-1}\Rightarrow u_n=768\cdot0{,}96^{n-1}$.
			\item Sau $10$ năm, ước tính giá trị của ô tô còn lại là $u_{10}=768\cdot0{,}96^9\approx 531{,}87$ triệu đồng.
		\end{enumEX} 
	}
\end{bt}

\begin{bt}%[1D2V3-8]
% nguon: Archive/2024/TOAN_11_XE/TOAN_11_XE/data4/11C2/KNTT-2.3.tex #559
%[Dự án Tex hóa SGK 11, Cánh Diều, 2023]%[Tiêu Khánh Văn]
	\immini
	{
		Một người nhảy bungee (một trò chơi mạo hiểm mà người chơi nhảy từ một nơi có địa thế cao xuống với dây đai an toàn buộc xung quanh người) từ một cây cầu và căng một sợi dây dài $100$ m. Sau mỗi lần rơi xuống, nhờ sự đàn hồi của dây, người nhảy được kéo lên một quãng đường có độ dài bằng $75$\% so với lần rơi trước đó và lại bị rơi xuống đúng bằng quãng đường vừa được kéo lên. Tính tổng quãng đường người đó đi được sau $10$ lần kéo lên và lại rơi xuống. 
	}
	{
		\begin{tikzpicture}[xscale=.5, font=\footnotesize, line join=round, line cap=round, >=stealth,yscale=1]
			\def\a{-0.12} % Hệ số a phải khác 0
			\def\b{0.86}
			\def\c{0}
			\def\m{-0.12} % Hệ số a phải khác 0
			\def\n{0.86}
			\def\p{-0.3}
			\clip (-2,-2)rectangle(9,3);
			\fill[red!40] (-2,1.8)--(9,1.8)--(9,1.55)--(-2,1.55)--cycle;%Mặt phẳng của cây cầu
			\fill[green!50] (-2,1.55)--(0,0)--(1.5,-1)--(-1,-2)--(-2,-2)--cycle;
			\fill[green!50] (9,1.55)--(7,0)--(5.5,-1)--(8,-2)--(9,-2)--cycle;
			\fill[blue!40] (-1,-2)--(8,-2)--(5.5,-1)--(1.5,-1)--cycle;
			\draw[color=blue!50,line width=2pt,<->] (2,-.5)--(2,1.5)node[left,midway]{$100$ m};	
			\draw[color=blue!50,line width=2pt,<->] (5,-.5)--(5,1)node[midway,right]{$0{,}75\cdot100$ m};
			\draw(3.5,1.55)--(3.5,-.25)node[rotate=200]{\faChild};
			\draw[color=white,<->] (3.5,1.8)--(3.5,1.55); 
		\end{tikzpicture}
	}
	\loigiai{
		Gọi $u_n$ là quãng đường người đó được kéo lên ở lần thứ $n$ được kéo lên và lại rơi xuống (đơn vị tính: mét). \\
		Ta có $u_1=0{,}75\cdot100=100\cdot1{,}5=75$ m và $u_n=0{,}75\cdot u_{n-1}$. \\
		Vậy $(u_n)$ là cấp số nhân với số hạng đầu $u_1=75$ và công bội $q=0{,}75$. \\
		Tổng quãng đường người đó đi được sau $10$ lần kéo lên và lại rơi xuống là 
		$$\begin{aligned}
			S&=100+2u_1+2u_2+\cdots+2u_{10}\\
			&=100+2S_{10}
			=100+2\cdot\dfrac{75\left(1-0{,}75^{10}\right)}{1-0{,}75}\\
			&\approx666{,}2 \text{ m}.
		\end{aligned}$$ 
	}
\end{bt}

\begin{bt}%[1D2N2-4]
% nguon: Archive/2024/TOAN_11_XE/TOAN_11_XE/data4/11C2/KNTT-2.2.tex #560

Cho cấp số $\left(u_n\right)$ có số hạng đầu $u_1=2$ và công sai $d=3$. Hãy viết năm số hạng đầu của cấp số cộng này.
\loigiai{Năm số hạng đầu của cấp số cộng này là:
\[u_1=2, u_2=u_1+3=5, u_3=u_2+3=8, u_4=u_3+3=11, u_5=u_4+3=14.\]}
\end{bt}

\begin{bt}%[1D2N2-2]
% nguon: Archive/2024/TOAN_11_XE/TOAN_11_XE/data4/11C2/KNTT-2.2.tex #561

Cho dãy số $\left(u_n\right)$ với $u_n=5n-1$. Chứng minh rằng $\left(u_n\right)$ là một cấp số cộng. Tìm số hạng đầu và công sai $d$ của nó.
\loigiai{Ta có
\[u_n-u_{n-1}=5n-1-\left(5(n-1)-1\right)=5,\]
với mọi $n\ge 2$. Do đó, $\left(u_n\right)$ là cấp số cộng có số hạng đầu là $u_1=5-1=4$ và công sai $d=5$.}
\end{bt}

\begin{bt}%[1D2N2-2]
% nguon: Archive/2024/TOAN_11_XE/TOAN_11_XE/data4/11C2/KNTT-2.2.tex #562

Cho dãy số $\left(u_n\right)$ với $u_n=-2n+3$. Chứng minh rằng $\left(u_n\right)$ là một cấp số cộng. Tìm số hạng đầu và công sai $d$ của cấp số cộng này.
\loigiai{
	Ta có 
	\[u_n-u_{n-1}=-2n+3-(-2(n-1)+3)=-2, \text{ với mọi } n\ge 2.\]
	Do đó, $(u_n)$ là cấp số cộng có số hạng đầu là $u_1=-2\cdot 1+3=1$ và công sai $d=-2$.
}
\end{bt}

\begin{bt}%[1D2H2-4]
% nguon: Archive/2024/TOAN_11_XE/TOAN_11_XE/data4/11C2/KNTT-2.2.tex #563

Tìm năm số hạng đầu và số hạng thứ $100$ của cấp số cộng $\left(u_n\right)\colon 10,5,\ldots$
\loigiai{Cấp số cộng này có số hạng đầu là $u_1=10$ và công sai là $d=-5$. Do đó,
\begin{itemize}
	\item Năm số hạng đầu của cấp số cộng là $10,5,0,-5,-10$.
	\item Số hạng thức $100$ là $u_{100}=u_1+99d=10+99\cdot (-5)=-485$.
\end{itemize}}
\end{bt}

\begin{bt}%[1D2H2-4]
% nguon: Archive/2024/TOAN_11_XE/TOAN_11_XE/data4/11C2/KNTT-2.2.tex #564

Số hạng thứ $10$ của một cấp số cộng $\left(u_n\right)$ bằng $48$ và số hạng thứ $18$ bằng $88$. Tìm số hạng thứ $100$ của cấp số cộng đó.
\loigiai{Giả sử $u_1$, $d$ lần lượt là số hạng đầu tiên, công sai của cấp số cộng thỏa mãn yêu cầu bài toán. Khi đó, theo giả thiết, ta có
\[\heva{&u_{10}=u_1+9d=48\\&u_{18}=u_1+17d=88.}\]
Giải hệ này ta được $u_1=3$ và $d=5$.\\
Vậy số hạng thức $100$ của cấp số cộng này là $u_{100}=u_1+99d=3+99\cdot 5=498$.}
\end{bt}

\begin{bt}%[1D2H2-3]
% nguon: Archive/2024/TOAN_11_XE/TOAN_11_XE/data4/11C2/KNTT-2.2.tex #565

Cho dãy số $\left(u_n\right)$ với $u_n=4n-3$. Chứng minh rằng $\left(u_n\right)$ là một cấp số cộng. Tìm số hạng đầu $u_1$ và công sai $d$ của cấp số cộng này. Từ đó viết số hạng tổng quát $u_n$ dưới dạng $u_n=u_1+(n-1)d$.
\loigiai{
	Ta có 
	\[u_n-u_{n-1}=4n-3-(4(n-1)-3)=4,\]
	với mọi $n\ge 2$. Do đó, $(u_n)$ là cấp số cộng có số hạng đầu là $u_1=4\cdot 1-3=1$ và công sai $d=4 $.\\
	Số hạng tổng quát $u_n=u_1+(n-1)d=1+(n-1)\cdot 4.$
}
\end{bt}

\begin{bt}%[1D2H2-6]
% nguon: Archive/2024/TOAN_11_XE/TOAN_11_XE/data4/11C2/KNTT-2.2.tex #566

Cần lấy tổng của bao nhiêu số hạng đầu của cấp số cộng $2$, $5$, $8$, $\ldots$ để được kết quả bằng $345$?
\loigiai{Cấp số cộng này có số hạng đầu là $u_1=2$ và công sai là $d=3$. Gọi $n$ là số các số hạng đầu của cấp số cộng cần lấy tổng, ta có
\[345=S_n=\dfrac{n}{2}\left(2u_1+(n-1)d\right)=\dfrac{n}{2}\left(2\cdot 2+(n-1)\cdot 3\right)=\dfrac{n}{2}\left(3n+1\right).\]
Do đó $3n^2+n-390=0$. Giải phương trình bậc hai này ta được $n=-\dfrac{230}{15}\text{(loại)}$ và $n=15$. Vậy phải lấy tổng của $15$ số hạng đầu của cấp số cộng đã cho để được tổng bằng $345$.}
\end{bt}

\begin{bt}%[1D2H2-7]
% nguon: Archive/2024/TOAN_11_XE/TOAN_11_XE/data4/11C2/KNTT-2.2.tex #567

Anh Nam được nhận vào làm việc ở một công ty về công nghệ với mức lương khởi điểm là $100$ triệu đồng một năm. Công ty sẽ tăng thêm lương cho anh Nam mỗi năm là $20$ triệu đồng. Tính tổng số tiền lương mà anh Nam nhận được sau $10$ năm làm việc cho công ty đó.
\loigiai{
	Số tiền lương của anh Nam qua các năm lập thành một cấp số cộng với số hạng đầu $u_1=100$ và công sai $d=20$.\\
	 Khi đó $u_{10}=u_1+9d=100+9\cdot 20=280$. \\
	  Vậy sau 10 năm làm việc cho công ty số tiền lương mà anh Nam sẽ nhận là $280$ triệu đồng.
}
\end{bt}

\begin{bt}%[1D2N2-4]
% nguon: Archive/2024/TOAN_11_XE/TOAN_11_XE/data4/11C2/KNTT-2.2.tex #568
%[DCHT Toán 11 - KNTT -Lê Hải Phụng] 
	Dãy số hữu hạn nào là một cấp số cộng? Vì sao?
	\begin{listEX}[2]
		\item  $-2$, $1$, $4$, $7$, $10$, $13$, $16$.
		\item  $ 1 $, $ -2 $, $ -4 $, $ -6 $, $ -8 $.
	\end{listEX}
	\dapso{ Dãy số 1 là một cấp số cộng, dãy số 2 không là một cấp số cộng.}
	\loigiai{
		\begin{enumerate}
			\item Ta thấy $ u_2=u_1+3 $  do $ 1=(-2)+3 $.\\
			Vì $ u_k=u_{k-1}+d,\ \forall k\geq2$ $\left(\ 1=\left(-2\right)+3;4=1+3;7=4+3;10=7+3;13=10+3;16=13+3\right) $ nên dãy số đã cho là cấp số cộng. 
			\item Ta thấy $ u_2=u_1+(-3) $  do $-2=1+(-3)$.\\
			Vì $ {u_3\neq u}_2+(-3) $ bởi $ \left(\ -4\neq-2+\left(-3\right)\right)\ $ nên dãy số đã cho không là cấp số cộng.
		\end{enumerate}
	}
\end{bt}

\begin{bt}%[1D2H2-4]
% nguon: Archive/2024/TOAN_11_XE/TOAN_11_XE/data4/11C2/KNTT-2.2.tex #569
%[DCHT Toán 11 - KNTT -Lê Hải Phụng] 
	Trong các dãy số dưới đây, dãy số nào là cấp số cộng? 
	\begin{listEX}[2]
		\item  Dãy số $\left({a_n}\right)$ với ${a_n}=4n-3$;
		\item  Dãy số $\left({c_n}\right)$ với ${c_n}={2018^n}$.
	\end{listEX}
	\dapso{Dãy số 1 là một cấp số cộng, dãy số 2 không là một cấp số cộng.}
	\loigiai{	
		\begin{enumerate}
			\item Ta có $a_{n+1}=4(n+1)-3=4n+1$ nên $a_{n+1}-a_n=(4n+1)-(4n-3)=4$,$\forall n\ge 1.$.\\
			Do đó $(a_n)$ là cấp số cộng.
			\item Ta có $c_{n+1}=2018^{n+1}$ nên $c_{n+1}-c_n=2018^{n+1}-2018^n=2017\cdot 2018^n$ (phụ thuộc vào giá trị của $n$).\\ 
			Suy ra $(c_n)$ không phải là một cấp số cộng.
		\end{enumerate}	
	}
\end{bt}

\begin{bt}%[1D2N2-3]
% nguon: Archive/2024/TOAN_11_XE/TOAN_11_XE/data4/11C2/KNTT-2.2.tex #570
%[DCHT Toán 11 - KNTT -Lê Hải Phụng] 
	Cho cấp số cộng $(u_n)$  có công thức số hạng tổng quát $u_n=3n+1$, $n\in\mathbb{N}^\ast$ . Tìm số hạng đầu $u_1$ và công sai $d$?
	\dapso{$u_1=4 $, $d=3$.}
	\loigiai{
		Từ công thức số hạng tổng quát, ta có $ u_1=4 $, $u_2=7$ suy ra $d=u_2-u_1=3$.
	}
\end{bt}

\begin{bt}%[1D2H2-4]
% nguon: Archive/2024/TOAN_11_XE/TOAN_11_XE/data4/11C2/KNTT-2.2.tex #571
%[DCHT Toán 11 - KNTT -Lê Hải Phụng] 
	Cho cấp số cộng $(u_n)$ với $u_1=3$, $u_2=9$. Công sai của cấp số cộng đã cho bằng bao nhiêu?
	\dapso{$ d=6 $}
	\loigiai{
		Cấp số cộng $(u_n)$ có số hạng tổng quát là $u_n=u_1+(n-1)d$ với $n \ge 2$.\\
		Suy ra $u_2=u_1+d \Leftrightarrow 9=3+d \Leftrightarrow d=6$.\\
		Vậy công sai của cấp số cộng đã cho là $6$.
	}
\end{bt}

\begin{bt}%[1D2V2-4]
% nguon: Archive/2024/TOAN_11_XE/TOAN_11_XE/data4/11C2/KNTT-2.2.tex #572
%[DCHT Toán 11 - KNTT -Lê Hải Phụng] 
	Tính số hạng đầu $u_1$ và công sai $d$ của một cấp số cộng biết $u_4=10$ và $u_7=19$.
	\dapso{$ u_1=1 $, $ d=3 $.}
	\loigiai{Ta có $ \heva{& u_4=10 \\ & u_7=19} \Leftrightarrow \heva{& u_1+3d=10 \\ & u_1+6d=19} \Leftrightarrow \heva{& u_1=1 \\ & d=3.}$}
\end{bt}

\begin{bt}%[1D2H2-3]
% nguon: Archive/2024/TOAN_11_XE/TOAN_11_XE/data4/11C2/KNTT-2.2.tex #573
%[Dự án DCHT-11-KNTT]%[Dao-V- Thuy]
	Xác định số hạng tổng quát của cấp số cộng $(u_n),$ biết $\heva{&u_7=8\\ &d=2.}$
	\dapso{$u_n=2n-6$}
	\loigiai{
		Ta có
		\begin{equation*}
			\heva{&u_7=8\\ &d=2} \Leftrightarrow \heva{&u_1+6d=8\\&d=2} \Leftrightarrow \heva{&u_1=-4\\ &d=2.}
		\end{equation*}
		Vậy công thức tổng quát của cấp số cộng
		\begin{center}
			$u_n=-4+(n-1)2 \Leftrightarrow u_n=2n-6 $ với $n \geq 2.$
		\end{center}	
	}
\end{bt}

\begin{bt}%[1D2H2-4]
% nguon: Archive/2024/TOAN_11_XE/TOAN_11_XE/data4/11C2/KNTT-2.2.tex #574
%[Dự án DCHT-11-KNTT]%[Dao-V- Thuy]
	Tìm số hạng đầu và công sai của cấp số cộng $(u_n)$, biết $\heva{&u_1+u_5-u_3=10\\ &u_1+u_6=17.}$
	\dapso{$u_1=16$, $d=-3$}
	\loigiai{
		Ta có
		\begin{align*}
			\heva{&u_1+u_5-u_3=10\\ &u_1+u_6=17} &\Leftrightarrow \heva{&u_1+u_1+4d-(u_1+2d)=10\\ &u_1+u_1+5d=17}\\ & \Leftrightarrow\heva{&u_1+2d=10 \\ &2u_1+5d=17} \Leftrightarrow \heva{&u_1=16 \\ &d=-3.}
		\end{align*}
		Vậy $u_1=16$, $d=-3$.
	}
\end{bt}

\begin{bt}%[1D2H2-3]
% nguon: Archive/2024/TOAN_11_XE/TOAN_11_XE/data4/11C2/KNTT-2.2.tex #575
%[Dự án DCHT-11-KNTT]%[Dao-V- Thuy]
	Cho cấp số cộng $(u_n)$ với $\heva{&u_1=-9\\ &u_{n-1}=u_n-5}$. Tìm số hạng tổng quát của cấp số cộng $(u_n)$.
	\dapso{$u_n= 5n-14$}
	\loigiai{
		Từ công thức $u_{n-1}=u_n-5 \Leftrightarrow u_n= u_{n-1}+5$, suy ra $d=5$.\\
		Vậy công thức tổng quát của cấp số cộng $(u_n)$ là $u_n=-9 + 5(n-1)= 5n-14$.
	}
\end{bt}

\begin{bt}%[1D2H2-3]
% nguon: Archive/2024/TOAN_11_XE/TOAN_11_XE/data4/11C2/KNTT-2.2.tex #576
%[Dự án DCHT-11-KNTT]%[Dao-V- Thuy]
	Cho cấp số cộng $(u_n)$ có $u_{20}=-52$ và $u_{51}=-145$. Hãy tìm số hạng tổng quát của cấp số cộng đó.
	\dapso{$u_n= -3n+8$}
	\loigiai{
		Ta có
		\begin{eqnarray*}
			\heva{&u_{20}=-52 \\&u_{51}=-145} &\Leftrightarrow& \heva{&u_1+19d=-52\\ &u_1+50d=-145} \Leftrightarrow \heva{&u_1=5 \\ &d= -3.}
		\end{eqnarray*}
		Vậy số hạng tổng quát cần tìm là $u_n= u_1+ (n-1)d= 5+(n-1) \cdot (-3)= -3n+8$.
	}
\end{bt}

\begin{bt}%[1D2H2-4]
% nguon: Archive/2024/TOAN_11_XE/TOAN_11_XE/data4/11C2/KNTT-2.2.tex #577
%[Dự án DCHT-11-KNTT]%[Dao-V- Thuy]
	Tìm số hạng đầu và công sai của cấp số cộng $(u_n)$, biết
	\begin{listEX}[2]
		\item $\heva{&u_9=5u_2\\ &u_{13}=2u_6+5.}$
		\item $\heva{&u_1-u_3+u_5=10\\ &u_1+u_6=7.}$
	\end{listEX}
	\dapso{$u_1=3$, $d=4$; $u_1=36$, $d=-13$}
	\loigiai
	{
		\begin{enumerate}
			\item Ta có
			\begin{eqnarray*}
				\heva{&u_9=5u_2\\ &u_{13}=2u_6+5} &\Leftrightarrow& \heva{&u_1+8d= 5 \left( u_1+d \right) \\ &u_1+12d = 2 \left( u_1+5d\right) + 5}\\
				&\Leftrightarrow& \heva{&-4u_1+3d=0\\ &-u_1+2d=5} \Leftrightarrow \heva{&u_1=3 \\ &d=4.}
			\end{eqnarray*}
			Vậy $u_1=3$, $d=4$.
			\item Ta có 
			\begin{eqnarray*}
				\heva{&u_1-u_3+u_5=10\\ &u_1+u_6=7} &\Leftrightarrow& \heva{&u_1-\left( u_1+2d\right) + \left( u_1+4d\right) = 10\\ &u_1+ \left( u_1+5d\right) = 7}\\
				&\Leftrightarrow& \heva{&u_1+2d=10\\ &2u_1+5d=7} \Leftrightarrow \heva{&u_1=36 \\ &d=-13.}
			\end{eqnarray*}
			Vậy $u_1=36$, $d=-13$.
		\end{enumerate}
	}
\end{bt}

\begin{bt}%[1D2V2-4]
% nguon: Archive/2024/TOAN_11_XE/TOAN_11_XE/data4/11C2/KNTT-2.2.tex #578
%[Dự án DCHT-11-KNTT]%[Dao-V- Thuy]
	Tìm số hạng đầu và công sai của cấp số cộng $(u_n)$, biết
	\begin{listEX}[2]
		\item $\heva{&-u_3+u_7=8\\ &u_2u_7=75.}$
		\item $\heva{&u_5=4u_3\\ &u_2u_6=-11.}$
	\end{listEX}
	\dapso{$\heva{&u_1=3\\ &d=2} \text{ hoặc } \heva{&u_1=-17\\ &d=2}$; $\heva{&u_1=-4\\ &d=3}$ hoặc $\heva{&u_1=4\\ &d=-3}$}
	\loigiai{
		\begin{enumerate}
			\item Ta có
			\begin{eqnarray*}
				\heva{&-u_3+u_7=8\\ &u_2u_7=75} &\Leftrightarrow& \heva{&-\left( u_1+2d\right) + \left( u_1+6d\right) = 8\\ &\left( u_1+d\right) \left( u_1+6d\right) = 75}\\
				&\Leftrightarrow& \heva{&4d=8\\ &u_1^2+7u_1d+6d^2=75}\\
				&\Leftrightarrow& \heva{&d=2\\ &u_1^2+14u_1-51=0}\\
				&\Leftrightarrow& \heva{&u_1=3\\ &d=2} \text{ hoặc } \heva{&u_1=-17\\ &d=2.}
			\end{eqnarray*}
			Vậy $\heva{&u_1=3\\ &d=2} \text{ hoặc } \heva{&u_1=-17\\ &d=2.}$
			\item Ta có 
			\begin{eqnarray*}
				\heva{&u_5=4u_3\\ &u_2u_6=-11} &\Leftrightarrow& \heva{&u_1+4d=4 \left( u_1+2d\right)\\ &\left( u_1+d\right) \left( u_1+5d\right)= -11}\\
				&\Leftrightarrow& \heva{&3u_1+4d=0 &(1)\\ &u_1^2+6du_1+5d^2=-11 &(2)}
			\end{eqnarray*}
			Từ $(1)$ suy ra $3u_1=-4d$. Thay vào $(2)$ ta được
			\begin{eqnarray*}
				9u_1^2+54du_1+45d^2=-99 &\Leftrightarrow& 16d^2 -72d^2+45d^2=-99\\
				&\Leftrightarrow& -11d^2=-99 \Leftrightarrow \hoac{&d=3\\ &d=-3.}
			\end{eqnarray*}
			Với $d=3$, ta có $u_1=-4$.\\
			Với $d=-3$, ta có $u_1=4$.\\
			Vậy $\heva{&u_1=-4\\ &d=3}$ hoặc $\heva{&u_1=4\\ &d=-3.}$
		\end{enumerate}
	}
\end{bt}

\begin{bt}%[1D2N2-4]
% nguon: Archive/2024/TOAN_11_XE/TOAN_11_XE/data4/11C2/KNTT-2.2.tex #579
%[Dự án DCHT-11-KNTT]%[Dao-V- Thuy]
	Cho cấp số cộng $(u_n)$, biết
	\begin{center}
		$\heva{u_1&=-15\\d&=18.}$
	\end{center}
	\begin{enumerate}
		\item Tìm $u_5$, $u_{10}$, $u_{15}$, $u_{20}$, $u_{25}.$
		\item Số $1209$ là số hạng thứ bao nhiêu ?
	\end{enumerate}
	\dapso{$u_5=57$, $u_{10}=147$, $u_{15}=237$, $u_{20}=327$, $u_{25}=417$;
		$1209$ là số hạng thứ $69$}
	\loigiai{
		\begin{enumerate}
			\item Áp dụng công thức $u_n=u_1+(n-1)d$ ta có \\
			$u_5=u_1+4d=-15+4.18=57.$\\ 
			$u_{10}=u_1+9d=-15+9.18=147.$\\
			$u_{15}=u_1+14d=-15+14.18=237.$\\
			$u_{20}=u_1+19d=-15+19.18=327.$\\
			$u_{25}=u_1+24d=-15+24.18=417.$
			\item Gọi $u_k$ là số hạng của số $1209$. Ta có
			\begin{center}
				$u_k=1209=u_1+(k-1)d= -15+18(k-1)= 18k-33 \Leftrightarrow k=69.$
			\end{center}
			Vậy $1209$ là số hạng thứ $69$.
		\end{enumerate}	
	}
\end{bt}

\begin{bt}%[1D2H2-4]
% nguon: Archive/2024/TOAN_11_XE/TOAN_11_XE/data4/11C2/KNTT-2.2.tex #580
%[Dự án DCHT-11-KNTT]%[Dao-V- Thuy]
	Tìm sáu số trong khoảng $(7; 35)$ để được một cấp số cộng gồm tám số hạng với $u_1=7$, $u_8=35$.
	\dapso{$11; \, 15; \,  19; \, 23; \, 27; \, 31$}
	\loigiai
	{
		Ta có
		\begin{eqnarray*}
			\heva{&u_1=7\\ &u_8=-35} \Leftrightarrow \heva{&u_1=7 \\ &u_1+7d=35} \Leftrightarrow \heva{&u_1=7 \\ &d=4.}
		\end{eqnarray*}
		Vậy sáu số đặt giữa các số $7$ và $35$ để được một cấp số cộng là
		\[
		11; \, 15; \,  19; \, 23; \, 27; \, 31.
		\]
	}
\end{bt}

\begin{bt}%[1D2H2-4]
% nguon: Archive/2024/TOAN_11_XE/TOAN_11_XE/data4/11C2/KNTT-2.2.tex #581
%[Dự án DCHT-11-KNTT]%[Dao-V- Thuy]
	Một cấp số cộng có năm số hạng mà tổng số hạng đầu và số hạng thứ ba bằng $28$, tổng của số hạng thứ ba và số hạng cuối bằng $40$. Hãy tìm cấp số cộng đó.
	\dapso{$11; 14; 17; 20; 23$}
	\loigiai
	{
		Ta có
		\begin{eqnarray*}
			\heva{&u_1+u_3=28\\ &u_3+u_5=40} \Leftrightarrow \heva{&u_1+ \left( u_1+2d\right)= 28\\ &\left( u_1+2d\right) + \left( u_1+4d\right)= 40}
			\Leftrightarrow \heva{&2u_1+2d=28\\ &2u_1+6d=40} \Leftrightarrow \heva{&u_1=11 \\ &d=3.}
		\end{eqnarray*}
		Vậy cấp số cộng cần tìm là
		\begin{center}
			11; 14; 17; 20; 23.
		\end{center}
	}
\end{bt}

\begin{bt}%[1D2H2-4]
% nguon: Archive/2024/TOAN_11_XE/TOAN_11_XE/data4/11C2/KNTT-2.2.tex #582
%[Dự án DCHT-11-KNTT]%[Dao-V- Thuy]
	Xác định $4$ góc của một tứ giác lồi, biết rằng $4$ góc hợp thành cấp số cộng và góc lớn nhất bằng $5$ lần góc nhỏ nhất.
	\dapso{$36^\circ; \, 72^\circ; \, 108^\circ; \, 144^\circ$}
	\loigiai
	{
		Gọi số đo bốn góc cần tìm là $u_1$, $u_2$, $u_3$, $u_4$. Ta có
		\begin{eqnarray*}
			\heva{&u_1+u_2+u_3+u_4=360\\ &u_5=5u_1} \Leftrightarrow \heva{&4u_1+6d=360\\ &4d=4u_1} \Leftrightarrow \heva{&u_1=36\\ &d=36.}
		\end{eqnarray*}
		Vậy số đo bốn góc cần tìm là
		\[
		36^\circ; \, 72^\circ; \, 108^\circ; \, 144^\circ.
		\]
	}
\end{bt}

\begin{bt}%[1D2H2-4]
% nguon: Archive/2024/TOAN_11_XE/TOAN_11_XE/data4/11C2/KNTT-2.2.tex #583
%[Dự án DCHT-11-KNTT]%[Dao-V- Thuy]
	Cho cấp số cộng $(u_n)$ với $\heva{&u_5=-43\\ &u_{21}=-171.}$
	\begin{enumerate}
		\item Tìm $d$ và $u_1$.
		\item Tìm $u_{29}$.
		\item $-16187$ là số hạng thứ bao nhiêu của cấp số cộng trên?
		\item $-35$ có thuộc cấp số cộng trên hay không?		
	\end{enumerate}
	\dapso{$\heva{&u_1=-11 \\ &d=-8}$; $u_{29}= -235$; $-16187$ là số hạng thứ $2023$; $-35$ thuộc cấp số cộng đã cho}
	\loigiai
	{
		\begin{enumerate}
			\item Ta có
			\begin{eqnarray*}
				\heva{&u_5=-43\\ &u_{21}=-171} \Leftrightarrow \heva{&u_1+4d=-43\\ &u_1+20d=-171} \Leftrightarrow \heva{&u_1=-11 \\ &d=-8.}
			\end{eqnarray*}
			\item Ta có $u_{29}= u_1+28d= -11+28 \cdot (-8)= -235$.
			\item Giả sử $u_k= -16187$, với $k \in \mathbb{N}^*$. Ta có
			\[
			u_1+ (k-1)d= -16187 \Leftrightarrow -11 + (k-1) \cdot (-8)= -16187 \Leftrightarrow k= 2023.
			\]
			Vậy $-16187$ là số hạng thứ $2023$ của cấp số cộng.
			\item Giả sử $u_m= -35$, với $m\in \mathbb{N}^*$. Ta có
			\[
			u_1+ (m-1) d= -35 \Leftrightarrow -11+ (m-1) \cdot (-8)= -35 \Leftrightarrow m=4. \quad (\text{thoả mãn})
			\]
			Vậy $-35$ thuộc cấp số cộng đã cho.
		\end{enumerate}
	}
\end{bt}

\begin{bt}%[1D2H2-7]
% nguon: Archive/2024/TOAN_11_XE/TOAN_11_XE/data4/11C2/KNTT-2.2.tex #584
%[DCHT Toán 11 - KNTT - Nguyễn Hữu Đức] 
	Một người có một khoản tiền gửi ngân hàng với lãi suất 10\% năm. Nếu sau $ 5 $ năm người đó nhận được tổng số tiền là $ 550 $ triệu đồng thì số tiền gửi ban đầu của người đó là bao nhiêu?
	\dapso{$366{,}67$ triệu đồng.}
	\loigiai{
		Gọi $x$ là số tiền gửi ban đầu của người đó $ (x>0) $.\\
		Sau 5 năm, số tiền nhận được bằng số tiền gốc cộng với lãi suất:
		$$
		x + 0{,}1x \times 5 = 1{,}5x.
		$$
		Theo đề bài, tổng số tiền nhận được sau 5 năm là $550$ triệu đồng, do đó ta có phương trình:
		$$
		1{,}5x = 550.
		$$
		Giải phương trình ta có:
		$$
		x = \frac{550}{1{,}5} \approx 366{,}67.
		$$
		Vậy số tiền gửi ban đầu của người đó là $366{,}67$ triệu đồng.
	}
\end{bt}

\begin{bt}%[1D2H2-7]
% nguon: Archive/2024/TOAN_11_XE/TOAN_11_XE/data4/11C2/KNTT-2.2.tex #585
%[DCHT Toán 11 - KNTT - Nguyễn Hữu Đức] 
	Bạn An muốn mua một món quà tặng mẹ nhân ngày mùng $8/3$. Bạn quyết định tiết kiệm từ ngày $1/2/2017$ đến hết ngày $6/3/2017$. Ngày đầu An có $5\,000$ đồng, kể từ ngày thứ hai số tiền An tiết kiệm được ngày sau cao hơn ngày trước mỗi ngày $1\,000$ đồng. Tính số tiền An tiết kiệm được để mua quà tặng mẹ.
	\dapso{$731\,000$ đồng.}
	\loigiai{
		Tính số ngày mà An tiết kiệm được từ ngày $1/2/2017$ đến hết ngày $6/3/2017$:\\
		Số ngày từ ngày $1/2/2017$ đến hết ngày $28/2/2017$ là $28$ ngày.\\
		Số ngày từ ngày $1/3/2017$ đến hết ngày $6/3/2017$ là $6$ ngày.\\
		Vậy An tiết kiệm được $28+6=34$ ngày.\\
		Gọi $x$ là số tiền An tiết kiệm được vào ngày đầu tiên, do ngày sau An tiết kiệm được nhiều hơn ngày trước mỗi ngày $1\,000$ đồng, nên số tiền An tiết kiệm được vào ngày thứ $n$ là $x + 1\,000(n-1)$.\\
		Vậy An tiết kiệm được tổng số tiền là:
		$$
		S=x + (x+1\,000) + (x+2\,000) + \cdots + (x+33 \times 1\,000).
		$$
		Ta có công thức của tổng số hạng cấp số cộng:
		$$
		S_n = \frac{n(u_1 + u_n)}{2}=\frac{2u_1+(n-1)d}{2}\cdot n
		$$
		Áp dụng công thức này, ta có:
		\begin{align*}
			&S= \frac{2u_1+(n-1)d}{2}\cdot n \\
			&= \frac{2\cdot 5000+33\cdot 1000}{2}\cdot 34. \\
			&=731\,000.
		\end{align*}
		Vậy số tiền An tiết kiệm được để mua quà tặng mẹ là $731\,000$ đồng.
	}
\end{bt}

\begin{bt}%[1D2H3-8]
% nguon: Archive/2024/TOAN_11_XE/TOAN_11_XE/data4/11C2/KNTT-2.2.tex #586
%[DCHT Toán 11 - KNTT - Nguyễn Hữu Đức] 
	Một hội đồng quản trị quyết định tăng lương cho nhân viên hàng năm theo tỷ lệ cố định. Ví dụ, lương của một nhân viên được tăng thêm $ 5 $\% so với năm trước. Hỏi nếu lương của một nhân viên là $ 10 $ triệu đồng/năm vào năm nay, thì lương của nhân viên đó sẽ là bao nhiêu vào năm thứ $ 5 $?
	\dapso{$12{,}1550625 $ triệu đồng/năm.}
	\loigiai{		
		Theo giả thiết, lương của nhân viên được tăng thêm $ 5 $ \% so với năm trước đó.
		\begin{itemize}
			\item Vậy lương của nhân viên vào năm thứ $ 2 $ sẽ là $ 10\cdot(1+0{,}05)=10{,}5 $ triệu đồng/năm.
			\item Tương tự, lương của nhân viên vào năm thứ $ 3 $ sẽ là $ 10{,}5 \cdot(1+0{,}05)=11{,}025 $ triệu đồng/năm.
			\item Lương của nhân viên vào năm thứ $ 4 $ sẽ là $ 11{,}025\cdot (1+0{,}05)=11{,}57625 $ triệu đồng/năm.
			\item Cuối cùng, lương của nhân viên vào năm thứ $ 5 $ sẽ là $ 11{,}57625\cdot(1+0{,}05)=12{,}1550625 $ triệu đồng/năm.
		\end{itemize}
		Vậy lương của nhân viên đó vào năm thứ $ 5 $ sẽ là $ 12{,}1550625 $ triệu đồng/năm.\\
		Chú ý: Lương của nhân viên đó vào năm thứ $ 5 $ sẽ là $ u_5=u_1+4d=10+4\cdot 10\cdot 0{,}05=12 $ triệu đồng chỉ đúng trong trường hợp lương của một nhân viên được tăng thêm $ 5 $\% so với năm đầu tiên.
	}
\end{bt}

\begin{bt}%[1D2N3-8]
% nguon: Archive/2024/TOAN_11_XE/TOAN_11_XE/data4/11C2/KNTT-2.2.tex #588
%[DCHT Toán 11 - KNTT - Nguyễn Hữu Đức]
	Hàng tháng ông An gửi vào ngân hàng một số tiền như nhau là $5\,000\,000$ đồng (vào ngày đầu mỗi tháng) với lãi suất $0{,}5\%$ một tháng, biết tiền lãi của tháng trước được nhập vào tiền gốc của tháng sau. Hỏi sau $36$ tháng ông An nhận được số tiền vốn và lãi là bao nhiêu? (làm tròn đến hàng đơn vị).
	\dapso{ $ 197\,663\,927 $ đồng.}
	\loigiai{
		Gọi $a$ là số tiền ông An gửi vào hàng tháng, $r$ là lãi suất trên một tháng và $P_n$ là số tiền vốn và lãi ông An nhận được sau $n$ tháng.
		\begin{itemize}
			\item Sau một tháng, ông An có số tiền là $P_1=a+ar=a(1+r)$.
			\item Đầu tháng thứ hai, ông An có số tiền là $P_1+a=a(1+r)+a$.
			\item Sau hai tháng, ông An có số tiền là $P_2=a(1+r)+a+\left[a(1+r)+a\right]r=a\left[(1+r)^2+(1+r)\right]$.
			\item Cuối tháng thứ $36$, ông An có số tiền là
			\begin{align*}
				P_{36}&=a\left[(1+r)^{36}+(1+r)^{35}+\ldots+(1+r)\right]\\
				&=a(1+r)\dfrac{(1+r)^{36}-1}{r}\\
				&=5000000\cdot (1+0{,}005)\cdot\dfrac{(1+0{,}005)^{36}-1}{0{,}005}\\
				&\approx 197\,663\,927 \quad \text{(đồng)}.
			\end{align*}
		\end{itemize}
	}
\end{bt}

\begin{bt}%[1D2N2-7]
% nguon: Archive/2024/TOAN_11_XE/TOAN_11_XE/data4/11C2/KNTT-2.2.tex #589
%[DCHT Toán 11 - KNTT - Nguyễn Hữu Đức] 
	Một xưởng có đăng tuyển công nhân với đãi ngộ về lương như sau: Trong quý đầu tiên thì xưởng trả là $ 6 $ triệu đồng/quý và kể từ quý thứ $ 2 $ sẽ tăng lên $ 0{,}5 $ triệu cho $ 1 $ quý. Hỏi với đãi ngộ trên thì sau $ 5 $ năm làm việc tại xưởng, tổng số lương của công nhân đó là bao nhiêu?
	\dapso{ $ 215 $ triệu đồng.}
	\loigiai{
		Giả sử công nhân làm cho xưởng $ n $ quý thì mức lương khi đó kí hiệu $ (u_n) $ (triệu đồng).\\
		Theo đề:\\
		Quý đầu: $ u_1 = 6 $ triệu.\\
		Các quý tiếp theo: $ u_{n+1} = u_{n} + 0,5 $ với $\forall n \ge 1$.\\
		Mức lương của công nhân mỗi quý là $ 1 $ số hạng của dãy số $ u_n $. Mặt khác, lương của quý sau hơn lương quý trước là $ 0,5 $ triệu nên dãy số $ u_n $ là một cấp số cộng với công sai $ d = 0{,}5 $.\\
		Ta biết $ 1 $ năm sẽ có $ 4 $ quý nên $ 5 $ năm sẽ có $ 5\cdot 4 = 20 $ quý. Theo yêu cầu của đề bài ta cần tính tổng của $ 20 $ số hạng đầu tiên của cấp số cộng ($ u_n $).\\
		Lương tháng quý $ 20 $ của công nhân: $ u_{20} = 6 + (20 - 1)\cdot 0{,}5 = 15{,}5 $ triệu đồng.\\
		Tổng số lương của công nhân nhận được sau $ 5 $ năm làm việc tại xưởng: $ S_{12}=20\cdot (6+15,5)2=215 $ (triệu đồng).
	}
\end{bt}

\begin{bt}%[1D2H2-3]
% nguon: Archive/2024/TOAN_11_XE/TOAN_11_XE/data4/11C2/KNTT-2.2.tex #590

	Xác định công sai, số hạng thứ $5$, số hạng tổng quát và số hạng thứ $100$ của mỗi cấp số cộng sau
	\begin{multicols}{2}
		\begin{enumerate}
			\item $4,9,14,19,\ldots;$
			\item $1,-1-3,-5,\ldots.$
		\end{enumerate}
	\end{multicols}	
	\loigiai{\begin{enumerate}
		\item Cấp số cộng với $u_1=4$ công sai $d=5$.\\
		Số hạng thứ $5$ là $u_5= u_1+4d=4+4\cdot 5=24 $.\\
		Số hạng tổng quát $u_n=u_1+(n-1)\cdot d=4+5(n-1).$\\
		Số hạng thứ $100$ là $u_{100}=4+99\cdot 5= 499$.
		\item Cấp số cộng với $u_1=1$ công sai $d=-2$.\\
		Số hạng thứ $5$ là $u_5= 1+4\cdot (-2)=-7$. \\
		Số hạng tổng quát là $u_n= 1-2(n-1)$.\\
		Số hạng thứ $100$ là $u_{100}=1-2\cdot 99=-197.$
	\end{enumerate}
	}
\end{bt}

\begin{bt}%[1D2H2-2]
% nguon: Archive/2024/TOAN_11_XE/TOAN_11_XE/data4/11C2/KNTT-2.2.tex #591

	Viết năm số hạng đầu của mỗi dãy số $(u_n)$ sau và xét xem nó có phải là cấp số cộng không. Nếu dãy số đó là cấp số cộng, hãy tìm công sai $d$ và viết số hạng tổng quát của nó dưới dạng $u_n=u_1+(n-1)d$.
	\begin{multicols}{2}
		\begin{enumerate}
			\item $u_n=3+5n$;
			\item $u_n=6n-4$;
			\item $u_1=2,u_n=u_{n-1}+n$;
			\item $u_1=2,u_{n-1}+3$.
		\end{enumerate}
	\end{multicols}	
	\loigiai{\begin{enumerate}
		\item Năm số hạng đầu $ 8,13,18,23,28.$ Đây là một cấp số cộng. Công sai $d=5$.\\
		 Số hạng tổng quát $u_n=8+(n-1)\cdot 5.$
		\item Năm số hạng đầu $2,8,14,20,26.$ Đây là một cấp số cộng với công sai $d=6$.\\
		 Số hạng tổng quát $u_n=2+(n-1)\cdot 6$.
		\item Năm số hạng đầu $2,4,7,11,16$. Đây không phải là cấp số cộng.
		\item Năm số hạng đầu $2,5,8,11,14.$ Đây là cấp số cộng với công sai $d=3$.\\
		 Số hạng tổng quát $u_n=2+(n-1)\cdot 3.$
		\end{enumerate}
	}
\end{bt}

\begin{bt}%[1D2H2-4]
% nguon: Archive/2024/TOAN_11_XE/TOAN_11_XE/data4/11C2/KNTT-2.2.tex #592

	Một cấp số cộng có số hạng thứ $5$ bằng $18$ và số hạng thứ $12$ bằng $32$. Tìm số hạng thứ $50$ của cấp số cộng này.
	\loigiai{
		Giả sử $u_1$, $d$ lần lượt là số hạng đầu tiên, công sai của cấp số cộng thỏa mãn yêu cầu bài toán. Khi đó, theo giả thiết, ta có
		\[\heva{&u_{5}=u_1+4d=18\\&u_{12}=u_1+11d=32.}\]
		Giải hệ này ta được $u_1=10$ và $d=2$.\\
		Vậy số hạng thức $50$ của cấp số cộng này là $u_{50}=u_1+49d=10+49\cdot 2=108$. }
\end{bt}

\begin{bt}%[1D2H2-6]
% nguon: Archive/2024/TOAN_11_XE/TOAN_11_XE/data4/11C2/KNTT-2.2.tex #593

	Một cấp số cộng có số hạng đầu bằng $5$ và công sai bằng $2$. Hỏi phải lấy tổng của bao nhiêu số hạng đầu của cấp số cộng này để có tổng bằng $2700$?
	\loigiai{
		Cấp số cộng này có số hạng đầu là $u_1=5$ và công sai là $d=2$. Gọi $n$ là số các số hạng đầu của cấp số cộng cần lấy tổng, ta có
		\[2700=S_n=\dfrac{n}{2}\left(2u_1+(n-1)d\right)=\dfrac{n}{2}\left(2\cdot 5+(n-1)\cdot 2\right)=\dfrac{n}{2}\left(2n+8\right).\]
		Do đó $2n^2+8n-5400=0$. Giải phương trình bậc hai này ta được $n=-54 \text{(loại)}$ và $n=50$. Vậy phải lấy tổng của $50$ số hạng đầu của cấp số cộng đã cho để được tổng bằng $2700$.}
\end{bt}

\begin{bt}%[1D2H2-7]
% nguon: Archive/2024/TOAN_11_XE/TOAN_11_XE/data4/11C2/KNTT-2.2.tex #594

	Giá của một chiếc xe ô tô lúc mới mua là $680$ triệu đồng. Cứ sau mỗi năm sử dụng, giá của chiếc ô tô đó giảm $55$ triệu đồng. Tính giá còn lại của chiếc xe sau $5$ năm sử dụng.
	\loigiai{
		Sau mỗi năm giá của chiếc ô tô giảm $55$ triệu đồng, do đó ta có thể xem giá của chiếc ô tô đó là một cấp số cộng với $u_1=680$ và công sai $d=-55$. \\
		Vậy giá của chiếc ô tô sau $5$ năm sử dụng là $ u_5=680+4\cdot (-55)=460$ triệu đồng.	
	}
\end{bt}

\begin{bt}%[1D2H2-7]
% nguon: Archive/2024/TOAN_11_XE/TOAN_11_XE/data4/11C2/KNTT-2.2.tex #595

	Một kiến trúc sư thiết kế một hội trường với $15$ ghế ngồi ở hàng thứ nhất, $18$ ghế ngồi ở hàng thứ hai, $21$ ghế ngồi ở hàng thứ ba, và cứ như vậy (số ghế ở hàng sau nhiều hơn 3 ghế so với số ghế ở hàng liền trước nó). Nếu muốn hội trường đó có sức chứa ít nhất $870$ ghế ngồi thì kiến trúc sư đó phải thiết kế tối thiểu bao nhiêu hàng ghế?
	\loigiai{
		Số ghế ở mỗi hàng của nhà hát lập thành một cấp số cộng với số hạng đầu $u_1=15$ và công sai $d=3$. Gọi $n$ là số hàng ghế trong hội trường. Tổng các số hạng này là
		\[S_{n}=\dfrac{n}{2}\left[2u_1+(n-1)d\right]=\dfrac{n}{2}\left[2\cdot 15+(n-1)\cdot 3\right]=\dfrac{n}{2}(3n+27)\ge 870.\]
		Do đó $3n^2+27n-1740\ge0$.\\
		 Giải bất phương trình bậc hai này ta được $n\le-29 \text{\ (loại)}$ và $n\ge 20$.\\
		  Vậy phải thiết kế tối thiểu 20 hàng ghế.
	}
\end{bt}

\begin{bt}%[1D2H2-7]
% nguon: Archive/2024/TOAN_11_XE/TOAN_11_XE/data4/11C2/KNTT-2.2.tex #596

	Vào năm $2020$, dân số của một thành phố là khoảng $1,2$ triệu người. Giả sử mỗi năm, dân số của thành phố này tăng thêm khoảng $30$ nghìn người. Hãy ước tính dân số của thành phố này vào năm $2030$.
	\loigiai{
		Ta thấy dân số của thành phố này qua các năm lập thành cấp số cộng với số hạng đầu $u_1=1200000$ và $d=30000$.\\
		 Vào năm 2030 dân số thành phố này tương ứng với năm thứ 11 kể từ năm 2020.\\
		  Do đó $u_{11}=1200000+10\cdot 30000=1500000$.\\
		Vậy dân số thành phố này vào năm 2030 là 1,5 triệu người.
	}
\end{bt}

\begin{bt}%[1D2H2-7]
% nguon: Archive/2024/TOAN_11_XE/TOAN_11_XE/data4/11C2/KNTT-2.2.tex #598
%[DCHT Toán 11 - KNTT - Nguyễn Hữu Đức] 
	Sinh nhật bạn của An vào ngày $ 01 $ tháng năm. An muốn mua một món quà sinh nhật cho bạn nên quyết định bỏ ống heo $ 100 $ đồng vào ngày $ 01 $ tháng $ 01 $ năm $ 2016 $, sau đó cứ liên tục ngày sau hơn ngày trước $ 100 $ đồng. Hỏi đến ngày sinh nhật của bạn, An đã tích lũy được bao nhiêu tiền? (thời gian bỏ ống heo tính từ ngày $ 01 $ tháng $ 01 $ năm $ 2016 $ đến ngày $ 30 $ tháng $ 04 $  năm $ 2016 $).\dapso{$ 738\,100 $ đồng.}
	\loigiai{
		Từ ngày $1$ tháng $1$ năm $2016$ đến ngày $30$ tháng $4$ năm $2016$ có tổng cộng $31+29+31+30=121$ ngày.\\
		Gọi $S$ là số tiền An tích lũy được vào ngày sinh nhật của bạn.\\
		Do An bỏ được $100$ đồng vào ngày đầu tiên nên số tiền An tích lũy được vào ngày thứ $n$ là 
		$$S= 100 + 100(n-1).$$
		Vậy tổng số tiền An tích lũy được là:
		$$
		S=100 + 200 + \cdots + 12\,100 = \frac{121(100 + 12\,100)}{2} = 738\,100
		.$$
		Vậy An đã tích lũy được $738\,100$ đồng vào ngày sinh nhật của bạn.}
\end{bt}

\begin{bt}%[1D2N2-7]
% nguon: Archive/2024/TOAN_11_XE/TOAN_11_XE/data4/11C2/KNTT-2.2.tex #599
%[DCHT Toán 11 - KNTT - Nguyễn Hữu Đức] 
	Người ta trồng $ 3\,003 $ cây theo dạng một hình tam giác như sau: hàng thứ nhất trồng $ 1 $ cây, hàng thứ hai trồng $ 2 $ cây, hàng thứ ba trồng $ 3 $ cây... cứ tiếp tục trồng như thế cho đến khi hết số cây. Số hàng cây được trồng là bao nhiêu?
	\dapso{$ 77 $ hàng.}
	\loigiai{
		Tổng số cây trồng được là $1 + 2 + 3 + \cdots + n$, nghĩa là tổng của $n$ số tự nhiên đầu tiên. Ta cần tìm số $n$ để tổng này bằng $3003$.\\
		Ta có công thức tổng của $n$ số tự nhiên đầu tiên là:
		$$
		1 + 2 + 3 + \cdots + n = \frac{n(n+1)}{2}.
		$$
		Giải phương trình:
		$$
		\frac{n(n+1)}{2} = 3\,003.
		$$
		Ta có:
		$$
		n(n+1) = 6\,006
		\Rightarrow n=77.$$
		Vậy số hàng cây được trồng là $77$.
	}
\end{bt}

\begin{bt}%[1D2H2-7]
% nguon: Archive/2024/TOAN_11_XE/TOAN_11_XE/data4/11C2/KNTT-2.2.tex #600
%[DCHT Toán 11 - KNTT - Nguyễn Hữu Đức] 
	Một công ty định mức sản phẩm hàng tháng theo cấp số cộng. Ví dụ, sản lượng hàng tháng của một công ty được tăng thêm $10$ sản phẩm so với tháng trước. Nếu công ty sản xuất được $ 100 $ sản phẩm trong tháng này, hỏi công ty sẽ sản xuất được bao nhiêu sản phẩm trong tháng thứ $ 12 $?
	\dapso{$ 210 $ sản phẩm.}
	\loigiai{
		Công thức cấp số cộng được sử dụng để tính sản lượng hàng tháng của công ty. Nếu công ty sản xuất được $100$ sản phẩm trong tháng này và sản lượng hàng tháng được tăng thêm $10$ sản phẩm so với tháng trước, ta có thể sử dụng công thức sau để tính sản lượng hàng tháng của công ty trong tháng thứ $12$:
		$$
		a_n = a_1 + (n-1)d.
		$$
		Trong đó $a_1$ là sản lượng hàng tháng ban đầu, $d$ là công sai và $n$ là số tháng.\\
		Với bài toán này, ta có: $a_1 = 100$, $d = 10$, $n = 12$.\\
		Sản lượng hàng tháng của công ty trong tháng thứ $12$ là:
		$$
		a_{12} = a_1 + (n-1)d = 100 + (12-1) \times 10 = 210.
		$$
		Vậy công ty sẽ sản xuất được $210$ sản phẩm trong tháng thứ $12$.
	}
\end{bt}

\begin{bt}%[1D2H2-4]
% nguon: Archive/2024/TOAN_11_XE/TOAN_11_XE/data4/11C2/KNTT-2.2.tex #601
%[Dự án DCHT-11-KNTT]%[Dao-V- Thuy]
	Giữa các số $10$ và $64$ hãy đặt thêm $17$ số nữa để được một cấp số cộng.
	\dapso{$13; 16; 19; 22; 25; 28; 31; 34; 37; 40; 43; 46; 49; 52; 55; 58; 61$}
	\loigiai{
		Ta có
		\begin{equation*}
			\heva{&u_1=10\\ &u_{19}=64} \Leftrightarrow \heva{&u_1=10\\ &u_1+18d=64} \Leftrightarrow \heva{&u_1=10\\ &d=3.}
		\end{equation*}
		Vậy $17$ số đặt thêm giữa các số $10$ và $64$ để được một cấp số cộng là
		\begin{center}
			13; 16; 19; 22; 25; 28; 31; 34; 37; 40; 43; 46; 49; 52; 55; 58; 61.
		\end{center} 
	}
\end{bt}

\begin{bt}%[1D2H2-4]
% nguon: Archive/2024/TOAN_11_XE/TOAN_11_XE/data4/11C2/KNTT-2.2.tex #602
%[Dự án DCHT-11-KNTT]%[Dao-V- Thuy]
	Tổng ba số hạng liên tiếp của một cấp số cộng bằng $2$ và tổng các bình phương của ba số đó bằng $\dfrac{14}{9}$. Xác định ba số đó và tính công sai của cấp số cộng.
	\dapso{$1;\dfrac{2}{3};\dfrac{1}{3}$ ứng với $d=-\dfrac{1}{3}$ hoặc $\dfrac{1}{3};\dfrac{2}{3};1$ ứng với $d=\dfrac{1}{3}$}
	\loigiai{
		Ta có hệ
		\begin{align*}
			&\quad \heva{&u_k+u_{k+1}+u_{k+2}=2\\ &u^2_k+u^2_{k+1}+u^2_{k+2}=\dfrac{14}{9}}
			\Leftrightarrow \heva{&u_k+u_k+d+u_k+2d=2\\ &u^2_k+\left(u_k+d\right)^2 +\left(u_k+2d\right)^2=\dfrac{14}{9}} \\
			&\Leftrightarrow \heva{&3u_k+3d=2\\ &3u^2_k+6u_kd+5d^2=\dfrac{14}{9}}
			\Leftrightarrow \heva{&u_k=1\\ &d=-\dfrac{1}{3}} \text{ hoặc } \heva{&u_k=\dfrac{1}{3}\\ &d=\dfrac{1}{3}.}
		\end{align*}
		Vậy ba số hạng liên tiếp của cấp số cộng thỏa yêu cầu bài toán: \\
		$1;\dfrac{2}{3};\dfrac{1}{3}$ ứng với $d=-\dfrac{1}{3}$ hoặc $\dfrac{1}{3};\dfrac{2}{3};1$ ứng với $d=\dfrac{1}{3}.$
	}
\end{bt}

\begin{bt}%[1D2H2-4]
% nguon: Archive/2024/TOAN_11_XE/TOAN_11_XE/data4/11C2/KNTT-2.2.tex #603
%[Dự án DCHT-11-KNTT]%[Dao-V- Thuy]
	Một cấp số cộng có $7$ số hạng với công sai $d$ dương và số hạng thứ tư bằng $11$. Hãy tìm các số hạng còn lại của cấp số cộng đó, biết hiệu của số hạng thứ ba và số hạng thứ năm bằng $6$.
	\dapso{$u_1=2$; $ u_2=5$; $u_4=11$; $u_6=17$; $u_7=20$}
	\loigiai{
		Gọi số hạng đầu của cấp số cộng là $u_1$, công sai $d$.
		Vì số hạng thứ tư của cấp số cộng bằng $11$ nên ta có $u_4=11$.\\
		Do $d$ dương nên $ u_5>u_3$.\\
		Vì hiệu của số hạng thứ ba và số hạng thứ năm bằng $6$ nên ta có $ u_5-u_3=6$.\\
		Ta có \begin{align*}
			\heva{&u_4=11\\&u_5-u_3=6 }
			\Leftrightarrow \heva{&u_1+3d=11\\&(u_1+4d)-(u_1+2d)=6 }
			\Leftrightarrow \heva{&u_1+3\cdot 3=11\\&d=3 }
			\Leftrightarrow \heva{&u_1=2\\&d=3.}
		\end{align*}
		Vậy các số  hạng còn lại của cấp số cộng là $u_1=2$; $ u_2=5$; $u_4=11$; $u_6=17$; $u_7=20$.
	}
\end{bt}

\begin{bt}%[1D2V2-4]
% nguon: Archive/2024/TOAN_11_XE/TOAN_11_XE/data4/11C2/KNTT-2.2.tex #604
%[Dự án DCHT-11-KNTT]%[Dao-V- Thuy]
	Tìm bốn số hạng liên tiếp của một cấp số cộng, biết rằng:
	\begin{enumerate}
		\item Tổng của chúng bằng $10$ và tổng bình phương bằng $70$.
		\item Tổng của chúng bằng $22$ và tổng bình phương bằng $66$.
		\item  Tổng của chúng bằng $36$ và tổng bình phương bằng $504$.
		\item  Chúng có tổng bằng $20$ và tích của chúng bằng $384$.
		\item  Tổng của chúng bằng $ 20$, tổng nghịch đảo của chúng bằng $ \dfrac{25}{24}$ và các số này là những số nguyên.
		\item  Nó là số đo của một tứ giác lồi và góc lớn nhất gấp $5$ lần góc nhỏ nhất.
	\end{enumerate}
	\dapso{$-2$; $ 1$; $ 4$; $7$.} 
	\dapso{không tồn tại bốn số hạng liên tiếp của cấp số cộng thỏa mãn yêu cầu đề bài. $0$; $ 6$; $ 12$; $18$.}
	\dapso{$2$; $ 4$; $ 6$; $8$ hoặc $5-\sqrt{241}$; $ \dfrac{15-\sqrt{241}}{3}$; $ \dfrac{15+\sqrt{241}}{3}$; $5+\sqrt{241}$.} 
	\dapso{$30^\circ$; $70^\circ$; $ 110^\circ$; $150^\circ$.}
	\loigiai{
		\begin{enumerate}
			\item Gọi bốn số hạng liên tiếp của cấp số cộng là $x-3d$; $x-d$; $x+d$, $x+3d$ trong đó $2d$ là công sai.\\
			Theo đề bài ta có 
			\begin{align*}
				&\quad \heva{& (x-3d)+(x-d)+(x+d)+(x+3d)=10\\& (x-3d)^2+(x-d)^2+(x+d)^2+(x+3d)^2=70}
				\Leftrightarrow \heva{& 4x=10 \\& 4x^2+20d^2=70}\\
				&\Leftrightarrow  \heva{&x=\dfrac{5}{2}\\& 4\cdot \left( \dfrac{5}{2}\right) ^2+20d^2=70}
				\Leftrightarrow  \heva{&x=\dfrac{5}{2}\\& d^2=\dfrac{9}{4}}
				\Leftrightarrow  \heva{&x=\dfrac{5}{2}\\& d=\pm \dfrac{3}{2}.}
			\end{align*}
			Vậy bốn số hạng liên tiếp của cấp số cộng là $-2$; $ 1$; $ 4$; $7$.
			\item Gọi bốn số hạng liên tiếp của cấp số cộng là $x-3d$; $x-d$; $x+d$; $x+3d$ trong đó $2d$ là công sai.\\
			Theo đề bài ta có 
			\begin{align*}
				&\quad \heva{& (x-3d)+(x-d)+(x+d)+(x+3d)=22\\& (x-3d)^2+(x-d)^2+(x+d)^2+(x+3d)^2=66}
				\Leftrightarrow \heva{& 4x=22 \\& 4x^2+20d^2=66}\\
				&\Leftrightarrow  \heva{&x=\dfrac{11}{2}\\& 4\cdot \left( \dfrac{11}{2}\right) ^2+20d^2=66}
				\Leftrightarrow  \heva{&x=\dfrac{11}{2}\\& d^2=\dfrac{-11}{4} 
					\ (\text{loại}).}\\
			\end{align*}
			Vậy không tồn tại bốn số hạng liên tiếp của cấp số cộng thỏa mãn yêu cầu đề bài.
			\item Gọi bốn số hạng liên tiếp của cấp số cộng là $x-3d$; $x-d$; $x+d$; $x+3d$ trong đó $2d$ là công sai.\\
			Theo đề bài ta có 
			\begin{align*}
				&\quad \heva{& (x-3d)+(x-d)+(x+d)+(x+3d)=36\\& (x-3d)^2+(x-d)^2+(x+d)^2+(x+3d)^2=504}
				\Leftrightarrow \heva{& 4x=36 \\& 4x^2+20d^2=504}\\
				&\Leftrightarrow  \heva{&x=9\\& 4\cdot 9^2+20d^2=504}
				\Leftrightarrow  \heva{&x=9\\& d^2=9}
				\Leftrightarrow  \heva{&x=9\\& d=\pm 3.}
			\end{align*}
			Vậy  bốn  số hạng liên tiếp của cấp số cộng là $0$; $ 6$; $ 12$; $18$.
			\item Gọi bốn số hạng liên tiếp của cấp số cộng là $x-3d$; $x-d$; $x+d$; $x+3d$ trong đó $2d$ là công sai.\\
			Theo đề bài ta có 
			\begin{align*}
				&\quad \heva{& (x-3d)+(x-d)+(x+d)+(x+3d)=20\\& (x-3d)(x-d)(x+d)(x+3d)=384}
				\Leftrightarrow  \heva{&x=5\\& (x^2-d^2)(x^2-9d^2)=384}\\
				&\Leftrightarrow  \heva{&x=5\\& (25-d^2)(25-9d^2)=384}
				\Leftrightarrow  \heva{&x=5\\& 9d^4-250d^2+241=0}
				\Leftrightarrow  \heva{&x=5\\& \hoac{&d^2=1\\& d^2=\dfrac{241}{9}}}
				\Leftrightarrow  \heva{&x=5\\& \hoac{&d=\pm 1\\& d=\pm \dfrac{\sqrt{241}}{3}.}}
			\end{align*}
			Vậy bốn số hạng liên tiếp của cấp số cộng là $2$; $ 4$; $ 6$; $8$ hoặc $5-\sqrt{241}$; $ \dfrac{15-\sqrt{241}}{3}$; $ \dfrac{15+\sqrt{241}}{3}$; $5+\sqrt{241}$.
			\item Gọi bốn số hạng liên tiếp của cấp số cộng là $x-3d$; $x-d$; $x+d$; $x+3d$ trong đó $2d$ là công sai trong đó $ 2d \in \mathbb{Z}$.\\
			Theo đề bài ta có 
			\begin{align*}
				&\quad \heva{& (x-3d)+(x-d)+(x+d)+(x+3d)=20\\& \dfrac{1}{x-3d}+\dfrac{1}{x-d}+\dfrac{1}{x+d}+\dfrac{1}{x+3d}=\dfrac{25}{24}}
				\Leftrightarrow \heva{& 4x=20 \\& \dfrac{1}{5-3d}+\dfrac{1}{5-d}+\dfrac{1}{5+d}+\dfrac{1}{5+3d}=\dfrac{25}{24}}\\
				&\Leftrightarrow  \heva{&x=5\\& \dfrac{10}{25-9d^2}+\dfrac{10}{25-d^2}=\dfrac{25}{24}}
				\Leftrightarrow  \heva{&x=5\\& 9d^4-250d^2+241=0}\\
				&\Leftrightarrow  \heva{&x=5\\& \hoac{&d^2=1 \\& d^2=\dfrac{241}{9} }}
				\Leftrightarrow  \heva{&x=5\\& \hoac{&d= \pm 1 \ (\text{thỏa mãn})\\& d= \pm \dfrac{\sqrt{241}}{3} \,\,(\text{loại vì} \,  2d \in \mathbb{Z}).}}
			\end{align*}
			Vậy bốn  số hạng  nguyên liên tiếp của cấp số cộng là $2$; $ 4$; $ 6$; $8$.
			\item Gọi bốn số hạng liên tiếp của cấp số cộng xếp theo thứ tự tăng dần  là $x-3d$; $x-d$; $x+d$; $x+3d$ trong đó $2d>0$ là công sai.\\
			Theo đề bài ta có 
			\begin{align*}
				&\quad \heva{& (x-3d)+(x-d)+(x+d)+(x+3d)=360^\circ\\& x+3d=5(x-3d)}\\
				&\Leftrightarrow \heva{& 4x=360^\circ\\& 4x=18d}
				\Leftrightarrow  \heva{&x=90^\circ\\& 4 \cdot 90^\circ=18d}
				\Leftrightarrow  \heva{&x=90^\circ\\& d=20^\circ.}
			\end{align*}
			Vậy bốn  góc của tứ giác lồi lần lượt là  $30^\circ$; $70^\circ$; $ 110^\circ$; $150^\circ$.
		\end{enumerate}
	}
\end{bt}

\begin{bt}%[1D2H2-3]
% nguon: Archive/2024/TOAN_11_XE/TOAN_11_XE/data4/11C2/KNTT-2.2.tex #605
%[Dự án DCHT-11-KNTT]%[Dao-V- Thuy]
	Xác định công thức tổng quát của cấp số cộng $(u_n)$, biết $\heva{&u_{11}=5\\ &d=-6.}$
	\loigiai{
		Ta có
		\begin{equation*}
			\heva{&u_{11}=5\\ &d=-6} \Leftrightarrow \heva{&u_1+10d=5\\ &d=-6} \Leftrightarrow \heva{&u_1=65\\ &d=-6.}
		\end{equation*}
		Vậy công thức tổng quát của cấp số cộng:
		\begin{center}
			$u_n=65+(n-1).(-6) \Leftrightarrow u_n=-6n+71$  với $n \geq 2.$
		\end{center}	
	}
\end{bt}

\begin{bt}%[1D2H2-4]
% nguon: Archive/2024/TOAN_11_XE/TOAN_11_XE/data4/11C2/KNTT-2.2.tex #606
%[Dự án DCHT-11-KNTT]%[Dao-V- Thuy]
	Tìm số hạng đầu và công sai của cấp số cộng $(u_n),$ biết $\heva{&u_2+u_5-u_3=10\\ &u_4+u_6=26.}$
	\loigiai{
		Ta có
		\begin{align*}
			\heva{&u_2+u_5-u_3=10\\ &u_4+u_6=26} 
			&\Leftrightarrow \heva{&u_1+d+u_1+4d-(u_1+2d)=10\\ &u_1+3d+u_1+5d=26}\\ 
			& \Leftrightarrow\heva{&u_1+3d=10 \\ &2u_1+8d=26}
			\Leftrightarrow \heva{&u_1=1 \\&d=3.}
		\end{align*}
		Vậy $u_1=1$, $d=3$.
	}
\end{bt}

\begin{bt}%[1D2H2-4]
% nguon: Archive/2024/TOAN_11_XE/TOAN_11_XE/data4/11C2/KNTT-2.2.tex #607
%[Dự án DCHT-11-KNTT]%[Dao-V- Thuy]
	Tìm số hạng đầu và công sai của cấp số cộng, biết
	\begin{listEX}[3]
		\item $\heva{&u_7 = 27\\&u_{15} = 59.}$
		\item $\heva{&u_9 = 5u_2\\&u_{13} = 2u_6 + 5.}$
		\item $\heva{&u_2 + u_4 - u_6 = -7\\&u_8 - u_7 = 2u_4.}$
		\item $\heva{&u_3 - u_7 = -8\\&u_2 \cdot u_7 = 75.}$
		\item $\heva{&u_6 + u_7 = 60\\&u_4^2 + u_{12}^2 = 1170.}$
	\end{listEX}
	\loigiai{
		\begin{enumerate}
			\item Ta có $\heva{&u_7 = 27\\&u_{15} = 59} \Leftrightarrow \heva{&u_1 + 6d = 27\\&u_1 + 14d = 59} \Leftrightarrow \heva{&u_1 = 3\\&d = 4.}$\\
			Vậy số hạng đầu của cấp số cộng là $u_1 = 3$, công sai là $d = 4$.
			\item Ta có $\heva{&u_9 = 5u_2\\&u_{13} = 2u_6 + 5} \Leftrightarrow \heva{&u_1 + 8d = 5u_1 + 5d\\&u_1 + 12d = 2u_1 + 10d + 5} \Leftrightarrow \heva{&4u_1 - 3d = 0\\&-u_1 + 2d = 5} \Leftrightarrow \heva{&u_1 = 3\\&d = 4.}$\\
			Vậy số hạng đầu của cấp số cộng là $u_1 = 3$, công sai là $d = 4$.
			\item Ta có $\heva{&u_2 + u_4 - u_6 = -7\\&u_8 - u_7 = 2u_4} \Leftrightarrow \heva{&u_1 + d + u_1 + 3d - u_1 - 5d = -7\\&u_1 + 7d - u_1 - 6d = 2u_1 + 6d} \Leftrightarrow \heva{&u_1 - d = -7\\&2u_1 + 5d = 0} \Leftrightarrow \heva{&u_1 = -5\\&d = 2.}$\\
			Vậy số hạng đầu của cấp số cộng là $u_1 = -5$, công sai là $d = 2$.
			\item Ta có $\heva{&u_3 - u_7 = -8\\&u_2 \cdot u_7 = 75} \Leftrightarrow \heva{&u_1 + 2d -u_1 - 6d = -8\\&(u_1 + d)(u_1 + 6d) = 75} \Leftrightarrow \heva{&d = 2\\&u_1^2 + 14u_1 - 51 = 0} \Leftrightarrow \heva{&d = 2\\&\hoac{&u_1 = 3\\&u_1 = -17.}}$\\
			Vậy số hạng đầu của cấp số cộng là $u_1 = 3$, công sai là $d = 2$ hoặc $u_1 = -17$, $d = 2$.
			\item Ta có $\heva{&u_6 + u_7 = 60\\&u_4^2 + u_{12}^2 = 1170} \Leftrightarrow \heva{&2u_6 + d = 60&(1)\\&(u_6 - 2d)^2 + (u_6 + 6d)^2 = 1170.&(2)}$\\
			Từ (1), suy ra $d = 60 - 2u_6$, thay vào (2), ta có
			$$(5u_6 - 120)^2 + (360 - 11u_6)^2 = 1170 \Leftrightarrow 146u_6^2 - 9120u_6 + 142830 = 0 \,\, (\text{vô nghiệm}).$$ 
			Vậy không tồn tại cấp số cộng thỏa yêu cầu bài toán.
		\end{enumerate}
	}
\end{bt}

\begin{bt}%[1D2N1-2]
% nguon: Archive/2024/TOAN_11_XE/TOAN_11_XE/data4/11C2/KNTT-2.1.tex #610
%[DCHT Toán 11 - KNTT -Tên GV]
	Xác định số hạng đầu là số hạng tổng quát của dãy số $(u_n)$ các số tự nhiên lẻ $1, 3, 5, 7, \ldots $
	\dapso{$u_n=2n-1$}	
	\loigiai{Dãy $(u_n)$ có số hạng đầu $u_1=1$ và số hạng tổng quát $u_n=2n-1$.}
\end{bt}

\begin{bt}%[1D2N1-2]
% nguon: Archive/2024/TOAN_11_XE/TOAN_11_XE/data4/11C2/KNTT-2.1.tex #611
%[DCHT Toán 11 - KNTT -Tên GV]
	Xác định số hạng đầu là số hạng tổng quát của dãy số $(v_n)$ các số nguyên dương chia hết cho $5$: $5,10,15,20,\ldots$
	\dapso{$v_n=5n$}
	\loigiai{Dãy $(v_n)$ có số hạng đầu $v_1=5$ và số hạng tổng quát $v_n=5n$.}
\end{bt}

\begin{bt}%[1D2N1-2]
% nguon: Archive/2024/TOAN_11_XE/TOAN_11_XE/data4/11C2/KNTT-2.1.tex #612
%[DCHT Toán 11 - KNTT -Tên GV]
	Viết năm số hạng đầu và số hạng thứ $100$ của dãy số $(u_n)$ có số hạng tổng quát $u_n=3n-2$.
	\dapso{$u_{100}=298$}
	\loigiai{Năm số hạng đầu của dãy số là $1,4,7,10,13$.\\
		Số hạng thứ $100$ của dãy là $u_{100}=3\cdot100-2=298$.}
\end{bt}

\begin{bt}%[1D2N1-3]
% nguon: Archive/2024/TOAN_11_XE/TOAN_11_XE/data4/11C2/KNTT-2.1.tex #613
%[DCHT Toán 11 - KNTT -Tên GV]
	Cho dãy số xác định bằng hệ thức truy hồi: $u_1=1, u_n=3u_{n-1}+2$ với $n\ge 2$. Viết ba số hạng đầu của dãy số này.
	\dapso{$u_1=1, u_2=5, u_3=17$}
	\loigiai{Ta có $u_1=1, u_2=3u_1+2=5, u_3=3u_2+2=17$.}
\end{bt}

\begin{bt}%[1D2N1-2]
% nguon: Archive/2024/TOAN_11_XE/TOAN_11_XE/data4/11C2/KNTT-2.1.tex #614
%[DCHT Toán 11 - KNTT -Tên GV]
	Dãy số $(u_n)$ cho bởi hệ thức truy hồi: $u_1=1, u_n=n \cdot u_{n-1}$ với $n \ge 2$. Viết năm số hạng đầu của dãy số và dự đoán công thức tổng quát $u_n$.
	\dapso{$u_n=n!$}
	\loigiai{Năm số hạng đầu của dãy là
		$$u_1=1, u_2=2\cdot u_1=2, u_3=3\cdot u_2=6, u_4=4 \cdot u_3= 24, u_5=5 \cdot u_4=124.$$
		Số hạng tổng quát\\
		Ta có $ u_2=2\cdot 1, u_3=6=3\cdot 2\cdot 1, u_4=24=4\cdot 3\cdot 2\cdot 1, u_5=124= 5\cdot4\cdot3\cdot2\cdot1  $.\\
		Vậy số hạng tổng quát $u_n=n!$.}
\end{bt}

\begin{bt}%[1D2N1-3]
% nguon: Archive/2024/TOAN_11_XE/TOAN_11_XE/data4/11C2/KNTT-2.1.tex #615

	Cho dãy số $u_n$ bao gồm các số nguyên tố. Tìm số hạng thứ $5$ của dãy số.
	\dapso{$u_5=11$}
	\loigiai{ 
		Ta có
		$u_1=2,u_2=3,u_3=5,u_4=7,u_5=11$. \\
		Vậy số hạng thứ $5$ của dãy số là $11$.
	}
\end{bt}

\begin{bt}%[1D2H1-3]
% nguon: Archive/2024/TOAN_11_XE/TOAN_11_XE/data4/11C2/KNTT-2.1.tex #616
%[VD 5 SGK KNTT]
	Cho dãy số xác định bằng hệ thức truy hồi
	$$
	u_1=1, u_n=3 u_{n-1}+2 \text { với } n \geq 2
	$$
	Viết ba số hạng đầu của dãy số này.
	\dapso{$u_5=17$}
	\loigiai{
		Ta có: $u_1=1, u_2=3 u_1+2=3 \cdot 1+2=5, u_3=3 u_2+2=3 \cdot 5+2=17$.
	}
\end{bt}

\begin{bt}%[1D2H1-3]
% nguon: Archive/2024/TOAN_11_XE/TOAN_11_XE/data4/11C2/KNTT-2.1.tex #617

	Cho dãy số $\left(u_n\right)\colon\heva{&u_1=5 \\ &u_{n+1}=u_n+n}$. Số $20$ là số hạng thứ mấy trong dãy?
	\dapso{số hạng thứ $6$}
	\loigiai{
		Ta có $u_1=5, u_2=6, u_3=8, u_4=11, u_5=16, u_6=20$.\\
		Vậy số $20$ là số hạng thứ $6$.}
\end{bt}

\begin{bt}%[1D2N1-4]
% nguon: Archive/2024/TOAN_11_XE/TOAN_11_XE/data4/11C2/KNTT-2.1.tex #618

	Xét sự tăng giảm của dãy số $(u_n)$ với $u_n=(-1)^n$.
	\dapso{dãy không tăng không giảm}
	\loigiai{
		Ta có:\\ $u_1=(-1)^1=-1,\,
		u_2=(-1)^2=1,\,
		u_3=(-1)^3=-1.$\\
		Vậy $(u_n)$ là dãy không tăng không giảm.
	}
\end{bt}

\begin{bt}%[1D2N1-4]
% nguon: Archive/2024/TOAN_11_XE/TOAN_11_XE/data4/11C2/KNTT-2.1.tex #619

	Xét tính tăng giảm của dãy số sau $(u_n)$ với $u_n=\dfrac{2n+1}{n+1}$.
	\dapso{dãy số tăng}
	\loigiai
	{
		Ta có: $u_n=\dfrac{2n+1}{n+1}=2-\dfrac{1}{n+1}$.\\
		$u_{n+1}-u_n=\left(2-\dfrac{1}{n+1+1}\right)-\left(2-\dfrac{1}{n+1}\right)=\dfrac{1}{n+1}-\dfrac{1}{n+2}>0, \forall n \in \mathbb{N}^\ast$.\\
		Vậy dãy số $(u_n)$ là dãy số tăng.
	}
\end{bt}

\begin{bt}%[1D2H1-4]
% nguon: Archive/2024/TOAN_11_XE/TOAN_11_XE/data4/11C2/KNTT-2.1.tex #620

	Xét tính tăng giảm của dãy số $(u_n)$ với $u_n=\sqrt{n}-\sqrt{n+2}$.
	\dapso{dãy số tăng}
	\loigiai{
		Ta có $u_n=\sqrt{n}-\sqrt{n+2}=\dfrac{-2}{\sqrt{n}+\sqrt{n+2}}$.\\
		Xét hiệu\\ 
		$\begin{aligned}
			u_{n+1}-u_n&=\dfrac{-2}{\sqrt{n+1}+\sqrt{n+3}}-\dfrac{-2}{\sqrt{n}+\sqrt{n+2}}\\
			&=\dfrac{2}{\sqrt{n}+\sqrt{n+2}}-\dfrac{2}{\sqrt{n+1}+\sqrt{n+3}}>0, \forall n\in \mathbb{N}^\ast.
		\end{aligned}$\\
		Vậy $(u_n)$ là dãy số tăng.
	}
\end{bt}

\begin{bt}%[1D2H1-4]
% nguon: Archive/2024/TOAN_11_XE/TOAN_11_XE/data4/11C2/KNTT-2.1.tex #621

	Xét tính tăng giảm của dãy số $(u_n)$ với $u_n=\dfrac{n}{3^n}$.
	\dapso{dãy số giảm}
	\loigiai{
		Ta có $u_n=\dfrac{n}{3^n}>0, \forall n \in \mathbb{N}^\ast$.\\
		Xét thương $\dfrac{u_{n+1}}{u_n}=\dfrac{n+1}{3^{n+1}}:\dfrac{n}{3^n}=\dfrac{n+1}{3.n}<1, \forall  n \in \mathbb{N}^\ast$.\\
		Vậy $(u_n)$ là dãy số giảm.
	}
\end{bt}

\begin{bt}%[1D2V1-4]
% nguon: Archive/2024/TOAN_11_XE/TOAN_11_XE/data4/11C2/KNTT-2.1.tex #622

	Xét tính tăng giảm của dãy số $(u_n)$ với $ \heva{& u_1=2\\
		& u_{n+1}=\dfrac{3u_n+1}{u_n+1}, n\in \mathbb{N}^\ast.}$
	\dapso{dãy số tăng}
	\loigiai{
		Giả sử $u_{n+1}>u_n , \forall n \in \mathbb{N}^\ast. \qquad (*)$\\
		Ta chứng minh $(*)$ bằng phương pháp quy nạp.
		\begin{itemize}
			\item Với $n=1, u_2=\dfrac{3.2+1}{2+1}=\dfrac{6}{3}=\dfrac{7}{3}>u_1=2.$
			\item Giả sử $(*)$ đúng khi $n=k, k\in \mathbb{N}^\ast$, tức là $u_{k+1}>u_k$.\\
			Ta sẽ chứng minh $(*)$ đúng với $n=k+1$, tức là
			$u_{k+2}>u_{k+1}$.\\
			Thật vậy\\ $u_{k+2}-u_{k+1}=\left(3-\dfrac{2}{u_{k+1}+1}\right)-\left(3-\dfrac{2}{u_k+1}\right)=\dfrac{2}{u_k+1}-\dfrac{2}{u_{k+1}+1}.$\\
			Theo giả thiết quy nạp ta có:\\ $u_{k+1}>u_k \Rightarrow u_{k+1}+1>u_k+1 \Rightarrow \dfrac{2}{u_k+1}>\dfrac{2}{u_{k+1}+1}$.\\
			Vậy $u_{k+2}-u_{k+1}>0$.\\
			Do đó, $(*)$ đúng với mọi số nguyên dương $n$.
		\end{itemize}
		Vậy $(u_n)$ là dãy số tăng.}
\end{bt}

\begin{bt}%[1D2N1-5]
% nguon: Archive/2024/TOAN_11_XE/TOAN_11_XE/data4/11C2/KNTT-2.1.tex #623
%[Trương Đăng Khoa]
	Chứng minh rằng dãy số $(u_n)$ với $u_n=\dfrac{3n}{n^2+9}$ bị chặn trên bởi $\dfrac{1}{2}$.
	\dapso{dãy số đã cho bị chặn trên bởi $\dfrac{1}{2}$}
	\loigiai{
		Với mọi $n\ge 1$, ta có $\dfrac{3n}{n^2+9}\le\dfrac{1}{2}\Leftrightarrow n^2+9\le 6n\Leftrightarrow(n-3)^2\le 0$ (đúng).\\
		Vậy dãy số đã cho bị chặn trên bởi $\dfrac{1}{2}$.
	}
\end{bt}

\begin{bt}%[1D2N1-5]
% nguon: Archive/2024/TOAN_11_XE/TOAN_11_XE/data4/11C2/KNTT-2.1.tex #624
%[Trương Đăng Khoa]
	Chứng minh rằng dãy số $(u_n)$ xác đinh bởi $u_n=\dfrac{8n+3}{3n+5}$ là một dãy số bị chặn.
	\dapso{dãy số bị chặn}
	\loigiai{
		Ta có $u_n>0$, $\forall n\ge 1$. Suy ra dãy số bị chặn dưới.\\
		Mặt khác $u_n=\dfrac{8n+3}{3n+5}<\dfrac{8n+3}{3n}=\dfrac{8}{3}+\dfrac{1}{n}<\dfrac{8}{3}+1=\dfrac{11}{3}$. Do đó dãy số bị chặn trên bởi $\dfrac{11}{3}$.\\
		Vậy dãy số đã cho bị chặn.
	}
\end{bt}

\begin{bt}%[1D2H1-5]
% nguon: Archive/2024/TOAN_11_XE/TOAN_11_XE/data4/11C2/KNTT-2.1.tex #625
%[Trương Đăng Khoa]
	Xét tính bị chặn của dãy số $\left(u_n\right)$ với $u_n=\dfrac{3n+1}{n+3}$.
	\dapso{dãy số bị chặn}
	\loigiai{
		Với $n\in \mathbb{N}^\ast$ ta có $u_n=\dfrac{3n+1}{n+3}>0$.\\
		Nên dãy $\left(u_n\right)$ bị chặn dưới bởi $0$.\\
		Mặt khác $u_n=\dfrac{3n+1}{n+3}=\dfrac{3n+9-8}{n+3}=3-\dfrac{8}{n+3}<3$, $\forall n\in\mathbb{N}^\ast$.\\
		Nên dãy $\left(u_n\right)$ bị chặn trên bởi $3$. Vậy dãy số $\left(u_n\right)$ bị chặn.
	}
\end{bt}

\begin{bt}%[1D2V1-5]
% nguon: Archive/2024/TOAN_11_XE/TOAN_11_XE/data4/11C2/KNTT-2.1.tex #626
%[Trương Đăng Khoa]
	Cho dãy số $(u_n)$ xác định bởi $u_1=1$ và $u_{n+1}=\dfrac{u_n+2}{u_n+1}$, $\forall n\ge 1$. Chứng minh rằng dãy $(u_n)$ bị chặn trên bởi sô $\dfrac{3}{2}$ và bị chặn dưới bởi số $1$.	
	\loigiai{
		Ta chứng minh $1\le u_n\le\dfrac{3}{2},\forall n\ge 1$ bằng phương pháp quy nạp.
		\begin{itemize}
			\item Với $n=1$ ta có $1\le u_1\le\dfrac{3}{2}$.
			\item Giả sử $1\le u_n\le\dfrac{3}{2}$ với mọi $n=k\ge 1$, tức là $1\le u_k\le\dfrac{3}{2}$. Ta cần chứng minh $1\le u_{k+1}\le\dfrac{3}{2}$.
		\end{itemize}
		Thật vậy 
		$u_{k+1}=1+\dfrac{1}{u_k+1}$.\\
		Vì $u_k+1>0$ nên $u_{k+1}=1+\dfrac{1}{u_k+1}>1$.\\
		Vì $u_k+1\ge 2$ nên $u_{k+1}=1+\dfrac{1}{u_k+1}\le 1+\dfrac{1}{2}=\dfrac{3}{2}$.\\
		Vậy $1\le u_n\le\dfrac{3}{2}$, $\forall n\ge 1$ hay dãy $(u_n)$ bị chặn trên bởi số $\dfrac{3}{2}$ và bị chặn dưới bởi số $1$.
	}
\end{bt}

\begin{bt}%[1D2V1-5]
% nguon: Archive/2024/TOAN_11_XE/TOAN_11_XE/data4/11C2/KNTT-2.1.tex #627
%[Trương Đăng Khoa]
	Xét tính bị chặn của dãy số $\left(u_n\right)$ với $u_n=\sin n+ \cos n$.
	\dapso{dãy số bị chặn}
	\loigiai{Ta có 
		\begin{eqnarray*}
            \sin n+\cos n &=& \sqrt{2}\left(\dfrac{1}{\sqrt{2}}\sin n+\dfrac{1}{\sqrt{2}}\cos n\right)\\
			              &=& \sqrt{2}\left(\sin n\cdot\cos \dfrac{\pi}{4}+\cos n\cdot\sin \dfrac{\pi}{4}\right) =\sqrt{2}\sin \left(n+\dfrac{\pi}{4}\right).
		\end{eqnarray*}
		\begin{eqnarray*}
			-1\leq \sqrt{2}\sin \left(n+\dfrac{\pi}{4}\right) \leq 1 &\Leftrightarrow&-\sqrt{2}\leq \sqrt{2} \sin \left(n+\dfrac{\pi}{4}\right)\leq \sqrt{2}\\
			&\Leftrightarrow& -\sqrt{2}\leq \sin n+\cos n \leq \sqrt{2},\ \forall n\in\mathbb{N}^\ast\\
				&\Leftrightarrow&-\sqrt{2}\leq u_n \leq \sqrt{2},\ \forall n\in\mathbb{N}^\ast.
		\end{eqnarray*}
		Vậy dãy số $\left(u_n\right)$ là dãy số bị chặn.}
\end{bt}

\begin{bt}%[1D2H2-7]
% nguon: Archive/2024/TOAN_11_XE/TOAN_11_XE/data4/11C2/KNTT-2.1.tex #628
%[Trương Đăng Khoa]%Ví dụ 1
	Một chồng cột gỗ được xếp thành các lớp, hai lớp liên tiếp hơn kém nhau một cột gỗ.
	\begin{center}
		\begin{tikzpicture}[font=\footnotesize, line join=round, line cap=round, >=stealth,scale=0.6]
			\def\r{0.2}
			\def\n{25}
			\def\g{110}
			\fill[teal!50!green](-6*\r,-3*\r)rectangle(3.5*\n*\r,0.5*\n*\r);
			\fill[teal!50!green,opacity=0.25](3.5*\n*\r,3*\r)rectangle(-6*\r,2*\n*\r);
			\foreach \j in {0,...,12}{
				\pgfmathsetmacro{\m}{\n-\j}
				\foreach \i in{0,...,\m}{
					\fill[left color=orange, right color=teal!30,draw=brown](2*\i*\r,0)++(60:2*\j*\r)++(\g:\r)--++(\g-90:6)arc(\g:-60:\r)--++(\g-270:6)--cycle;
					\fill[orange!20!brown!40,draw=teal](2*\i*\r,0)++(60:2*\j*\r)circle(\r);
				}
			}
		\end{tikzpicture}
	\end{center}
	\begin{enumerate}
		\item  Gọi $u_1=25$ là số cột gỗ có ở hàng dưới cùng của chồng cột gỗ, $u_n$ là số cột gỗ có ở hàng thứ $n$ tính từ dưới lên trên. Xét tính tăng, giảm của dãy số này.
		\item  Gọi $v_1=14$ là số cột gỗ có ở hàng trên cùng của chồng cột gỗ, $v_n$ là số cột gỗ có ở hàng thứ $n$ tính từ trên xuống dưới. Xét tinh tăng, giảm của dãy số này.
	\end{enumerate}
	\loigiai{
		\begin{enumerate}
			\item Ta có $u_n=26-n>u_{n+1}=26-n-1=25-n$.\\
			Vậy dãy số $(u_n)$ là dãy số giảm.
			\item Ta có $v_n=13+n<v_{n+1}=13+n+1=14+n$.\\
			Vậy dãy số $(u_n)$ là dãy số tăng
		\end{enumerate}
	}
\end{bt}

\begin{bt}%[1D2H1-6]
% nguon: Archive/2024/TOAN_11_XE/TOAN_11_XE/data4/11C2/KNTT-2.1.tex #629
%[Trương Đăng Khoa]%Ví dụ 2
	Trên lưới ô vuông, mỗi ô cạnh $1$ đơn vị, người ta vẽ $8$ hình vuông và tô màu khác nhau như hình vẽ. Tìm dãy số biểu diễn độ dài cạnh của $8$ hình vuông đó từ nhỏ đến lớn. Có nhận xét gì về dãy số trên?
	\begin{center}
		\begin{tikzpicture}[scale=0.8]
			\def\r{21}
			\def\hv(#1){
				\ifnum #1= 1\else
				\pgfmathsetmacro{\R}{250*rnd}
				\pgfmathsetmacro{\G}{250*rnd}
				\pgfmathsetmacro{\B}{250*rnd}
				\definecolor{mau}{RGB}{\R,\G,\B}
				\fill[mau!30](0,0)rectangle(\r,\r);
				\draw[red,line width=1pt] (0,0) arc(180:90:\r)(0,0)rectangle(\r,\r);
				\pgfmathtruncatemacro{\k}{#1-1}
				\begin{scope}[shift={(45:\r*sqrt(2))},rotate=-90,scale={(sqrt(5)-1)/2}]
					\hv(\k)
					\pgfmathsetmacro{\n}{int((1/(sqrt(5))*(((1+sqrt(5))/2)^(\k)-(1-(sqrt(5))/2)^(\k)+1)}
					\ifnum \k>1
					\path (\r/2,\r/2)node[scale=1.75]{\color{red}$\n$};
					\else
					\fi
				\end{scope}
				\fi
			}
			\begin{scope}[scale=0.35]
				\hv(9)
				\draw[teal](0,0)grid(34,21);
				\path(21/2,21/2)node[scale=2]{21};
			\end{scope}
		\end{tikzpicture}
	\end{center}
	\loigiai{
		\begin{multicols}{4}
			\begin{itemize}
				\item $u_1=1$.
				\item $u_2=1$.
				\item $u_3=2$.
				\item $u_4=3$.
				\item $u_5=5$.
				\item $u_6=8$.
				\item $u_7=13$.
				\item $u_8=21$.
			\end{itemize}
		\end{multicols}
		Ta có dãy số $\left(u_n\right)\colon\heva{& u_1=1\\ &u_2=1\\ &u_n=u_{n-1}-u_{n-2}.}$
	}
\end{bt}

\begin{bt}%[1D2V3-8]
% nguon: Archive/2024/TOAN_11_XE/TOAN_11_XE/data4/11C2/KNTT-2.1.tex #630
%[Trương Đăng Khoa]% Ví dụ 3
	Chị Mai gửi tiền tiết kiệm vào ngân hàng theo thể thức lãi kép như sau. Lần đầu chị gửi $100$ triệu đồng. Sau đó, cứ hết $1$ tháng chị lại gửi thêm vào ngân hàng $6$ triệu đồng. Biết lãi suất của ngân hàng là $0{,}5\%$ một tháng. Gọi $P_n$ (triệu đồng) là số tiền chị có trong ngân hàng sau $n$ tháng.
	\begin{enumerate}
		\item Tính số tiền chị có trong ngân hàng sau $1$ tháng.
		\item Tính số tiền chị có trong ngân hàng sau $3$ tháng.
		\item Dự đoán công thức của $P_n$ tính theo $n$.
	\end{enumerate}
	\loigiai{
		\begin{enumerate}
			\item Số tiền chị có trong ngân hàng sau $1$ tháng là $P_1=+100+100\cdot0{,}5\%+6=100{,}5+6$ (triệu đồng).
			\item Số tiền chị có trong ngân hàng sau 2 tháng là 
			\begin{eqnarray*}
				P_2&=&100{,}5+6+(100{,}5+6)\cdot 0{,}5\%+6\\
				&=&(100{,}5+6)(1+0{,}5\%)+6\\
				&=& 100{,}5(1+0{,}5\%)+6\cdot(1+0{,}5\%)+6\, (\text{triệu đồng}).
			\end{eqnarray*}
			Số tiền chị có trong ngân hàng sau $3$ tháng là
			\begin{eqnarray*}
				P_3&=&(100{,}5+6)(1+0{,}5 \%)+6+[(100{,}5+6)(1+0{,}5 \%)+6] \cdot 0{,}5 \%+6\\
				&=& 100{,}5 \cdot(1+0{,}5 \%)^2+6(1+0{,}5 \%)^2+6 \cdot(1+0{,}5 \%)+6 \text{(triệu đồng)}.
			\end{eqnarray*}
			\item Số tiền chị có trong ngân hàng sau $4$ tháng là
			\begin{eqnarray*}
				P_4&=&(100{,}5+6)(1+0{,}5 \%)^2+6 \cdot(1+0{,}5 \%)+6+\left[(100{,}5+6)(1+0{,}5 \%)^2+6 \cdot(1+0{,}5 \%)+6\right]\cdot 0{,}5 \%+6\\ 
				&=&100{,}5 \cdot(1+0{,}5 \%)^3+6 \cdot(1+0{,}5 \%)^3+6\cdot(1+0{,}5 \%)^2+6 \cdot(1+0{,}5 \%)+6\, (\text{triệu đồng}).
			\end{eqnarray*}
			Số tiền chị có trong ngân hàng sau $n$ tháng là
			$$P_n=100{,}5 \cdot(1+0{,}5 \%)^{n-1}+6\cdot(1+0{,}5 \%)^{n-1}+6\cdot(1+0{,}5 \%)^{n-2}+6 \cdot(1+0{,}5 \%)^{n-3}+\ldots+6$$ với mọi $n \in\mathbb{N}^\ast$.
		\end{enumerate}
	}
\end{bt}

\begin{bt}%[1D2V2-7]
% nguon: Archive/2024/TOAN_11_XE/TOAN_11_XE/data4/11C2/KNTT-2.1.tex #631
%[Trương Đăng Khoa]% Ví dụ 4
	Anh Thanh vừa được tuyển dụng vào một công ty công nghệ, được cam kết lương năm đầu sẽ là $200$ triệu đồng và lương mỗi năm tiếp theo sẽ được tăng thêm $25$ triệu đồng. Gọi $s_n$ (triệu đồng) là lương vào năm thứ $n$ mà anh Thanh làm việc cho công ty đó. Khi đó ta có
	$$s_1=200,\, s_n=s_{n-1}+25\, \text{với}\, n \ge 2.$$
	\begin{enumerate}
		\item Tính lương của anh Thanh vào năm thứ $5$ làm việc cho công ty.
		\item Chứng minh $(s_n)$ là dãy số tăng. Giải thích ý nghĩa thực tế của kết quả này. 
	\end{enumerate}
	\loigiai{
		\begin{enumerate}
			\item Ta có \begin{eqnarray*}
				s_2&=&s_1+25=200+25=225\\
				s_3&=&s_2+25=225+25=250\\
				s_4&=&s_3+25=250+25=275\\
				s_5&=&s_4+25=275+25=300. 
			\end{eqnarray*}
			Vậy lương của anh Thanh vào năm thứ $5$ làm việc cho công ty là $300$ triệu đồng.
			\item  Ta có $s_n=s_{n-1}+25\Leftrightarrow s_n-s_{n-1}=25>0$ với mọi $n\ge 2$, $n\in\mathbb{N}^\ast$.\\
			Tức là $s_n>s_{n-1}$ với mọi $n\ge 2$, $n\in\mathbb{N}^\ast$.\\
			Vậy $(s_n)$ là dãy số tăng.\\
			Điều này có nghĩa là mức lương hàng năm của anh Thanh tăng dần theo thời gian làm việc.
		\end{enumerate}
	}
\end{bt}

\begin{bt}%[1D2C1-6]
% nguon: Archive/2024/TOAN_11_XE/TOAN_11_XE/data4/11C2/KNTT-2.1.tex #633
%[Trương Đăng Khoa]%Ví dụ 6
	Chị Hương vay trả góp một khoản tiền $100$ triệu đồng và đồng ý trả dần $2$ triệu đồng mỗi tháng với lãi suất $0{,}8\%$ số tiền còn lại của mỗi tháng.
	Gọi $A_n$, ($n\in\mathbb{N}$) là số tiền còn nợ (triệu đồng) của chị Hương sau $n$ tháng.
	\begin{enumerate}
		\item Tìm lần lượt $A_0$, $A_1$, $A_2$, $A_3$, $A_4$, $A_5$, $A_6$ đễ tính số tiền còn nợ của chị Hương sau $6$ tháng.
		\item  Dự đoán hệ thức truy hồi đối với dãy số $(A_n)$.
	\end{enumerate}
	\loigiai{
		\begin{enumerate}
			\item  Ta có $A_0=100$ (triệu đồng).
			\begin{itemize}
				\item Tiền lãi chị Hương phải trả sau $1$ tháng là $100\cdot 0{,}8\%=0{,}8$ (triệu đồng).\\
				Do đó, số tiền gốc chị Hương trả được sau $1$ tháng là $2-0{,}8=1{,}2$ (triệu đồng).\\
				Khi đó, số tiền còn nợ của chị Hương sau $1$ tháng là 
				$A_1=100-1{,}2=98{,}8$ (triệu đồng).
				\item  Tiền lãi chị Hương phải trả sau $2$ tháng là $98{,}8\cdot 0{,}8\%=0{,}7904$ (triệu đồng).\\
				Do đó, số tiền gốc chị Hương trả được sau $2$ tháng là $2-0{,}7904=1{,}2096$ (triệu đồng).\\
				Khi đó, số tiền còn nợ của chị Hương sau $2$ tháng là 
				$A_2=98{,}8-1{,}2096=97{,}5904$ (triệu đồng).
				\item Tiền lãi chị Hương phải trả sau $3$ tháng là $97{,}5904\cdot 0{,}8\%=0{,}7807232$ (triệu đồng).\\
				Do đó, số tiền gốc chị Hương trả được sau $3$ tháng là $2-0{,}7807232=1{,}2192768$ (triệu đồng).\\
				Khi đó, số tiền còn nợ của chị Hương sau $3$ tháng là 
				$A_3=97{,}5904-1{,}2192768=96{,}3711232$ (triệu đồng).
				\item Tiền lãi chị Hương phải trả sau $4$ tháng là $96{,}3711232\cdot 0{,}8\%\approx 0{,}77097$ (triệu đồng).\\
				Do đó, số tiền gốc chị Hương trả được sau $4$ tháng là $2-0{,}77097=1{,}22903$ (triệu đồng).\\
				Khi đó, số tiền còn nợ của chị Hương sau $4$ tháng là 
				$A_4=96{,}3711232-1{,}22903=95{,}1420932$ (triệu đồng).
				\item Tiền lãi chị Hương phải trả sau $5$ tháng là $95{,}1420932\cdot 0{,}8\%\approx 0{,}76114$ (triệu đồng).\\
				Do đó, số tiền gốc chị Hương trả được sau $5$ tháng là $2-0{,}76114=1{,}23886$ (triệu đồng).\\
				Khi đó, số tiền còn nợ của chị Hương sau $5$ tháng là $A_5=95{,}1420932-1{,}23886=93{,}9032332$ (triệu đồng).
				\item Tiền lãi chị Hương phải trả sau $6$ tháng là $93{,}9032332\cdot 0{,}8\%\approx 0{,}75123$ (triệu đồng).\\
				Do đó, số tiền gốc chị Hương trả được sau $6$ tháng là $2-0{,}75123=1{,}24877$ (triệu đồng).\\
				Khi đó, số tiền còn nợ của chị Hương sau $6$ tháng là $A_6=93{,}9032332-1{,}24877=92{,}6544632$ (triệu đồng).
			\end{itemize}
			\item Dự đoán hệ thức truy hồi đối với dãy số $(A_n)$ là 
			\[A_0=100,\, A_n=A_{n-1}-\left(2-A_{n-1} \cdot 0{,}8 \%\right)=1{,}008 A_{n-1}-2\]
		\end{enumerate}
	}
\end{bt}

\begin{bt}%[1D2C3-4]
% nguon: Archive/SP_KhoLatex/ID6_old/1-GK1-KNTT-De15-NH23-24.tex #2712

	Tìm $4$ số hạng đầu của một cấp số nhân biết tổng $3$ số hạng đầu bằng $\dfrac{148}{9}$, đồng thời theo thứ tự chúng là số hạng thứ $1$, thứ $4$, thứ $8$ của một cấp số cộng có công sai khác $0$.
	\loigiai{
		Gọi $4$ số hạng đầu của cấp số nhân đã cho là $u_1$, $u_2$, $u_3$, $u_4$; công bội của cấp số nhân là $ q $, công sai của cấp số cộng là $ d $ $\left(d\neq0\right)$.\\
		Tổng $3$ số hạng đầu của cấp số nhân bằng $\dfrac{148}{9}$ nên $u_1+u_2+u_3=\dfrac{148}{9} \Leftrightarrow u_1+u_1\cdot q+u_1\cdot q^2=\dfrac{148}{9}\left(1\right)$.\\
		Do $u_1$, $u_2$, $u_3$ theo thứ tự chúng là số hạng thứ $1$, thứ $4$, thứ $8$ của một cấp số cộng có công sai $d\neq0$ nên\\
		$\heva{&u_1\cdot q=u_1+3d\left(2\right)\\&u_1\cdot q^2=u_1+7d\left(3\right)}$.\\
		Nhân phương trình $\left(2\right)$ với $ 7 $ và nhân phương trình $\left(3\right)$ với $ 3 $, sau đó trừ hai phương trình theo vế ta được $u_1\left(3q^2-7q+4\right)=0\left(4\right)$.\\
		Từ phương trình $\left(1\right)$ ta có $u_1\neq0$. Khi đó $\left(3\right)\Leftrightarrow 3q^2-7q+4=0 \Leftrightarrow \hoac{&q=1\\&q=\dfrac{4}{3}}$.
		\begin{itemize}
			\item[+)] Với $q=1$, thay vào $\left(1\right)$ suy ra $u_1=u_2=u_3=\dfrac{148}{27}$ (loại do $u_1$, $_2$, $u_3$ theo thứ tự chúng là số
			hạng thứ $1$, thứ $4$, thứ $8$ của một cấp số cộng có công sai $d\neq0$).
			\item[+)] Với $q=\dfrac{4}{3}$, thay vào $\left(1\right)$ suy ra $u_1=4$, $u_2=\dfrac{16}{3}$, $u_3=\dfrac{64}{9}$, $u_4=\dfrac{256}{27}$.
		\end{itemize}
			Vậy $4$ số hạng đầu của cấp số nhân là $u_1=4$, $u_2=\dfrac{16}{3}$, $u_3=\dfrac{64}{9}$, $u_4=\dfrac{256}{27}$.
	}
\end{bt}

\begin{bt}%[1D2H2-7]
% nguon: Archive/SP_KhoLatex/ID6_old/1-GK1-CanhDieu-De5-NH23-24.tex #2745
%[Pj17-0-CK2-NH22-23--TeamTeXHoa--VoVanTu]
	Đầu năm mới, anh Hùng muốn tiết kiệm để cuối năm mua xe máy mới với giá $60$ triệu. Bắt đầu từ tháng $1$ năm $2024$ anh tiết kiệm bằng cách gửi ngân hàng mỗi tháng $5$ triệu theo hình thức gửi góp, biết lãi suất ngân hàng là $0,5 \% /$ tháng và tiền lãi của tháng trước sẽ nhập vào gốc của tháng sau để tiếp tục tính lãi.
	\begin{enumerate}
		\item Em hãy lập công thức tính số tiền tiết kiệm được sau $n$ tháng $\left(n \in \mathbb{N}^{*}\right)$.
		\item Hỏi sau đúng $11$ tháng gửi tiết kiệm anh Hùng còn thiếu bao nhiêu tiền để mua được xe máy.
	\end{enumerate}
	\loigiai{
	\begin{enumerate}		
		\item Đặt $r$ là lãi suất hàng tháng. Khi đó số tiền thu được
		\begin{itemize}
			\item Sau tháng thứ 1: $P_{1}=5+5 r=5(1+r)$.
			\item Sau tháng thứ 2: $P_{2}=\left(P_{1}+5\right)(1+r)=5(1+r)^{2}+5(1+r)$.
			\item Sau tháng thứ 3: $P_{3}=\left(P_{2}+5\right)(1+r)=5(1+r)^{3}+5(1+r)^{2}+5(1+r)$.
			\item Sau tháng thứ $n: P_{n}=5(1+r)^{n}+5(1+r)^{n-1}+\ldots+5(1+r)=5(1+r) \dfrac{(1+r)^{n}-1}{r}$.
		\end{itemize}
		\item Tổng số tiền thu được sau đúng 11 tháng là $$P_{11}=5(1+0{,}5 \%) \dfrac{(1+0{,}5 \%)^{11}-1}{0{,}5 \%} \approx 56{,}6 ~(\text{triệu đồng}).$$
		Như vậy anh còn thiếu khoảng $3{,}34$ triệu nữa là đủ tiền mua xe.
	\end{enumerate} 
}
\end{bt}

\begin{bt}%[1D2V2-7]
% nguon: Archive/SP_KhoLatex/ID6_old/1-GK1-CTST-De10-NH23-24.tex #2755
%[Câu 37]
	Giá tiền công khoan giếng tại một cơ sở được tính như sau: Giá của mét khoan đầu tiên là $60000$ đồng và kể từ mét khoan thứ hai, giá của mỗi mét sau tăng thêm $7\%$ giá của mét khoan ngay trước đó. Hỏi nếu khoan $50$m thì hết bao nhiêu tiền?
\loigiai{
	Gọi $u_n$ là giá tiền của mét khoan thứ $n$.\\
	Khi đó ta có: $u_n=u_{n-1}+u_{n-1}\cdot 7\%=u_{n-1}\left(1+7\%\right)=1{,}07\cdot u_{n-1}\left(\forall n\ge 2,\,n\in \mathbb{N}\right)$.\\
	Suy ra dãy số $\left(u_n\right)$ lập thành cấp số nhân có số hạng đầu $u_1=60000$ và công bội $q=1{,}07$.\\
	Vậy số tiền phải trả khi khoan $50$m là\\
	$S_{50}=u_1+u_2+\cdots +u_{50}=\dfrac{u_1 \left(q^{50}-1\right)}{q-1} =\dfrac{60000\left(1{,}07^{50}-1\right)}{1{,}07-1}\approx 24391736$ đồng.}
\end{bt}

\begin{bt}%[1D2H3-6]
% nguon: Archive/SP_KhoLatex/ID6_old/1K2-ot.tex #2943
 
    Một hình vuông màu vàng có cạnh 1 đơn vị dài được chia thành chín hình vuông nhỏ hơn và hình vuông ở chính giữa được tô màu xanh như hình. Mỗi hình vuông màu vàng nhỏ hơn lại được chia thành chín hình vuông con, và mỗi hình vuông con ở chính giữa lại được tô màu xanh. Nếu quá trình này được tiếp tục lặp lại năm lần, thì tồng diện tích các hình vuông được tô màu xanh bao nhiêu?
   \begin{center}
   	\begin{center}
   		\begin{tikzpicture}[>=stealth,thick,scale=0.7]
   			\def\n{1}
   			\def\a{3}
   			\pgfmathsetmacro{\m}{int(3^(\n))}
   			\def\hv#1{
   				\ifnum#1>0
   				\fill[blue!50] (-\a/3,-\a/3) rectangle (\a/3,\a/3);
   				\pgfmathtruncatemacro{\k}{#1-1}
   				\foreach \i in {0,...,3}{\begin{scope}[shift={(90*\i:2)},scale=1/3]\hv{\k}\end{scope}}
   				\foreach \i in {0,...,3}{\begin{scope}[shift={(45+90*\i:{4/sqrt(2)})},scale=1/3]\hv{\k}\end{scope}}
   				\fi
   			}
   			\fill[orange!50] (-\a,-\a) rectangle (\a,\a);
   			\hv{\n}
   			\foreach \i in {0,1,...,\m}{
   				\draw[blue!50] 
   				({-\a+2*\i *\a/\m},\a)--++(270:2*\a)
   				(\a,{-\a+2*\i *\a/\m})--++(180:2*\a)
   				;
   			}
   		\end{tikzpicture}
   		\begin{tikzpicture}[>=stealth,thick,scale=0.7]
   			\def\n{2}
   			\def\a{3}
   			\pgfmathsetmacro{\m}{int(3^(\n))}
   			\def\hv#1{
   				\ifnum#1>0
   				\fill[blue!50] (-\a/3,-\a/3) rectangle (\a/3,\a/3);
   				\pgfmathtruncatemacro{\k}{#1-1}
   				\foreach \i in {0,...,3}{\begin{scope}[shift={(90*\i:2)},scale=1/3]\hv{\k}\end{scope}}
   				\foreach \i in {0,...,3}{\begin{scope}[shift={(45+90*\i:{4/sqrt(2)})},scale=1/3]\hv{\k}\end{scope}}
   				\fi
   			}
   			\fill[orange!50] (-\a,-\a) rectangle (\a,\a);
   			\hv{\n}
   			\foreach \i in {0,1,...,\m}{
   				\draw[blue!50] 
   				({-\a+2*\i *\a/\m},\a)--++(270:2*\a)
   				(\a,{-\a+2*\i *\a/\m})--++(180:2*\a)
   				;
   			}
   		\end{tikzpicture}
   		\begin{tikzpicture}[>=stealth,thick,scale=0.7]
   			\def\n{4}
   			\def\a{3}
   			\def\hv#1{
   				\ifnum#1>0
   				\fill[blue!50] (-\a/3,-\a/3) rectangle (\a/3,\a/3);
   				\pgfmathtruncatemacro{\k}{#1-1}
   				\foreach \i in {0,...,3}{\begin{scope}[shift={(90*\i:2)},scale=1/3]\hv{\k}\end{scope}}
   				\foreach \i in {0,...,3}{\begin{scope}[shift={(45+90*\i:{4/sqrt(2)})},scale=1/3]\hv{\k}\end{scope}}
   				\fi
   			}
   			\fill[orange!50] (-\a,-\a) rectangle (\a,\a);
   			\hv{\n}
   		\end{tikzpicture}
   	\end{center}
   \end{center}
    \loigiai{
        \begin{itemize}
        \item Diện tích ô vuông màu xanh sau lần phân chia thứ nhất là $\dfrac{1}{9}$, số ô vuông màu xanh được tạo thêm là $8^0$.
        \item Diện tích ô vuông màu xanh sau lần phân chia thứ hai là $\dfrac{1}{9^2}$, số ô vuông màu xanh được tạo thêm là $8^1$.
        \item $ \ldots $
        \item Diện tích ô vuông màu xanh sau lần phân chia thứ năm là $\dfrac{1}{9^5}$, số ô vuông màu xanh được tạo thêm là $8^4$.\\
        \item Tổng diện tích các ô vuông màu xanh là: $\dfrac{1}{9}+\dfrac{1}{9^2}\cdot 8^1+\ldots +\dfrac{1}{9^5}\cdot {8^4}=0{,}445$.
        \end{itemize}
    }
\end{bt}

\begin{bt}%[1D2H2-5]
% nguon: Archive/SP_KhoLatex/ID6_old/1K2-ot.tex #2944

    Tìm giá trị $ x $ dương nhỏ nhất thỏa mãn ba số $\sin x$, $\sin 2x$, $\sqrt{3}\cos x$ lập thành cấp số cộng.
    \loigiai{
        Theo bài ra, ba số $\sin x$, $\sin 2x$, $\sqrt{3}\cos x$ lập thành cấp số cộng nên suy ra
        \allowdisplaybreaks
        \begin{eqnarray*}
        && \sin x+\sqrt{3}\cos x=2\sin 2x\\
        &\Leftrightarrow& \dfrac{1}{2}\sin x+\dfrac{\sqrt{3}}{2}\cos x=\sin x\cos \dfrac{\pi }{3}+\sin \dfrac{\pi }{3}\cdot\cos x=\sin 2x\\
        &\Leftrightarrow& \sin \left( x+\dfrac{\pi }{3} \right)=\sin 2x\Leftrightarrow \hoac{
            & 2x=x+\dfrac{\pi }{3}+k2\pi \\
            & 2x=\pi -x-\dfrac{\pi }{3}+k2\pi \\
        }\Leftrightarrow \hoac{
            & x=\dfrac{\pi }{3}+k2\pi \\
            & x=\dfrac{2\pi }{9}+\dfrac{k2\pi }{3} \\
        },\ \left( k\in \mathbb{Z} \right).
    \end{eqnarray*}
        Nghiệm dương $ x $ nhỏ nhất sẽ ứng với $k=0$. Vậy $x=\dfrac{\pi }{3}$ hoặc $x=\dfrac{2\pi }{9}$.}
\end{bt}

\begin{bt}%[1D2V2-2]
% nguon: Archive/SP_KhoLatex/ID6_old/1K2-ot.tex #2945

    Chứng minh rằng ba số dương $a$, $b$, $c$ theo thứ tự lập thành một cấp số cộng khi và chỉ khi các số $\dfrac{1}{\sqrt{b}+\sqrt{c}}$, $\dfrac{1}{\sqrt{c}+\sqrt{a}}$, $\dfrac{1}{\sqrt{a}+\sqrt{b}}$ theo thứ tự lập thành cấp số cộng.
    \loigiai{
        Ba số $\dfrac{1}{\sqrt{b}+\sqrt{c}}$, $\dfrac{1}{\sqrt{c}+\sqrt{a}}$, $\dfrac{1}{\sqrt{a}+\sqrt{b}}$ lập thành cấp số cộng khi và chỉ khi
        \allowdisplaybreaks
        \begin{eqnarray*}
        && \dfrac{1}{\sqrt{c}+\sqrt{a}}-\dfrac{1}{\sqrt{b}+\sqrt{c}}=\dfrac{1}{\sqrt{a}+\sqrt{b}}-\dfrac{1}{\sqrt{b}+\sqrt{c}}\\
        &\Leftrightarrow& \dfrac{\sqrt{b}-\sqrt{a}}{\left( \sqrt{c}+\sqrt{a} \right)\left( \sqrt{b}+\sqrt{c} \right)}=\dfrac{\sqrt{c}-\sqrt{b}}{\left( \sqrt{a}+\sqrt{b} \right)\left( \sqrt{c}+\sqrt{a} \right)}\\
        &\Leftrightarrow& \left( \sqrt{b}-\sqrt{a} \right)\left( \sqrt{b}+\sqrt{a} \right)=\left( \sqrt{c}-\sqrt{b} \right)\left( \sqrt{c}+\sqrt{b} \right)\\
        &\Leftrightarrow& b-a=c-b\Leftrightarrow a,\ b,\ c \text{ lập thành cấp số cộng.}
        \end{eqnarray*}}
\end{bt}

\begin{bt}%[1D2H2-4]
% nguon: Archive/SP_KhoLatex/ID6_old/1K2-ot.tex #2946
 :
    Chu vi của một đa giác là $ 45 $ cm, số đo các cạnh của nó lập thành một cấp số cộng với công sai $d=3$ cm. Biết cạnh lớn nhất là $ 15 $ cm, tính số cạnh của đa giác đó.
    \loigiai{
        Gọi cạnh nhỏ nhất của đa giác là $u_1$ và số cạnh của đa giác là $ n $.\\
        Ta có $15=u_1+\left( n-1 \right)\cdot 3$ hay $u_1=18-3n>0\Rightarrow  n<6$.\\
        Tổng các cạnh là $ 45 $ cm, ta có $45=\dfrac{n\left( 15+18-3n \right)}{2}$ hay $3n^2-33n+90=0$.\\
        Giải phương trình với $n\in \mathbb{N}^*;n<6$, ta được $n=5$.}
\end{bt}

\begin{bt}%[1D2H2-4]
% nguon: Archive/SP_KhoLatex/ID6_old/1K2-ot.tex #2947

    Tìm ba số hạng liên tiếp của một cấp số cộng, biết rằng:
    \begin{enumerate}
    \item Tổng của chúng bằng $ 15 $ và tích của chúng bằng $ 105 $.
    \item Tổng của chúng bằng $ 21 $ và tổng bình phương của chúng bằng $ 155 $.
    \end{enumerate}
    \loigiai{
        Gọi ba số hạng liên tiếp của cấp số cộng là $a,b,c\Rightarrow a+c=2b\left( * \right)$.
        \begin{enumerate}
        \item Theo bài ra, ta có $\heva{
            & a+b+c=15 \\
            & abc=105 \\
        }$, kết hợp với $\left( * \right)$, ta được 
        $$\heva{& a+b+c=15 \\& a+c=2b \\& abc=105}
        \Leftrightarrow \heva{& 3b=15 \\& a+c=2b \\& abc=105}
        \Leftrightarrow \heva{& b=5 \\& c=10-a \\& 5a\left( 10-a \right)=105}
        \Leftrightarrow \heva{& a=3 \\& b=5 \\& c=7} \text{ hoặc } \heva{& a=7 \\& b=5 \\& c=3.}$$
        \item Theo bài ra, ta có $\heva{& a+b+c=21 \\& a^2+b^2+c^2=155}$, kết hợp với $\left( * \right)$, ta được 
        $\heva{& a+b+c=21 \\& a+c=2b \\& a^2+b^2+c^2=155}$
        $$\Leftrightarrow \heva{& 3b=21 \\& a+c=2b \\& a^2+b^2+c^2=155}
        \Leftrightarrow \heva{& b=7 \\& c=14-a \\& a^2+\left( 14-a \right)^2+7^2=155}
        \Leftrightarrow \heva{& a=5 \\& b=7 \\& c=9} \text{ hoặc } \heva{& a=9 \\& b=7 \\& c=5.}$$
        \end{enumerate}}
\end{bt}

\begin{bt}%[1D2V2-2]
% nguon: Archive/SP_KhoLatex/ID6_old/1K2-ot.tex #2948
 
    Cho ba số $a$, $b$, $c$ theo thứ tự là ba số hạng liên tiếp của một cấp số cộng. Chứng minh rằng $a^2+2bc=c^2+2ab$ và $a^2+8bc=\left( 2b+c \right)^2$.
    \loigiai{
        Vì $a$, $b$, $c$ theo thứ tự lập thành cấp số cộng suy ra $a+c=2b$.\\
        Ta có $$a^2+2bc-c^2-2ab=\left( a-c \right)\left( a+c \right)+2b\left( c-a \right)=\left( a-c \right)\left( a+c-2b \right)=0.$$
        Suy ra $a^2+2bc-c^2-2ab=0\Leftrightarrow {a^2}+2bc=c^2+2ab $ điều phải chứng minh.\\
        Lại có $$a^2+8bc-\left( 2b+c \right)^2=a^2+4\left( a+c \right)c-\left( a+2c \right)^2=a^2+4ac+4c^2-a^2-4ac-4c^2.$$
        Suy ra $a^2+8bc-\left( 2b+c \right)^2=0\Leftrightarrow {a^2}+8bc=\left( 2b+c \right)^2 $ điều phải chứng minh.}
\end{bt}

\begin{bt}%[1D2H3-6]
% nguon: Archive/SP_KhoLatex/ID6_old/1K2-ot.tex #2950

    Cho cấp số nhân có $\heva{& u_1+u_5=51 \\& u_2+u_6=102.}$
    \begin{enumerate}
    \item Tìm số hạng đầu tiên và công bội.
    \item Tính tổng của $ 10 $ số hạng đầu tiên.
    \item Tổng của bao nhiêu số hạng đầu sẽ bằng $ 765 $.
    \item Số $ 12\ 288 $ là số hạng thứ mấy?
    \end{enumerate}
    \loigiai{
        Ta có $\heva{
            & u_1+u_5=51 \\
            & u_2+u_6=102}
        \Leftrightarrow \heva{
            & u_1+u_1 q^4=51 \\
            & u_1q+u_1 q^5=102}
        \Leftrightarrow \heva{
            & u_1\left( 1+q^4 \right)=51 \\
            & u_1q\left( 1+q^4 \right)=102}
        \Leftrightarrow \heva{& u_1=3 \\& q=2.}$
        \begin{enumerate}
        \item Vậy số hạng đầu $u_1=3$ và công bội $q=2$.
        \item Tổng của $ 10 $ số hạng đầu tiên $S_{10}=u_1\cdot\dfrac{1-q^{10}}{1-q}=3\cdot\dfrac{1-2^{10}}{1-2}=3069$.
        \item Ta có $S_n=u_1\cdot\dfrac{1-q^n}{1-q}=3\cdot\dfrac{1-2^n}{1-2}=765\Leftrightarrow n=8$.\\ 
        Vậy tổng của $ 8 $ số hạng đầu tiên bằng $ 765 $.
        \item Giả sử $u_n=12\ 288$. Theo công thức tổng quát của cấp số nhân, ta có
        $$u_n=u_1\cdot q^{n-1}\Leftrightarrow 12\ 288=3\cdot 2^{n-1}\Leftrightarrow n=13.$$
        Vậy $12\ 288$ là số hạng thứ $ 13 $ của cấp số nhân.
        \end{enumerate}}
\end{bt}

\begin{bt}%[1D2H3-2]
% nguon: Archive/SP_KhoLatex/ID6_old/1K2-ot.tex #2951

    Tìm số hạng đầu và công bội của cấp sống nhân $\left( u_n \right)$, biết
    \begin{enumEX}{2}
    \item $\heva{& u_5-u_1=15 \\& u_4-u_2=6.}$  
    \item $\heva{& {u_{20}}=8{u_{17}} \\& u_3+u_5=240.}$
    \item $\heva{& u_1-u_3+u_5=65 \\& u_1+u_7=325.}$        
    \item $\heva{& u_2-u_4+u_5=10 \\& u_3-u_5+u_6=20.}$
    \end{enumEX}
    \loigiai{
    \begin{enumEX}{1}
        \item Ta có $\heva{
            & u_5-u_1=15 \\
            & u_4-u_2=6}
        \Leftrightarrow \heva{
            & u_1q^4-u_1=15 \\
            & u_1q^3-u_1q=6}
        \Leftrightarrow \heva{
            u_1\left( q^4-1 \right)=15 & \left( 1 \right) \\
            u_1q\left( {q^2}-1 \right)=6. & \left( 2 \right)}$\\
        Lấy $\left( 1 \right)$ chia $\left( 2 \right)$, ta được $\dfrac{q^2+1}{q}=\dfrac{15}{6}\Leftrightarrow 2q^2-5q+2=0\Leftrightarrow q=2$ hoặc $q=\dfrac{1}{2}$. \\
        Vậy $\heva{& u_1=1 \\& q=2}$ hoặc $\heva{& u_1=-16 \\& q=\dfrac{1}{2}.}$\\
        \item Ta có $\heva{
            & u_{20}=8u_{17} \\
            & u_3+u_5=240}
        \Leftrightarrow \heva{
            & u_1{q^{19}}=8u_1{q^{16}} \\
            & u_1{q^2}+u_1q^4=240}
        \Leftrightarrow \heva{
            & q^3=8 \\
            & u_1{q^2}+u_1q^4=240}
        \Leftrightarrow \heva{
            & u_1=12 \\
            & q=2.}$\\
        \item Ta có $\heva{
            & u_1-u_3+u_5=65 \\
            & u_1+u_7=325}
        \Leftrightarrow \heva{
            & u_1-u_1{q^2}+u_1q^4=65 \\
            & u_1+u_1{q^6}=325}
        \Leftrightarrow \heva{
            u_1\left( 1-q^2+q^4 \right)=65 & \left( 1 \right) \\
            u_1\left( 1+q^6 \right)=325. & \left( 2 \right)}$\\
        Lấy $\left( 2 \right)$ chia $\left( 1 \right)$, ta được $\dfrac{1+q^6}{1-q^2+q^4}=\dfrac{325}{65}\Leftrightarrow 1+q^2=5\Leftrightarrow q=\pm 2$. \\
        Vậy $\heva{& u_1=5 \\& q=2}$ hoặc $\heva{& u_1=5 \\& q=-2.}$
        \item Ta có $\heva{
            & u_2-u_4+u_5=10 \\
            & u_3-u_5+u_6=20}
        \Leftrightarrow \heva{
            & u_1q-u_1q^3+u_1q^4=10 \\
            & u_1{q^2}-u_1q^4+u_1{q^5}=20}
        \Leftrightarrow \heva{
            & u_1q\left( 1-q^2+q^3 \right)=10 & \left( 1 \right) \\
            & u_1{q^2}\left( 1-q^2+q^3 \right)=20. & \left( 2 \right)}$\\
        Lấy $\left( 2 \right)$ chia $\left( 1 \right)$, ta được $q=2$. Vậy $\heva{
            & u_1=1 \\
            & q=2.}$
        \end{enumEX}}
\end{bt}

\begin{bt}%[1D2H3-6]
% nguon: Archive/SP_KhoLatex/ID6_old/1K2-ot.tex #2952

    Tìm số hạng đầu và công bội của cấp sống nhân $\left( u_n \right)$, biết
    \begin{enumEX}{2}
    \item $\heva{& u_1=1 \\ & S_8=\dfrac{3^8-1}{2}.}$
    \item $\heva{& {S_4}=40 \\& S_8=680.}$
    \end{enumEX}
    \loigiai{
    \begin{enumEX}{1}
        \item Ta có $\heva{& u_1=1 \\& {S_8}=\dfrac{3^8-1}{2}}
        \Leftrightarrow \heva{& u_1=1 \\& u_1\cdot\dfrac{q^8-1}{q-1}=\dfrac{3^8-1}{2}}
        \Leftrightarrow \heva{& u_1=1 \\& q=3.}$\\
        \item Ta có $\heva{& S_4=40 \\& S_8=680}
        \Leftrightarrow \heva{
            & u_1\cdot\dfrac{q^4-1}{q-1}=40 & \left( 1 \right) \\
            & u_1\cdot\dfrac{q^8-1}{q-1}=680. & \left( 2 \right)}$\\
        Lấy $\left( 2 \right)$ chia cho $\left( 1 \right)$, ta được $q^4+1=17\Leftrightarrow q^4=16\Leftrightarrow q=\pm 2$.\\
        Vậy $\heva{& u_1=\dfrac{8}{3} \\& q=2}$ hoặc $\heva{& u_1=-8 \\& q=-2.}$
    \end{enumEX}}
\end{bt}

\begin{bt}%[1D2H3-7]
% nguon: Archive/SP_KhoLatex/ID6_old/1K2-ot.tex #2953

    Tìm ba số khác nhau tạo thành cấp số cộng có tổng bằng $ 6 $, biết rằng nếu hoán đổi vị trí số hạng thứ nhất và số hạng thứ hai đồng thời giữ nguyên số hạng thứ ba ta được cấp số nhân.
    \loigiai{
        Gọi ba số cần tìm là $u_1$, $u_2$, $u_3$ với $u_1\ne u_2\ne u_3\ne u_1$.\\
        $u_1$, $u_2$, $u_3$ tạo thành cấp số cộng với công sai $d\ne 0$ nên $u_1$, $u_2=u_1+d$, $u_3=u_1+2d$.\\
        Hơn nữa, $$u_1+u_2+u_3=6\Leftrightarrow u_1+\left( u_1+d \right)+\left( u_1+2d \right)=6\Leftrightarrow u_1+d=2.$$
        $u_2$, $u_1$, $u_3$ tạo thành cấp số nhân hay $u_1+d$, $u_1$, $u_1+2d$ tạo thành cấp số nhân 
        \begin{eqnarray*}
        &\Leftrightarrow& \left( u_1+d \right)\left( u_1+2d \right)=u_1^2
        \Leftrightarrow \left( u_1+d \right)\left( u_1+d+d \right)=u_1^2\\
        &\Leftrightarrow& 2\left( 2+2-u_1 \right)=u_1^2\Leftrightarrow u_1^2+2u_1-8=0\Leftrightarrow \hoac{
                & u_1=2 \\
                & u_1=-4.}
        \end{eqnarray*}
        Với $u_1=2$, suy ra $d=0$ không thỏa mãn.\\
        Với $u_1=-4$, suy ra $d=6$. \\
        Vậy ba số cần tìm là $-4$, $2$, $8$.}
\end{bt}

\begin{bt}%[1D2N2-4]
% nguon: Archive/SP_KhoLatex/ID6_old/1K2-2.tex #3050
[NB]%[DCHT Toán 11 - KNTT -Lê Hải Phụng] 
	Trong các dãy số sau, dãy số nào là một cấp số cộng?
	\begin{listEX}[1]
		\item $1$, $-3$, $-7$, $-11$, $-15$, $\ldots$;
		\item $1$, $-2$, $-4$, $-6$, $-8,$ $\ldots$.
		\item $ \dfrac{1}{2} $, $0$, $-\dfrac{1}{2}$, $-1$, $-\dfrac{3}{2}$, $\ldots$
	\end{listEX}
	\dapso{1) và 3) là cấp số cộng.}
	\loigiai{Ta lần lượt đi kiểm tra: $ u_2-u_1=u_3-u_2=u_4-u_3=\ldots $?\\
		Xét từng dãy số thì ta thấy 1) và 3) là cấp số cộng. 
	}
\end{bt}

\begin{bt}%[1D2N2-4]
% nguon: Archive/SP_KhoLatex/ID6_old/1K2-2.tex #3051
[NB]%[DCHT Toán 11 - KNTT -Lê Hải Phụng] 
	Trong các dãy số sau, dãy nào là cấp số cộng. Tìm số hạng đầu và công sai của cấp số cộng đó.
	\begin{listEX}[2]
		\item Dãy số $ (u_n) $ với $ u_n=19n-5 $;
		\item Dãy số $ (u_n) $ với $ u_n=n^2+n+1 $. 
	\end{listEX}
	\dapso{Dãy số 1) $ (u_n) $ là một cấp số cộng với số hạng đầu là $ u_1=19\cdot1-5=14 $ và công sai $ d=19 $. Dãy số 2) không là một cấp số cộng.}
	\loigiai{
		\begin{enumerate}
			\item Dãy số $ (u_n) $ với $ u_n=19n-5 $.\\
			Ta có $ u_{n+1}-u_n=19(n+1)-5-(19n-5)=19 $. Vậy $ (u_n) $ là một cấp số cộng với số hạng đầu là $ u_1=19\cdot1-5=14 $ và công sai $ d=19 $.
			\item Dãy số $ (u_n) $ với $ u_n=n^2+n+1 $.\\
			Ta có $ u_{n+1}-u_n=(n+1)^2+(n+1)+1-(n^2+n+1)=2n+2$ phụ thuộc vào $ n $. Vậy $ (u_n) $ không là một cấp số cộng.
	\end{enumerate}}
\end{bt}

\begin{bt}%[1D2H2-4]
% nguon: Archive/SP_KhoLatex/ID6_old/1K2-2.tex #3052
[TH]%[DCHT Toán 11 - KNTT -Lê Hải Phụng] 
	Cho cấp số cộng $\left(u_n\right)$ với $u_1=3$, $u_2=9$. Công sai của cấp số cộng đã cho bằng bao nhiêu?
	\dapso{Công sai của cấp số cộng đã cho là 6.}
	\loigiai{Cấp số cộng $(u_n)$ có số hạng tổng quát là
		$u_n=u_1+\left(n-1\right)d$ với $n \ge 2$
		(số hạng đầu $u_1$ và công sai $d$)\\
		Suy ra $ u_2=u_1+d\Leftrightarrow9=3+d\Leftrightarrow d=6 $.\\
		Vậy công sai của cấp số cộng đã cho là 6.
	}
\end{bt}

\begin{bt}%[1D2H2-4]
% nguon: Archive/SP_KhoLatex/ID6_old/1K2-2.tex #3053
[TH]%[DCHT Toán 11 - KNTT -Lê Hải Phụng] 
	Tìm số hạng đầu tiên, công sai của cấp số cộng sau $ \heva{&u_5=19\\&u_9=35.}$
	
	\dapso{Số hạng đầu tiên $ u_1=3 $, công sai $ d=4 $.}
	\loigiai{Áp dụng công thức $ u_n=u_1+(n-1)d $ ta có $\heva{&u_5=19\\&u_9=35} \Leftrightarrow \heva{&u_1+4d=19\\&u_1+8d=35} \Leftrightarrow \heva{&u_1=3\\&d=4.}$}
	
	Vậy số hạng đầu tiên $ u_1=3 $, công sai $ d=4 $.
\end{bt}

\begin{bt}%[1D2V2-4]
% nguon: Archive/SP_KhoLatex/ID6_old/1K2-2.tex #3055
[VD]%[DCHT Toán 11 - KNTT -Lê Hải Phụng] 
Cho cấp số cộng $ (u_n) $ thỏa mãn $ \heva{&u_2-u_3+u_5=10\\&u_4+u_6=26} $. Xác định số hạng đầu $ u_1 $ và công sai $ d $ cấp số cộng.         
\dapso{$ u_1=1 $, $ d=3 $.}
\loigiai{Ta có $ \heva{&u_2-u_3+u_5=10\\&u_4+u_6=26} \Leftrightarrow \heva{& u_1+d-(u_1+2d)+u_1+4d=10 \\ & u_1+3d+u_1+5d=26} \Leftrightarrow \heva{& u_1+3d=10 \\ & u_1+4d=13} \Leftrightarrow \heva{u_1=1\\d=3.}$}
\end{bt}

\begin{bt}%[1D2C2-4]
% nguon: Archive/SP_KhoLatex/ID6_old/1K2-2.tex #3056
[VDC]%[DCHT Toán 11 - KNTT -Lê Hải Phụng] 
Tính số hạng đầu $ u_1 $ và công sai $d$ của một cấp số cộng biết $ \heva{&u_1+u_2+u_3=27\\&u_1^2+u_2^2+u_3^2=275} $

\dapso{$ u_1=5 $, $ d=4 $ hoặc $ u_1=13 $, $ d=-4 $.}
\loigiai{Ta có $ \heva{&u_1+u_2+u_3=27\\&u_1^2+u_2^2+u_3^2=275} \Leftrightarrow \heva{&u_2-d+u_2+u_2+d=27\\&(u_2-d)^2+u_2^2+(u_2+d)^2=275}\Leftrightarrow \heva{&u_2=9\\&3u_2^2+2d^2=275.}$}

Thay $ u_2=9 $ vào $ 3u_2^2+2d^2=275 $ ta được $ d=4 $ hay $ d=-4 $.

Vậy $ u_1=5 $, $ d=4 $ hoặc $ u_1=13 $, $ d=-4 $.
\end{bt}

\begin{bt}%[1D2N3-2]
% nguon: Archive/SP_KhoLatex/ID6_old/1K2-3.tex #3117
[NB]%[DCHT Toán 11 - KNTT -Tên Huỳnh Thanh Chí]
	Dãy số $25$; $5$; $1$; $\dfrac{1}{5}$; $\ldots$ có phải là một cấp số nhân không? Nếu có hãy tìm công bội của cấp số nhân đó.
	\dapso{Dãy số $25$; $5$; $1$; $\dfrac{1}{5}$; $\ldots$ là một cấp số nhân với công bội $ q=\dfrac{1}{5} $.}
	\loigiai{
	Ta có $ \dfrac{u_2}{u_1}=\dfrac{u_3}{u_2}=\ldots=\dfrac{1}{5} $ là một số không đổi.\\
	Do đó dãy số $25$; $5$; $1$; $\dfrac{1}{5}$; $\ldots$ là một cấp số nhân với công bội $ q=\dfrac{1}{5} $.
	}
\end{bt}

\begin{bt}%[1D2H3-2]
% nguon: Archive/SP_KhoLatex/ID6_old/1K2-3.tex #3118
[NB]%[DCHT Toán 11 - KNTT -Tên Huỳnh Thanh Chí]
	Dãy số $1$; $n$; $n^2$; $n^3$; $ n^4 $; $\ldots$ (với $ n>1 $) có phải là một cấp số nhân không? Nếu có hãy tìm công bội của cấp số nhân đó.
	\dapso{Dãy số $1$; $n$; $n^2$; $n^3$; $ n^4 $; $\ldots$ (với $ n>1 $) là một cấp số nhân với công bội $ q=n $.}
	\loigiai{
		Ta có $ \dfrac{u_2}{u_1}=\dfrac{u_3}{u_2}=\ldots=n $ (với $ n>1 $) là một số không đổi.\\
		Do đó dãy số $1$; $n$; $n^2$; $n^3$; $ n^4 $; $\ldots$ (với $ n>1 $) là một cấp số nhân với công bội $ q=n $.
	}
\end{bt}

\begin{bt}%[1D2H3-4]
% nguon: Archive/SP_KhoLatex/ID6_old/1K2-3.tex #3119
[TH]%[DCHT Toán 11 - KNTT -Tên Huỳnh Thanh Chí]
	Cho dãy số $ (u_n) $ được xác định bởi $ \heva{& u_1=2\\ & u_{n+1}=u_n^2} $. Hỏi dãy số $ (u_n) $ có là một cấp số nhân hay không?
	\dapso{Dãy số $ (u_n) $ không là một cấp số nhân.}
	\loigiai{
	Ta có $ u_2=u_1^2=4,u_3=u_2^2=16,u_4=u_3^2=256 $.\\
	Suy ra $ \dfrac{u_2}{u_1}=2 $; $ \dfrac{u_3}{u_2}=4 $ và $ \dfrac{u_4}{u_3}=16 $. Vì $ \dfrac{u_2}{u_1}\ne \dfrac{u_3}{u_2}\ne \dfrac{u_4}{u_3} $ nên $ (u_n) $ không là một cấp số nhân	
	}
\end{bt}

\begin{bt}%[1D2H3-2]
% nguon: Archive/SP_KhoLatex/ID6_old/1K2-3.tex #3120
[TH]%[DCHT Toán 11 - KNTT -Tên Huỳnh Thanh Chí]
	Cho dãy số $ (u_n) $, biết $ u_1=2 $ và $ u_{n+1}=\dfrac{1}{3}u_n $. Chứng minh $ (u_n) $ là một cấp số nhân và tìm số hạng $ u_3 $.
	\dapso{}
	\loigiai{
	Ta có $ u_{n+1}=\dfrac{1}{3}u_n\Rightarrow \dfrac{u_{n+1}}{u_n}=\dfrac{1}{3} $ là một số không đổi nên $ (u_n) $ là một cấp số nhân với công bội là $ q=\dfrac{1}{3} $.\\
	Do đó $ u_3=u_2\cdot q=u_1\cdot q^2=2\cdot \dfrac{1}{3^2}=\dfrac{2}{9} $.
	}
\end{bt}

\begin{bt}%[1D2H3-2]
% nguon: Archive/SP_KhoLatex/ID6_old/1K2-3.tex #3121
[TH]%[DCHT Toán 11 - KNTT -Tên Huỳnh Thanh Chí]
	Cho $ (u_n) $ là cấp số nhân có công bội $ q\ne 0,u_1\ne 0 $. Chứng minh rằng dãy số $ (v_n) $ với $ v_n=\dfrac{u_nu_{2n+1}}{4} $ cũng là một cấp số nhân.
	\dapso{$ (v_n) $ là một cấp số nhân với công bội là $ q^3 $.}
	\loigiai{
		Ta có $ \dfrac{v_n}{v_{n-1}}=\dfrac{\dfrac{u_nu_{2n+1}}{4}}{\dfrac{u_{n-1}u_{2(n-1)+1}}{4}}=\dfrac{u_1q^{n-1}\cdot u_1q^{2n}}{u_1q^{n-2}\cdot u_1q^{2n-2}}=q^3 $. Do đó $ (v_n) $ là một cấp số nhân với công bội là $ q^3 $.}
\end{bt}

\begin{bt}%[1D2V3-2]
% nguon: Archive/SP_KhoLatex/ID6_old/1K2-3.tex #3122
[VD]%[DCHT Toán 11 - KNTT -Tên Huỳnh Thanh Chí]
	Cho dãy số $ (u_n) $ được xác định bởi $ \heva{& u_1=3\\ & u_{n+1}=2u_n-2},\forall n\in\mathbb{N}^* $. Chứng minh rằng dãy số $ (v_n) $ xác định bởi $ v_n=2u_n-4,\forall n\in\mathbb{N}^* $ là một cấp số nhân. Hãy xác định số hạng đầu và công bội của cấp số nhân đó.
	\dapso{$ (v_n) $ là cấp số nhân với công bội $ q=2 $ và số hạng đầu $ v_1=2 $.}
	\loigiai{
	Ta có $ v_n=2u_n-4 $ (1) và $ v_{n+1}=2u_{n+1}-4 $ (2).\\
	Theo đề ta có $ u_{n+1}=2u_n-2 \Rightarrow 2u_{n+1}-4=2(2u_n-4) $ (3).\\
	Thay (1) và (2)	vào (3) ta được $ v_{n+1}=2v_n \Rightarrow \dfrac{v_{n+1}}{v_n}=2,\forall n\in\mathbb{N}^* $.\\
	Suy ra $ (v_n) $ là cấp số nhân với công bội $ q=2 $ và số hàng đầu $ v_1=2u_1-4=2\cdot 3-4=2 $.
	}
\end{bt}

\begin{bt}%[1D2N3-4]
% nguon: Archive/SP_KhoLatex/ID6_old/1K2-3.tex #3128
[NB]%[DCHT Toán 11 - KNTT -Tên Huỳnh Thanh Chí]
	Tìm số hạng thứ $ 100 $ của cấp số nhân $ 8;-4;2;-1;\ldots $
	\dapso{$ u_{100}=-\dfrac{1}{2^{96}} $.}
	\loigiai{
	Cấp số nhân này có số hạng đầu $ u_1=8 $ và công bội $ q=\dfrac{-4}{8}=-\dfrac{1}{2} $.\\
	Do đó số hạng tổng quát $ u_n=8\cdot \left(-\dfrac{1}{2}\right)^{n-1} $.\\
	Vậy $ u_{100}=8\cdot \left(-\dfrac{1}{2}\right)^{100-1}=8\cdot \left(-\dfrac{1}{2}\right)^{99}=-\dfrac{1}{2^{96}} $.
	}
\end{bt}

\begin{bt}%[1D2H3-3]
% nguon: Archive/SP_KhoLatex/ID6_old/1K2-3.tex #3129
[NB]%[DCHT Toán 11 - KNTT -Tên Huỳnh Thanh Chí]
	Tìm số hạng tổng quát của dãy số $ 3;12;48;192;\ldots $, biết dãy $ (u_n) $ là một cấp số nhân.  
	\dapso{$ u_n=3\cdot 4^{n-1} $.}
	\loigiai{
		Vì dãy số $ (u_n) $ là một cấp số nhân nên $ q=\dfrac{u_2}{u_1}=\dfrac{12}{3}=4 $ và số hạng đầu $ u_1=3 $.\\
		Do đó dãy số $ 3;12;48;192;\ldots $ là một cấp số nhân có số hạng tổng quát là $ u_n=u_1q^{n-1}=3\cdot 4^{n-1} $.}
\end{bt}

\begin{bt}%[1D2H3-3]
% nguon: Archive/SP_KhoLatex/ID6_old/1K2-3.tex #3130
[TH]%[DCHT Toán 11 - KNTT -Tên Huỳnh Thanh Chí]
	Tìm số hạng tổng quát của cấp số nhân, biết $ \heva{& u_1+u_3=51\\ & u_2+u_4=153.} $
	\dapso{$ u_n=\dfrac{51}{10}\cdot 3^{n-1} $.}
	\loigiai{
		Vì $ (u_n) $ là một cấp số nhân nên $ u_n=u_1\cdot q^{n-1} $.\\
		Ta có $ \heva{& u_1+u_3=51\\ & u_2+u_4=153}\Leftrightarrow\heva{& u_1+u_1q^2=51\\ & u_1q+u_1q^3=153}\Leftrightarrow\heva{& u_1(1+q^2)=51 \qquad (1)\\ & u_1q(1+q^2)=153 \qquad (2).} $\\
		Chia từng vế của $ (2) $ cho $ (1) $ ta được $ \dfrac{u_1q(1+q^2)}{u_1(1+q^2)}=\dfrac{153}{51} \Leftrightarrow q=3 $.\\
		Suy ra $ u_1=\dfrac{51}{1+q^2}=\dfrac{51}{10} $, $ u_n=u_1\cdot q^{n-1}=\dfrac{51}{10}\cdot 3^{n-1} $.\\
		Vậy số hạng tổng quát $ u_n=\dfrac{51}{10}\cdot 3^{n-1} $.
	}
\end{bt}

\begin{bt}%[1D2H3-3]
% nguon: Archive/SP_KhoLatex/ID6_old/1K2-3.tex #3131
[TH]%[DCHT Toán 11 - KNTT -Tên Huỳnh Thanh Chí]
	Tìm số hạng đầu, công bội và số hạng tổng quát của cấp số nhân, biết $ \heva{& u_3=15\\ & u_6=120.} $
	\dapso{$ u_1=\dfrac{15}{4} $, $ q=2 $ và $ u_n=15\cdot 3^{n-3} $.}
	\loigiai{
		Vì $ (u_n) $ là một cấp số nhân nên $ u_n=u_1\cdot q^{n-1} $.\\	
		Ta có $ \heva{& u_3=15\\ & u_6=120}\Leftrightarrow \heva{& u_1q^2=15 \qquad (1)\\ & u_1q^5=120\qquad (2).} $\\
		Chia từng vế của $ (2) $ cho $ (1) $ ta được $ \dfrac{u_1q^5}{u_1q^2}=\dfrac{120}{15} \Leftrightarrow q^3=8 \Leftrightarrow q=2 $.\\
		Suy ra $ u_1=\dfrac{15}{q^2}=\dfrac{15}{2^2}=\dfrac{15}{4} $, $ u_n=u_1\cdot q^{n-1}=\dfrac{15}{4}\cdot 2^{n-1}=15\cdot 2^{n-3} $.\\
		Vậy $ u_1=\dfrac{15}{4} $, $ q=2 $ và $ u_n=15\cdot 3^{n-3} $.
	}
\end{bt}

\begin{bt}%[1D2H3-3]
% nguon: Archive/SP_KhoLatex/ID6_old/1K2-3.tex #3132
[TH]%[DCHT Toán 11 - KNTT -Tên Huỳnh Thanh Chí]
	Tìm số hạng tổng quát của cấp số nhân, biết $ \heva{& u_4=35\\ & u_8=560.} $
	\dapso{$ u_n=35\cdot 2^{n-4} $ với $ q=2 $ hoặc $ u_n=35\cdot (-2)^{n-4} $ với $ q=-2 $.}
	\loigiai{
		Vì $ (u_n) $ là một cấp số nhân nên $ u_n=u_1\cdot q^{n-1} $.\\	
		Ta có $ \heva{& u_4=35\\ & u_8=560}\Leftrightarrow \heva{& u_1q^3=35 \qquad (1)\\ & u_1q^7=560\qquad (2).} $\\
		Chia từng vế của $ (2) $ cho $ (1) $ ta được $ \dfrac{u_1q^7}{u_1q^3}=\dfrac{560}{35} \Leftrightarrow q^4=16 \Leftrightarrow \hoac{& q=2\\ & q=-2.} $\\
		Với $ q=2 $. Suy ra $ u_1=\dfrac{35}{q^3}=\dfrac{35}{8} $, $ u_n=u_1\cdot q^{n-1}=\dfrac{35}{8}\cdot 2^{n-1}=35\cdot 2^{n-4} $.\\
		Với $ q=-2 $. Suy ra $ u_1=\dfrac{35}{q^3}=-\dfrac{35}{8} $, $ u_n=u_1\cdot q^{n-1}=-\dfrac{35}{8}\cdot (-2)^{n-1}=35\cdot (-2)^{n-4} $.\\
		Vậy $ u_n=35\cdot 2^{n-4} $ với $ q=2 $ hoặc $ u_n=35\cdot (-2)^{n-4} $ với $ q=-2 $.
	}
\end{bt}

\begin{bt}%[1D2V3-8]
% nguon: Archive/SP_KhoLatex/ID6_old/1K2-3.tex #3133
[VD]%[DCHT Toán 11 - KNTT -Tên Huỳnh Thanh Chí]
	Trong một lọ nuôi cấy vi khuẩn, ban đầu có $ 5\ 000 $ con vi khuẩn và số lượng vi khuẩn tăng lên thêm $ 8\% $ mỗi giờ. Hỏi sau $ 5 $ giờ thì số lượng vi khuẩn là bao nhiêu?
	\dapso{}
	\loigiai{
	Ta có $ A=5\ 000 $ là số lượng vi khuẩn ban đầu, $ r=8\%=0{,}08 $ là tỉ lệ gia tăng vi khuẩn sau một giờ.\\
	Vì sau mỗi giờ, số lượng vi khuẩn tăng lên bằng số lượng vi khuẩn trước đó cộng cho số lượng vi khuẩn tăng lên. Nên ta có:
	\begin{itemize}
	\item Tại thời điểm ban đầu: $ T_0=5\ 000 $.
	\item Tại thời điểm sau $ 1 $ giờ: $ T_1=5\ 000+5\ 000\cdot 0{,}08= 5\ 000\cdot(1{,}08)$.
	\item Tại thời điểm sau $ 2 $ giờ: $ T_2=T_1+T_1\cdot\cdot 0{,}08=T_1\cdot (1{,}08)=5\ 000\cdot(1{,}08)^2 $.
	\end{itemize}
	Do đó ta có thể nhận thấy rằng, số lượng vi khuẩn ở thời gian $ n $ giờ là một cấp số nhân có số hạng đầu $ u_1=5\ 000 $ và công bội $ q=1{,}08 $.\\
	Suy ra số hạng tổng quát $ u_n=5\ 000\cdot (1{,}08)^{n-1} $.\\
	Vậy số lượng vi khuẩn sau $ 5 $ giờ là $ u_5=5\ 000\cdot (1{,}08)^{5-1}\approx 6802 $ (vi khuẩn).
	}
\end{bt}

\begin{bt}%[1D2H3-2]
% nguon: Archive/SP_KhoLatex/ID6_old/1K2-3.tex #3139
[TH]
	Cho cấp số nhân $(u_n)$ biết $\heva{&u_4-u_2=25 \\ &u_3-u_1=50.}$
	\begin{enumEX}{1}
		\item Tìm số hạng đầu và công bội của cấp số nhân $(u_n)$.
		\item Tìm số hạng thứ $8$ của cấp số nhân $(u_n)$.
	\end{enumEX}
	\dapso{$\heva{&q=\dfrac{1}{2} \\ &u_1=-200}$, $u_8=-\dfrac{25}{16}$}
	\loigiai{
		\begin{enumerate}
			\item Ta có $\heva{&u_4-u_2=25 \\ &u_3-u_1=50} 
			\Leftrightarrow \heva{&u_1(q^3-q)=25 \\ &u_1(q^2-q)=50} 
			\Rightarrow \heva{&q=\dfrac{1}{2} \\ &u_1=-200.}$
			\item Ta có $u_8=u_1\cdot q^7=-200\cdot \dfrac{1}{2^7}=-\dfrac{25}{16}$.
		\end{enumerate}
	}
\end{bt}

\begin{bt}%[1D2H3-4]
% nguon: Archive/SP_KhoLatex/ID6_old/1K2-3.tex #3140
[TH]
	Tìm số hạng thứ $10$ của cấp số nhân $(u_n)$ biết $\heva{&u_4-u_2=72 \\ &u_5-u_3=144.}$	
	\dapso{$u_{10}=6144$}
	\loigiai{
		Ta có $\heva{&u_4-u_2=72 \\ &u_5-u_3=144}\Leftrightarrow \heva{&u_4-u_2=72 \\ &q(u_4-u_2)=144}
		\Rightarrow \heva{&q=2 \\ &u_1(q^3-q)=72} \Leftrightarrow \heva{&q=2 \\ &u_1=12.}$ \\
		Khi đó $u_{10}=u_1\cdot q^9=6144$.
	}
\end{bt}

\begin{bt}%[1D2H3-4]
% nguon: Archive/SP_KhoLatex/ID6_old/1K2-3.tex #3141
[TH]
	Cho một cấp số nhân có $5$ số hạng biết $2$ số hạng đầu là số dương, tích số hạng đầu và số hạng thứ $3$ là $1$, tích số hạng thứ $3$ và số hạng cuối là $\dfrac{1}{16}$. Tìm cấp số nhân này.	
	\dapso{$2; 1; \dfrac{1}{2}; \dfrac{1}{4}; \dfrac{1}{8}$}
	\loigiai{
		Gọi $5$ số hạng cần tìm có dạng $\dfrac{x}{q^2}$; $\dfrac{x}{q}$; $x$; $xq$; $xq^2$.\\
		Theo đề ra ta có $\heva{&\dfrac{x}{q^2}\cdot x=1 \\ &x\cdot xq^2=\dfrac{1}{16}} 
		\Leftrightarrow \heva{&x=\dfrac{1}{2} \\ &q=\dfrac{1}{2}}$ (do hai số hạng đầu dương nên $q>0$). \\
		Vậy $5$ số hạng cần tìm là $2; 1; \dfrac{1}{2}; \dfrac{1}{4}; \dfrac{1}{8}$.
	}
\end{bt}

\begin{bt}%[1D2H3-2]
% nguon: Archive/SP_KhoLatex/ID6_old/1K2-3.tex #3142
[TH]
	Tìm số hạng đầu và công bội của cấp số nhân $(u_n)$ biết $\heva{&u_2+u_5-u_4=10 \\ &u_3+u_6-u_5=20.}$
	\dapso{$\heva{&q=2 \\ &u_1=1}$}
	\loigiai{
		Ta có $\heva{&u_2+u_5-u_4=10 \\ &u_3+u_6-u_5=20} 
		\Leftrightarrow \heva{&u_1(q+q^4-q^3)=10 \\ &u_1(q^2+q^5-q^4)=20}
		\Leftrightarrow \heva{&q=2 \\ &u_1=1.}$
	}
\end{bt}

\begin{bt}%[1D2H3-4]
% nguon: Archive/SP_KhoLatex/ID6_old/1K2-3.tex #3143
[TH]
	Tìm $5$ số lập thành một cấp số nhân có công bội bằng $\dfrac{1}{4}$ số thứ nhất và tổng $2$ số đầu là $\dfrac{5}{4}$.
	\dapso{$1; \dfrac{1}{4}; \dfrac{1}{16}; \dfrac{1}{64}; \dfrac{1}{128}$ hoặc $-5; -\dfrac{5}{4}; -\dfrac{5}{16}; -\dfrac{5}{64}; -\dfrac{1}{128}$}
	\loigiai{
		Theo đề, ta có $\heva{&q=\dfrac{1}{4}u_1 \\ &u_1+u_2=\dfrac{5}{4}}
		\Leftrightarrow \heva{&q=\dfrac{1}{4}u_1 \\ &u_1+u_1\cdot q=\dfrac{5}{4}}
		\Leftrightarrow \heva{&q=\dfrac{1}{4}u_1 \\ &u_1^2+4u_1-5=0}
		\Leftrightarrow \heva{&q=\dfrac{1}{4} \\ &u_1=1}$ hoặc $\heva{&q=-\dfrac{5}{4} \\ &u_1=-5.}$\\
		Vậy có hai CSN là $1; \dfrac{1}{4}; \dfrac{1}{16}; \dfrac{1}{64}; \dfrac{1}{128}$ và $-5; -\dfrac{5}{4}; -\dfrac{5}{16}; -\dfrac{5}{64}; -\dfrac{1}{128}$.
	}
\end{bt}

\begin{bt}%[1D2H3-4]
% nguon: Archive/SP_KhoLatex/ID6_old/1K2-3.tex #3144
[TH]
	Tìm $3$ số lập thành một cấp số nhân có tổng là $63$ và tích là $1728$.
	\dapso{$3; 12; 48$}
	\loigiai{
		Gọi ba số cần tìm là $\dfrac{x}{q}; x; xq$. Theo đề ra, ta có $x^3=1728\Rightarrow x=12$. \\
		Mặt khác $\dfrac{x}{q}+x+xq=63\Leftrightarrow 12q+12+\dfrac{12}{q}=63
		\Leftrightarrow 12q^2-51q+12=0 \Leftrightarrow \hoac{&q=4 \\ &q=\dfrac{1}{4}\cdot}$ \\
		Vậy CSN cần tìm là $3; 12; 48$.
	}
\end{bt}

\begin{bt}%[1D2H3-5]
% nguon: Archive/SP_KhoLatex/ID6_old/1K2-3.tex #3150
[TH]
	Xác định $x$ dương để $2x-3$; $x$; $2x+3$ lập thành cấp số nhân.
	\dapso{$x=\sqrt{3}$}
	\loigiai{
		Ba số $2x-3$; $x$; $2x+3$ lập thành cấp số nhân khi $x^2=(2x-3)(2x+3) \Leftrightarrow x=\pm \sqrt{3}$.\\
		Do $x>0$ nên chọn $x=\sqrt{3}$.
	}
\end{bt}

\begin{bt}%[1D2H3-4]
% nguon: Archive/SP_KhoLatex/ID6_old/1K2-3.tex #3151
[TH]
	Cho cấp số nhân $x, 12, y, 192$. Tìm $x$ và $y$.
	\dapso{$\heva{&x=-3 \\ &y=-48}$}	
	\loigiai{
		Bốn số $x, 12, y, 192$ lập thành CSN khi $\heva{&xy=12^2 \\ &y^2=12\cdot 192}
		\Leftrightarrow \heva{&x=3 \\ &y=48}$ hoặc $\heva{&x=-3 \\ &y=-48.}$
	}
\end{bt}

\begin{bt}%[1D2H3-5]
% nguon: Archive/SP_KhoLatex/ID6_old/1K2-3.tex #3152
[TH]
	Tìm $x$ để dãy số $1$, $x^2$, $6-x^2$ lập thành cấp số nhân.
	\dapso{$x=\pm \sqrt{2}$}
	\loigiai{
		Ta có $1, x^2, 6-x^2$ lập thành cấp số nhân $\Leftrightarrow x^4=6-x^2 \Leftrightarrow x= \pm \sqrt{2}$.
	}
\end{bt}

\begin{bt}%[1D2H3-4]
% nguon: Archive/SP_KhoLatex/ID6_old/1K2-3.tex #3153
[TH]
	Viết $6$ số xen giữa hai số $-2$ và $256$ để được một cấp số nhân có $8$ số hạng. Tìm cấp số nhân này.
	\dapso{$-2; 4; -8; 16; -32; 64; -128; 256$}
	\loigiai{
		Theo đề ra, ta có $\heva{&u_1=-2 \\ &u_8=256} \Leftrightarrow \heva{&u_1=-2 \\ &u_1\cdot q^7=256}
		\Leftrightarrow \heva{&u_1=-2 \\ &q=-2.}$ \\
		Cấp số nhân cần tìm là $-2; 4; -8; 16; -32; 64; -128; 256$.
	}
\end{bt}

\begin{bt}%[1D2V3-5]
% nguon: Archive/SP_KhoLatex/ID6_old/1K2-3.tex #3154
[VD]%
	Bốn góc của một tứ giác lồi lập thành một cấp số nhân, góc lớn nhất gấp $8$ lần góc nhỏ nhất. Tìm $4$ góc đó.
	\dapso{$24^\circ; 48^\circ; 96^\circ; 192^\circ$}	
	\loigiai{
		Giả sử $4$ góc của tứ giác là $A\leq B\leq C\leq D$. Suy ra $A+B+C+D=360^\circ$.\\
		Theo đề, ta có $D=8A \Leftrightarrow Aq^3=8A \Leftrightarrow q=2$. Khi đó, ta được
		$$A(1+q+q^2+q^3)=360^\circ \Rightarrow A=24^\circ.$$
		Vậy $4$ góc của tứ giác lần lượt là $24^\circ; 48^\circ; 96^\circ; 192^\circ$.
	}
\end{bt}

\begin{bt}%[1D2V3-5]
% nguon: Archive/SP_KhoLatex/ID6_old/1K2-3.tex #3155
[VD]%
	Tìm tất cả các giá trị của tham số $m$ để phương trình sau có ba nghiệm phân biệt lập thành một cấp số nhân $x^3-7mx^2+2(m^2+6 m)x-64=0$.
	\dapso{$m-8$}
	\loigiai{
		+ Điều kiện cần: \\
		Giả sử phương trình đã cho có ba nghiệm phân biệt $x_1; x_2; x_3$ lập thành một cấp số nhân. \\
		Theo định lý Vi-ét, ta có $x_1 \cdot x_2 \cdot x_3=64$. \\
		Theo tính chất của cấp số nhân, ta có $x_1\cdot x_3=x_2^2$. Suy ra ta có $x_2^3=64 \Leftrightarrow x_2=4$. \\
		Thay $x=4$ vào phương trình đã cho ta được 
		$$4^3-7m\cdot 4^2+2(m^2+6m) \cdot 4-64=0
		\Leftrightarrow m^2-8m=0
		\Leftrightarrow \hoac{&m=0 \\ &m=8.}$$
		+ Điều kiện đủ: \\
		Với $m=0$ thay vào phương trình đã cho ta được: $x^3-64=0$ hay $x=4$
		(nghiệm kép-loại). \\
		Với $m=8$ thay vào phương trình đã cho nên ta có phương trình $x^3-56x^2+224x-64=0$. \\
		Phương trình này có $3$ nghiệm phân biệt lập thành cấp số nhân. \\
		Vậy $m=8$ là giá trị cần tìm.
	}
\end{bt}

\begin{bt}%[1D2N3-6]
% nguon: Archive/SP_KhoLatex/ID6_old/1K2-3.tex #3161
%[DCHT Toán 11 - KNTT -Đỗ Chí Tâm] 
	Một cấp số nhân có số hạng đầu $u_1=3$ và công bội $q=2 $. Tính tổng $8$ số hạng đầu của cấp số nhân.
	\dapso{$765$}		
	\loigiai{
		Ta có $S_8=\dfrac{u_1\left(1-q^8\right)}{1-q}=\dfrac{3 \left(1-2^8\right)}{1-2}=765$.	
	}	
\end{bt}

\begin{bt}%[1D2N3-6]
% nguon: Archive/SP_KhoLatex/ID6_old/1K2-3.tex #3162
%[DCHT Toán 11 - KNTT -Đỗ Chí Tâm]
	Một cấp số nhân có số hạng đầu $u_1=1$ và công bội $q=3$. Tính $S_{10}.$
	\dapso{$29524$}	
	\loigiai{
		Ta có $S_{10}=\dfrac{u_1\left(1-q^{10}\right)}{1-q}=1.\dfrac{1-3^{10}}{1-3}=29524$.	
	}	
\end{bt}

\begin{bt}%[1D2N3-6]
% nguon: Archive/SP_KhoLatex/ID6_old/1K2-3.tex #3163
%[DCHT Toán 11 - KNTT -Đỗ Chí Tâm] 
	Một cấp số nhân $(u_n)$ có $u_1=4$ và công bội $q=2$. Tính $S_{20}$.
	\dapso{$4194300$}	
	\loigiai{
		Ta có $S_{20}=\dfrac{u_1\left(1-q^{20}\right)}{1-q}=4194300$.}
\end{bt}

\begin{bt}%[1D3H1-5]
% nguon: Archive/SP_KhoLatex/ID6_old/1K2-3.tex #3164
%[DCHT Toán 11 - KNTT -Đỗ Chí Tâm] 	
	Tính tổng $S=1+\dfrac{1}{3}+\dfrac{1}{3^2}+\cdots+\dfrac{1}{3^n}+\cdots$.
	\dapso{$\dfrac{3}{2}$}			
	\loigiai{
		Đây là tổng của một cấp số nhân lùi vô hạn với $u_1=1, q=\dfrac{1}{3}$\\
		Suy ra   $S=\dfrac{u_1}{1-q}=\dfrac{1}{1-\dfrac{1}{3}}=\dfrac{3}{2}.$
	}
\end{bt}

\begin{bt}%[1D2H3-6]
% nguon: Archive/SP_KhoLatex/ID6_old/1K2-3.tex #3165
%[DCHT Toán 11 - KNTT -Đỗ Chí Tâm] 	
	Cho cấp số nhân có $q=-3, S_6=730$. Tính $u_1$.
	\dapso{$4$}		
	\loigiai{
		$S_6=u_1.\dfrac{1-q^6}{1-q}\Rightarrow u_1=S_6\cdot\dfrac{1-q}{1-q^6}=730\cdot \dfrac{1-(-3)}{1-(-3)^6}=4$.
	}
\end{bt}

\begin{bt}%[1D2H3-6]
% nguon: Archive/SP_KhoLatex/ID6_old/1K2-3.tex #3166
%[DCHT Toán 11 - KNTT -Đỗ Chí Tâm] 
	Một cấp số nhân $(u_n)$ có $u_1=-5$, $u_2=10$.Tính tổng của $15$ số hạng đầu của cấp số nhân đó.
	\dapso{$-54615$}	
	\loigiai{
		Công bội của cấp số nhân đã cho là: $q=\dfrac{u_2}{u_1}=\dfrac{10}{-5}=-2$.\\
		Tổng của $15$ số hạng đầu của cấp số nhân đó là $S_{15}=-5\cdot \dfrac{1-(-2)^{15}}{1-(-2)}=-54615$.
	}
\end{bt}

\begin{bt}%[1D2H3-6]
% nguon: Archive/SP_KhoLatex/ID6_old/1K2-3.tex #3167
%[DCHT Toán 11 - KNTT -Đỗ Chí Tâm] 
	Một cấp số nhân $(u_n)$ có $u_1=2$, $u_2=-2$.Tính tổng của $9$ số hạng đầu của cấp số nhân đó.
	\dapso{$2$}	
	\loigiai{
		Công bội của cấp số nhân đã cho là: $q=\dfrac{u_2}{u_1}=\dfrac{-2}{2}=-1$.\\
		Tổng của $9$ số hạng đầu của cấp số nhân đó là: $S_9=2\cdot \dfrac{1-(-1)^9}{1-(-1)}=2$.
	}
\end{bt}

\begin{bt}%[1D2V3-6]
% nguon: Archive/SP_KhoLatex/ID6_old/1K2-3.tex #3168
%[DCHT Toán 11 - KNTT -Đỗ Chí Tâm] 
	Một cấp số nhân $(u_n)$ có $u_3=8$, $u_5=32$ và công bội $q>0$. Tính tổng của $10$ số hạng đầu tiên của cấp số nhân.
	\dapso{$2046$}	
	\loigiai{
		$\heva{u_3&=8\\u_5&=32}\Leftrightarrow \heva{u_1.q^2&=8\\u_1.q^4&=32}\Rightarrow q^2=\dfrac{32}{8}=4\Rightarrow q=2, u_1=2$\\
		\hspace*{5.2cm}$\Rightarrow S_{10}=u_1.\dfrac{1-q^{10}}{1-q}=2.\dfrac{2-2^{10}}{1-2}=2046$.
	}
\end{bt}

\begin{bt}%[1D2V3-6]
% nguon: Archive/SP_KhoLatex/ID6_old/1K2-3.tex #3169
%[DCHT Toán 11 - KNTT -Đỗ Chí Tâm] 	
	Tính tổng $S=2+6+18+...+13122$.
	\dapso{$19682$}
	\loigiai{
		Xét cấp số nhân có $u_1=2, q=3$. Khi đó $13122=u_1.q^{n-1}\Leftrightarrow 13122=2.3^{n-1}\Leftrightarrow n=9$\\
		Vậy $S=S_9=u_1\dfrac{1-q^9}{1-q}=2.\dfrac{1-3^9}{1-3}=19682$
	}
\end{bt}

\begin{bt}%[1D2V3-6]
% nguon: Archive/SP_KhoLatex/ID6_old/1K2-3.tex #3170
%[DCHT Toán 11 - KNTT -Đỗ Chí Tâm] 	
	Tính tổng $S=1+2+4+8+\cdots+1024$.	
	\dapso{$2047$}
	\loigiai{
		Xét cấp số nhân có $u_1=1, q=2$. Khi đó $1024=u_1.q^{n-1}\Leftrightarrow 1024=1.2^{n-1}\Leftrightarrow n=11$.\\
		Vậy $S=S_{11}=u_1\dfrac{1-q^{11}}{1-q}=1.\dfrac{1-2^{11}}{1-2}=2047$.
	}
\end{bt}

\begin{bt}%[1D2V3-6]
% nguon: Archive/SP_KhoLatex/ID6_old/1K2-3.tex #3171
%[DCHT Toán 11 - KNTT -Đỗ Chí Tâm] 	
	Một cấp số nhân có $u_1=1, q=3$, biết $S_n=3280$. Tìm $n$.	
	\dapso{$8$}
	\loigiai{
		$S_n=u_1\dfrac{1-q^n}{1-q}=1.\dfrac{1-3^n}{1-3}=3280\Rightarrow n=8$.
	}
\end{bt}

\begin{bt}%[1D2V3-6]
% nguon: Archive/SP_KhoLatex/ID6_old/1K2-3.tex #3172
%[DCHT Toán 11 - KNTT -Đỗ Chí Tâm] 	
	Một cấp số nhân $(u_n)$ có $u_4+u_6=-540, u_3+u_5=180$. Tính $S_5$.
	\dapso{$122$}
	\loigiai{
		$\heva{u_4+u_6&=-540\\u_3+u_5&=180}\Leftrightarrow \heva{u_1.q^3+u_1.q^5&=-540\\u_1.q^2+u_1.q^4&=180} \Leftrightarrow \heva{u_1q^3(1+q^2)&=-540\\u_1q^2(1+q^2)&=180}\Leftrightarrow \heva{u_1&=2\\ q&=-3}$.\\
		Vậy $S_5=u_1\dfrac{1-q^5}{1-q}=2.\dfrac{1-(-3)^5}{1+3}=122$.
	}
\end{bt}

\begin{bt}%[1D2V3-6]
% nguon: Archive/SP_KhoLatex/ID6_old/1K2-3.tex #3173
%[DCHT Toán 11 - KNTT -Đỗ Chí Tâm] 
	Bốn số hạng liên tiếp của một cấp số nhân, trong đó số hạng thứ hai nhỏ hơn số hạng thứ nhất $35$, còn số hạng thứ ba lớn hơn số hạng thứ tư $560$. Tìm tổng của bốn số hạng trên, biết công bội mang giá trị dương.
	\dapso{$-\dfrac{2975}{3}$}
	\loigiai{
		Theo đề ta có $\heva{u_1-u_2&=35\\u_3-u_4&=560}\Leftrightarrow \heva{&u_1-u_1q=35\\&u_1q^2-u_1q^3=560}\Leftrightarrow \heva{&u_1(1-q)=35\,\,\,\,\,\,\, (1)\\&u_1q^2(1-q)=560\,\,\,\, (2)}$\\
		Thay $(1)$ vào $(2)$ ta được $q^2=16\Leftrightarrow q=\pm 4$.\\
		Với $q=4$ thay vào $(1)$ ta được $u_1=-\dfrac{35}{3}$.	\\
		$S_4=u_1.\dfrac{1-q^4}{1-q}=-\dfrac{2975}{3}$.
	}	
\end{bt}

\begin{bt}%[1D2C3-6]
% nguon: Archive/SP_KhoLatex/ID6_old/1K2-3.tex #3174
%[DCHT Toán 11 - KNTT -Đỗ Chí Tâm] 
	Tổng của một cấp số nhân lùi vô hạn bằng $ \dfrac{1}{4} $, tổng ba số hạng đầu tiên của cấp số nhân đó bằng $ \dfrac{7}{27} $. Tổng của số hạng đầu và công bội của cấp số nhân đó bằng
	\dapso{$S=0$}	
	\loigiai{Gọi $ u_1 $ và $ q $ với ($ |q|<1 $) lần lượt là số hạng đầu và cộng bội của cấp số nhân lùi vô hạn. Theo giả thiết, ta có $$\begin{cases}
			\dfrac{u_1}{q-1}=\dfrac{1}{4}\\
			u_1+u_1q+u_1q^2=\dfrac{7}{27}
		\end{cases} \Leftrightarrow \begin{cases}
			\dfrac{u_1}{q-1}=\dfrac{1}{4}\\
			u_1 (1-q^3)=\dfrac{7}{27}(1-q)
		\end{cases} \Leftrightarrow \begin{cases}\dfrac{u_1}{1-q}=\dfrac{1}{4}\\ q^3=-\dfrac{1}{27} \end{cases}\Leftrightarrow
		\heva{& u_1=\dfrac{1}{3}\\& q=-\dfrac{1}{3}.}$$ 
		Vậy $ u_1+q=0 $.
	}
\end{bt}

\begin{bt}%[1D2C3-8]
% nguon: Archive/SP_KhoLatex/ID6_old/1K2-3.tex #3175
%[DCHT Toán 11 - KNTT -Đỗ Chí Tâm] 
	Một du khách vào trường đua ngựa đặt cược, lần đầu đặt $20.000$ đồng, mỗi lần sau tiền đặt gấp đôi số tiền lần đặt trước. Người đó thua $10$ lần liên tiếp và thắng ở lần thứ $11$. Hỏi du khách trên thắng hay thua bao nhiêu tiền?
	\dapso{$20000$ đồng}
	\loigiai{
		Số tiền du khách đặt trong mỗi lần (kể từ lần đầu) là một cấp số nhân có $u_1=20.000$ và công bội $q=2$.\\
		Du khách thua trong 10 lần liên tiếp đầu tiên nên tổng số tiền thua là
		$$S_{10}=\dfrac{u_1(1-q^{10})}{1-q}=\dfrac{20000(1-2^{10})}{1-2}=20000(2^{10}-1) \text{(đồng).}$$ 
		Số tiền du khách thắng trong lần thứ 11 là $u_{11}=u_1q^{10}=20000.2^{10}$ (đồng).\\
		Ta có $u_{11}-S_{10}=20000>0$. Vậy du khách thắng $20000$ đồng.
	}
\end{bt}

\begin{bt}%[1D2V3-7]
% nguon: Archive/SP_KhoLatex/ID6_old/1K2-3.tex #3181
%[DCHT Toán 11 - KNTT -Đỗ Chí Tâm]
	Biết $x, y, x+4$ theo thứ tự lập thành cấp số cộng và $x+1, y+1, 2y+2$ theo thứ tự lập thành cấp số nhân với $x, y$ là số thực dương. Tính $x+y$.
	\dapso{$4$}
	\loigiai{
		Theo giả thiết ta có:\\
		$\heva{&x+(x+4)=2y\\&(x+1)(2y+2)=(y+1)^2}\Leftrightarrow \heva{&y=x+2\\&(x+1)(2x+6)=(x+3)^2}\Leftrightarrow \heva{&x=1\Rightarrow y=3\\&x=-3\Rightarrow y=-1}$\\
		Do đó $x+y=4$.
	}	
\end{bt}

\begin{bt}%[1D2V3-7]
% nguon: Archive/SP_KhoLatex/ID6_old/1K2-3.tex #3182
%[DCHT Toán 11 - KNTT -Đỗ Chí Tâm] 
	Cho $3$ số $a, b, c$ theo thứ tự tạo thành một cấp số nhân với công bội khác $1$. Biết cũng theo thứ tự đó chúng lần lượt là số hạng thứ nhất, thứ tư và thứ tám của một cấp số cộng với công sai $d\ne 0$. Tính $\dfrac{a}{d}$.
	\dapso{$9$}
	\loigiai{
		$a, b, c$ lần lượt là số hạng thứ nhất, thứ tư, thứ tám của một CSC với công sai $d$\\
		ta có $\heva{b &=a+3d\\c&=a+7d}$.\\
		Mặt khác $a, b, c$ là $3$ số hạng liên tiếp của CSN nên\\ $a.c=b^2\Leftrightarrow a(a+7d)=(a+3d)^2\Leftrightarrow a^2+7ad=a^2+6ad+9d^2\Leftrightarrow 9d^2=ad\Leftrightarrow \dfrac{a}{d}=9$.
	}	
\end{bt}

\begin{bt}%[1D2V3-7]
% nguon: Archive/SP_KhoLatex/ID6_old/1K2-3.tex #3183
%[DCHT Toán 11 - KNTT -Đỗ Chí Tâm] 
	Tìm tích các số dương $a$ và $b$ sao cho $a, a + 2b, 2a + b$ lập thành một cấp số cộng và $(b + 1)^2, ab + 5,(a + 1)^2$ lập thành một cấp số nhân.
	\dapso{$3$}
	\loigiai{
		Theo tính chất CSC ta có $a+(2a+b)=2(a+2b)\,\,\,\, (1)$\\
		Theo tính chất CSN ta có $(b+1)^2.(a+1)^2=(ab+5)^2\,\,\,\, (2)$\\
		Từ $(1)$ ta được $a=3b$, thay vào $(2)$ ta được $(b+1)^2(3b+1)^2=(3b^2+5)^2$\\
		$\Leftrightarrow \hoac{&(b+1)(3b+1)=(3b^2+5)\\& (b+1)(3b+1)=-(3b^2+5) \,\, (\text{Vô nghiêm})}\Leftrightarrow b=1, a=3\Rightarrow ab=3.$	
	}	
\end{bt}

\begin{bt}%[1D2V3-7]
% nguon: Archive/SP_KhoLatex/ID6_old/1K2-3.tex #3184
%[DCHT Toán 11 - KNTT -Đỗ Chí Tâm] 
	$a,b,c\,(a\ne b\ne c)$ là ba số hạng liên tiếp của một cấp số cộng và $b,c,a$ là ba số hạng liên tiếp của một cấp số nhân, đồng thời $a.b.c=125$. Tìm $a,b,c$.	
	\dapso{$5$}
	\loigiai{
		$a,b, c$ là ba số hạng liên tiếp của cấp số cộng, nên có $a+c=2b$.\\
		$b,c,a$ là ba số hạng liên tiếp của một cấp số nhân, nên có $b.a=c^2$.\\
		Ta có hệ $\heva{&a+c=2b\,\,\,\, (1)\\&b.a=c^2\,\,\,\,\,\,\,\,(2) \\&a.b.c=125\,\,\,\, (3)}$\\
		Thay $(2)$ vào $(3)$ ta được $c^3=125\Rightarrow c=5$\\
		Thay $c=5$ vào $(1), (2)$ ta được hệ $\heva{a+5&=2b\\ab&=25}\Leftrightarrow \heva{&a=2b-5\\&2b^2-5b-25=0}\Leftrightarrow \hoac{&b=5\Rightarrow a=5\\&b=-\dfrac{5}{2}\Rightarrow a=-10}$\\
		Vậy $a=-10, b=-\dfrac{5}{2}, c=5$.	
	}	
\end{bt}

\begin{bt}%[1D2V3-7]
% nguon: Archive/SP_KhoLatex/ID6_old/1K2-3.tex #3185
%[DCHT Toán 11 - KNTT -Đỗ Chí Tâm] 
	Một cấp số cộng và một cấp số nhân đều là các dãy tăng các số hạng thứ nhất của hai dãy số đều bằng $3$, các số hạng thứ hai bằng nhau. Tỉ số giữa các số hạng thứ ba của CSN và CSC là $\dfrac{9}{5}$. Tìm tích ba số hạng của cấp số cộng thỏa mãn tính chất trên.
	\dapso{$405$}
	\loigiai{
		Gọi $u_1, u_2, u_3$ là $3$ số hạng liên tiếp của CSC.\\
		Gọi $a_1, a_2, a_3$ là $3$ số hạng liên tiếp của CSN.\\
		Theo đề ta có hệ $\heva{&u_1=a_1=3\\&u_2=a_2\\&a_3=\dfrac{9}{5}u_3}\Leftrightarrow \heva{&u_1=a_1=3\\&3+d=3q\\&5(3q^2)=9(3+2d)}\Rightarrow q=3 \vee q=\dfrac{3}{5}$\\
		Chọn $q=3$ vì dãy tăng, khi đó $d=6$. \\
		Vậy $3$ số hạng của cấp số cộng là $3; 9; 15\Rightarrow 3\cdot9\cdot15=405$	
	}	
\end{bt}

\begin{bt}%[1D2V3-7]
% nguon: Archive/SP_KhoLatex/ID6_old/1K2-3.tex #3186
%[DCHT Toán 11 - KNTT -Đỗ Chí Tâm] 
	Một CSC và CSN đều có số hạng đầu tiên là bằng 5, số hạng thứ hai của CSC lớn hơn số hạng thứ hai của CSN là 10, còn các số hạng thứ 3 của hai cấp số thì bằng nhau. Tìm tổng các số hạng của cấp số cộng biết công bội của cấp số nhân không âm.
	\dapso{$75$}
	\loigiai{
		Gọi $u_1, u_2, u_3$ là $3$ số hạng liên tiếp của CSC với công sai $d$.\\
		Gọi $a_1, a_2, a_3$ là $3$ số hạng liên tiếp của CSN với công bội $q$.\\
		Theo đề bài ta có: $\heva{&u_1=a_1=5\\&u_2-a_2=10\\&u_3=a_3}\Leftrightarrow \heva{&u_1=a_1=5\\&u_1+d-a_1q=10\\&u_1+2d=a_1q^2}\Leftrightarrow \heva{&u_1=a_1=5\\&d=5+5q\\&5+2d=5q^2}\Rightarrow q=3\vee q=-1 (L)$\\
		Với $q=3\Rightarrow d=20$. Vậy CSC là $5;25;45 \Rightarrow S=5+25+45=75$
	}	
\end{bt}

\begin{bt}%[1D2C3-7]
% nguon: Archive/SP_KhoLatex/ID6_old/1K2-3.tex #3187
%[DCHT Toán 11 - KNTT -Đỗ Chí Tâm] 	
	Ba số khác nhau có tổng bằng $114$ có thể coi là ba số hạng liên tiếp của một CSN, hoặc coi là số hạng thứ nhất, thứ tư và thứ hai mươi lăm của một CSC. Tìm ba số đó.
	\dapso{$2; 14; 98$}
	\loigiai{
		Gọi $u_1, u_2, u_3$ là $3$ số hạng liên tiếp của CSN với công bội $q$.\\
		Theo đề $u_1=a_1, u_2=a_4, u_3=a_{25}$ với $a_1,a_4, a_{25}$ là 3 số hạng của CSC với công sai $d$.\\
		Ta có $\heva{&a_4=a_1+3d\\&a_{25}=a_1+24d}\Rightarrow 8a_4-a_{25}=7a_1\Leftrightarrow 8u_2-u_3=7u_1\Leftrightarrow 8u_1q-u_1q^2=7u_1$\\
		\hspace*{4.2cm}$\Leftrightarrow q^2-8q+7=0 \Leftrightarrow q=1 (L) \vee q=7 (N)$\\
		Theo đề ta cũng có $u_1+u_2+u_3=114\Leftrightarrow u_1+u_1q+u_1q^2=114\Rightarrow u_1=2$\\
		Vậy $3$ số cần tìm là $2; 14; 98$.
	}	
\end{bt}

\begin{bt}%[1D2C3-7]
% nguon: Archive/SP_KhoLatex/ID6_old/1K2-3.tex #3188
%[DCHT Toán 11 - KNTT -Đỗ Chí Tâm] 
	Ba số khác nhau có tổng là $217$ có thể coi là các số hạng liên tiếp của một CSN hoặc là các số hạng thứ $2$ thứ $9$ và thứ $44$ của một CSC. Tìm 3 số đó. 
	\dapso{$7; 35; 175$}
	\loigiai{
		Gọi $u_1, u_2, u_3$ là $3$ số hạng liên tiếp của CSN với công bội $q$.\\
		Theo đề $u_1=a_2, u_2=a_9, u_3=a_{44}$ với $a_2,a_9, a_{44}$ là 3 số hạng của CSC với công sai $d$.\\
		Ta có $\heva{&a_9=a_2+7d\\&a_{44}=a_2+42d}\Rightarrow 6a_9-a_{44}=5a_2\Leftrightarrow 6u_2-u_3=5u_1\Leftrightarrow 6u_1q-u_1q^2=5u_1$\\
		\hspace*{4.2cm}$\Leftrightarrow q^2-6q+5=0 \Leftrightarrow q=1 (L) \vee q=5 (N)$\\
		Theo đề ta cũng có $u_1+u_2+u_3=217\Leftrightarrow u_1+u_1q+u_1q^2=217\Rightarrow u_1=7$\\
		Vậy $3$ số cần tìm là $7; 35; 175$.
	}	
\end{bt}

\begin{bt}%[1D2V3-7]
% nguon: Archive/SP_KhoLatex/ID6_old/1K2-3.tex #3189
%[DCHT Toán 11 - KNTT -Đỗ Chí Tâm] 
	$a,b,c$ là ba số hạng liên tiếp của một CSN và $a,b+2,c+9$ là ba số hạng liên tiếp của một CSC, đồng thời $a,b+2,c$ là ba số hạng liên tiếp của một CSN khác. Tìm $a$.
	\dapso{$m=\dfrac{-7\pm3\sqrt{5}}{2}$}
	\loigiai{
		Vì $a,b,c$ là ba số hạng liên tiếp của CSN, ta có $ac=b^2\,\,\,\, (1)$\\
		Vì $a,b+2,c+9$ là ba số hạng liên tiếp của CSC, ta có $a+(c+9)=2(b+2)\,\,\,\, (2)$\\
		Vì $a,b+2,c$ là ba số hạng liên tiếp của CSN, ta có $a.c=(b+2)^2\,\,\,\, (3)$\\
		Thế $(1)$ vào $(3)$, ta được $b^2=(b+2)^2\Leftrightarrow b=-1$\\
		Thay $b=-1$ vào $(1), (2)$, ta được $\heva{&ac=1\\&a+c=-7}\Leftrightarrow a=\dfrac{-7+3\sqrt{5}}{2}\vee a=\dfrac{-7-3\sqrt{5}}{2}$
	}	
\end{bt}

\begin{bt}%[1D2C3-7]
% nguon: Archive/SP_KhoLatex/ID6_old/1K2-3.tex #3190
%[DCHT Toán 11 - KNTT -Đỗ Chí Tâm] 
	Một CSC và CSN có cùng các số hạng thứ $m+1$, thứ $n+1$, thứ $p+1$ và $3$ số hạng này là $3$ số dương $a, b, c$. Tính $T=a^{b-c}.b^{c-a}.c^{a-b}$.
	\dapso{$m=1$}
	\loigiai{
		$a=u_1+md=q^m.v_1$\\
		$b=u_1+nd=q^n.v_1$\\
		$c=u_1+pd=q^p.v_1$\\
		Suy ra $T=a^{b-c}.b^{c-a}.c^{a-b}=\left(q^mv_1 \right)^{(n-p)d}. \left(q^nv_1 \right)^{(p-m)d}. \left(q^pv_1 \right)^{(m-n)d}=1$
	}	
\end{bt}

\begin{bt}%[1D2C3-5]
% nguon: Archive/SP_KhoLatex/ID6_old/1K2-3.tex #3191
%[DCHT Toán 11 - KNTT -Đỗ Chí Tâm] 
	Tìm $m$ dương để phương trình $x^3+(5-m)x^2+(6-5m)x-6m=0 \,\,\,(*)$ có $3$ nghiệm phân biệt lập thành cấp số nhân.
	\dapso{$m=\sqrt{6}$}
	\loigiai{
		$(*)\Leftrightarrow (x+2) \left(x^2+(3-m)x-3m \right)=0\Leftrightarrow x=-2 \vee x=-3 \vee x=m$.\\
		Để $(*)$ có 3 nghiệm phân biệt thì $m\ne -3$ và $m\ne -2$.\\
		Do $3$ nghiệm này lập thành cấp số nhân, ta sắp xếp các nghiệm này theo thứ tự tăng dần được các dãy số sau
		\begin{itemize}
			\item $-3;-2; m$ lập thành CSN $\Leftrightarrow$ $-3m=(-2)^2\Leftrightarrow m=-\dfrac{4}{3}$
			\item $-3; m; -2$ lập thành CSN $\Leftrightarrow -3(-2)=m^2\Leftrightarrow m=\pm \sqrt{6}$
			\item $m; -3; -2$ lập thành CSN $\Leftrightarrow m(-2)=(-3)^2\Leftrightarrow m=-\dfrac{9}{2}$
		\end{itemize}	
		So với điều kiện thì $m=\sqrt{6}$ thỏa yêu cầu bài toán.
	}	
\end{bt}

\begin{bt}%[1D2C2-5]
% nguon: Archive/SP_KhoLatex/ID6_old/1K2-3.tex #3192
%[DCHT Toán 11 - KNTT -Đỗ Chí Tâm] 
	Tìm tham số $m$ để phương trình $x^3-(2m+1)x^2+2mx=0 \,\, (*)$ có $3$ nghiệm phân biệt lập thành một cấp số cộng, biết $m<0$.
	\dapso{$m=-\dfrac{1}{2}$}
	\loigiai{
		$(*)\Leftrightarrow x.\left(x^2-(2m+1)x+2m \right)=0\Leftrightarrow x=0 \vee x=1 \vee x=2m$.\\
		Để $(*)$ có 3 nghiệm phân biệt thì $m\ne 0$ và $m\ne \dfrac{1}{2}$.\\
		Do $3$ nghiệm này lập thành cấp số cộng, ta sắp xếp các nghiệm này theo thứ tự tăng dần được các dãy số sau
		\begin{itemize}
			\item $2m;0; 1$ lập thành CSC $\Leftrightarrow$ $2m+1=2.0\Leftrightarrow m=-\dfrac{1}{2}$
			\item $0; 2m; 1$ lập thành CSC $\Leftrightarrow 0+1=4m\Leftrightarrow m=\dfrac{1}{4}$
			\item $0; 1; 2m$ lập thành CSC $\Leftrightarrow 0+2m=2.1\Leftrightarrow m=1$
		\end{itemize}	
		Vậy $m=-\dfrac{1}{2}$ là giá trị cần tìm.
	}	
\end{bt}

\begin{bt}%[1D2H3-8]
% nguon: Archive/SP_KhoLatex/ID6_old/1K2-3.tex #3200
[TH] %[DCHT Toán 11 - KNTT - Dung Phuong] 
	Vào năm 2020, dân số của một quốc gia là khoảng $97$ triệu người và tốc độ tăng trưởng dân số là $0{,}91 \%$. Nếu tốc độ tăng trưởng dân số này được giữ nguyên hằng năm, hãy ước tính dân số của quốc gia đó vào năm 2030.
	\dapso{$106{,}1973784$}
	\loigiai{
		Dân số năm 2021 tăng lên so với năm 2020 là $97 \cdot 0{,}91 \% $ triệu người.\\
		Dân số năm 2021 là 
		\begin{center}
			$97 + 97 \cdot 0{,}91 \% = 97\cdot (1+0{,}91 \%)$ triệu người.
		\end{center}
		Dân số năm 2022 tăng lên so với năm 2021 là $97\cdot (1+0{,}91 \%)\cdot 0{,}91 \% $ triệu người.\\
		Dân số năm 2022 là 
		\begin{center}
			$97\cdot (1+0{,}91 \%) + 97\cdot (1+0{,}91 \%) \cdot 0{,}91 \% = 97\cdot (1+0{,}91 \%)^2$ triệu người.
		\end{center}
		Dân số năm 2023 tăng lên so với năm 2021 là $97\cdot (1+0{,}91 \%)^2\cdot 0{,}91 \% $ triệu người.\\
		Dân số năm 2023 là
		\begin{center}
			$97\cdot (1+0{,}91 \%)^2 + 97\cdot (1+0{,}91 \%)^2\cdot 0{,}91 \% = 97\cdot (1+0{,}91 \%)^3$ triệu người.
		\end{center}
		Tương tự vậy ta có dân số năm 2030 là $97\cdot (1+0{,}91 \%)^{10} = 106{,}1973784$ triệu người.
	}
\end{bt}

\begin{bt}%[1D2H3-8]
% nguon: Archive/SP_KhoLatex/ID6_old/1K2-3.tex #3202
[TH] %[DCHT Toán 11 - KNTT - Dung Phuong] 
	Giả sử một thành phố có dân số năm 2022 là khoảng $2{,}1$ triệu người và tốc độ gia tăng dân số trung bình mỗi năm là $0{,}75 \%$.
	\begin{listEX}[1]
		\item [a)]  Dự đoán dân số của thành phố đó vào năm $2032$; \dapso{$\approx2262924$ (người)}
		\item [b)]  Nếu tốc độ gia tăng dân số vẫn giữ nguyên như trên thì uớc tính vào năm nào dân số của thành phố đó sẽ tăng gấp đôi so với năm $2022$?  \dapso{$2116$}	\end{listEX}
	
	\loigiai{
		\begin{listEX}
			\item [a)] 	Giả sử dân số năm $2022$ là $u_1=2{,}1\cdot 10^6$ thì dân số năm 
			$2023$ 
			là 
			$u_2=u_1+ 0{,}0075u_1=1{,}0075u_1$.\\
			Tương tự dân số năm $2024$ là $u_3=1{,}0075u_2$.\\
			Do đó dân số của thành phố qua các năm lập thành một cấp số nhân với 
			$u_1=2{,}1\cdot10^6$; $q=1{,}0075$.\\
			Vậy dân số năm $2032$ tương ứng với $u_{11}=u_1\cdot q^{10}=2,1\cdot 
			10^6\cdot1{,}0075^{10}\approx2262924$ (người).
			\item [b)] Giả sử đến năm thứ $n$ thì dân số gấp đôi năm $2022$. \\
			Suy ra 
			$u_n=2u_1 \Leftrightarrow q^{n-1}=2\Leftrightarrow  1{,}0075^{n-1}=2 
			\Leftrightarrow n \approx 93{,}7.$\\
			Vậy $94$ năm sau tức là năm $2116$ thì dân số thành phố sẽ gấp đôi năm $2022$.
	\end{listEX}}
\end{bt}

\begin{bt}%[1D2H3-8]
% nguon: Archive/SP_KhoLatex/ID6_old/1K2-3.tex #3203 
% KIEM_TRA: cần xem lại đáp số/dạng
[TH] %[DCHT Toán 11 - KNTT - Dung Phuong] 
	Giả sử anh Tuấn kí hợp đồng lao động trong $10$ năm với điều khoản về tiền lương như sau: Năm thứ nhất, tiền lương của anh Tuấn là $60$ triệu. Kể từ năm thứ hai trở đi, mỗi năm tiền lương của anh Tuấn được tăng lên $8 \%$. Tính tổng số tiền lương anh Tuấn lĩnh được trong $10$ năm đi làm (đơn vị: triệu đồng, làm tròn đến hàng phần nghìn).
	\dapso{$\approx 869{,}194$ triệu người}
	\loigiai{
		Gọi $u_n$ là số tiền lương (triệu đồng) anh Tuấn được lĩnh ở năm làm việc thứ $n$. Ta có: $u_1=60$;
		\[u_n=u_{n-1}+u_{n-1} \cdot 0{,}08=u_{n-1} \cdot(1+0{,}08)=u_{n-1} \cdot 1{,}08. \]
		Do đó, $\left(u_n\right)$ là cấp số nhân có số hạng đầu $u_1=60$, công bội $q=1{,}08$. Áp dụng công thức tính tổng $S_n$, ta có tổng số tiền lương anh Tuấn lĩnh được trong $10$ năm đi làm là
		\[S_{10}=\dfrac{60\cdot\left(1-1{,}08^{10}\right)}{1-1{,}08} \approx 869{,}194\ (\text{triệu người}). \]
	}
\end{bt}

\begin{bt}%[1D2H3-8]
% nguon: Archive/SP_KhoLatex/ID6_old/1K2-3.tex #3204
[TH] %[DCHT Toán 11 - KNTT - Dung Phuong] 
	Một công ty xây dựng mua một chiếc máy ủi với giá $3$ tỉ đồng. Cứ sau mỗi năm sử dụng, giá trị của chiếc máy ủi này lại giảm $20 \%$ so với giá trị của nó trong năm liền trước đó. Tìm giá trị còn lại của chiếc máy ủi đó sau $5$ năm sử dụng.
	\dapso{$480$ triệu đồng}
	\loigiai{
		Gọi $(u_n)$ là dãy số biểu diễn giá trị của chiếc máy ủi theo từng năm. \\
		Dãy số này là một cấp số nhân có $u_1=3$, $q=0{,}2$.\\
		Số hạng tổng quát của cấp số nhân này là $u_n=3\cdot 0{,}2^{n-1}$.\\
		Ta có $u_5=3\cdot 0{,}2^4=\dfrac{3}{325}=4{,}8\cdot10^{-3}$.\\
		Tương ứng giá trị của chiếc máy ủi sau $5$ năm là $480$ triệu đồng.
	}
\end{bt}

\begin{bt}%[1D2H3-8]
% nguon: Archive/SP_KhoLatex/ID6_old/1K2-3.tex #3205
[TH] %[DCHT Toán 11 - KNTT - Dung Phuong] 
	Một gia đình mua một chiếc ô tô giá $800$ triệu đồng. Trung bình sau mỗi năm sử dụng, giá trị còn lại của ô tô giảm đi $4 \%$ (so với năm trước đó).
	\begin{enumEX}[a)]{1} 
		\item Viết công thức tính giá trị của ô tô sau $1$ năm, $2$ năm sử dụng.
		\item Viết công thức tính giá trị của ô tô sau $n$ năm sử dụng.
		\item Sau $10$ năm, giá trị của ô tô ước tính còn bao nhiêu triệu đồng? 
	\end{enumEX} 
	\dapso{$\approx 531{,}87$ triệu đồng}
	\loigiai{
		Gọi $u_n$ là giá trị còn lại của ô tô sau $n$ năm sử dụng. 
		\begin{enumEX}[a)]{1} 
			\item Giá trị của ô tô sau $1$ năm sử dụng là $u_1=800-800\cdot0{,}04=800\cdot0{,}96=768$ triệu đồng.\\
			Giá trị của ô tô sau $2$ năm sử dụng là $u_2=u_1-u_1\cdot0{,}04=u_1\cdot0{,}96=737{,}28$ triệu đồng.
			\item Ta có $u_n=u_{n-1}-u_{n-1}\cdot0{,}04=u_{n-1}\cdot0{,}96$. \\
			Do đó, $(u_n)$ là cấp số nhân với số hạng đầu $u_1=768$ và công bội $q=0{,}96$. \\
			Vậy sau $n$ năm sử dụng, giá trị còn lại của chiếc ô tô là $u_n=u_1q^{n-1}\Rightarrow u_n=768\cdot0{,}96^{n-1}$.
			\item Sau $10$ năm, ước tính giá trị của ô tô còn lại là $u_{10}=768\cdot0{,}96^9\approx 531{,}87$ triệu đồng.
		\end{enumEX} 
	}
\end{bt}

\begin{bt}%[1D2V3-8]
% nguon: Archive/SP_KhoLatex/ID6_old/1K2-3.tex #3206
 [VD] %[DCHT Toán 11 - KNTT - Dung Phuong] 
	Ông An vay ngân hàng $1$ tỉ đồng với lãi suất $12\%/$năm. Ông đã trả nợ theo cách: Bắt đầu từ tháng thứ nhất sau khi vay, cuối tháng ông trả ngân hàng số tiền là $a$ (đồng) và đã trả hết nợ sau đúng $2$ năm kể từ ngày vay. Hỏi số tiền mỗi tháng mà ông An phải trả là bao nhiêu đồng (làm tròn kết quả đến hàng nghìn)?
	\dapso{$47073500$}
	\loigiai{
		Do lãi suất là $12\%$/năm tương đương với lãi là $1\%$/tháng.\\
		Sau $1$ tháng, ông An còn nợ là: $10^9.(1+1\%)-a=10^9.(1,01)-S_1$.\\
		Sau $2$ tháng, ông An còn nợ là: $10^9.(1.01)^2-a.(1.01)-a=10^9(1,01)^2-S_2$.\\
		Sau $3$ tháng, ông An còn nợ là: $10^9.(1.01)^3-a(1.01)^2-a(1.01)-a=10^9.(1.01)^3-S_3$.\\
		Sau $24$ tháng, ông An còn nợ là: $10^9.(1.01)^{24}-S_{24}=0$.\\
		Do đó $S_{24}=10^9.(1.01)^{24}  \Leftrightarrow a.\dfrac{1-(1.01)^{24}}{1-(1.01)}=10^9.(1.01)^{24} \Leftrightarrow a =\dfrac{10^9.(1.01)^{24}.0.01}{(1.01)^{24}-1}\approx 47073472,22$.\\
		Vậy mỗi tháng ông An phải trả $47073500$.
	}
\end{bt}

\begin{bt}%[1D2H3-8]
% nguon: Archive/SP_KhoLatex/ID6_old/1K2-3.tex #3208
 [TH] %[DCHT Toán 11 - KNTT - Dung Phuong] 
	Một cái tháp có $11$ tầng. Diện tích của mặt sàn tầng $2$ bằng nửa diện tích của mặt đáy tháp và diện tích của mặt sàn mỗi tầng bằng nửa diện tích của mặt sàn mỗi tầng ngay bên dưới. Biết mặt đáy tháp có diện tích là $12 288m^2$. Tính diện tích của mặt sàn tầng trên cùng của tháp theo đơn vị mét vuông.
	\dapso{$12m^2$}
	\loigiai{
		Do diện tích của mặt sàn tính từ tầng một lập thành một cấp số nhân với $u_2=\dfrac{1}{2}.12288=6144$ và $q=\dfrac{1}{2}$.\\
		Ta có $\heva{u_2&=6144 \\ q&=\dfrac{1}{2}}  \Leftrightarrow \heva{u_1&=12288 \\ q&=\dfrac{1}{2}}$.\\
		Ta có $u_{11}=u_1.q^{10}=12288.\dfrac{1}{12^{10}}=12m^2$.
		Vậy diện tích của mặt sàn tầng trên cùng là	$12m^2$.
	}
\end{bt}

\begin{bt}%[1D2H3-8]
% nguon: Archive/SP_KhoLatex/ID6_old/1K2-3.tex #3210
[TH] %[DCHT Toán 11 - KNTT - Dung Phuong] 
	Một cây đàn organ có tần số âm thanh các phím liên tiếp tạo thành một cấp số nhân. Cho biết tần số phím La trung là $400$ Hz và tần số của phím La cao cao hơn $12$ phím là $800$ Hz (nguồn: https://vi.wikipedia.org/wikiOrgan). Tìm công bội của cấp số nhân nói trên (làm tròn kết quả đến hàng phần nghìn).
	\dapso{$q = \pm 1{,}414$}
	\loigiai{
		Theo đề ta có $\heva{&u_1=400\\&u_{13}=800} \Leftrightarrow \heva{&u_1=400\\&u_1q^{12}=800} \Rightarrow q^{12} = 2 \Rightarrow q = \pm 1{,}414$.
	}
\end{bt}

\begin{bt}%[1D2V3-8]
% nguon: Archive/SP_KhoLatex/ID6_old/1K2-3.tex #3211
[VD]%[DCHT Toán 11 - KNTT - Dung Phuong]  
	Một loại thuốc được dùng mỗi ngày một lần. Lúc đầu nồng độ thuốc trong máu của bệnh nhân tăng nhanh, nhưng mỗi liều kế tiếp có tác dụng ít hơn liều trước đó. Lượng thuốc trong máu ở ngày thứ nhất là $50 \,\mathrm{mg}$, và mỗi ngày sau đó giảm chỉ còn một nửa so với ngày kề trước đó. Tính tổng lượng thuốc (tính bằng $\mathrm{mg}$) trong máu của bệnh nhân sau khi dùng thuốc $10$ ngày liên tiếp.
	\dapso{$99{,}902$ mg.}
	\loigiai{
		Gọi $(u_n)$ là dãy số biểu diễn giá trị của lượng thuốc trong máu của bệnh nhân theo từng ngày. \\
		Dãy số này là một cấp số nhân có $u_1=50$, $q=\dfrac{1}{2}$.\\
		Tổng của $n$ số hạng đầu tiên của cấp số nhân là $S_n=u_1\dfrac{1-q^n}{1-q}$.\\
		Theo bài toán, ta có $S_{10}=50 \cdot\dfrac{1-\left(\dfrac{1}{2}\right)^{10}}{1-\dfrac{1}{2}} \approx 99{,}902$.\\
		Vậy tổng lượng thuốc trong máu của bệnh nhân sau khi dùng thuốc $10$ ngày liên tiếp là $99{,}902$ mg.	
	}
\end{bt}

\begin{bt}%[1D2V3-7]
% nguon: Archive/SP_KhoLatex/ID6_old/1-GK1-CTST-De11-NH23-24.tex #3220
%Bài 2
\begin{listEX}
\item Cho cấp số cộng $\left(u_n\right)$ có $u_2+u_5=42$ và $u_4+u_9=66$. Tìm số hạng đầu $u_1$ và công sai $d$.
\item Cho cấp số nhân $\left(v_n\right)$ có $v_1+v_2+v_3=7$ và $v_5-v_2=14$. Tìm số hạng đầu $u_1$ và công bội $q$.
\item Ông Sơn trồng cây trên một mãnh đất hình tam giác theo quy luật: ở hàng thứ nhất có $1$ cây, ở hàng thứ hai có $2$ cây, ở hàng thứ ba có $3$ cây,..., ở hàng thứ $n$ có $n$ cây. Biết rằng ông đã trồng hết $11325$ cây. Hỏi số hàng cây được trồng theo cách trên lă bao nhiêu?
\end{listEX}
\loigiai{
\begin{listEX}
\item Áp dụng công thức số hạng tổng quát của cấp số cộng ta có
$$\heva{&u_{2}+u_{5}=42\\&u_{4}+u_{9}=66}\Leftrightarrow\heva{&2u_{1}+5d=42\\&2 u_{2}+11d=66}\Leftrightarrow\heva{&u_1=11\\&d=4.}$$
Vậy số hạng đầu là $u_1=11$ và công sai là $d=4$.
\item Áp dụng công thức số hạng tổng quát của cấp số nhân ta có
$$\heva{&v_{1}+v_{2}+v_{3}=7\\&v_{5}-v_{2}=14}\Leftrightarrow\heva{&v_{1}+v_{1}q+v_{1}q^{2}=7\\&v_{1}q^{4}-v_{1}q=14}\Leftrightarrow\heva{&v_1\left(1+q+q^2\right)=7\\&v_1q(q-1)\left(q^2+q+1\right)=14.}$$
Lấy $(2)$ chia $(1)$ ta được
$$q(q-1)=2 \Leftrightarrow q^2-q=2 \Leftrightarrow q^2-q-2=0 \Leftrightarrow\hoac{&q=2\\&q=-1.}$$
Với $q=2$, thay vào (1) ta được $v_1=1$.\\
Với $q=-1$, thay vào (1) ta được $v_1=7$.\\
Vậy số hạng đầu $v_1=1$ và công bội $q=2$ hoặc số hạng đầu $v_1=7$ và công bội $q=-1$.
\item Đặt $u_1=1$, $u_2=2$, $u_3=3$, $\ldots$, $u_n=n$. Dãy số $\left(u_n\right)$ là một cấp số cộng với công sai $d=1$.\\
Ta có $S_n=11325$.\\
$S_n=n u_1+\dfrac{n(n-1)}{2} \cdot d$ \\
$\Rightarrow 11325=n \cdot 1+\dfrac{n(n-1)}{2} \cdot 1 \Rightarrow n^2+n-22650=0 \Rightarrow\hoac{&n=150\\&n=-151.}$\\
Vậy có 150 hàng cây.
\end{listEX}
}
\end{bt}

\begin{bt}%[1D2N1-3]
% nguon: Archive/SP_KhoLatex/ID6_old/1K2-1.tex #3228
[NB]%[DCHT Toán 11 - KNTT -Tên GV]
	Xét dãy số hữu hạn gồm các số tự nhiên lẻ nhỏ hơn 20, sắp xếp theo thứ tự từ bé đến lớn. Liệt kê tất cả các số hạng của dãy số này, tìm số hạng đầu và số hạng cuối của dãy. 
	\dapso{$u_1=1$, $u_{11}=19$}
	\loigiai{Các số hạng của dãy là $1,3,5,7,9,10,11,13,15,17,19$.\\
		Số hạng đầu của dãy là $u_1=1$.\\
		Số hạng cuối của dãy là $u_{11}=19$.}
\end{bt}

\begin{bt}%[1D2N1-2]
% nguon: Archive/SP_KhoLatex/ID6_old/1K2-1.tex #3229
[TH]%[DCHT Toán 11 - KNTT -Tên GV]
	Xét dãy số gồm tất cả các số tự nhiên chia cho $5$ dư $1$. Xác định số hạng tổng quát của dãy số.
	\dapso{$u_n=5n+1$}
	\loigiai{Các số tự nhiên chia cho $5$ dư $1$ gồm các số sau:
		$6,11,16,21, \ldots $\\
		Số hạng tổng quát $u_n=5n+1$.}
\end{bt}

\begin{bt}%[1D2N1-2]
% nguon: Archive/SP_KhoLatex/ID6_old/1K2-1.tex #3230
[NB]%[DCHT Toán 11 - KNTT -Tên GV] 
	Tìm năm số hạng đầu và số hạng thứ $100$ của dãy $(u_n)$ có số hạng tổng quát $u_n= \dfrac{(-1)^n}{n}$.
	\dapso{$u_1=-1, \dfrac{1}{2},-\dfrac{1}{3}, \dfrac{1}{4}, -\dfrac{1}{5} $, $u_{100}=\dfrac{1}{100}$}
	\loigiai{
		Năm số hạng đầu của dãy là $u_1=-1, \dfrac{1}{2},-\dfrac{1}{3}, \dfrac{1}{4}, -\dfrac{1}{5} $.\\
		Số hạng thứ $100$ là $u_{100}=\dfrac{1}{100}$.}
\end{bt}

\begin{bt}%[1D2N1-3]
% nguon: Archive/SP_KhoLatex/ID6_old/1K2-1.tex #3231
[NB]%[DCHT Toán 11 - KNTT -Tên GV]
	Viết năm số hạng đầu của dãy số gồm các số nguyên tố theo thứ tự tăng dần.
	\dapso{$2,3,5,7,11$}
	\loigiai {Năm số hạng đầu của dãy số trên là $2,3,5,7,11$.} 
\end{bt}

\begin{bt}%[1D2N1-2]
% nguon: Archive/SP_KhoLatex/ID6_old/1K2-1.tex #3232
[NB]%[DCHT Toán 11 - KNTT-Tên GV]
	Viết năm số hạng đầu của dãy $(u_n)$ với số hạng tổng quát là $u_n=n!$.
	\dapso{$1,2,6,24,120$}
	\loigiai{Năm số hạng đầu của dãy trên là $1,2,6,24,120$.}
\end{bt}

\begin{bt}%[1D2N1-3]
% nguon: Archive/SP_KhoLatex/ID6_old/1K2-1.tex #3236
[NB]
	Cho dãy số $u_n=\dfrac{1}{\sqrt{n}+1}$. Tìm số hạng $u_4$.	
	\dapso{$u_4=\dfrac{1}{3}$}
	\loigiai{ Ta có
		$u_4=\dfrac{1}{\sqrt{4}}+1=\dfrac{1}{3}.$		
	}
\end{bt}

\begin{bt}%[1D2N1-2]
% nguon: Archive/SP_KhoLatex/ID6_old/1K2-1.tex #3237
[NB]
	Cho dãy số $(u_n)$ có số hạng tổng quát: $u_n=2 n+\sqrt{n^2+4}$. Tìm số hạng thứ $6$ của dãy số.
	\dapso{$u_6=12+2\sqrt{10}$}
	\loigiai{
		Ta có $u_6=12+2 \sqrt{10}$.
	}
\end{bt}

\begin{bt}%[1D2N1-3]
% nguon: Archive/SP_KhoLatex/ID6_old/1K2-1.tex #3238
[NB]
	Cho dãy số $(u_n)$ xác định bởi: $\heva{&u_1=-1 ; u_2=3 \\&u_{n+1}=5 u_n-6 u_{n-1} \forall n \geq 2}.$ Tìm số hạng thứ $7$ của dãy.
	\dapso{$3261$}
	\loigiai{
		Ta có
		$$
		u_3=5 u_2-6 u_1=21 ;~ u_4=5 u_3-6 u_2=87 ;~ u_5=309 ;~ u_6=1023 ;~ u_7=3261
		$$
		Vậy số hạng thứ $7$ của dãy là $3261$.
	}
\end{bt}

\begin{bt}%[1D2N1-3]
% nguon: Archive/SP_KhoLatex/ID6_old/1K2-1.tex #3239
[NB]
	Viết năm số hạng đầu của dãy số Fibonacci $\left(F_n\right)$ cho bởi hệ thức truy hồi
	$$
	\heva{
		&F_1=1, F_2=1 \\
		&F_n=F_{n-1}+F_{n-2}~(n \geq 3) .
	}
	$$
	\dapso{$F_3=2,~F_4=3,~F_5=5$}
	\loigiai{
		Ta có $F_3=2,~F_4=3,~F_5=5.$
	}
\end{bt}

\begin{bt}%[1D2H1-6]
% nguon: Archive/SP_KhoLatex/ID6_old/1K2-1.tex #3240
[NB]
	Người ta nuôi cấy $5$ con vi khuẩn ecoli trong môi trường nhân tạo. Cứ $30$ phút thì vi khuẩn ecoli sẽ nhân đôi 1 lần. Tính số lượng vi khuẩn thu được sau $1,2,3$ lần nhân đôi.
	\dapso{$u_2=10, u_3=20, u_4=40$}
	\loigiai{
		Đặt $u_1=5$, gọi số vi khuẩn sau $n$ lần phân chia là $u_{n+1}$, khi đó ta có dãy số $(u_n)$ thỏa mãn $$u_1=5, \; u_{n+1}=2u_n$$
		Ta có $u_2=10, u_3=20, u_4=40$.
	}	
\end{bt}

\begin{bt}%[1D2H1-3]
% nguon: Archive/SP_KhoLatex/ID6_old/1K2-1.tex #3241
[TH]
	Cho dãy số $(u_n)$ được xác định bởi $u_n=\dfrac{n^2+3n+7}{n+1}$.
	\begin{listEX}
		\item Viết năm số hạng đầu của dãy.
		\item Dãy số có bao nhiêu số hạng nhận giá trị nguyên.
	\end{listEX}
	\dapso{$u_1=\dfrac{11}{2}$; $u_2=\dfrac{17}{3}$; $u_3=\dfrac{25}{4}$; $u_4=7$; $u_5=\dfrac{47}{6}$. $u_4=7 $}
	\loigiai{		
		\begin{listEX}
			\item Ta có năm số hạng đầu của dãy
			$u_1=\dfrac{1^2+3.1+7}{1+1}=\dfrac{11}{2}$; $u_2=\dfrac{17}{3}$; $u_3=\dfrac{25}{4}$; $u_4=7$; $u_5=\dfrac{47}{6}$.
			\item Ta có: $u_n=n+2+\dfrac{5}{n+1}$, do đó $u_n$ nguyên khi và chỉ khi $ \dfrac{5}{n+1}$ nguyên hay $ n+1 $ là ước của 5. Điều đó xảy ra khi $ n+1=5\Leftrightarrow n=4 $. Vậy dãy số có duy nhất một số hạng nguyên là $u_4=7 $.
		\end{listEX}		
	}
\end{bt}

\begin{bt}%[1D2V1-3]
% nguon: Archive/SP_KhoLatex/ID6_old/1K2-1.tex #3242
[VD]
	Cho dãy số $\left(x_n\right)$ thỏa mãn điều kiện $x_1=1, x_{n+1}-x_n=\dfrac{1}{n(n+1)}, n=1,2,3, \ldots$. Số hạng $x_{2023}$ bằng
	\dapso{$x_{2023}=\dfrac{4045}{2023}$}
	\loigiai{
		Ta có
		$$
		\begin{aligned}
			x_{n+1}-x_n=\dfrac{1}{n(n+1)}=\dfrac{1}{n}-\dfrac{1}{n+1} & \Leftrightarrow \sum_{k=1}^{n-1}\left(x_{k+1}-x_k\right)=\sum_{k=1}^{n-1}\left(\dfrac{1}{k}-\dfrac{1}{k+1}\right) \\
			& \Leftrightarrow x_n-x_1=1-\dfrac{1}{n} \\
			& \Leftrightarrow x_n=\dfrac{2n-1}{n} .
		\end{aligned}
		$$
	}
\end{bt}

\begin{bt}%[1D2C1-3]
% nguon: Archive/SP_KhoLatex/ID6_old/1K2-1.tex #3243
[VDC]
	Cho dãy số $\left(u_n\right)$ biết $\heva{&u_1=99 \\&u_{n+1}=u_n-2 n-1, n \geq 1}$. Hỏi số $-861$ là số hạng thứ mấy?
	\dapso{$-861$ là số hạng thứ $31$}
	\loigiai{
		Ta có
		$$
		\begin{aligned}
			&u_n &=& &u_{n-1}-2 n+1 \\
			&u_{n-1} & = & &u_{n-2}-2 n+3 \\
			&\vdots &\vdots&  &\vdots \\
			&u_3 & = & &u_2-2 n+2 n-5 \\
			&u_2 & = & &u_1-2 n+2 n-3
		\end{aligned}
		$$
		Suy ra
		$$
		\begin{aligned}
			& u_n=u_1-2 n \cdot(n-1)+1+3+5+\cdots+(2 n-5)+(2 n-3) \\
			& u_n=99-2 n^2+2 n+\dfrac{n-1}{2}\cdot[2 \cdot 1+(n-2) \cdot 2]=100-n^2
		\end{aligned}
		$$
		Giả sử $u_n=-861 \Rightarrow n^2=961 \Rightarrow n=31$ (vì $n \in \mathbb{N}$).
		Vậy số $-861$ là số hạng thứ $31$ .}
\end{bt}

\begin{bt}%[1D2N1-4]
% nguon: Archive/SP_KhoLatex/ID6_old/1K2-1.tex #3249
[NB]
	Xét tính tăng giảm của dãy số $(u_n)$ với $u_n=\dfrac{\sqrt{2}}{3^n}$.
	\dapso{dãy số giảm}
	\loigiai{
		Ta có $u_n>0, \forall n \in \mathbb{N}^\ast$.\\
		Xét thương $$\dfrac{u_{n+1}}{u_n}=\dfrac{\sqrt{2}}{3^{n+1}}: \dfrac{\sqrt{2}}{\sqrt{3^2}}=\dfrac{3^n}{3^{n+1}}=\dfrac{1}{3}<1.$$
		Vậy $\left(u_n\right)$ là dãy số giảm.
	}
\end{bt}

\begin{bt}%[1D2N1-4]
% nguon: Archive/SP_KhoLatex/ID6_old/1K2-1.tex #3250
[NB]
	Xét tính tăng giảm của dãy số $\left(u_n\right)$ với $u_n=\dfrac{1}{n(n+1)}$.
	\dapso{dãy số tăng}
	\loigiai{
		Ta có $u_n=\dfrac{1}{n(n+1)}=\dfrac{1}{n}-\dfrac{1}{n+1}$.
		Xét hiệu:
		$$
		\begin{aligned}
			u_{n+1}-u_n & =\left(\dfrac{1}{n}-\dfrac{1}{n+1}\right)-\left(\dfrac{1}{n+1}-\dfrac{1}{n+2}\right) \\
			& =\dfrac{1}{n}-\dfrac{1}{n+2}>0, \forall n \in \mathbb{N}^\ast
		\end{aligned}
		$$
		Vậy  $\left(u_n\right)$ là dãy số tăng.
	}
\end{bt}

\begin{bt}%[1D2H1-4]
% nguon: Archive/SP_KhoLatex/ID6_old/1K2-1.tex #3251
[TH]
	Xét tính tăng giảm của dãy số $\left(u_n\right)$ với $u_n=n+\cos ^2 n$.
	\dapso{dãy số tăng}
	\loigiai{
		Xét hiệu
		$$
		\begin{aligned}
			u_{n+1}-u_n & =\left(n+1+\cos ^2(n+1)\right)-\left(n+\cos ^2 n\right) \\
			& =1+\cos ^2(n+1)-\cos ^2 n \\
			& =\cos ^2(n+1)+\sin ^2 n>0, \forall n \in \mathbb{N}^\ast .
		\end{aligned}
		$$
		Vậy $\left(u_n\right)$ là dãy số tăng.
	}
\end{bt}

\begin{bt}%[1D2H1-4]
% nguon: Archive/SP_KhoLatex/ID6_old/1K2-1.tex #3252
{TH}
	Xét tính tăng giảm của dãy số $(u_n)$ với $u_n=\dfrac{1}{n+1}+\dfrac{1}{n+2}+\ldots+\dfrac{1}{2n}$.
	\dapso{dãy số giảm}
	\loigiai{
		Xét hiệu\\
		$\begin{aligned}
			u_{n+1}-u_n&=\left(\dfrac{1}{n+2}+\dfrac{1}{n+3}+\ldots+\dfrac{1}{2(n+1)}\right)-\left(\dfrac{1}{n+1}+\dfrac{1}{n+2}+\ldots+\dfrac{1}{2n}\right)\\
			&=\dfrac{1}{n+2}-\dfrac{1}{2n+1}-\dfrac{1}{2n+2}\\
			&=\dfrac{1}{2n+2}-\dfrac{1}{2n+1}<0, \forall n\in \mathbb{N}^\ast.
		\end{aligned}$\\
		Vậy $(u_n)$ là dãy số giảm.
	}
	
\end{bt}

\begin{bt}%[1D2V1-4]
% nguon: Archive/SP_KhoLatex/ID6_old/1K2-1.tex #3254
[VD]
	Xét tính tăng giảm của dãy số $\left(u_n\right)$ cho bởi
	$$
	\left(u_n\right)\colon\heva{
		&u_1=1 ; u_2=2 \\
		&u_{n+1}=\sqrt{u_n}+\sqrt{u_{n-1}} \forall n \geq 2
	}
	$$
	\dapso{dãy số tăng}
	\loigiai{
		Ta chứng minh dãy $\left(u_n\right)$ là dãy tăng bằng phương pháp quy nạp.\\
		Dễ thấy $u_1<u_2<u_3$.\\
		Giả sử $u_{k-1}<u_k ~\forall k \geq 2$, ta chứng minh $u_{k+1}>u_k$.\\
		Thật vậy ta có $u_{k+1}=\sqrt{u_k}+\sqrt{u_{k-1}}>\sqrt{u_{k-1}}+\sqrt{u_{n-2}}=u_k$.\\ Vậy $\left(u_n\right)$ là dãy tăng.
	}
	
\end{bt}

\begin{bt}%[1D2V1-4]
% nguon: Archive/SP_KhoLatex/ID6_old/1K2-1.tex #3255
[VD]
	Cho dãy số $\left(u_n\right)$ biết $u_n=\dfrac{b \cdot2 n^2+1}{n^2+3}$ và $b \in \mathbb{R}$. Hãy xác định $b$ để
	\begin{listEX}[2]
		\item $\left(u_n\right)$ là dãy số giảm.
		\item $\left(u_n\right)$ là dãy số tăng.
	\end{listEX}
	\dapso{$b<\dfrac{1}{6}$ dãy số giảm; $b>\dfrac{1}{6}$ dãy số tăng}
	\loigiai{
		Ta có
		$$
		u_n=2 b+\dfrac{1-6 b}{n^2+3}
		$$
		Xét hiệu $$u_{n+1}-u_n=\dfrac{1-6 b}{(n+1)^2+3}-\dfrac{1-6 b}{n^2+3}=(1-6 b) \cdot\left(\dfrac{1}{(n+1)^2+3}-\dfrac{1}{n^2+3}\right)=A_n.$$
		\begin{listEX}
			\item Để $\left(u_n\right)$ là dãy sỗ giảm thì $A_n<0, \forall n \in \mathbb{N}^\ast$.
			$$
			A_n<0 \Leftrightarrow 1-6 b>0 \Leftrightarrow b<\dfrac{1}{6}
			$$
			\item Để $\left(u_n\right)$ là dãy số tăng thì $A_n>0, \forall n \in \mathbb{N}^\ast$.
			$$
			A_n>0 \Leftrightarrow 1-6 b<0 \Leftrightarrow b>\dfrac{1}{6}.
			$$
		\end{listEX}
		
	}
\end{bt}

\begin{bt}%[1D2C1-4]
% nguon: Archive/SP_KhoLatex/ID6_old/1K2-1.tex #3256
[VDC]
	Xét tính tăng giảm của dãy số $\left(u_n\right)$ với $u_n=\sin n+\cos n$.
	\dapso{dãy khōng tăng, không giảm}
	\loigiai{
		Ta có: $u_n=\sin n+\cos n=\sqrt{2} \sin \left(n+\dfrac{\pi}{4}\right)$.
		Xét hiệu
		$$
		\begin{aligned}
			u_{n+1}-u_n&=\sqrt{2} \sin \left(n+1+\dfrac{\pi}{4}\right)-\sqrt{2} \sin \left(n+\dfrac{\pi}{4}\right) \\
			&=2 \sqrt{2} \cdot \cos \left(2 n+\dfrac{1}{2}+\dfrac{\pi}{4}\right) \cdot \sin \dfrac{1}{2}=A_n . \\
		\end{aligned}
		$$
		Với  $n=1, A_1>0$. Với  $n=100, A_{100}<100$ . \\
		Vậy $\left(u_n\right)$ là dãy khōng tăng, không giảm.
	}
\end{bt}

\begin{bt}%[1D2H1-5]
% nguon: Archive/SP_KhoLatex/ID6_old/1K2-1.tex #3263
[TH]%[Trương Đăng Khoa]
	Xét tính bị chặn của các dãy số sau 
		\begin{listEX}[3]
			\item $u_n=\dfrac{1}{2n^2-1}$.
			\item 
			$u_n=3\cdot\cos\dfrac{n x}{3}$. 
			\item  $u_n=2n^3+1$.
			\item  $u_n=\dfrac{n^2+2n}{n^2+n+1}$.
			\item  $u_n=n+\dfrac{1}{n}$.
		\end{listEX}
	\loigiai{
		\begin{enumerate}
			\item  $u_n=\dfrac{1}{2n^2-1}$.\\
			Ta có $2n^2-1\ge 1\Rightarrow u_n=\dfrac{1}{2n^2-1}\le 1$, $\forall n\ge 1$.\\
			Vậy dãy số bị chặn trên bởi $1$.\\
			\item $u_n=3\cdot\cos\dfrac{n x}{3}$ có $-1\le\cos\dfrac{n x}{3}\le 1\Rightarrow-3\le 3\cdot\cos\dfrac{n x}{3}\le 3$.\\
			Vậy dãy số bị chặn dưới bởi $-3$ và chặn trên bởi $3$.
			\item  $u_n=2n^3+1$ có $2n^3+1\ge 3$, $\forall n\ge 1$.\\
			Vậy dãy số bị chặn dưới bởi $3$.
			\item $u_n=\dfrac{n^2+2n}{n^2+n+1}$ có $u_n=\dfrac{n^2+2n}{n^2+n+1}=1+\dfrac{n-1}{n^2+n+1}\ge 1$, $\forall n\ge 1$.\\
			Vậy dãy số bị chặn dưới bởi $1$.
			\item  $u_n=n+\dfrac{1}{n}$ có $u_n=n+\dfrac{1}{n}\ge 2\sqrt{n\cdot\dfrac{1}{n}}=2$, $\forall n>0$.\\
			Vậy dãy số bị chặn bởi $2$.
		\end{enumerate}
	}
\end{bt}

\begin{bt}%[1D2V1-5]
% nguon: Archive/SP_KhoLatex/ID6_old/1K2-1.tex #3264
[VD]%[Trương Đăng Khoa]
	Xét tính bị chặn của dãy số $(u_n)$ với:
	\begin{listEX}[3]	 
		\item $u_{n}=\dfrac{4}{n}-5$.
		\item $u_{n}=\dfrac{n+4}{n+2}$.
		\item $u_{n}=\dfrac{5}{n^2+1}+\dfrac{n+2}{n+1}+\cos n$.
	\end{listEX}
	\loigiai{
		\begin{enumerate} 
			\item $u_{n}=\dfrac{4}{n}-5$.\\
			Ta có $u_n=\dfrac{4}{n}-5 \le \dfrac{4}{1}-5=-1$, $\forall n \in \mathbb{N}^{*}$ suy ra dãy $(u_n)$ bị chặn trên bởi $-1$.\\
			Mặt khác $u_n=\dfrac{4}{n}-5 \ge -5 \,\, \forall n \in \mathbb{N}^{*}$ suy ra dãy $(u_n)$ bị chặn dưới bởi $-5$.\\
			Vậy dãy $(u_n)$ bị chặn.	
			\item $u_{n}=\dfrac{n+4}{n+2}$.\\
			Ta có $u_n=	\dfrac{n+4}{n+2}=1+\dfrac{2}{n+2}> 1$, $\forall n \in \mathbb{N}^{*}$ suy ra dãy $(u_n)$ bị chặn dưới bởi $1$.\\
			Mặt khác $u_n=\dfrac{n+4}{n+2}=1+\dfrac{2}{n+2} \le 1+\dfrac{2}{1+2}=\dfrac{3}{5}$, $\forall n \in \mathbb{N}^{*}$ suy ra dãy $(u_n)$ bị chặn trên bởi $\dfrac{3}{5}$.\\
			Vậy dãy $(u_n)$ bị chặn.
			\item $u_{n}=\dfrac{5}{n^2+1}+\dfrac{n+2}{n+1}+\cos n$.\\
			Ta có $u_n=	\dfrac{5}{n^2+1}+\dfrac{n+2}{n+1}+\cos n=\dfrac{5}{n^2+1}+1+\dfrac{1}{n+1}+\cos n<5$, $\forall n \in \mathbb{N}^{*}$.\\
			Suy ra dãy $(u_n)$ bị chặn trên bởi $5$.\\
			Mặt khác $u_n=	\dfrac{5}{n^2+1}+\dfrac{n+2}{n+1}+\cos n=\dfrac{5}{n^2+1}+1+\dfrac{1}{n+1}+\cos n>0$, $\forall n \in \mathbb{N}^{*}$.\\
			Suy ra dãy $(u_n)$ bị chặn trên bởi $0$.\\
			Vậy dãy $(u_n)$ bị chặn.
		\end{enumerate}			
	}		
\end{bt}

\begin{bt}%[1D2C1-5]
% nguon: Archive/SP_KhoLatex/ID6_old/1K2-1.tex #3265
[VDC]%[Trương Đăng Khoa]
	Xét tính bị chặn của dãy số $u_n=\left(1+\dfrac{1}{n}\right)^n$, $n\in N^\ast$.
	\loigiai{
		Ta có $u_n=\left(1+\dfrac{1}{n}\right)^n>0$, $\forall n\in N^\ast$ nên $(u_n)$ bị chặn dưới $(1)$.\\
		Lại có $\begin{aligned}[t]
			u_n&\ =\left(1+\dfrac{1}{n}\right)^n=\displaystyle\sum\limits_{k=0}^n C_n^k\left(\dfrac{1}{n}\right)^k\\
			&\ =\displaystyle\sum\limits_{k=0}^n\left[\dfrac{n!}{k!\cdot(n-k)!\cdot n^k}\right]\\
			&\ =\displaystyle\sum\limits_{k=0}^n\left[\dfrac{1}{k!}\cdot\dfrac{(n-k+1)}{n}\cdot\dfrac{(n-k+2)}{n}\ldots\dfrac{(n-k+k)}{n}\right]\le\displaystyle\sum\limits_{k=0}^n\dfrac{1}{k!},\, n\in \mathbb{N}^\ast
		\end{aligned}$\\
		Mà $\begin{aligned}[t]
			\displaystyle\sum\limits_{k=0}^n\dfrac{1}{k!}&\ \le 1+1+\dfrac{1}{1\cdot2}+\dfrac{1}{2\cdot3}+\dfrac{1}{3\cdot4}+\ldots+\dfrac{1}{(n-1)\cdot n}\\
			&\ =2+\left(1-\dfrac{1}{2}\right)+\left(\dfrac{1}{2}-\dfrac{1}{3}\right)+\ldots+\left(\dfrac{1}{n-1}-\dfrac{1}{n}\right)\\
			&\ 
			=3-\dfrac{1}{n}<3,\, \forall n\in \mathbb{N}^\ast.
		\end{aligned}$\\
		Suy ra $u_n<3$, $\forall n\in \mathbb{N}^\ast$ nên dãy số $(u_n)$ bị chặn trên $(2)$.\\
		Từ $(1)$ và $(2)$  suy ra dãy số $(u_n)$ bị chặn.}
\end{bt}

\begin{bt}%[1D2V1-5]
% nguon: Archive/SP_KhoLatex/ID6_old/1K2-1.tex #3266
[VD]%[Trương Đăng Khoa]
	Cho dãy số $(u_n)$ xác định bởi $u_1=0$ và $u_{n+1}=\dfrac{1}{2}u_n+4$, $ \forall n\geq 1$.
	\begin{enumerate}
		\item Chứng minh dãy $(u_n)$ bị chặn trên bởi số $8$.
		\item Chứng minh dãy $(u_n)$ tăng, từ đó suy ra dãy $(u_n)$ bị chặn.
	\end{enumerate}
	\loigiai{
		\begin{enumerate}
			\item Ta chứng minh $u_n\leq 8$ với mọi $n\geq 1$.
			\begin{itemize}
				\item Khi $n=1$, ta có $u_1=0 <8$.
				\item Giả sử $u_n\leq 8$ với $n=k\geq 1$, tức là $u_k\leq 8$.\\
				Ta cần chứng minh $u_{k+1}\leq 8$.\\
				Thật vậy, $u_{k+1}=\dfrac{1}{2}u_k+4\leq \dfrac{1}{2}\cdot 8+4\leq 8$.
			\end{itemize}
			Vậy $u_n\leq 8$ với mọi $n\geq 1$, hay $(u_n)$ bị chặn trên bởi $8$.
			\item Với mọi $n\geq 1$, ta có $u_{n+1}-u_n=4-\dfrac{1}{2}u_n$. Mà $u_n\leq 8$ nên $u_{n+1}-u_n\geq 0$.\\
			Suy ra $u_n$ là dãy số tăng. Do đó $(u_n)$ bị chặn dưới bởi $u_1=0$.\\
			Kết hợp với câu a, ta được dãy số $(u_n)$ bị chặn.
		\end{enumerate}
	}
\end{bt}

\begin{bt}%[1D2V1-5]
% nguon: Archive/SP_KhoLatex/ID6_old/1K2-1.tex #3267
[VD]%[Trương Đăng Khoa]
	Trong các dãy số $(u_n)$ sau, dãy số nào bị chặn trên, bị chặn dưới và bị chặn?
	\begin{listEX}[3]
		\item $u_n=n^2+5$.
		\item $u_n=\dfrac{3n+1}{2n+5}$.
		\item $u_n=(-1)^n\cos \dfrac{\pi}{2n}$.
		\item $u_n=\dfrac{n^2+2n}{n^2+n+1}$.
		\item $u_n=\dfrac{n}{\sqrt{n^2+2n}+n}$.
	\end{listEX}
	\loigiai{
		\begin{enumerate}
			\item Dãy số bị chặn dưới bởi $6$, không bị chặn trên.
			\item Dãy $(u_n)$ bị chặn dưới bởi $0$. Vì $u_n<\dfrac{3n+1}{2n}=\dfrac{3}{2}+\dfrac{1}{2n}<\dfrac{3}{2}+1=\dfrac{5}{2}$ nên dãy số bị chặn trên bởi $\dfrac{5}{2}$. Vậy dãy số bị chặn.
			\item Ta có $|u_n|\leq 1$ nên dãy số bị chặn trên bởi 1, bị chặn dưới bởi $-1$.
			\item Dãy số bị chặn dưới bởi $0$. Vì $u_n<\dfrac{n^2+2n}{n^2}=1+\dfrac{2}{n}\leq 3$ nên dãy số bị chặn trên. Vậy dãy số bị chặn.
			\item Ta có $0<u_n\leq 1$ vậy dãy số bị chặn.
		\end{enumerate}
	}
\end{bt}

\begin{bt}%[1D2H2-7]
% nguon: Archive/SP_KhoLatex/ID6_old/1K2-1.tex #3268
%[Trương Đăng Khoa]%Ví dụ 1
	Một chồng cột gỗ được xếp thành các lớp, hai lớp liên tiếp hơn kém nhau một cột gỗ.
	\begin{center}
		\begin{tikzpicture}[font=\footnotesize, line join=round, line cap=round, >=stealth,scale=0.8]
			\def\r{0.2}
			\def\n{25}
			\def\g{110}
			\fill[teal!50!green](-6*\r,-3*\r)rectangle(3.5*\n*\r,0.5*\n*\r);
			\fill[teal!50!green,opacity=0.25](3.5*\n*\r,3*\r)rectangle(-6*\r,2*\n*\r);
			\foreach \j in {0,...,12}{
				\pgfmathsetmacro{\m}{\n-\j}
				\foreach \i in{0,...,\m}{
					\fill[left color=orange, right color=teal!30,draw=brown](2*\i*\r,0)++(60:2*\j*\r)++(\g:\r)--++(\g-90:6)arc(\g:-60:\r)--++(\g-270:6)--cycle;
					\fill[orange!20!brown!40,draw=teal](2*\i*\r,0)++(60:2*\j*\r)circle(\r);
				}
			}
		\end{tikzpicture}
	\end{center}
	\begin{enumerate}
		\item  Gọi $u_1=25$ là số cột gỗ có ở hàng dưới cùng của chồng cột gỗ, $u_n$ là số cột gỗ có ở hàng thứ $n$ tính từ dưới lên trên. Xét tính tăng, giảm của dãy số này.
		\item  Gọi $v_1=14$ là số cột gỗ có ở hàng trên cùng của chồng cột gỗ, $v_n$ là số cột gỗ có ở hàng thứ $n$ tính từ trên xuống dưới. Xét tinh tăng, giảm của dãy số này.
	\end{enumerate}
	\loigiai{
		\begin{enumerate}
			\item Ta có $u_n=26-n>u_{n+1}=26-n-1=25-n$.\\
			Vậy dãy số $(u_n)$ là dãy số giảm.
			\item Ta có $v_n=13+n<v_{n+1}=13+n+1=14+n$.\\
			Vậy dãy số $(u_n)$ là dãy số tăng
		\end{enumerate}
	}
\end{bt}

\begin{bt}%[1D2N1-2]
% nguon: Archive/SP_KhoLatex/ID6_old/1K2-5.tex #3373
%[Dự án 2 - TLDH - TeamTeXHoa - Nguyễn Cường]
	Cho dãy số $(u_n)$ xác định bởi $u_n=\dfrac{n+(-1)^n}{2n+1}$. Tìm $5$ số hạng đầu tiên của dãy số.
	\loigiai 
	{
	Ta có $u_1=0$; $u_2=\dfrac{3}{5}$; $u_3=\dfrac{2}{7}$; $u_4=\dfrac{5}{9}$; $u_5=\dfrac{4}{11}$.
}
\end{bt}

\begin{bt}%[1D2H1-2]
% nguon: Archive/SP_KhoLatex/ID6_old/1K2-5.tex #3374
%[Dự án 2 - TLDH - TeamTeXHoa - Nguyễn Cường]
	Cho dãy số $(u_n)$, từ đó dự đoán $u_n$
	\begin{multicols}{2}
		\begin{enumerate}
			\item $(u_n)\colon \heva{&u_1=5\\&u_{n+1}=u_n+3,\, n\ge 1}$;
			\item $(u_n)\colon \heva{&u_1=3\\&u_{n+1}=4u_n,\, n\ge 1}$.
		\end{enumerate}
	\end{multicols}
	\loigiai 
	{
	\begin{enumerate}
		\item Ta có 
		\begin{itemize}
						\begin{multicols}{2}
			\item $u_1=5$; 
			\item $u_2=5+1\cdot 3$; 
			\item $u_3=5+2\cdot 3$;
			\item $\ldots$
			\item Dự đoán $u_n=5+(n-1)\cdot 3$.
				\end{multicols}
		\end{itemize}
	Bằng phương pháp chứng minh quy nạp, ta chứng minh được $$u_n=5+(n-1)\cdot 3,\,\forall n \in \mathbb{N^*}.$$
	Vậy dãy số $(u_n)$ có số hạng tổng quát là $u_n=3n-2$.
		\item Ta có 				
		\begin{itemize}
			\begin{multicols}{2}
			\item $u_1=3$; 
			\item $u_2=3\cdot 4$; 
			\item $u_3=3\cdot 4^2$;
			\item $\ldots$
			\item $u_n=3\cdot 4^{n-1}$.
				\end{multicols}
		\end{itemize}		
	\end{enumerate}	
	}
\end{bt}

\begin{bt}%[1D2H1-2]
% nguon: Archive/SP_KhoLatex/ID6_old/1K2-5.tex #3375
%[Dự án 2 - TLDH - TeamTeXHoa - Nguyễn Cường]
	Cho dãy số $(u_n)$, từ đó dự đoán $u_n$.
	\begin{multicols}{2}
		\begin{enumerate}
			\item $(u_n)\colon \heva{&u_1=1\\&u_{n+1}=2u_n+3}$;
			\item $(u_n)\colon \heva{&u_1=3\\&u_{n+1}=\sqrt{1+u_n^2}}$.
		\end{enumerate}
	\end{multicols}
	\loigiai 
	{
			\begin{enumerate}
				\item Ta có 
				\begin{itemize}
						\begin{multicols}{2}
					\item $u_1=1=2^2-3$; 
					\item $u_2=5=2^3-3$; 
					\item $u_3=13=2^4-3$;
					\item $\ldots$
					\item  Dự đoán $u_n=2^{n+1}-3$.
						\end{multicols}
				\end{itemize}
				\item Ta có 
				\begin{itemize}
						\begin{multicols}{2}
					\item $u_1=3=\sqrt{3^2+0}$; 
					\item $u_2=\sqrt{10}=\sqrt{3^2+1}$; 
					\item $u_3=\sqrt{11}=\sqrt{3^2+2}$; 
					\item $\ldots$
					\item Dự đoán $u_n=\sqrt{3^2+n-1}$; 
						\end{multicols}
				\end{itemize}
			\end{enumerate}
	
	}
\end{bt}

\begin{bt}%[1D2H1-4]
% nguon: Archive/SP_KhoLatex/ID6_old/1K2-5.tex #3376
%[Dự án 2 - TLDH - TeamTeXHoa - Nguyễn Cường]
	Xét tính đơn điệu của các dãy số sau
	\begin{multicols}{2}
		\begin{enumerate}
			\item $u_n=2n+3$;
			\item $u_n=\dfrac{\sqrt{n}}{2^n}$.
		\end{enumerate}
	\end{multicols}
	\loigiai 
	{
	\begin{enumerate}
		\item Ta có $u_n=2n+3$ và $u_{n+1}=2(n+1)+3=2n+5$.\\
		Do đó, với mọi $n\in \mathbb{N}^*$, ta có $u_{n+1}-u_n=2n+5-2n-3=2>0$.\\
		Suy ra dãy số đã cho là dãy số tăng.
		\item Ta có $u_n=\dfrac{\sqrt{n}}{2^n}$ và $u_{n+1}=\dfrac{\sqrt{n+1}}{2^{n+1}}$.\\
		Do đó, với mọi $n\in \mathbb{N}^*$, ta có $\dfrac{u_{n+1}}{u_n}=\dfrac{\sqrt{n+1}}{2^{n+1}}\cdot \dfrac{2^n}{\sqrt{n}}=\dfrac{1}{2}\sqrt{\dfrac{n+1}{n}}$.\\
		Giả sử $\dfrac{1}{2}\sqrt{\dfrac{n+1}{n}}>1\Leftrightarrow\dfrac{1}{4}\cdot\dfrac{n+1}{n}>1\Leftrightarrow n+1>4n\Leftrightarrow 3n<1$ (vô lý).\\
		Vậy $\dfrac{u_{n+1}}{u_n}<1$. Suy ra dãy số đã cho là dãy số giảm.
	\end{enumerate}
	}
\end{bt}

\begin{bt}%[1D2H1-4]
% nguon: Archive/SP_KhoLatex/ID6_old/1K2-5.tex #3377
%[Dự án 2 - TLDH - TeamTeXHoa - Nguyễn Cường]
	Xét tính đơn điệu của các dãy số sau
	\begin{multicols}{2}
		\begin{enumerate}
			\item $u_n=\dfrac{n}{n^2+1}$;
			\item $u_n=\dfrac{\sqrt{n+1}-n}{n}$.
		\end{enumerate}
	\end{multicols}
	\loigiai 
	{
		\begin{enumerate}
			\item Ta có $u_n=\dfrac{n}{n^2+1}$ và $u_{n+1}=\dfrac{n+1}{(n+1)^2+1}=\dfrac{n+1}{n^2+2n+2}$.\\
			Do đó, với mọi $n\in \mathbb{N}^*$, ta có 
			\allowdisplaybreaks
			\begin{eqnarray*}
				u_{n+1}-u_n&=&\dfrac{n+1}{n^2+2n+2}-\dfrac{n}{n^2+1},\,\forall n\in \mathbb{N}^*.\\
				&=&\dfrac{(n+1)(n^2+1)-n(n^2+2n+2)}{(n^2+1)(n^2+2n+2)},\,\forall n\in \mathbb{N}^*.\\
				&=&\dfrac{n^3+n^2+n+1-n^3-2n^2-2n}{(n^2+1)(n^2+2n+2)},\,\forall n\in \mathbb{N}^*.\\
				&=&\dfrac{-n^2-n+1}{(n^2+1)(n^2+2n+2)}<0,\,\forall n\in \mathbb{N}^*.
			\end{eqnarray*}
			Suy ra dãy số đã cho là dãy số giảm.
			\item Ta có $u_n=\dfrac{\sqrt{n+1}-n}{n}=\dfrac{\sqrt{n+1}}{n}-1$ và $u_{n+1}=\dfrac{\sqrt{n+2}}{n+1}-1$.\\
			Do đó, với mọi $n\in \mathbb{N}^*$, ta có $$u_{n+1}-u_n=\dfrac{\sqrt{n+2}}{n+1}-\dfrac{\sqrt{n+1}}{n}=\dfrac{n\sqrt{n+2}-(n+1)\sqrt{n+1}}{n(n+1)}.$$
			Giả sử \allowdisplaybreaks
			\begin{eqnarray*}
				u_{n+1}-u_n>0&\Leftrightarrow&n\sqrt{n+2}-(n+1)\sqrt{n+1}>0,\,\forall n\in \mathbb{N}^*.\\
				&\Leftrightarrow&n\sqrt{n+2}>(n+1)\sqrt{n+1},\,\forall n\in \mathbb{N}^*.\\
				&\Leftrightarrow&n^2(n+2)>(n+1)^3,\,\forall n\in \mathbb{N}^*.\\
				&\Leftrightarrow&n^3+2n^2>n^3+3n^2+3n+1,\,\forall n\in \mathbb{N}^*.\\
				&\Leftrightarrow&-n^2-3n-1>0 (\text{vô lý}).
			\end{eqnarray*}
			Vậy $u_{n+1}-u_n<0$. Suy ra dãy số đã cho là dãy số giảm.
		\end{enumerate}
	}
\end{bt}

\begin{bt}%[1D2H1-4]
% nguon: Archive/SP_KhoLatex/ID6_old/1K2-5.tex #3378
%[Dự án 2 - TLDH - TeamTeXHoa - Nguyễn Cường]
	Xét tính đơn điệu của các dãy số sau
	\begin{multicols}{2}
		\begin{enumerate}
			\item $u_n=\dfrac{1}{n}-2$;
			\item $u_n=\dfrac{n-1}{n+1}$.
		\end{enumerate}
	\end{multicols}
	\loigiai 
	{
	\begin{enumerate}
		\item Ta có $u_n=\dfrac{1}{n}-2$ và $u_{n+1}=\dfrac{1}{n+1}-2$.\\
		Do đó, với mọi $n\in \mathbb{N}^*$, ta có $u_{n+1}-u_n=\dfrac{1}{n+1}-\dfrac{1}{n}=-\dfrac{1}{n(n+1)}<0$.\\
		Suy ra dãy số đã cho là dãy số giảm.
		\item Ta có $u_n=\dfrac{n-1}{n+1}=1-\dfrac{2}{n+1}$ và $u_{n+1}=1-\dfrac{2}{n+2}$.\\
		Do đó, với mọi $n\in \mathbb{N}^*$, ta có $u_{n+1}-u_n=-\dfrac{2}{n+2}+\dfrac{2}{n+1}=\dfrac{2}{(n+2)(n+1)}>0$.\\
		Suy ra dãy số đã cho là dãy số tăng.
	\end{enumerate}
	}
\end{bt}

\begin{bt}%[1D2H1-4]
% nguon: Archive/SP_KhoLatex/ID6_old/1K2-5.tex #3379
%[Dự án 2 - TLDH - TeamTeXHoa - Nguyễn Cường]
	Xét tính đơn điệu của các dãy số sau
	\begin{multicols}{2}
		\begin{enumerate}
			\item $u_n=\dfrac{2n+1}{5n+2}$;
			\item $u_n=2n^2+5$.
		\end{enumerate}
	\end{multicols}
	\loigiai 
	{
		\begin{enumerate}
			\item Ta có $u_n=\dfrac{2n+1}{5n+2}=\dfrac{2}{5}+\dfrac{1}{5(5n+2)}$ và $u_{n+1}=\dfrac{2}{5}+\dfrac{1}{5(5n+7)}$.\\
			Do đó, với mọi $n\in \mathbb{N}^*$, ta có $$u_{n+1}-u_n=\dfrac{1}{5(5n+7)}-\dfrac{1}{5(5n+2)}=-\dfrac{1}{(5n+2)(5n+7)}<0.$$
			Suy ra dãy số đã cho là dãy số giảm.
			\item Ta có $u_n=2n^2+5$ và $u_{n+1}=2(n+1)^2+5$.\\
			Do đó, với mọi $n\in \mathbb{N}^*$, ta có $u_{n+1}-u_n=2(n+1)^2+5-2n^2-5=4n+2>0$.\\
			Suy ra dãy số đã cho là dãy số tăng.
		\end{enumerate}
	}
\end{bt}

\begin{bt}%[1D2H1-4]
% nguon: Archive/SP_KhoLatex/ID6_old/1K2-5.tex #3380
%[Dự án 2 - TLDH - TeamTeXHoa - Nguyễn Cường]
	Xét tính đơn điệu của các dãy số sau
	\begin{multicols}{2}
		\begin{enumerate}
			\item $u_n=\dfrac{2n^2-1}{n^2+1}$;
			\item $u_n=\sqrt{n+1}-\sqrt{n}$.
		\end{enumerate}
	\end{multicols}
	\loigiai 
	{
		\begin{enumerate}
			\item Ta có $u_n=\dfrac{2n^2-1}{n^2+1}=2-\dfrac{3}{n^2+1}$ và $u_{n+1}=2-\dfrac{3}{(n+1)^2+1}$.\\
			Với mọi $n\in \mathbb{N}^*$, ta có
			\allowdisplaybreaks
			\begin{eqnarray*}
				(n+1)^2>n^2&\Leftrightarrow&\dfrac{1}{(n+1)^2+1}<\dfrac{1}{n^2+1}\\
				&\Leftrightarrow&2-\dfrac{3}{(n+1)^2+1}>2-\dfrac{3}{n^2+1}\\
				&\Leftrightarrow& u_{n+1}>u_n.
			\end{eqnarray*}
			Suy ra dãy số đã cho là dãy số tăng.
			\item Ta có $u_n=\sqrt{n+1}-\sqrt{n}=\dfrac{1}{\sqrt{n+1}+\sqrt{n}}$ và $u_{n+1}=\dfrac{1}{\sqrt{n+2}+\sqrt{n+1}}$.\\
			Với mọi $n\in \mathbb{N}^*$, ta có
			\allowdisplaybreaks
			\begin{eqnarray*}
				(\sqrt{n+2}+\sqrt{n+1}>\sqrt{n+1}+\sqrt{n}&\Leftrightarrow&\dfrac{1}{\sqrt{n+2}+\sqrt{n+1}}<\dfrac{1}{\sqrt{n+1}+\sqrt{n}}\\
				&\Leftrightarrow& u_{n+1}<u_n.
			\end{eqnarray*}
			Suy ra dãy số đã cho là dãy số giảm.
		\end{enumerate}
	}
\end{bt}

\begin{bt}%[1D2H1-4]
% nguon: Archive/SP_KhoLatex/ID6_old/1K2-5.tex #3381
%[Dự án 2 - TLDH - TeamTeXHoa - Nguyễn Cường]
	Xét tính đơn điệu của các dãy số sau
	\begin{multicols}{2}
		\begin{enumerate}
			\item $u_n=\dfrac{3n^2-2n+1}{n+1}$;
			\item $u_n=\dfrac{\sqrt{n+1}-1}{n}$.
		\end{enumerate}
	\end{multicols}
	\loigiai 
	{
		\begin{enumerate}
			\item Ta có $u_n=\dfrac{3n^2-2n+1}{n+1}=3n-5+\dfrac{6}{n+1}$ và $u_{n+1}=3n-2+\dfrac{6}{n+2}$.\\
			Do đó, với mọi $n\in \mathbb{N}^*$, ta có $$u_{n+1}-u_n=3-\dfrac{6}{(n+1)(n+2)}.$$
			Với $n\in \mathbb{N}^*$, $(n+1)(n+2)\ge 6\Leftrightarrow \dfrac{6}{(n+1)(n+2)}\le 1\Leftrightarrow 3-\dfrac{6}{(n+1)(n+2)}>0$.\\
			Suy ra, $u_{n+1}-u_n>0$ hay dãy số đã cho là dãy số tăng.
			\item Ta có $u_n=\dfrac{\sqrt{n+1}-1}{n}=\dfrac{n}{n\left(\sqrt{n+1}+1\right)}=\dfrac{1}{\sqrt{n+1}+1}$ và $u_{n+1}=\dfrac{1}{\sqrt{n+2}+1}$.\\
			Do đó, với mọi $n\in \mathbb{N}^*$, ta có $u_{n+1}-u_n<0$.\\
			Suy ra dãy số đã cho là dãy số giảm.
		\end{enumerate}
	}
\end{bt}

\begin{bt}%[1D2H1-4]
% nguon: Archive/SP_KhoLatex/ID6_old/1K2-5.tex #3382
%[Dự án 2 - TLDH - TeamTeXHoa - Nguyễn Cường]
Xét tính tăng giảm của dãy số $(u_n)$ với $u_n=\dfrac{3^n}{2^{n+1}}$.
	\loigiai 
	{
	Ta có $u_{n+1}=\dfrac{3^{n+1}}{2^{n+2}}$, suy ra $\dfrac{u_{n+1}}{u_n}=\dfrac{3^{n+1}}{2^{n+2}}\cdot \dfrac{2^{n+1}}{3^n}=\dfrac{3}{2}>1$.\\
	Do $u_n>0,\, \forall n\in\mathbb{N}^*$ nên $u_{n+1}>u_n, \, \forall n\in\mathbb{N}^*$ hay dãy số đã cho là dãy số tăng.
	}
\end{bt}

\begin{bt}%[1D2H1-4]
% nguon: Archive/SP_KhoLatex/ID6_old/1K2-5.tex #3383
%[Dự án 2 - TLDH - TeamTeXHoa - Nguyễn Cường]
	Xét tính tăng giảm của dãy số $(u_n)$ với $u_n=\dfrac{\sqrt{n}}{2^n}$.
	\loigiai 
	{
		Ta có $u_{n+1}=\dfrac{\sqrt{n+1}}{2^{n+1}}$, suy ra $\dfrac{u_{n+1}}{u_n}=\dfrac{\sqrt{n+1}}{2^{n+1}}\cdot \dfrac{2^n}{\sqrt{n}}=\dfrac{1}{2}\cdot\sqrt{\dfrac{n+1}{n}}=\dfrac{1}{2}\cdot\sqrt{1+\dfrac{1}{n}}$.\\
		Với $\forall n\in\mathbb{N}^*$, suy ra $n\ge 1\Rightarrow \dfrac{1}{n}\le 1\Rightarrow\dfrac{u_{n+1}}{u_n}\le \dfrac{1}{2}\cdot\sqrt{2}<1$.\\
		Do $u_n>0,\, \forall n\in\mathbb{N}^*$ nên $u_{n+1}<u_n, \, \forall n\in\mathbb{N}^*$ hay dãy số đã cho là dãy số giảm.
	}
\end{bt}

\begin{bt}%[1D2H1-4]
% nguon: Archive/SP_KhoLatex/ID6_old/1K2-5.tex #3384
%[Dự án 2 - TLDH - TeamTeXHoa - Nguyễn Cường]
	Xét tính tăng giảm của dãy số $(u_n)$ với $u_n=\dfrac{3^n}{n^2}$.
	\loigiai 
	{
		Ta có $u_{n+1}=\dfrac{3^{n+1}}{(n+1)^2}$, suy ra $\dfrac{u_{n+1}}{u_n}=\dfrac{3^{n+1}}{(n+1)^2}\cdot \dfrac{n^2}{3^n}=3\left(\dfrac{n}{n+1}\right)^2\Rightarrow \dfrac{u_n}{u_{n+1}}=\dfrac{1}{3}\left(1+\dfrac{1}{n}\right)^2$.\\
		Khi đó, $\dfrac{u_n}{u_{n+1}}>1\Leftrightarrow 1+\dfrac{1}{n}>\sqrt{3}\Leftrightarrow \dfrac{1}{n}<\sqrt{3}-1\Leftrightarrow n>\dfrac{1}{\sqrt{3}-1}$, mà $n\in\mathbb{N}^*\Rightarrow n\ge 2$.\\
		Hơn nữa $u_n>0,\, \forall n\in\mathbb{N}^*$ nên $\heva{&u_{n+1}<u_n\Leftrightarrow n=1\\&u_{n+1}>u_n\Leftrightarrow n\ge 2.}$\\
		Do đó, dãy số đã cho không tăng không giảm.
	}
\end{bt}

\begin{bt}%[1D2H1-4]
% nguon: Archive/SP_KhoLatex/ID6_old/1K2-5.tex #3385
%[Dự án 2 - TLDH - TeamTeXHoa - Nguyễn Cường]
	Xét tính tăng giảm của dãy số $(u_n)$ với $u_n=\sqrt{n}-\sqrt{n-1}$.
	\loigiai 
	{
		Ta có $u_{n+1}=\sqrt{n+1}-\sqrt{n}$, suy ra $u_{n+1}-u_n=\sqrt{n+1}-2\sqrt{n}+\sqrt{n-1}$.\\
		Lại có $$\left(\sqrt{n+1}+\sqrt{n-1}\right)^2-\left(2\sqrt{n}\right)^2=2n+2\sqrt{n^2-1}-4n=2\left(\sqrt{n^2-1}-n\right)<0,\,\forall n\in\mathbb{N}^*.$$
		Do đó, $u_{n+1}-u_n=\sqrt{n+1}-2\sqrt{n}+\sqrt{n-1}<0,\,\forall n\in\mathbb{N}^*$.\\
		Vậy dãy số đã cho là dãy số giảm.
	}
\end{bt}

\begin{bt}%[1D2V1-4]
% nguon: Archive/SP_KhoLatex/ID6_old/1K2-5.tex #3386
%[Dự án 2 - TLDH - TeamTeXHoa - Nguyễn Cường]
	Với giá trị nào của $a$ thì dãy số $(u_n)$ với $u_n=\dfrac{na+2}{n+1}$
	\begin{multicols}{2}
		\begin{enumerate}
			\item là dãy số tăng.
			\item là dãy số giảm.
		\end{enumerate}
	\end{multicols}
	\loigiai 
	{
		Ta có $u_n=\dfrac{na+2}{n+1}=a+\dfrac{2-a}{n+1}$, suy ra $u_{n+1}=a+\dfrac{2-a}{n+2}$.\\
		Khi đó, $u_{n+1}-u_n=\dfrac{a-2}{(n+1)(n+2)}$.
	\begin{enumerate}
		\item $(u_n)$ là dãy số tăng khi và chỉ khi $u_{n+1}-u_n=\dfrac{a-2}{(n+1)(n+2)}>0\Leftrightarrow a>2$.
		\item $(u_n)$ là dãy số giảm khi và chỉ khi $u_{n+1}-u_n=\dfrac{a-2}{(n+1)(n+2)}<0\Leftrightarrow a<2$.
	\end{enumerate}
	}
\end{bt}

\begin{bt}%[1D2H1-5]
% nguon: Archive/SP_KhoLatex/ID6_old/1K2-5.tex #3387
%[Dự án 2 - TLDH - TeamTeXHoa - Nguyễn Cường]
Xét tính bị chặn của các dãy số sau
\begin{multicols}{2}
	\begin{enumerate}
		\item $u_n=\dfrac{n^2+1}{2n^2-3}$;
		\item $u_n=\dfrac{7n+5}{5n+7}$.
	\end{enumerate}
\end{multicols}
	\loigiai 
	{
		\begin{enumerate}
			\item Ta có $u_n=\dfrac{n^2+1}{2n^2-3}=\dfrac{1}{2}+\dfrac{5}{2(2n^2-3)}$.\\
			Với $\heva{&n=1\Rightarrow u_1=-2\\&\forall n\ge 2\Rightarrow 2n^2-3>0\Rightarrow u_n>\dfrac{1}{2}}\Rightarrow u_n\ge -2\, \forall n\in\mathbb{N}^*$.\\
			Xét $\dfrac{u_{n+1}}{u_n}=\dfrac{(n+1)^2+1}{2(n+1)^2-3}\cdot\dfrac{2n^2-3}{n^2+1}$.\\
			Nhận thấy, $\forall u_n>0$ thì
			\allowdisplaybreaks
			\begin{eqnarray*}
				\dfrac{u_{n+1}}{u_n}<1&\Leftrightarrow&(n^2+2n+2)(2n^2-3)<(n^2+1)(2n^2+4n-1)\\
				&\Leftrightarrow& 4n^4-3n^2+4n^3-6n+4n^2-6<4n^4+4n^3-n^2+2n^2+4n-1\\
				&\Leftrightarrow& n^2-6n-6<n^2+4n-1\\
				&\Leftrightarrow& 0<10n+5,\,\forall n\in\mathbb{N}^*.
			\end{eqnarray*}
		Do đó, $u_{n+1}<u_n<\ldots<u_2=1$.\\
		Vậy $-2<u_n<1,\,\forall n\in\mathbb{N}^*$, suy ra dãy số đã cho bị chặn.
			\item Ta có $u_n=\dfrac{7n+5}{5n+7}=\dfrac{7}{5}-\dfrac{24}{5(5n+7)}<\dfrac{7}{5},\,\forall n\in\mathbb{N}^*$.\\
			Mặt khác $u_n\ge 1,\,\forall n\in\mathbb{N}^*$.\\
			Do đó, $1\le u_n<\dfrac{7}{5},\,\forall n\in\mathbb{N}^*$.\\ Vậy dãy số đã cho bị chặn.
		\end{enumerate}
	}
\end{bt}

\begin{bt}%[1D2H1-5]
% nguon: Archive/SP_KhoLatex/ID6_old/1K2-5.tex #3388
%[Dự án 2 - TLDH - TeamTeXHoa - Đào Trung Kiên]
	Xét tính bị chặn của các dãy số sau
	\begin{multicols}{2}
		\begin{enumerate}
			\item  $u_n=\dfrac{1}{2 n^2-3}$.
			\item  $u_n=\dfrac{1}{n(n+1)}$.
		\end{enumerate}
	\end{multicols}
	\loigiai{
		\begin{enumerate}
			\item Với $n=1 \Rightarrow u_1=-1 \quad \quad \Rightarrow u_1 \geq -1$.\\
			Với $\forall n \geq 2\Rightarrow 2 n^2-3>0,\Rightarrow u_n>0$.\\ Xét $\dfrac{u_{n-1}}{u_n}=\dfrac{2 n^2-3}{2(n+1)^2-3}<1 \Leftrightarrow n<n+1$.\\
			Do đó, suy ra: $u_n<u_{n-1}<\ldots<u_2=\dfrac{1}{5}$.\\
			Vậy $-1 \leq u_n<\dfrac{1}{5} \Rightarrow\left(u_n\right)$ bị chặn.
			\item  Ta dễ dàng thấy\\
			+ $u_n>0, \forall n$ do đó nó bị chặn dưới.\\
			+ Vì $n(n+1) \geq 2 \Leftrightarrow u_n \leq \dfrac{1}{2}$ do đó nó bị chặn trên.\\
			Vậy ta được $0<u_n \leq \dfrac{1}{2},\,\forall n \in \mathbb{N^*}$, do đó nó bị chặn.	
		\end{enumerate}
	}
\end{bt}

\begin{bt}%[1D2H1-5]
% nguon: Archive/SP_KhoLatex/ID6_old/1K2-5.tex #3389
%[Dự án 2 - TLDH - TeamTeXHoa - Đào Trung Kiên]
	Xét tính bị chặn của các dãy số sau
	\begin{multicols}{2}
		\begin{enumerate}
			\item  $u_n=\dfrac{1}{2 n^2-1}$.
			\item  $u_n=\dfrac{n-1}{\sqrt{n^2+1}}$.
		\end{enumerate}
	\end{multicols}
	\loigiai{
		\begin{enumerate}
			\item  Với $\forall n \in \mathbb{N}^*: 2 n^2-1 > 0$ nên $u_n>0, \forall n\in \mathbb{N}^*$\\
			Xét $\dfrac{u_{n+1}}{u_n}=\dfrac{2 n^2-1}{2(n+1)^2-1}<1 \Leftrightarrow n<n+1$.\\
			Do đó $u_n<u_{n-1}<\ldots<u_2<u_1=1$.\\
		Suy ra $-1<u_n<1,\,\forall n\in \mathbb{N}^*$.\\Vậy dãy số $\left(u_n\right)$ bị chặn.
			\item  Với $\forall n \in \mathbb{N}^*$, ta có $n-1 \geq 0$ và $\sqrt{n^2+1}>0$ nên $u_n \geq 0, \forall n$.\\
			Và $\dfrac{n-1}{\sqrt{n^2+1}} \leq 1,\,\forall n \in \mathbb{N}$.\\
		Suy ra $-1 \leq u_n \leq 1,\,\forall n\in \mathbb{N}^*$.\\ Vậy dãy số $\left(u_n\right)$ bị chặn.
		\end{enumerate}
	}
\end{bt}

\begin{bt}%[1D2H1-5]
% nguon: Archive/SP_KhoLatex/ID6_old/1K2-5.tex #3390
%[Dự án 2 - TLDH - TeamTeXHoa - Đào Trung Kiên]
	Xét tính bị chặn của các dãy số sau
	\begin{multicols}{2}
		\begin{enumerate}
			\item  $u_n=\dfrac{2 n^2}{n^2+1}$.
			\item  $u_n=\dfrac{2 n^2+2 n+1}{n^2+n+4}$.
		\end{enumerate}
	\end{multicols}
	\loigiai{
		\begin{enumerate}
			\item 	Vì $\left\{\begin{array}{l}2 n^2 \geq 0 \\ n^2+1>0\end{array},\, \forall n \in \mathbb{N}^*\Rightarrow u_n \geq 0.\right.$\\
			Mặt khác, $u_n=\dfrac{2\left(n^2+1\right)-2}{n^2+1}=2-\dfrac{2}{n^2+1}<2,\, \forall n \in \mathbb{N}^*$.\\
		Suy ra  $0\le u_n<2,\,\forall n\in \mathbb{N}^*$.\\
		Vậy dãy số $\left(u_n\right)$ bị chặn.
			\item  Vì $\left\{\begin{array}{l}2 n^2+2 n+1=2\left(n^2+1\right)+1>0 \\ n^2+n+4=n(n+1)+4>0\end{array}, \forall n \in \mathbb{N}^*\Rightarrow u_n>0\right.$\\
			Mặt khác, $u_n=\dfrac{2 n^2+2 n+1}{n^2+n+4}=\dfrac{2\left(n^2+n+4\right)-7}{n^2+n+4}=2-\dfrac{7}{n^2+n+4}<2$.\\
			Suy ra  $0<u_n<2,\,\forall n\in \mathbb{N}^*$.\\
			Vậy dãy số $\left(u_n\right)$ bị chặn.
		\end{enumerate}
	}
\end{bt}

\begin{bt}%[1D2H1-5]
% nguon: Archive/SP_KhoLatex/ID6_old/1K2-5.tex #3391 
% KIEM_TRA: cần xem lại đáp số/dạng
%[Dự án 2 - TLDH - TeamTeXHoa - Đào Trung Kiên]
	Cho dãy số $\left(u_n\right)$, với $u_n=\dfrac{3 n+(-1)^n}{4 n+(-1)^{n+1}}$.
	\begin{enumerate}
		\item  Tính $6$ số hạng đầu tiên của dãy, nêu nhận xét về tính đơn điệu của dãy số.
		\item  Tính $u_{2 n}$ và $u_{2 n-1}$. Chứng minh rằng $0<u_n \leq \dfrac{3 n+4}{4 n-1}$.
	\end{enumerate}
	\loigiai{
		\begin{enumerate}
			\item Ta có: $u_1=\dfrac{2}{5} ; u_2=1 ; u_3=\dfrac{8}{13} ; u_4=\dfrac{13}{15} ; u_5=\dfrac{16}{21} ; u_6=\dfrac{19}{23}$.\\
			Nhận thấy dãy số không tăng cũng không giảm.
			\item Ta có $\heva{&u_{2n}=\dfrac{6 n+1}{8 n-1} \\& u_{2n-1}=\dfrac{6 n+2}{8 n+5}.} (n\in \mathbb{N}^*)$\\
			Tổng quát, với $n=2k, (k \geq 1, k \in \mathbb{N}^*)\Rightarrow u_n=\dfrac{3 n+1}{4 n-1} \Rightarrow 0<u_n \leq \dfrac{3 n+1}{4 n-1}$.\\
			Với $n=2k+1, (k \in \mathbb{N}) \Rightarrow u_n=\dfrac{3 n-1}{4 n+1} \Rightarrow \heva{&u_n>0 \\& u_n=\dfrac{3 n-1}{4 n+1}<\dfrac{3 n+4}{4 n+1}<\dfrac{3 n+4}{4 n-1}.}$\\
			$ \Rightarrow 0<u_n \leq \dfrac{3 n+4}{4 n+1}.$\\
			Vậy với mọi $n\in \mathbb{N}^*$ thì $0<u_n \leq \dfrac{3 n+4}{4 n-1}$.	
		\end{enumerate}
	}
\end{bt}

\begin{bt}%[1D2H1-5]
% nguon: Archive/SP_KhoLatex/ID6_old/1K2-5.tex #3392
%[Dự án 2 - TLDH - TeamTeXHoa - Đào Trung Kiên]
	Xét tính bị chặn trên, bị chặn dưới, bị chặn của các dãy số $\left(u_n\right)$ cho bởi
	\begin{multicols}{2}
		\begin{enumerate}
			\item  $u_n=\dfrac{2 n+3}{n+2}$.
			\item  $u_n=\dfrac{1}{n(n+1)}$.
		\end{enumerate}
	\end{multicols}
	\loigiai{
		\begin{enumerate}
			\item  $u_{n+1}-u_n=\dfrac{2 n+5}{n+3}-\dfrac{2 n+3}{n+2}=\dfrac{1}{(n+3)(n+2)}>0$ nên dãy là dãy tăng.\\
			Hơn nữa $u_n=\dfrac{2 n+3}{n+2}=\dfrac{2(n+2)-1}{n+2}=2-\dfrac{1}{n+2}<1 \Rightarrow u_n$ bị chặn trên bởi 2, chặn dưới bởi $u_1=\dfrac{5}{3}$.\\
			Vậy dãy đã cho bị chặn.
			\item Dễ thấy $u_n>0$ và $\dfrac{u_{n-1}}{u_n}=\dfrac{n(n+1)}{(n+1)(n+2)}=\dfrac{n}{n+2}<1$ suy ra dãy số đã cho là dãy giảm và bị chặn trên bởi $u_1=\dfrac{1}{2}$, bị chặn dưới bởi $0$ nên dãy số bị chặn.
		\end{enumerate}
	}
\end{bt}

\begin{bt}%[1D2V1-5]
% nguon: Archive/SP_KhoLatex/ID6_old/1K2-5.tex #3393
%[Dự án 2 - TLDH - TeamTeXHoa - Đào Trung Kiên]
	Xét tính bị chặn trên, bị chặn dưới, bị chặn của các dãy số $\left(u_n\right)$ cho bởi
	\begin{multicols}{2}
		\begin{enumerate}
			\item $u_n=\dfrac{n^2+2 n}{n^2+n+1}$.
			\item $u_n=\dfrac{n}{\sqrt{n^2+2 n}+n}$.
		\end{enumerate}
	\end{multicols}
	\loigiai{
		\begin{enumerate}
			\item  $u_{n+1}-u_n=\dfrac{n^2+2 n+1+2 n+2}{n^2+2 n+1+n+1+1}-\dfrac{n^2+2 n}{n^2+n+1}=\dfrac{n^2+4 n+3}{n^2+3 n+3}-\dfrac{n^2+2 n}{n^2+n+1}>0$ và
			$$
			u_n=\dfrac{n^2+2 n}{n^2+n+1}=\dfrac{n^2+n+1+n}{n^2+n+1}=1+\dfrac{n}{n^2+n+1}>1.
			$$
			Nên dãy đã cho là dãy tăng, bị chặn dưới bởi $1$.\\
			Dễ thấy $u_n<2, \forall n\in \mathbb{N}^*$ nên dãy số bị chặn.
			\item Ta có $u_n=\dfrac{n}{\sqrt{n^2+2 n}+n}=\dfrac{n\left(\sqrt{n^2+2 n}-n\right)}{2 n}=\dfrac{\sqrt{n^2+2 n}-n}{2}>0$.\\
			Lại có
			$$
			\begin{aligned}
				& \dfrac{u_{n-1}}{u_n}=\dfrac{\sqrt{n^2+4 n+3}-n-1}{\sqrt{n^2+2 n}-n}>1 \Leftrightarrow \sqrt{n^2+4 n+3}>\sqrt{n^2+2 n}+1 \\
				& \Leftrightarrow n^2+4 n+3>n^2+2 n+1+2 \sqrt{n^2+2 n}\\
				& \Leftrightarrow n+1>\sqrt{n^2+2 n} \Leftrightarrow n^2+2 n+1>n^2+2 n \quad (*).
			\end{aligned}
			$$
			Do $(*)$ hiển nhiên đúng nên ta có dãy đã cho là dãy tăng, và bị chặn dưới bởi $u_1=\dfrac{1}{\sqrt{3}+1}$.\\
			Hơn nữa $u_n=\dfrac{n}{\sqrt{n^2+2 n}+n}<\dfrac{n}{n}=1 \Rightarrow u_n$ bị chặn trên bởi $1$.\\
			Vậy dãy đã cho bị chặn.	
		\end{enumerate}
	}
\end{bt}

\begin{bt}%[1D2H1-4]
% nguon: Archive/SP_KhoLatex/ID6_old/1K2-5.tex #3394
%[Dự án 2 - TLDH - TeamTeXHoa - Đào Trung Kiên]
	Chứng minh rằng dãy số $u_n=\dfrac{n+3}{n+1}$ giảm và bị chặn.
	\loigiai{
		Xét \begin{eqnarray*}
			u_{n+1}-u_n&=&\dfrac{n+4}{n+2}-\dfrac{n+3}{n+1}=\dfrac{(n+4)(n+1)-(n+2)(n+3)}{(n+2)(n+1)}\\
			&=&\dfrac{n^2+5 n+4-n^2-5 n-6}{(n+2)(n+1)}=-\dfrac{2}{(n+2)(n+1)},\,\forall n\in \mathbb{N}^*.
		\end{eqnarray*} 
		Nhận thấy $u_{n+1}-u_n<0=u_{n+1}<u_n$, do đó, dãy số $u_n$ giảm.\\
		Viết lại $u_n$ dưới dạng $u_n=1+\dfrac{2}{n+1}>1 \Rightarrow\left(u_n\right)$ bị chặn dưới và bị chặn.}	
\end{bt}

\begin{bt}%[1D2V1-4]
% nguon: Archive/SP_KhoLatex/ID6_old/1K2-5.tex #3395
%[Dự án 2 - TLDH - TeamTeXHoa - Đào Trung Kiên]
	Chứng minh rằng dãy số $u_n=\dfrac{1}{1\cdot 2}+\dfrac{1}{2\cdot 3}+\dfrac{1}{3\cdot 4}+\ldots+\dfrac{1}{n(n+1)}$ tăng và bị chặn trên.
	\loigiai{
		Viết lại $u_n$ dưới dạng
		$$
		u_n=\dfrac{2-1}{1\cdot 2}+\dfrac{3-2}{2\cdot 3}+\dfrac{4-3}{3\cdot 4}+\ldots+\dfrac{(n+1)-n}{n(n+1)}=1-\dfrac{1}{2}+\dfrac{1}{2}-\dfrac{1}{3}+\dfrac{1}{3}-\dfrac{1}{4}+\ldots+\dfrac{1}{n}-\dfrac{1}{n+1}=1-\dfrac{1}{n+1}
		$$
		Xét hiệu: $u_{n+1}-u_n=1-\dfrac{1}{n+2}-\left(1-\dfrac{1}{n+1}\right)=\dfrac{1}{n+1}-\dfrac{1}{n+2}>0 \Rightarrow\left(u_n\right)$ tăng.\\
		Nhận thấy $u_n=1-\dfrac{1}{n+1}<1 \Rightarrow\left(u_n\right)$ bị chặn trên}
\end{bt}

\begin{bt}%[1D2V1-5]
% nguon: Archive/SP_KhoLatex/ID6_old/1K2-5.tex #3396
%[Dự án 2 - TLDH - TeamTeXHoa - Đào Trung Kiên]
	Chứng minh rằng dãy số $u_n=\dfrac{n^2+1}{2 n^2-3}$ là một dãy số bị chặn.
	\loigiai{
		Viết lại $u_n$ dưới dạng $u_n=\dfrac{n^2-\dfrac{3}{2}}{2 n^2-3}+\dfrac{5}{2\left(2 n^2-3\right)}=\dfrac{1}{2}+\dfrac{5}{2\left(2 n^2-3\right)}$.\\
		Ta có $u_1=-2$.\\
		Với $\forall n \geq 2\Rightarrow 2 n^2-3>0 \Rightarrow u_n>\dfrac{1}{2}$.\\
		Xét $\dfrac{u_{n+1}}{u_n}=\dfrac{(n+1)^2+1}{2(n+1)^2-3} \cdot \dfrac{2 n^2-3}{n^2+1}$.\\
		Nhận thấy với $\forall u_n>0$ thì 
		\begin{eqnarray*}
			\dfrac{u_{n-1}}{u_n}<1 &\Leftrightarrow&\left(n^2+2 n+2\right)\left(2 n^2-3\right)<\left(n^2+1\right)\left(2 n^2+4 n-1\right)\\
			& \Leftrightarrow& 4 n^4-3 n^2+4 n^3-6 n+4 n^2-6<4 n^4+4 n^3-n^2+2 n^2+4 n-1\\
			& \Leftrightarrow& n^2-6 n-6<n^2+4 n-1 \\
			& \Leftrightarrow& 0<10 n+5, \forall n \in \mathbb{N}^*.
		\end{eqnarray*}
		Do đó, $u_{n+1}<u_n<\ldots<u_2=1$.\\
		Vậy $-2<u_n<1 \Rightarrow\left(u_n\right)$ bị chặn	}
\end{bt}

\begin{bt}%[1D2V1-5]
% nguon: Archive/SP_KhoLatex/ID6_old/1K2-5.tex #3397
%[Dự án 2 - TLDH - TeamTeXHoa - Đào Trung Kiên]
	Cho dãy số $\left\{\begin{array}{l}u_1=0 \\ u_{n-1}=\dfrac{1}{2} u_n+4,,\, \forall n \ge2.\end{array}\right.$
	\begin{enumerate}
		\item  Chứng minh rằng $u_n<8$.
		\item  Xét tính bị chặn của dãy số.
	\end{enumerate}
	\loigiai{
		\begin{enumerate}
			\item Giả sử tồn tại $u_n \geq 8, (\forall n \in \mathbb{N}^*, n>1)$ thì $u_{n-1} \geq 8$.  Suy ra $u_{n-2}, u_{n-3} \cdots u_2, u_1 \geq 8$.\\
			Điều đó vô lí do $u_1=0<8$. Nên điều giả sử là sai. Suy ra $u_n<8$	
			\item  Xét $u_{n+1}-u_n=\dfrac{1}{2} u_n+4-u_n=4-\dfrac{u_n}{2}=\dfrac{8-u_n}{2}>0 \Rightarrow u_{n+1}>u_n$.\\
			Suy ra dãy tăng. Mà $u_n<8$ và $u_1 \geq 0 \Rightarrow u_n>0$. Suy ra dãy bị chặn dưới.\\
			Vậy dãy tăng và bị chặn.
		\end{enumerate}
	}
\end{bt}

\begin{bt}%[1D2V1-5]
% nguon: Archive/SP_KhoLatex/ID6_old/1K2-5.tex #3398
%[Dự án 2 - TLDH - TeamTeXHoa - Đào Trung Kiên]
	Cho dãy số $\heva{&u_1=1 \\ &u_{n+1}=\dfrac{u_n+2}{u_n+1},\,\forall n\geq 1.}$
	\begin{enumerate}
		\item  Tìm $5$ số hạng đầu tiên của dãy số.
		\item Chứng minh rằng dãy số bị chặn dưới bởi $1$ và bị chặn trên bởi $\dfrac{3}{2}$.
	\end{enumerate}
	\loigiai{
		\begin{enumerate}
			\item $u_1=1 ; u_2=\dfrac{1+2}{1+1}=\dfrac{3}{2} ; u_3=\dfrac{\dfrac{3}{2}+2}{\dfrac{3}{2}+1}=\dfrac{7}{5} ; u_4=\dfrac{\dfrac{7}{5}+2}{\dfrac{7}{5}+1}=\dfrac{17}{12} ; u_5=\dfrac{\dfrac{17}{12}+2}{\dfrac{17}{12}+1}=\dfrac{41}{29}.$
			\item $u_1=1>0\Rightarrow u_n>0$ suy ra $u_{n+1}=1+\dfrac{1}{u_n+1}>1$.\\
			Giả sử $u_n>\dfrac{3}{2}\Rightarrow u_{n-1}<1$ (vô lí).\\
			Vậy $u_n\leq \dfrac{3}{2},\,\forall n \in \mathbb{N^*}$.\\
			Vậy ta có điều phải chứng minh.
		\end{enumerate}
	}
\end{bt}

\begin{bt}%[1D2V1-4]
% nguon: Archive/SP_KhoLatex/ID6_old/1K2-5.tex #3399
%[Dự án 2 - TLDH - TeamTeXHoa - Đào Trung Kiên]
	Chứng minh rằng dãy số $\left\{\begin{array}{l}u_1=\sqrt{2} \\ u_{n+1}=\sqrt{u_n+2}\end{array}\right.$ tăng và bị chặn trên bởi $2$.
	\loigiai{
		Dễ thấy $u_n>1$.\\
		Giả sử tồn tại $u_n \geq 2 \Rightarrow \sqrt{u_{n-1}+2} \geq 2 \Rightarrow u_{n-1} \geq 2$.\\
		Như vậy, nếu tồn tại $u_n \geq 2$ thì suy ra $u_{n-1} \geq 2$, từ đó cũng suy ra được $$u_{n-2}, u_{n-3} \ldots u_2, u_1 \geq 2.$$
		Do $u_1=\sqrt{2}<2$ nên điều giả sử là sai.\\
		Suy ra $u_n<2$.\\
		Xét $u_{n+1}-u_n=\sqrt{u_n+2}-u_n=\dfrac{u_n+2-u_n^2}{\sqrt{u_n+2}+u_n}=\dfrac{\left(2-u_n\right)\left(1+u_n\right)}{\sqrt{u_n+2}+u_n}>0$.\\
		Suy ra $u_{n+1}>u_n$ nên dãy số $(u_n)$ à dãy số tăng.\\
		Vậy dãy đã cho tăng và bị chặn trên bởi $2$.}
\end{bt}

\begin{bt}%[1D2H1-2]
% nguon: Archive/SP_KhoLatex/ID6_old/1K2-5.tex #3400
%[Dự án 2 - TLDH - TeamTeXHoa - Đào Trung Kiên]%[ID6]
	Viết năm số hạng đầu và số hạng thứ 100 của các dãy số $\left(u_n\right)$ có số hạng tổng quát cho bởi
	\begin{multicols}{3}
		\begin{enumerate}
			\item  $u_n=3 n-2$;
			\item  $u_n=3 \cdot 2^n$;
			\item  $u_n=\left(1+\dfrac{1}{n}\right)^n$.
		\end{enumerate}
	\end{multicols}
	\loigiai{
		\begin{enumerate}
			\item  $u_1=1, u_2=4, u_3=7, u_4=10, u_5=13, u_{100}=298$.
			\item  $u_1=6, u_2=12, u_3=24, u_4=48, u_5=96, u_{100}=3,803 \times 10^{30}$.
			\item  $u_1=2, u_2=\dfrac{9}{4}, u_3=\dfrac{64}{27}, u_4=\dfrac{625}{256}, u_5=2,48832, u_{100}=2{,}7148$.
		\end{enumerate}
	}
\end{bt}

\begin{bt}%[1D2H1-4]
% nguon: Archive/SP_KhoLatex/ID6_old/1K2-5.tex #3402
%[Dự án 2 - TLDH - TeamTeXHoa - Đào Trung Kiên]%[ID6]
	Xét tính tăng, giảm của dãy số $\left(u_n\right)$, biết
	\begin{multicols}{3}
		\begin{enumerate}
			\item $u_n=2 n-1$;
			\item  $u_n=-3 n+2$;
			\item  $u_n=\dfrac{(-1)^{n-1}}{2^n}$.
		\end{enumerate}
	\end{multicols}
	\loigiai{
		\begin{enumerate}
			\item  Ta có $u_{n+1}-u_n=2>0,\forall n \in \mathbb{N}^*$ suy ra đây là dãy số tăng.
			\item  Ta có $u_{n+1}-u_n=-3<0,\forall n \in \mathbb{N}^*$ suy ra đây là dãy số giảm.
			\item  $u_3=\dfrac{1}{8}>u_2=\dfrac{-1}{4}<u_1=\dfrac{1}{2}$ suy ra đây là dãy số không tăng không giảm.
		\end{enumerate}
	}	
\end{bt}

\begin{bt}%[1D2H1-5]
% nguon: Archive/SP_KhoLatex/ID6_old/1K2-5.tex #3403
%[Dự án 2 - TLDH - TeamTeXHoa - Đào Trung Kiên]%[ID6]
	Trong các dãy số $\left(u_n\right)$ sau, dãy số nào bị chặn dưới, bị chặn trên, bị chặn?
	\begin{multicols}{2}
		\begin{enumerate}
			\item  $u_n=n-1$;
			\item  $u_n=\dfrac{n+1}{n+2}$;
			\item  $u_n=\sin n$;
			\item  $u_n=(-1)^{n-1} n^2$.
		\end{enumerate}
	\end{multicols}
	\loigiai{
		\begin{enumerate}
			\item  Ta có $u_n=n-1 \geq 0\left(\forall n \in \mathbb{N}^*\right)$ suy ra $u_n$ bị chặn dưới.
			\item  Ta có: $u_n=\dfrac{n+1}{n+2}=1-\dfrac{1}{n+2}<1 ; u_n=\dfrac{n+1}{n+2}=1-\dfrac{1}{n+2} \geq 0\left(\forall n \in \mathbb{N}^*\right)$.
			Suy ra $u_n$ bị chặn.
			\item  $u_n=\sin n$ do đó $-1 \leq u_n \leq 1\left(\forall n \in \mathbb{N}^*\right)$. Suy ra $u_n$ bị chặn.
			\item  Ta có $u_n=(-1)^{n-1} n^2>0$ nếu $\mathrm{n}$ là số tự nhiên lẻ.\\
			$u_n=(-1)^{n-1} n^2<0$ nếu $\mathrm{n}$ là số tự nhiên chẵn.\\
			Từ đó suy ra dây số không bị chặn trên, không bị chặn dưới và không bị chặn.
		\end{enumerate}
	}
\end{bt}

\begin{bt}%[1D2H1-2]
% nguon: Archive/SP_KhoLatex/ID6_old/1K2-5.tex #3404
%[Dự án 2 - TLDH - TeamTeXHoa - Đào Trung Kiên]%[ID6]
	Viết số hạng tổng quát của dãy số tăng gồm tất cả các số nguyên dương mà mỗi số hạng của nó
	\begin{enumerate}
		\item Đều chia hết cho $3$;
		\item Khi chia cho $4$ dư $1$.
	\end{enumerate}
	\loigiai{
		\begin{enumerate}
			\item $u_n=3n, \forall n \in \mathbb{N}^*$
			\item  $u_n=4n+1, \forall n \in \mathbb{N}^*$.
		\end{enumerate}
	}
\end{bt}

\begin{bt}%[1D2H2-4]
% nguon: Archive/SP_KhoLatex/ID6_old/1-GK1-CanhDieu-De1-NH23-24.tex #3407
%[0-GK1-NH22-23--TeamTeXHoa--Đào Trung Kiên]
		Cho cấp số cộng $\left(u_n\right)$ có $u_{10}=22; u_1+u_6=5$. Tìm số hạng thứ $5$ của cấp số cộng.
		\loigiai{Gọi số hạng đầu, công sai của cấp số cộng là $u_1$ và $d$.\\
			Ta có $\left\{\begin{array}{l}u_{10}=22\\ u_1+u_6=5\end{array}\Leftrightarrow\left\{\begin{array}{l}u_1+9 d=22\\ u_1+u_1+5 d=5\end{array}\Leftrightarrow\left\{\begin{array}{l}u_1+9 d=22\\ 2 u_1+5 d=5\end{array}\Leftrightarrow\left\{\begin{array}{l}u_1=-5\\ d=3.\end{array}\right.\right.\right.\right.$\\
			Số hạng thứ $5$ của cấp số cộng là $u_5=u_1+4d=-5+4\cdot 3=7$.}
	\end{bt}

\begin{bt}%[1D2V3-8]
% nguon: Archive/SP_KhoLatex/ID6_old/1-GK1-CanhDieu-De1-NH23-24.tex #3408
%[0-GK1-NH22-23--TeamTeXHoa--Đào Trung Kiên]
		Bạn Ngọc thả một quả bóng cao su từ độ cao $20$(m) so với mặt đất, mỗi lần chạm đất quả bóng lại nảy lên một độ cao bằng $\dfrac{4}{5}$ độ cao lần rơi trước. Biết rằng quả bóng luôn chuyển động vuông góc với mặt đất. Tính tổng quãng đường quả bóng đã di chuyển được (từ lúc thả bóng cho đến lúc bóng không nảy nữa).
		\loigiai{
			Ta có quãng đường bóng bay bằng tổng quãng đường bóng nảy lên và quãng đường bóng rơi xuống.\\
			Vì mỗi lần bóng nảy lên bằng $\dfrac{4}{5}$ lần nảy trước nên ta có tổng quãng đường bóng nảy lên là\\
			$$
			S_1=20\cdot\dfrac{4}{5}+20\cdot\left(\dfrac{4}{5}\right)^2+20\cdot\left(\dfrac{4}{5}\right)^3+\ldots+20\cdot\left(\dfrac{4}{5}\right)^n+\ldots
			$$\\
			Đây là tổng của cấp số nhân lùi vô hạn có số hạng đầu $u_1=20\cdot\dfrac{4}{5}=16$ và công bội $q=\dfrac{4}{5}$.\\
			Suy ra $S_1=\dfrac{16}{1-\tfrac{4}{5}}=80$.\\
			Tổng quãng đường bóng rơi xuống bằng khoảng cách độ cao ban đầu và tổng quãng đường bóng nảy lên nên là $S_2=20+20\cdot\left(\dfrac{4}{5}\right)+20\cdot\left(\dfrac{4}{5}\right)^2+\ldots+20\cdot\left(\dfrac{4}{5}\right)^n+\ldots$\\
			Đây là tổng của cấp số nhân lùi vô hạn với số hạng đầu $v_1=20$ và công bội $q=\dfrac{4}{5}$.\\
			Suy ra $S_2=\dfrac{20}{1-\tfrac{4}{5}}=100$.\\
			Vậy tổng quãng đường bóng bay là $S=S_1+S_2=180$(m).}
	\end{bt}

\begin{bt}%[1D2V3-2]
% nguon: Archive/SP_KhoLatex/ID6_old/1K2-7.tex #3420
%[Dự án 2 - TLDH - TeamTeXHoa - Le Hoang Anh]
	Cho dãy số $u_{n}$ xác định bởi $\heva{&u_{1}=2 \\& u_{n+1}=4 u_{n}+9, \quad \left(n \geq 1 \right).}$
	\begin{enumerate}
		\item Chứng minh dãy số $\left( v_n \right)$ với $v_{n}=u_{n}+3, \, n \geq 1$ là một cấp số nhân.
		\item Tìm công thức tổng quát của dãy số $u_n$.
	\end{enumerate}
	\loigiai{
		\begin{enumerate}
			\item Ta có $v_n=u_n+3$, suy ra $v_{n+1}=u_{n+1}+3= \left(4 u_n+9\right) +3$.\\
			Do đó $\dfrac{v_{n+1}}{v_n}=\dfrac{\left(4u_n+9\right) +3}{u_n+3}=\dfrac{4\left(u_n+3\right)}{u_n+3}=4$.\\
			Vậy $\left(v_{n}\right)$ là cấp số nhân với số hạng đầu $v_{1}=u_{1}+3=2+3=5$ và công bội $q=4$.
			\item  Do $\left(v_{n}\right)$ là cấp số nhân với $\heva{&v_{1}=5 \\& q=4}$ nên số hạng tổng quát của $v_{n}=v_{1} \cdot q^{n-1}=5 \cdot 4^{n-1}$.\\
			Suy ra công thức tổng quát của dãy số $\left(u_{n}\right)$ là $u_{n}=v_{n}-3=5 \cdot 4^{n-1}-3$.
		\end{enumerate}
	}
\end{bt}

\begin{bt}%[1D2V3-2]
% nguon: Archive/SP_KhoLatex/ID6_old/1K2-7.tex #3421
%[Dự án 2 - TLDH - TeamTeXHoa - Le Hoang Anh]
	Xét trong các dãy số số sau, dãy số nào là cấp số nhân, (nếu có) tìm công bội của cấp số nhân đó.
	\begin{multicols}{2}
		\begin{enumerate}
			\item $u_{n}=(-3)^{2 n+1}$.
			\item $u_{n}=(-1)^{n} \cdot 5^{3 n+2}$.
			\item $\heva{& u_{1}=2 \\ & u_{n+1}=u_{n}^{2}.}$
			\item $\heva{& u_{1}=3 \\& u_{n+1}=\dfrac{9}{u_n}.}$
		\end{enumerate}
	\end{multicols}
	\loigiai{
		\begin{enumerate}
			\item Ta có $\dfrac{u_{n+1}}{u_{n}}=\dfrac{(-3)^{2 n+3}}{(-3)^{2 n+1}}=(-3)^{2}=9$ (không đổi). Kết luận $\left(u_{n}\right)$ là cấp số nhân với công bội $q=9$.
			\item Ta có $\dfrac{u_{n+1}}{u_{n}}=\dfrac{(-1)^{n+1} \cdot 5^{3(n+1)+2}}{(-1)^{n} \cdot 5^{3 n+2}}=-1 \cdot 5^{3}=-125$ (không đổi). Kết luận $\left(u_{n}\right)$ là cấp số nhân với công bội $q=-125$.
			\item Ta có $u_2=u_{1}^{2}=4$, $u_3=u_{2}^{2}=16$, $u_4=u_{3}^{2}=256$, suy ra $\dfrac{u_2}{u_1}=\dfrac{4}{2}=2$ và $\dfrac{u_4}{u_3}=\dfrac{256}{16}=16$\\
			$\Rightarrow \dfrac{u_2}{u_1} \neq \dfrac{u_4}{u_3}$. Do đó $\left(u_n\right)$ không là cấp số nhân.
			\item $\dfrac{u_{n+1}}{u_n}=\dfrac{\frac{9}{u_n}}{\frac{9}{u_{n-1}}}=\dfrac{u_{n-1}}{u_n} \Rightarrow u_{n+1}=u_{n-1}, \forall n \geq 2$. Do đó có\\
			$u_1=u_3=u_5=\cdots=u_{2 n+1} \cdots$ (1)
			
			Và $u_2=u_4=u_6=\cdots=u_{2 n}=\cdots$ (2)
			
			Theo đề bài có $u_1=3 \Rightarrow u_2=\dfrac{9}{u_1}=3$. (3)
			
			Từ (1), (2), (3) suy ra $u_1=u_2=u_3=u_4=u_5=\cdots=u_{2n}=u_{2n+1} \cdots$ \\
			Kết luận $\left(u_n\right)$ là cấp số nhân với công bội $q=1$.
		\end{enumerate}
	}
\end{bt}

\begin{bt}%[1D2V3-5]
% nguon: Archive/SP_KhoLatex/ID6_old/1K2-7.tex #3423
%[Dự án 2 - TLDH - TeamTeXHoa - Le Hoang Anh]
	Cho $x$, $3$, $y$ theo thứ tự lập thành cấp số nhân và $x^{4}=y \sqrt{3}$. Tìm $x$, $y$.
	\loigiai{
		Ta có $x \cdot y=9 \Rightarrow y=\dfrac{9}{x}$.\\
		Thay vào $x^{4}=y \sqrt{3} \Leftrightarrow x^{4}=\dfrac{9}{x} \sqrt{3} \Leftrightarrow x^{5}=(\sqrt{3})^{4} \cdot \sqrt{3} \Leftrightarrow x^{5}=(\sqrt{3})^{5} \Leftrightarrow x=\sqrt{3}$\\
		$\Rightarrow y=\dfrac{9}{\sqrt{3}}=3 \sqrt{3}$. Kết luận $\heva{&x=\sqrt{3} \\& y=3 \sqrt{3}.}$
	}
\end{bt}

\begin{bt}%[1D2V3-4]
% nguon: Archive/SP_KhoLatex/ID6_old/1K2-7.tex #3424
%[Dự án 2 - TLDH - TeamTeXHoa - Le Hoang Anh]
	Tìm số hạng đầu và công bội của cấp số nhân, biết
	\begin{multicols}{3}
		\begin{enumerate}
			\item $ \heva{& u_1+u_5=51 \\& u_2+u_6=102;} $
			\item $ \heva{&u_1+u_2+u_3=135 \\& u_4+u_5+u_6=40;} $
			\item $ \heva{&u_2=6 \\& S_3=43.} $
		\end{enumerate}
	\end{multicols}
	\loigiai{
		\begin{enumerate}
			\item $ \heva{&u_{1}+u_{5}=51 \\& u_{2}+u_{6}=102} $ $ \Leftrightarrow $ 
			$ \heva{& u_{1}+u_{1} q^{4}=51 \\& u_{1} q+u_{1} q^{5}=102} $ $ \Leftrightarrow $ 
			$ \heva{&u_{1}\left(1+q^{4}\right)=51 & (*) \\& u_{1} q\left(1+q^{4}\right)=102. & (* *) }$
			
			Lấy $\dfrac{(* *)}{(*)} \Leftrightarrow \dfrac{u_{1} q\left(1+q^{4}\right)}{u_{1}\left(1+q^{4}\right)}=\dfrac{102}{51} \Leftrightarrow q=2 \Rightarrow u_{1}=\dfrac{51}{1+q^{4}}=\dfrac{51}{17}=3$.\\
			\item $\heva{&u_{1}+u_{2}+u_{3}=135 \\& u_{4}+u_{5}+u_{6}=40} $ $ \Leftrightarrow $
			$ \heva{& u_{1}+u_{1} q+u_{1} q^{2}=135 \\& u_{1} q^{3}+u_{1} q^{4}+u_{1} q^{5}=40} $ $ \Leftrightarrow $
			$ \heva{&u_{1}\left(1+q+q^{2}\right)=135 & (*) \\& u_{1} q^{3}\left(1+q+q^{2}\right)=40. & (* *)} $
			
			Lấy $\dfrac{(* *)}{(*)} \Leftrightarrow \dfrac{u_{1} q^{3}\left(1+q+q^{2}\right)}{u_{1}\left(1+q+q^{2}\right)}=\dfrac{40}{135} \Leftrightarrow q^{3}=\dfrac{8}{27} 
			\Leftrightarrow q=\dfrac{2}{3}$\\
			$\Rightarrow u_{1}=\dfrac{135}{1+q+q^{2}}=\dfrac{1215}{19}$.\\
			Kết luận có công bội $q=\dfrac{2}{3}$ và số hạng đầu tiên $u_{1}=\dfrac{1215}{19}$.
			\item $ \heva{&u_{2}=6 \\& S_{3}=43} $ $ \Leftrightarrow $
			$ \heva{&u_{1} q=6 \\& u_{1}+u_{2}+u_{3}=43} $ $ \Leftrightarrow $
			$ \heva{&u_{1} q=6 \\& u_{1}+u_{1} q+u_{1} q^{2}=43} $ $ \Leftrightarrow $
			$ \heva{&u_{1} q=6 &(*) \\& u_{1}\left(1+q+q^{2}\right)=43. &(* *)} $
			
			Lấy $\dfrac{(*)}{(* *)} \Leftrightarrow \dfrac{u_{1} q}{u_{1}\left(1+q+q^{2}\right)}=\dfrac{6}{43}$
			$\Leftrightarrow 43 q=6\left(1+q+q^{2}\right) \Leftrightarrow 6 q^{2}-37 q+6=0 \Leftrightarrow \hoac{&q=6 \\& q=\dfrac{1}{6}.}$\\
			Với $q=6 \Rightarrow u_{1}=1$.\\
			Với $q=\dfrac{1}{6} \Rightarrow u_{1}=36$.\\
			Kết luận $\heva{&q=6 \\& u_{1}=1}$ hoặc $\heva{&q=\frac{1}{6} \\& u_{1}=36.}$
			
		\end{enumerate}
	}
\end{bt}

\begin{bt}%[1D2H3-2]
% nguon: Archive/SP_KhoLatex/ID6_old/1K2-7.tex #3427
,[1K2K7-3]%[Dự án 2 - TLDH - TeamTeXHoa - Thinh Ngo]
	Tìm số hạng đầu và công bội của cấp số nhân, biết: 
	\begin{enumEX}{2}
		\item $\heva{& u_4-u_2=72 \\ & u_5-u_3=144;}$
		\item $\heva{& u_1-u_3+u_5=65 \\ &u_1+u_7=325;}$
		\item $\heva{& u_3+u_5=90 \\ & u_2-u_6=240;}$
		\item $\heva{& u_1+u_2+u_3=14 \\ & u_1\cdot u_2\cdot u_3=64;}$
		\item $\heva{& u_1+u_2+u_3=21 \\ & \dfrac{1}{u_1}+\dfrac{1}{u_2}+\dfrac{1}{u_3}=\dfrac{7}{12};}$
		\item $\heva{& u_1+u_2+u_3+u_4=30 \\ & u_1^2+u_2^2+u_3^2+u_4^2=340.}$
	\end{enumEX}	
	\loigiai{
		\begin{enumerate}
			\item $\heva{& u_4-u_2=72 \\ & u_5-u_3=144} \Leftrightarrow \heva{& u_1q^3-u_1q=72 \\ & u_1q^4-u_1q^2=144} \Leftrightarrow \heva{& u_1q\left(q^2-1\right)=72 & (1) \\ &u_1q^2\left(q^2-1\right)=144. & (2)}$\\
			Lấy $\dfrac{(2)}{(1)}$ ta được $q=2$, thay $q=2$ vào $(1)$ ta có $u_1=12$.
			\item $\heva{& u_1-u_3+u_5=65 \\ &u_1+u_7=325;} \Leftrightarrow \heva{& u_1- u_1q^2+u_1q^4=65 \\ & u_1+u_1q^6=325} \Leftrightarrow \heva{& u_1\left(1-q^2+q^4\right)=65 & (1) \\ &u_1\left(1+q^6\right)=325. & (2)}$\\
			Lấy $\dfrac{(2)}{(1)}$ ta được $\dfrac{1+q^6}{1-q^2+q^4}=5 \Leftrightarrow q^6 - 5q^4 + 5q^2 - 4 = 0 \Leftrightarrow q^2 = 4 \Leftrightarrow q = \pm 2$, thay $q=\pm2$ vào $(2)$ ta có $u_1=5$.
			\item $\heva{& u_3+u_5=90 \\ & u_2-u_6=240} \Leftrightarrow \heva{& u_1q^2+u_1q^4=90 \\ & u_1q-u_1q^5=240} \Leftrightarrow \heva{& u_1q^2\left(1+q^2\right)=90 & (1) \\ & u_1q\left(1-q^4\right)=240. &(2)}$ \\
			Lấy $ \dfrac{(2)}{(1)}$ ta được
			\begin{eqnarray*}
				\dfrac{u_1q\left(1-q^4\right)}{u_1q^2\left(1+q^2\right)}=\dfrac{240}{90}&\Leftrightarrow&\dfrac{\left(1-q^2\right)\left(1+q^2\right)}{q\left(1+q^2\right)}=\dfrac{8}{3}\Leftrightarrow\dfrac{1-q^2}{q}=\dfrac{8}{3}\\
				&\Leftrightarrow& 3q^2+8q-3=0\Leftrightarrow q=\dfrac{1}{3}\vee q=-3.
			\end{eqnarray*}
			\begin{itemize}
				\item Với $q=\dfrac{1}{3}$ thay vào $(1)$ ta có $u_1=729$.
				\item Với $q=-3$ thay vào $(1)$ ta có $u_1=1$.
			\end{itemize}
			\item $\heva{& u_1+u_2+u_3=14 \\ & u_1\cdot u_2\cdot u_3=64} \Leftrightarrow \heva{& u_1+u_1q+u_1q^2=14 \\ & u_1\cdot u_1\cdot q\cdot u_1\cdot q^2=64} \Leftrightarrow \heva{& u_1\left(1+q+q^2\right)=14 &(1) \\ & \left(u_1q\right)^3=64. & (2)} $\\
			$(2)\Leftrightarrow{u_1}q=4\Rightarrow{u_1}=\dfrac{4}{q}$ , thay vào (1) được\\ $$\dfrac{4}{q}\left(1+q+q^2\right)=14
			\Leftrightarrow 2q^2-5q+2=0\Leftrightarrow q=2\vee q=\dfrac{1}{2}.$$
			\begin{itemize}
				\item Với $q=2\Rightarrow u_1=2$.
				\item Với $q=\dfrac{1}{2}\Rightarrow u_1=8$.
			\end{itemize}
			\item $\heva{& u_1+u_2+u_3=21 \\ & \dfrac{1}{u_1}+\dfrac{1}{u_2}+\dfrac{1}{u_3}=\dfrac{7}{12};} \Leftrightarrow \heva{& u_1+u_1q+u_1q^2=21 \\ & \dfrac{1}{u_1}+\dfrac{1}{u_1q}+\dfrac{1}{u_1q^2}=\dfrac{7}{12}} \Leftrightarrow \heva{& u_1\left(1+q+q^2\right)=21 &(1) \\ & \dfrac{q^2+q+1}{u_1q^2}=\dfrac{7}{12}. & (2)} $\\
			$(1)\Rightarrow 1+q+q^2=\dfrac{21}{u_1}$ , thay vào (2) ta được $$\dfrac{21}{u_1}\cdot\dfrac{1}{u_1q^2}=\dfrac{7}{12}\Leftrightarrow{\left(u_1q\right)^2}=36\Leftrightarrow u_1q=\pm 6\Leftrightarrow u_1=\pm \dfrac{6}{q}.$$
			\begin{itemize}
				\item Với $u_1=\dfrac{6}{q}$ thay vào (1) ta có \\ $\dfrac{6}{q}\left(1+q+q^2\right)=21\Leftrightarrow 2q^2-5q+2=0\Leftrightarrow q=2 \vee q=\dfrac{1}{2}$.
				\begin{itemize}
					\item Nếu $q=2\Rightarrow{u_1}=3$.
					\item Nếu $q=\dfrac{1}{2}\Rightarrow{u_1}=12$.
				\end{itemize}
				\item Với $u_1=-\dfrac{6}{q}$ thay vào (1) ta có \\ $-\dfrac{6}{q}\left(1+q+q^2\right)=21\Leftrightarrow 2q^2+9q+2=0\Leftrightarrow q=\dfrac{-9+\sqrt{65}}{4}\vee q=\dfrac{-9-\sqrt{65}}{4}$.
				\begin{itemize}
					\item Nếu $q=\dfrac{-9+\sqrt{65}}{4}\Rightarrow{u_1}=\dfrac{27+3\sqrt{65}}{2}$.
					\item Nếu $q=\dfrac{-9-\sqrt{65}}{4}\Rightarrow{u_1}=\dfrac{27-3\sqrt{65}}{2}$. 
				\end{itemize}
				
			\end{itemize}
			\item \begin{eqnarray*}
				\heva{& u_1+u_2+u_3+u_4=30 \\ & u_1^2+u_2^2+u_3^2+u_4^2=340} &\Leftrightarrow& \heva{& u_1+u_1q+u_1q^2+u_1q^3=30 \\ & u_1^2+u_1^2q^2+u_1^2q^4+u_1^2q^6=340}\\& \Leftrightarrow& \heva{& u_1\left(1+q+q^2+q^3\right)=30 \\ & u_1^2\left(1+q^2+q^4+q^6\right)=340}\\ &\Leftrightarrow & \heva{& u_1\left(1+q\right)\left(1+q^2\right)=30 \\ & u_1^2\left(1+q^4\right)\left(1+q^2\right)=340}\\ &\Leftrightarrow& \heva{& u_1^2\left(1+q\right)^2\left(1+q^2\right)^2=900 & (1) \\ & u_1^2\left(1+q^4\right)\left(1+q^2\right)=340. & (2)}
			\end{eqnarray*}
			Lấy $\dfrac{(1)}{(2)}\Leftrightarrow\dfrac{\left(1+q\right)^2\left(1+q^2\right)}{1+q^4}=\dfrac{45}{17}$, quy đồng rút gọn ta được \\ $14q^4-17q^3-17q^2-17q+14=0$ $\Leftrightarrow 14q^2-17q-17-\dfrac{17}{q}+\dfrac{14}{q^2}=0$\\
			$\Leftrightarrow 14\left(q^2+\dfrac{1}{q^2}\right)-17\left(q+\dfrac{1}{q}\right)-45=0$. \quad $(3)$\\
			Đặt $t=q+\dfrac{1}{q}$, điều kiện $\left| t\right|\ge 2$.\\
			$(3)\Leftrightarrow 14t^2-17t-45=0\Leftrightarrow t=\dfrac{5}{2}\vee t=-\dfrac{9}{7}~ \text{(loại)}$.\\
			Khi $t=\dfrac{5}{2}\Rightarrow q+\dfrac{1}{q}=\dfrac{5}{2}\Leftrightarrow 2q^2-5q+2=0\Leftrightarrow q=2\vee q=\dfrac{1}{2}$.
			\begin{itemize}
				\item Với $q=2\Rightarrow u_1=2$. 
				\item Với $q=\dfrac{1}{2}\Rightarrow u_1=16$.
			\end{itemize}
		\end{enumerate}
	}
	
\end{bt}

\begin{bt}%[1D2V3-7]
% nguon: Archive/SP_KhoLatex/ID6_old/1K2-7.tex #3428
%[Dự án 2 - TLDH - TeamTeXHoa - Thinh Ngo]
	Tìm $a, b$ biết rằng $1, a, b$ là ba số hạng liên tiếp của cấp số cộng và $1, a^2, b^2$ là ba số hạng liên tiếp của một cấp số nhân.
	\loigiai{
		Theo đề bài ta có hệ phương trình $\heva{& 1+b=2a\\ &{b^2}=a^4\\} \Leftrightarrow \heva{& 1+b=2a &(1)\\ 
			& b=\pm{a^2}.}$
		\begin{itemize}
			\item Với $b=a^2$ thay vào $(1)$ ta được \\ $1+a^2=2a\Leftrightarrow{a^2}-2a+1=0\Leftrightarrow a=1\Rightarrow b=1$.
			\item Với $b=-a^2$ thay vào $(1)$ ta được \\ $1-a^2=2a\Leftrightarrow{a^2}+2a-1=0\Leftrightarrow a=-1+\sqrt{2} ~ \vee a=-1-\sqrt{2}$.
			\begin{itemize}
				\item Khi $a=-1+\sqrt{2}\Rightarrow b=-\left(-1+\sqrt{2}\right)^2\Leftrightarrow b=-3+2\sqrt{2}$.
				\item Khi $a=-1-\sqrt{2}\Rightarrow b=-\left(-1-\sqrt{2}\right)^2\Leftrightarrow b=-3-2\sqrt{2}$.
			\end{itemize}
			Kết luận $\heva{& a=1\\& b=1} \vee \heva{& a=-1+\sqrt{2}\\	& b=-3+2\sqrt{2}} \vee \heva{		& a=-1-\sqrt{2}\\& b=-3-2\sqrt{2}}$ thỏa yêu cầu đề bài.
		\end{itemize}
	}
	
\end{bt}

\begin{bt}%[1D2V3-4]
% nguon: Archive/SP_KhoLatex/ID6_old/1K2-7.tex #3429
%[Dự án 2 - TLDH - TeamTeXHoa - Thinh Ngo]
	Tìm số hạng đầu của cấp số nhân biết công bội bằng $3$, tổng số các số hạng là $728$ và số hạng cuối bằng $486$.
	\loigiai{
		Theo đề bài ta có $\heva{&S_n=728\\ 
			&u_n=486} \Leftrightarrow \heva{
			&\dfrac{u_1\left(1-q^n\right)}{1-q}=728\\ 
			&u_1\cdot q^{n-1}=486} \Leftrightarrow\heva{&u_1-u_1q^n=728(1-q)\\ &u_1q^n=486q}\\
		\Rightarrow u_1-486q=728(1-q)\Leftrightarrow u_1=2$.
	}
	
\end{bt}

\begin{bt}%[1D2V3-7]
% nguon: Archive/SP_KhoLatex/ID6_old/1K2-7.tex #3430
%[Dự án 2 - TLDH - TeamTeXHoa - Thinh Ngo]
	Cho $3$ số tạo thành một cấp số cộng có tổng $21$. Nếu thêm $2, 3, 9$ lần lượt vào số thứ nhất, số thứ hai, số thứ ba tạo thành một cấp số nhân. Tìm $3$ số đó.
	\loigiai{
		Gọi $u_1,u_2,u_3$ lập thành cấp số cộng.\\
		Theo đề bài $u_1+2; u_2+3; u_3+9$ là ba số liên tiếp tạo thành cấp số nhân.\\
		Theo đề bài ta có 
		\begin{eqnarray*}
			\heva{&u_1+u_2+u_3=21\\ 
				&u_1+u_3=2u_2\\ 
				&\left(u_1+2\right)\left(u_3+9\right)=\left(u_2+3\right)^2} &\Leftrightarrow&\heva{& 3u_2=21\\ 
				&u_1+u_3=2u_2\\ 
				&\left(u_1+2\right)\left(u_3+9\right)=\left(u_2+3\right)^2} \\
			&\Leftrightarrow&\heva{
				&u_2=7\\ 
				&u_1=14-u_3\\ 
				&\left(14-u_3+2\right)\left(u_3+9\right)=100. \quad(*)}
		\end{eqnarray*}
		Giải $(*)$ $\left(16-u_3\right)\left(u_3+9\right)=100 \Leftrightarrow-u_3^2+7u_3+44=0 \Leftrightarrow u_3=11 \vee u_3=-4$.
		\begin{itemize}
			\item Với $u_3=11\Rightarrow u_1=3$.
			\item Với $u_3=-4\Rightarrow u_1=18$.
		\end{itemize}
	}	
\end{bt}

\begin{bt}%[1D2V3-7]
% nguon: Archive/SP_KhoLatex/ID6_old/1K2-7.tex #3431
%[Dự án 2 - TLDH - TeamTeXHoa - Thinh Ngo]
	Cho $3$ số dương có tổng là $65$ lập thành một cấp số nhân tăng, nếu bớt một đơn vị ở số hạng thứ nhất và $19$ đơn vị ở số hạng thứ ba ta được một cấp số cộng. Tìm $3$ số đó.
	\loigiai{
		Gọi $u_1; u_2; u_3$ theo thứ tự đó lập thành một cấp số nhân.\\
		Theo đề bài ta có $u_1-1; u_2; u_3-19$ theo thứ tự đó lập thành một cấp số cộng.\\
		Ta có \begin{eqnarray*}
			\heva{&u_1+u_2+u_3=65\\ 
				&u_1-1+u_3-19=2u_2}
			&\Leftrightarrow&\heva{&u_1+u_2+u_3=65\\ 
				&u_1-2u_2+u_3=20}\\&\Leftrightarrow&\heva{&u_1+u_1\cdot q + u_1\cdot q^2=65\\&u_1-2u_1\cdot q + u_1\cdot q^2=20}\\ &\Leftrightarrow&\heva{&u_1\left(1+q+q^2\right)=65 &(1)\\ 
				&u_1\left(1-2q+q^2\right)=20. & (2)}
		\end{eqnarray*}
		\begin{eqnarray*}
			\text{Lấy}~ \dfrac{(1)}{(2)}&\Leftrightarrow&\dfrac{1+q+q^2}{1-2q+q^2}=\dfrac{65}{20}=\dfrac{13}{4} \Leftrightarrow 4\left(1+q+q^2\right)=13\left(1-2q+q^2\right)\\
			&\Leftrightarrow& 9q^2-30q+9=0 \Leftrightarrow q=3 \vee q=\dfrac{1}{3}.
		\end{eqnarray*}
		Vì $u_1; u_2; u_3$ theo thứ tự lập thành cấp số nhân tăng dần nên chọn $q=3\Rightarrow u_1=5$.\\
		Vậy $u_1=5; u_2=15; u_3=45$.
	}
	
	
\end{bt}

\begin{bt}%[1D2V3-4]
% nguon: Archive/SP_KhoLatex/ID6_old/1K2-7.tex #3432
%[Dự án 2 - TLDH - TeamTeXHoa - Thinh Ngo]
	Tìm $3$ số hạng liên tiếp của một cấp số nhân biết tổng của chúng là $19$ và tích là $216$.
	\loigiai{
		Gọi ba số hạng liên tiếp của cấp số nhân là $u_1,u_2,u_3$ với công bội là $q$. Theo đề bài ta có hệ phương trình $\heva{&u_1+u_2+u_3=19\\&u_1\cdot u_2\cdot u_3=216} \Leftrightarrow\heva{&u_1+u_1q+u_1q^2=19\\	&{\left(u_1q\right)^3}=6^3} 
		\Leftrightarrow\heva{&u_1\left(1+q+q^2\right)=19 \quad(*)\\& u_1q=6\Rightarrow u_1=\dfrac{6}{q}.}$\\
		Thay $u_1=\dfrac{6}{q}$ vào $(*)$ ta được $6q^2-13q+6=0\Leftrightarrow q=\dfrac{3}{2} \vee q=\dfrac{2}{3}$.
		\begin{itemize}
			\item Với $q=\dfrac{3}{2}\Rightarrow u_1=4,u_2=6,u_3=9$.
			\item Với $q=\dfrac{2}{3}\Rightarrow u_1=9,u_2=6,u_3=4$.
		\end{itemize}
		
	}
\end{bt}

\begin{bt}%[1D2V3-2]
% nguon: Archive/SP_KhoLatex/ID6_old/1K2-7.tex #3433
%[Dự án 2 - TLDH - TeamTeXHoa - Thinh Ngo]
	Tìm công bội của một cấp số nhân có số hạng đầu là $7$, số hạng cuối là $448$ và tổng số các số hạng là $889$.
	\loigiai{
		Theo đề bài ta có $\left\{\begin{aligned}
			&{S_n}=889\\ 
			&{u_n}=448\\ 
		\end{aligned}\right.\Leftrightarrow\left\{\begin{aligned}
			&\dfrac{u_1\left(q^n-1\right)}{q-1}=889\\ 
			&{u_1}{q^{n-1}}=448\\ 
		\end{aligned}\right.\Leftrightarrow\left\{\begin{aligned}
			&{u_1}{q^n}-u_1=889(q-1)&\quad(1)\\ 
			&{u_1}{q^n}=448q. &\quad (2)\\ 
		\end{aligned}\right.$\\
		Thay $(2)$ vào $(1)$ ta được $448q-7=889q-889\Leftrightarrow q=2$.	
	}
	
\end{bt}

\begin{bt}%[1D2V3-4]
% nguon: Archive/SP_KhoLatex/ID6_old/1K2-7.tex #3434
%[Dự án 2 - TLDH - TeamTeXHoa - Thinh Ngo]
	Tìm bốn số hạng liên tiếp của một cấp số nhân, trong đó số hạng thứ hai nhỏ hơn số hạng thứ nhất $35$, còn số hạng thứ ba lớn hơn số hạng thứ tư $560$.
	\loigiai{
		Theo đề bài ta có  $\left\{\begin{aligned}
			&{u_1}-u_2=35\\ 
			&{u_3}-u_4=560\\ 
		\end{aligned}\right.\Leftrightarrow\left\{\begin{aligned}
			&{u_1}-u_1q=35\\ 
			&{u_1}{q^2}-u_1q^3=560\\ 
		\end{aligned}\right.\Leftrightarrow\left\{\begin{aligned}
			&{u_1}(1-q)=35 &\quad (1)\\ 
			&{u_1}{q^2}(1-q)=560. & \quad (2)\\ 
		\end{aligned}\right.$\\
		Thay $(1)$ vào $(2)$ ta được $q^2=16\Leftrightarrow q=\pm 4$.
		\begin{itemize}
			\item Với $q=4$ thay vào $(1)$ ta được \\
			$u_1=-\dfrac{35}{3}, u_2=u_1q=-\dfrac{140}{3}, u_3=-\dfrac{560}{3}, u_4=-\dfrac{2240}{3}$.
			\item Với $q=-4$ thay vào $(1)$ ta được \\
			$u_1=7, u_2=u_1q=-28, u_3=112, u_4=-448$.
		\end{itemize}
		
	}
	
\end{bt}

\begin{bt}%[1D2V3-7]
% nguon: Archive/SP_KhoLatex/ID6_old/1K2-7.tex #3435
%[Dự án 2 - TLDH - TeamTeXHoa - Thinh Ngo]
	Tìm các số dương $a$ và $b$ sao cho $a, a+2b, 2a+b$ lập thành một cấp số cộng và $(b+1)^2, ab+5, (a+1)^2$ lập thành một cấp số nhân.
	\loigiai{
		Theo tính chất của cấp số cộng ta có $a+(2a+b)=2(a+2b)$. \hfill $(1)$\\
		Theo tính chất của cấp số nhân ta có $(b+1)^2\cdot(a+1)^2=(ab+5)^2$. \hfill $(2)$\\
		Từ $(1)$ khai triển rút gọn ta được $a=3b$, thay vào $(2)$ ta có\\
		$(b+1)^2(3b+1)^2=(3b^2+5)^2\Leftrightarrow (b+1)(3b+1)=\pm (3b^2+5)$.
		\begin{itemize}
			\item Với $\left(b+1\right)\left(3b+1\right)=3b^2+5\Leftrightarrow b=1\Rightarrow a=3$.
			\item Với $\left(b+1\right)\left(3b+1\right)=-3b^2-5\Leftrightarrow 6b^2+4b+6=0$ (vô nghiệm).
		\end{itemize}
		Kết luận $a=3, b=1$.
	}
	
\end{bt}

\begin{bt}%[1D2V3-6]
% nguon: Archive/SP_KhoLatex/ID6_old/1K2-7.tex #3436
%[Dự án 2 - TLDH - TeamTeXHoa - Thinh Ngo]
	Tính các tổng sau:
	\begin{enumerate}
		\item $S_n=2+2^2+2^3+\cdot\cdot\cdot+2^n$;
		\item $S_n=\dfrac{1}{2}+\dfrac{1}{2^2}+\dfrac{1}{2^3}+\cdot\cdot\cdot+\dfrac{1}{2^n}$;
		\item $S_n=\left(3+\dfrac{1}{3}\right)^2+\left(9+\dfrac{1}{9}\right)^2+\cdot\cdot\cdot+\left(3^n+\dfrac{1}{3^n}\right)^2$;
		\item $S_n=6+66+666+\cdot\cdot\cdot+\underbrace{666...6}_{n~\text{số}~ 6}$.
	\end{enumerate}
	\loigiai{
		\begin{enumerate}
			\item Ta có dãy số $2, 2^2, 2^3, \cdot\cdot\cdot ,2^n$ là một cấp số nhân với $n$ số hạng, có số hạng đầu $u_1=2$ và công bội $q=\dfrac{2^2}{2}=2$. Do đó $S_n=u_1\cdot\dfrac{1-q^n}{1-q}=2\cdot\dfrac{1-2^n}{1-2}=2\left(2^n-1\right)$.
			\item Ta có dãy số $\dfrac{1}{2}, \dfrac{1}{2^2}, \dfrac{1}{2^3}, \cdot\cdot\cdot, \dfrac{1}{2^n}$ là một cấp số nhân với $n$ số hạng, có số hạng đầu $u_1=\dfrac{1}{2}$ và công bội $q=\dfrac{\dfrac{1}{2^2}}{\dfrac{1}{2}}=\dfrac{1}{2}$. Do đó $S_n=u_1\cdot\dfrac{1-q^n}{1-q}=\dfrac{1}{2}\cdot\dfrac{1-\left(\dfrac{1}{2}\right)^n}{1-\dfrac{1}{2}}=1-\dfrac{1}{2^n}$.
			\item \begin{eqnarray*} S_n&=&\left(3+\dfrac{1}{3}\right)^2+\left(9+\dfrac{1}{9}\right)^2+\cdot\cdot\cdot+\left(3^n+\dfrac{1}{3^n}\right)^2\\
				&=&3^2+2+\dfrac{1}{3^2}+3^4+2+\dfrac{1}{3^4}+\cdot\cdot\cdot+3^{2n}+2+\dfrac{1}{3^{2n}}\\
				&=&\left(3^2+3^4+\cdot\cdot\cdot+3^{2n}\right)+\left(\dfrac{1}{3^2}+\dfrac{1}{3^4}+\cdot\cdot\cdot+\dfrac{1}{3^{2n}}\right)+\underbrace{2+2+2+\cdot\cdot\cdot+2.}_n
			\end{eqnarray*}
			\begin{itemize}
				\item Có dãy số $3^2,3^4,\cdot\cdot\cdot ,3^{2n}$ là cấp số nhân với $n$ số hạng, có số hạng đầu $u_1=3^2$ và công bội $q=\dfrac{3^4}{3^2}=9$. Do đó $S_1=u_1\cdot\dfrac{1-q^n}{1-q}=9\cdot\dfrac{1-9^n}{1-9}=\dfrac{9}{8}\left(9^n-1\right)$.
				\item Có dãy số $\dfrac{1}{3^2},\dfrac{1}{3^4},\cdot\cdot\cdot ,\dfrac{1}{3^{2n}}$ là cấp số nhân với $n$ số hạng, có số hạng đầu $u_1=\dfrac{1}{3^2}$ và công bội $q=\dfrac{1}{9}$. Do đó $S_2=u_1\cdot\dfrac{1-q^n}{1-q}=\dfrac{1}{9}\cdot\dfrac{1-\dfrac{1}{9^n}}{1-\dfrac{1}{9}}=\dfrac{1}{8}\left(1-\dfrac{1}{9^n}\right)=\dfrac{9^n-1}{8\cdot9^n}$.
			\end{itemize}	
			Vậy $S_n=\dfrac{9}{8}\left(9^n-1\right)+\dfrac{9^n-1}{8\cdot9^n}+2n=\dfrac{\left(9^n-1\right)\left(9^{n+1}+1\right)}{8\cdot9^n}+2n$.
			\item \begin{eqnarray*} S_n&=&6+66+666+\cdot\cdot\cdot+\underbrace{666...6}_{n~\text{số}~ 6}=\dfrac{6}{9}\left(9+99+999+\cdot\cdot\cdot+\underbrace{999...9}_{n~\text{số}~ 9}\right) \\
				&=&\dfrac{2}{3}\left[(10-1)+(100-1)+(1000-1)+\cdot\cdot\cdot+(10^n-1)\right]\\
				&=&\dfrac{2}{3}\left[10+10^2+10^3+\cdot\cdot\cdot+10^n-n\right]\\
				&=&\dfrac{2}{3}\left[10\cdot\dfrac{10^n-1}{10-1}-n\right]\\
				&=&\dfrac{20}{27}\left(10^n-1\right)-\dfrac{2n}{3}.
			\end{eqnarray*}
		\end{enumerate}
	}	
\end{bt}

\begin{bt}%[1D2V3-6]
% nguon: Archive/SP_KhoLatex/ID6_old/1K2-7.tex #3437
%[Dự án 2 - TLDH - TeamTeXHoa- Phạm Hoàng Đăng]
	Với hình vuông $A_1B_1C_1D_1$ như hình vẽ bên, cách tô màu như phần gạch sọc được gọi là cách tô màu \lq\lq đẹp\rq\rq. Một nhà thiết kế tiến hành tô màu cho một hình vuông như hình bên, theo quy trình sau
	\begin{center}
		\begin{tikzpicture}[>=stealth,smooth,line join=round,line cap=round,font=\footnotesize,scale=1]
			\path
			(0,0) coordinate (A_1)
			(0:4) coordinate (B_1)
			(-90:4) coordinate (D_1)
			($(B_1)+(D_1)-(A_1)$) coordinate (C_1)
			($(A_1)!1/3!(C_1)$) coordinate (A_2)
			($(A_1)!2/3!(C_1)$) coordinate (C_2)
			($(B_1)!1/3!(D_1)$) coordinate (B_2)
			($(B_1)!2/3!(D_1)$) coordinate (D_2)
			($(A_1)!1/3!(B_1)$) coordinate (ab1)
			($(A_1)!2/3!(B_1)$) coordinate (ab2)
			($(A_1)!1/3!(D_1)$) coordinate (ad1)
			($(A_1)!2/3!(D_1)$) coordinate (ad2)
			($(B_1)!1/3!(C_1)$) coordinate (bc1)
			($(B_1)!2/3!(C_1)$) coordinate (bc2)
			($(D_1)!1/3!(C_1)$) coordinate (dc1)
			($(D_1)!2/3!(C_1)$) coordinate (dc2)
			;
			\fill[pattern=north east lines] (A_1)--(A_2)--(D_2)--(C_2)--(C_1)--(D_1)--cycle
			;
			\draw[thick] (A_1)--(B_1)--(C_1)--(D_1)--cycle
			(A_1)--(A_2)--(D_2)--(C_2)--(C_1)
			;
			\draw[dashed] (A_2)--(C_2)
			(ab1)--(dc1) (ab2)--(dc2) (ad1)--(bc1)(ad2)--(bc2)
			;
			\foreach \x/\g in {A_1/135,B_1/40,D_1/-130,C_1/-40,A_2/45,B_2/40,D_2/-130,C_2/40} \fill[black] (\x) circle (1pt) +(\g:3mm) node {$\x$};
		\end{tikzpicture}
	\end{center}
	\begin{itemize}
		\item \textbf{Bước 1.} Tô màu \lq\lq đẹp\rq\rq\, cho hình vuông $A_1B_1C_1D_1$.
		\item \textbf{Bước 2.} Tô màu \lq\lq đẹp\rq\rq\, cho hình vuông $A_2B_2C_2D_2$  là hình vuông ở chính giữa khi chia hình vuông $A_1B_1C_1D_1$ thành $9$ phần bằng nhau như hình vẽ.
		\item \textbf{Bước 3.} Tô màu \lq\lq đẹp\rq\rq\, cho hình vuông $A_3B_3C_3D_3$ là hình vuông ở chính giữa khi chia hình vuông $A_2B_2C_2D_2$ thành $9$ phần bằng nhau. Cứ tiếp tục như vậy. Hỏi cần ít nhất bao nhiêu bước để tổng diện tích phần được tô màu chiếm $49{,}99\%$.
	\end{itemize}
	\loigiai{
		Gọi diện tích được tô màu ở mỗi bước là $u_n$, $n\in\mathbb{N}^*$. Dễ thấy dãy các giá trị $u_n$ là một cấp số nhân với số hạng đầu $u_1=\dfrac{4}{9}$ và công bội $q=\dfrac{1}{9}$.\\
		Gọi $S_k$ là tổng của $k$ số hạng đầu trong cấp số nhân đang xét thì $S_k=\dfrac{u_1(q^k-1)}{q-1}$.\\
		Để tổng diện tích phần được tô màu chiếm $49{,}99\%$ thì $$S_k=\dfrac{u_1(q^k-1)}{q-1}\geq 0{,}4999\Leftrightarrow k\geq 3{,}8.$$
		Vậy cần ít nhất $4$ bước.
	} 
\end{bt}

\begin{bt}%[1D2V3-6]
% nguon: Archive/SP_KhoLatex/ID6_old/1K2-7.tex #3438
%[Dự án 2 - TLDH - TeamTeXHoa- Phạm Hoàng Đăng]
	Cho hình vuông $(C_1)$ có cạnh bằng $a$. Người ta chia mỗi cạnh của hình vuông thành bốn phần bằng nhau và nối các điểm chia một cách thích hợp để có hình vuông $(C_2)$ (Hình vẽ).
	\begin{center}
		\begin{tikzpicture}
			\def\g{20}
			\def\a{6}
			\foreach \i in {0,...,5}{
				\pgfmathsetmacro{\b}{\a /(sin(\g)+cos(\g))^(\i)}
				\pgfmathsetmacro{\quay}{\i *\g}
				\draw[rotate=\quay] (-.5*\b,-.5*\b) rectangle (.5*\b,.5*\b);
			}
		\end{tikzpicture}
	\end{center}
	Từ hình vuông $(C_2)$ lại tiếp tục làm như trên ta nhận được dãy các hình vuông $C_1,C_2,C_3,\ldots,C_n$. Gọi $S_i\,(i\in\{1;2;3;\ldots\})$ là diện tích của hình vuông $C_i$. Đặt $T=S_1+S_2+S_3+\cdots+S_n+\cdots$. Biết $T=\dfrac{32}{2}$, tính $a$?
	\loigiai{
		Cạnh của hình vuông $C_2$ là $a_2=\sqrt{\left(\dfrac{3}{4}a\right)^2+\left(\dfrac{1}{4}a\right)^2}=\dfrac{a\sqrt{10}}{4}$.\\ 
		Do đó diện tích $S_2=\dfrac{5}{8}a^2=\dfrac{5}{8}S$.\\
		Cạnh của hình vuông $C_3$ là $a_3=\sqrt{\left(\dfrac{3}{4}a_2\right)^2+\left(\dfrac{1}{4}_2\right)^2}=\dfrac{a_2\sqrt{10}}{4}=a\left(\dfrac{\sqrt{10}}{4}\right)^2$.\\ 
		Do đó diện tích $S_3=\left(\dfrac{5}{8}\right)^2a^2=\dfrac{5}{8}S_2$.\\
		Lý luận tương tự ta có $S_1, S_2, S_3,\ldots, S_n, \ldots$ tạo thành một dãy cấp số nhân lùi vô hạn có $u_1=S_1=a^2$ và công bội $q=\dfrac{5}{8}$.\\
		Do đó $T=\dfrac{S_1}{1-q}=\dfrac{8a^2}{3}$.\\
		Với $T=\dfrac{32}{3}$ ta có $a^2=4\Leftrightarrow a=2$.  
	} 
\end{bt}

\begin{bt}%[1D2N3-3]
% nguon: Archive/SP_KhoLatex/ID6_old/1K2-7.tex #3439
%[Dự án 2 - TLDH - TeamTeXHoa- Phạm Hoàng Đăng]
	Xác định công bội, số hạng thứ $5$, số hạng tổng quát và số hạng thứ $100$ của mỗi cấp số nhân sau
	\begin{enumerate}
		\item $1,4,16,\ldots $; \item $2,-\dfrac{1}{2},\dfrac{1}{8},\ldots$
	\end{enumerate} 
	\loigiai{
		\begin{enumerate}
			\item Ta có: $4:1=4$, $16:4=4$. Do đó công bội của cấp số nhân là $q=4$.\\
			Số hạng tổng quát $u_n=4^{n-1}$.\\
			Số hạng thứ 5 là $u_5=4^{5-1}=256$.\\
			Số hạng thứ 100 là $u_{100}=4^{100-1}=4{,}017\cdot {10}^{59}$.
			\item Ta có $\left( -\dfrac{1}{2} \right):2=-\dfrac{1}{4}$; $\dfrac{1}{8}:\left( -\dfrac{1}{2}\right)=-\dfrac{1}{4}$. Do đó công bội $q=-\dfrac{1}{4}$.\\
			Số hạng tổng quát $u_n=2\cdot \left(-\dfrac{1}{4}\right)^{n-1}$.\\
			Số hạng thứ $5$ là $u_5=2\cdot\left(-\dfrac{1}{4}\right)^{5-1}=\dfrac{1}{128}$.\\
			Số hạng thứ $100$ lầ $u_{100}=2\cdot \left(-\dfrac{1}{4}\right)^{100-1}=-4{,}978\cdot 10^{-60}$.
		\end{enumerate}
	} 
\end{bt}

\begin{bt}%[1D2H2-2]
% nguon: Archive/SP_KhoLatex/ID6_old/1K2-7.tex #3440
%[Dự án 2 - TLDH - TeamTeXHoa- Phạm Hoàng Đăng]
	Viết năm số hạng đầu của mỗi dãy số $(u_n)$ sau và xem nó có phải là cấp số nhân không. Nếu nó là cấp số nhân, hãy tìm công bội $q$ và viết công thức số hạng tổng quát của nó dưới dạng $u_n=u_1\cdot q^{n-1}$
	\begin{multicols}{2}
		\begin{enumerate}
			\item $u_n=5n$; 
			\item $u_n=5^n$;
			\item $u_1=1,u_n=nu_{n-1}$;
			\item $u_1=1,u_n=5u_{n-1}$.
		\end{enumerate}
	\end{multicols}
	\loigiai{
		\begin{enumerate}
			\item Năm số hạng đầu của dãy là $5;10;15;20;25$.\\
			Ta có $10:5=2\ne 15:10=\dfrac{3}{2}$ suy ra $(u_n)$ không phải cấp số nhân. 
			\item Năm số hạng đầu của dãy là $5;25;125;625;3125$.\\
			Ta có $u_n=5^n$ nên $u_{n+1}=5^{n+1}\Rightarrow \dfrac{u_{n+1}}{u_n}=\dfrac{5^{n+1}}{5^n}=5\,(\forall n\ge 2)$.\\
			Do đó $(u_n)$ là cấp số nhân có công bội $q=5$.\\
			Số hạng tổng quát $u_n=5\cdot {5}^{n-1}$. 
			\item Năm số hạng đầu của dãy $1;2;6;24;120$.\\
			Ta có $2:1=2\ne 6:2=3$ nên $(u_n)$ không phải là cấp số nhân.
			\item Năm số hạng đầu của dãy $1;5;25;125;625$.\\
			Ta có $u_n=5u_{n-1}$ nên $\dfrac{u_n}{u_{n-1}}=5\,(\forall n\ge 2)$.\\
			Do đó $(u_n)$ là cấp số nhân với công sai $q=5$.\\
			Số hạng tổng quát $u_n=5^{n-1}$.
		\end{enumerate}
		
	} 
\end{bt}

\begin{bt}%[1D2H2-2]
% nguon: Archive/SP_KhoLatex/ID6_old/1K2-6.tex #3474
%[Dự án 1 - TLDH - TeamTeXHoa - Nguyễn Văn A]%[ID6]
Chứng minh các dãy số sau là cấp số cộng.
\begin{enumerate}
\item Dãy số $(u_n)$ với $u_n=2020n-2021$.
\item Dãy số $(u_n)$ với $u_n=-2n+5$.
\end{enumerate}
\loigiai{
\begin{enumerate}
\item Dãy số $(u_n)$ với $u_n=2020n-2021$.\\
Ta có $u_{n+1}-u_n=2020(n+1)-2021-(2020n-2021)=2020$.\\
Vậy $(u_n)$ là một cấp số cộng với công sai $d=2020$.
\item Dãy số $(u_n)$ với $u_n=-2n+5$.\\
Ta có $u_{n+1}-u_n=-2(n+1)+5-(-2n+5)=-2$.\\
Vậy $(u_n)$ là một cấp số cộng với công sai $d=-2$.
\end{enumerate}
}
\end{bt}

\begin{bt}%[1D2H2-4]
% nguon: Archive/SP_KhoLatex/ID6_old/1K2-6.tex #3476
%[Dự án 1 - TLDH - TeamTeXHoa - Diamond]
Cho cấp số cộng $(u_n)$ có $u_3=15$ và $d=-2$. Tìm $u_n$.
\loigiai{
Ta có $\heva{&15=u_3=u_1+2d\\&d=-2} \Leftrightarrow \heva{&u_1=19\\&d=-2}\Rightarrow u_n=u_1+(n-1)d=-2n+21$.
}
\end{bt}

\begin{bt}%[1D2H2-4]
% nguon: Archive/SP_KhoLatex/ID6_old/1K2-6.tex #3477
%[Dự án 1 - TLDH - TeamTeXHoa - Diamond]
Một cấp số cộng có $8$ số hạng. Số hạng đầu là $5$, số hạng thứ tám là $40$. Khi đó công sai $d$ của cấp số cộng đó là bao nhiêu?
\loigiai{
Ta có $\heva{&u_1=5\\&40=u_8=u_1+7d}\Rightarrow d=5$.
}
\end{bt}

\begin{bt}%[1D2H2-4]
% nguon: Archive/SP_KhoLatex/ID6_old/1K2-6.tex #3478
%[Dự án 1 - TLDH - TeamTeXHoa - Diamond]
Cho cấp số cộng $(u_n)$ có $u_1=123$ và $u_3-u_{15}=84$. Tìm số hạng $u_{17}$.
\loigiai{
Ta có công sai của cấp số cộng là $d=\dfrac{u_3-u_{15}}{3-15}=\dfrac{84}{-12}=-7$.\\
Suy ra $u_{17}=u_1+(17-1)d=11$.
}
\end{bt}

\begin{bt}%[1D2V2-6]
% nguon: Archive/SP_KhoLatex/ID6_old/1K2-6.tex #3479
%[Dự án 1 - TLDH - TeamTeXHoa - Diamond]
Cho cấp số cộng $(u_n)$ có $u_1+2u_5=0$ và $S_4=14$. Tính số hạng đầu $u_1$ và công sai $d$ của cấp số cộng.
\loigiai{
Ta có $u_1+2u_5=0\Leftrightarrow u_1+2(u_1+4d)=0\Leftrightarrow 3u_1+8d=0$.\\
Lại có $S_4=14\Leftrightarrow \dfrac{4(2u_1+3d)}{2}=14\Leftrightarrow 2u_1+3d=7$.\\
Ta có hệ phương trình $\heva{&3u_1+8d=0\\&2u_1+3d=7}\Leftrightarrow \heva{&u_1=8\\&d=-3.}$
}
\end{bt}

\begin{bt}%[1D2H2-6]
% nguon: Archive/SP_KhoLatex/ID6_old/1K2-6.tex #3480
%[Dự án 1 - TLDH - TeamTeXHoa - Diamond]
Cho cấp số cộng $(u_n)$ có $u_1=4$ và $d=-5$. Tính tổng $100$ số hạng đầu tiên của cấp số cộng.
\loigiai{
Ta có $S_n=nu_1+\dfrac{n(n-1)}{2}d\Rightarrow S_{100}=100u_1+\dfrac{100\cdot 99}{2}d=-24350$.
}
\end{bt}

\begin{bt}%[1D2V2-6]
% nguon: Archive/SP_KhoLatex/ID6_old/1K2-6.tex #3481
%[Dự án 1 - TLDH - TeamTeXHoa - Diamond]
Xét các số nguyên dương chia hết cho $3$. Tính tổng số $50$ số nguyên dương đầu tiên. 
\loigiai{
Số nguyên dương chia hết cho $3$ có dạng $3n$, $(n\in\mathbb{N}^*)$ nên chúng lập thành cấp số cộng \[u_n=3n\Rightarrow \heva{&u_1=3\\&u_{50}=150}\Rightarrow S_{50}=\dfrac{50}{2}\left(u_1+u_{50}\right)=3825.\]
}
\end{bt}

\begin{bt}%[1D2V2-6]
% nguon: Archive/SP_KhoLatex/ID6_old/1K2-6.tex #3482
%[Dự án 1 - TLDH - TeamTeXHoa - Diamond]
Tính tổng $S=1-2+3-4+5+\cdots+(2n-1)-2n$ với $n\geq 1$ và $n\in\mathbb{N}$.
\loigiai{
Với mọi $n\in\mathbb{N}^*$ thì $(2n-1)-2n=-1$.\\
Ta có $S=(1-2)+(3-4)+(5-6)+\cdots+((2n-1)-2n)$. Do đó ta xem $S$ là tổng của $n$ số hạng, mà mỗi số hạng đều bằng $-1$ nên $S=-n$.
}
\end{bt}

\begin{bt}%[1D2V2-6]
% nguon: Archive/SP_KhoLatex/ID6_old/1K2-6.tex #3483
%[Dự án 1 - TLDH - TeamTeXHoa - Diamond]
Cho cấp số cộng $(u_n)$ thỏa mãn $u_2+u_8+u_9+u_{15}=100$. Tính tổng $16$ số hạng đầu tiên của cấp số cộng đã cho.
\loigiai{
Ta có $u_2+u_8+u_9+u_{15}=100\Leftrightarrow 4u_1+30d=100\Leftrightarrow 2u_1+15d=50$.\\
Khi đó $S_{16}=\dfrac{16}{2}(u_1+u_{16})=8(2u_1+15d)=8\cdot 50=400$.
}
\end{bt}

\begin{bt}%[1D2V2-6]
% nguon: Archive/SP_KhoLatex/ID6_old/1K2-6.tex #3484
%[Dự án 1 - TLDH - TeamTeXHoa - Diamond]
Cho cấp số cộng $(u_n)$ có công sai $d=-3$ và $u_2^2+u_3^2+u_4^2$ đạt giá trị nhỏ nhất. Tính tổng $S_{100}$ của $100$ số hạng đầu tiên của cấp số cộng đó. 
\loigiai{
Đặt $a=u_1$ thì $u_2^2+u_3^2+u_4^2=(a+d)^2+(a+2d)^2+(a+3d)^2=3a^2-36a+126=3(a-6)^2+18\geq 18$ với mọi $a$. \\
Dấu bằng xảy ra khi $a-6=0\Leftrightarrow a=6$. Suy ra $u_1=6$. \\
Ta có $S_{100}=\dfrac{100\cdot [2u_1+(100-1)d]}{2}=-14250$.
}
\end{bt}

\begin{bt}%[1D2V2-6]
% nguon: Archive/SP_KhoLatex/ID6_old/1K2-6.tex #3485
%[Dự án 1 - TLDH - TeamTeXHoa - Diamond]
Biết $u_4+u_8+u_{12}+u_{16}=224$. Tính $S_{19}$.
\loigiai{
Ta có
\begin{eqnarray*}
u_4+u_8+u_{12}+u_{16}=224&\Leftrightarrow& u_1+3d+u_1+7d+u_1+11d+u_1+15d=224\\
&\Leftrightarrow& 4u_1+36d=224\\
&\Leftrightarrow& u_1+9d=56.
\end{eqnarray*}
Ta có $S_{19}=\dfrac{19}{2}(2u_1+18d)=19(u_1+9d)=19\cdot 56=1064$.
}
\end{bt}

\begin{bt}%[1D2H2-5]
% nguon: Archive/SP_KhoLatex/ID6_old/1K2-6.tex #3486
%[Dự án 1 - TLDH - TeamTeXHoa - Diamond]
Cho các số $-4$; $1$; $6$; $x$ theo thứ tự lập thành một cấp số cộng. Tìm $x$.
\loigiai{
Vì các số $-4$; $1$; $6$; $x$ theo thứ tự $u_1$, $u_2$, $u_3$, $u_4$ lập thành cấp số cộng nên
$$u_4-u_3=u_3-u_2\Rightarrow x-6=6-1\Leftrightarrow x=11.$$
}
\end{bt}

\begin{bt}%[1D2H2-5]
% nguon: Archive/SP_KhoLatex/ID6_old/1K2-6.tex #3487
%[Dự án 1 - TLDH - TeamTeXHoa - Diamond]
Nếu các số $5+m$; $7+2m$; $17+m$ theo thứ tự lập thành cấp số cộng thì $m$ bằng bao nhiêu?
\loigiai{
Ba số $5+m$; $7+2m$; $17+m$ theo thứ tự $u_1$, $u_2$, $u_3$ lập thành cấp số cộng nên
$$u_1+u_3=2u_2\Leftrightarrow (5+m)+(17+m)=2(7+2m)\Leftrightarrow m=4.$$}
\end{bt}

\begin{bt}%[1D2H2-5]
% nguon: Archive/SP_KhoLatex/ID6_old/1K2-6.tex #3488
%[Dự án 1 - TLDH - TeamTeXHoa - Diamond]
Với giá trị nào của $x$ và $y$ thì các số $-7$; $x$; $11$; $y$ theo thứ tự đó lập thành một cấp số cộng?
\loigiai{
Bốn số $-7$; $x$; $11$; $y$ theo thứ tự $u_1$, $u_2$, $u_3$, $u_4$ lập thành cấp số cộng nên
$$\heva{&u_4-u_3=u_3-u_2\\&u_4-u_3=u_2-u_1}\Leftrightarrow \heva{&y-11=11-x\\&y-11=x+7}\Leftrightarrow \heva{&x+y=22\\&-y=-18}\Leftrightarrow \heva{&x=2\\&y=20.}$$
}
\end{bt}

\begin{bt}%[1D2V2-2]
% nguon: Archive/SP_KhoLatex/ID6_old/1K2-6.tex #3490
%[Dự án 1 - TLDH - TeamTeXHoa - Diamond]
Cho $a$, $b$, $c$ là ba số hạng liên tiếp của một cấp số cộng. Chứng minh rằng
\begin{enumerate}
\item $a^2+2bc=c^2+2ab$.
\item $a^2+8bc=(2b+c)^2$.
\end{enumerate}
\loigiai{
Vì $a$, $b$, $c$ là ba số hạng liên tiếp của một cấp số cộng nên $a+c=2b\Leftrightarrow a=2b-c$.
\begin{enumerate}
\item Ta có $a^2-2ab=(2b-c)^2-2(2b-c)\cdot b=4b^2-4bc+c^2-4b^2+2bc=c^2-2bc$.\\
Vậy $a^2-2ab=c^2-2bc\Leftrightarrow a^2+2bc=c^2+2ab$.
\item Ta có $a^2+8bc=(2b-c)^2+8bc=4b^2-4bc+c^2+8bc=4b^2+4bc+c^2=(2b+c)^2$.
\end{enumerate}
}
\end{bt}

\begin{bt}%[1D2H2-5]
% nguon: Archive/SP_KhoLatex/ID6_old/1K2-6.tex #3491
%[Dự án 1 - TLDH - TeamTeXHoa - Diamond]
Một tam giác vuông có chu vi bằng $3a$ và ba cạnh lập thành một cấp số cộng. Tính độ dài ba cạnh của tam giác theo $a$.
\loigiai{
Gọi $x$, $y$, $z$ theo thứ tự là độ dài ba cạnh của tam giác $(x<y<z)$.\\
Chu vi của tam giác là $x+y+z=3a$.\hfill(1)\\
Theo tính chất của cấp số cộng, ta có $x+z=2y$.\hfill(2)\\
Tam giác đã cho vuông nên $x^2+y^2=z^2$.\hfill(3)\\
Thay (2) và (1), ta được $3y=3a\Leftrightarrow y=a$.\\
Thay $y=a$ vào (2), ta được $x+z=2a\Rightarrow x=2a-z$.\\
Thay $x=2a-z$ và $y=a$ vào (3), ta được
$$(2a-z)^2+a^2=z^2\Leftrightarrow 5a^2-4az=0\Leftrightarrow z=\dfrac{5a}{4}\Rightarrow x=\dfrac{3a}{4}.$$
Vậy độ dài ba cạnh của tam giác là $\dfrac{3a}{4}$, $a$, $\dfrac{5a}{4}$.
}
\end{bt}

\begin{bt}%[1D2C2-5]
% nguon: Archive/SP_KhoLatex/ID6_old/1K2-6.tex #3492
%[Dự án 1 - TLDH - TeamTeXHoa - Diamond]
Tìm tất cả các giá trị của tham số $m$ để phương trình sau có ba nghiệm phân biệt lập thành một cấp số cộng $$x^3-3mx^2+2m(m-4)x+9m^2-m=0.$$
\loigiai{
\begin{itemize}
\item Điều kiện cần.\\
Giả sử phương trình đã cho có ba nghiệm phân biệt $x_1$, $x_2$, $x_3$ lập thành một cấp số cộng. Theo định lý Vi-ét đối với phương trình bậc ba, ta có $x_1+x_2+x_3=3m$. Vì $x_1$, $x_2$, $x_3$ lập thành cấp số cộng nên $x_1+x_3=2x_2$.\\
Suy ra $3x_2=3m\Leftrightarrow x_2=m$.\\ Thay $x_2=m$ vào phương trình đã cho, ta được 
$$m^3-3m\cdot m^2+2m(m-4)\cdot m+9m^2-m=0\Leftrightarrow m^2-m=0\Leftrightarrow \hoac{&m=0\\&m=1.}$$
\item Điều kiện đủ.
\begin{itemize}
\item Với $m=0$ thì ta có phương trình $x^3=0\Leftrightarrow x=0$ (phương trình có nghiệm duy nhất). Do đó $m=0$ không phải giá trị cần tìm. 
\item Với $m=1$, ta có phương trình $x^3-3x^2-6x+8=0\Leftrightarrow \hoac{&x=1\\&x=-2\\&x=4.}$\\ 
Ba nghiệm $-2$; $1$; $4$ lập thành một cấp số cộng nên $m=1$ là giá trị cần tìm.
\end{itemize}
\end{itemize} 
}
\end{bt}

\begin{bt}%[1D2C2-5]
% nguon: Archive/SP_KhoLatex/ID6_old/1K2-6.tex #3493
%[Dự án 1 - TLDH - TeamTeXHoa - Diamond]
Tìm giá trị của tham số $m$ để phương trình sau có bốn nghiệm phân biệt lập thành một cấp số cộng $x^4-10x^2+2m^2+7m=0$.  $\hfill(1)$
\loigiai{
Đặt $t=x^2$, $(t\geq 0)$. Khi đó ta có phương trình $t^2-10t+2m^2+7m=0$. $\hfill(2)$\\
Phương trình $(1)$ đã cho có $4$ nghiệm phân biệt khi và chỉ khi phương trình $(2)$ có $2$ nghiệm dương phân biệt $\Leftrightarrow \heva{&5^2-(2m^2+7m)>0\\&2m^2+7m>0}\Leftrightarrow 0<2m^2+7m<25$. $\hfill(*)$\\
Với điều kiện trên thì $(2)$ có hai nghiệm dương phân biệt là $t_1$, $t_2$ $(t_1<t_2)$.\\
Khi đó phương trình đã cho có bốn nghiệm phân biệt là $-\sqrt{t_2}$; $-\sqrt{t_1}$; $\sqrt{t_1}$; $\sqrt{t_2}$.\\
Bốn nghiệm này lập thành một cấp số cộng khi $$-\sqrt{t_1}-\left(-\sqrt{t_2}\right)=\sqrt{t_1}-\left(-\sqrt{t_1}\right)=\sqrt{t_2}-\sqrt{t_1}\Leftrightarrow t_2=9t_1.$$ 
Theo định lý Vi-ét ta có $t_1+t_2=10$ và $t_1t_2=2m^2+7m$.\\
Suy ra ta có hệ phương trình $\heva{&t_2=9t_1\\&t_1+t_2=10\\&t_1t_2=2m^2+7m}\Leftrightarrow \heva{&t_1=1\\&t_2=9&2m^2+7m=9}\\
\Rightarrow \hoac{& m=1\\&m=-\dfrac{9}{2}} (\text{thỏa điều kiện} (*)).$
}
\end{bt}

\begin{bt}%[1D2H2-3]
% nguon: Archive/SP_KhoLatex/ID6_old/1K2-6.tex #3494
[Bài $2.8$ - Sách giáo khoa]%[Dự án 1 - TLDH - TeamTeXHoa - Diamond]
Xác định công sai, số hạng thứ $5$, số hạng tổng quát và số hạng thứ $100$ của mỗi cấp số cộng sau
\begin{enumerate}
\item $4$, $9$, $14$, $19,\ldots $
\item $1$, $-1$, $-3$, $-5,\ldots$
\end{enumerate}
\loigiai{
\begin{enumerate}
\item Ta có $9-4=5$, $14-9=5$. Suy ra cấp số cộng có $u_1=4$, công sai $d=5$.\\
Số hạng tổng quát của dãy số là $u_n=4+5(n-1)$.\\
Số hạng thứ $5$ là $u_5=4+5(5-1)=24$.\\
Số hạng thứ $100$ là $u_{100}=4+5(100-1)=499$.
\item Ta có $-1-1=-2$; $-3-(-1)=-2$. Suy ra cấp số cộng có $u_1=1$, công sai $d=-2$.\\
Số hạng tổng quát của dãy số là $u_n=1+(-2)(n-1)=3-2n$.\\
Số hạng thứ $5$ là $u_5=3-2\cdot 5=-7$.\\
Số hạng thứ $100$ là $u_{100}=3-2\cdot 100=-197$.
\end{enumerate}
}
\end{bt}

\begin{bt}%[1D2H2-2]
% nguon: Archive/SP_KhoLatex/ID6_old/1K2-6.tex #3495
[Bài $2.9$ - Sách giáo khoa]%[Trần Phú Hiếu]
	Viết năm số hạng đầu của mỗi dãy số $\left(u_n\right)$ sau và xét xem nó có phải là cấp số cộng không. Nếu dãy số đó là cấp số cộng, hãy tìm công sai $d$ và viết số hạng tổng quát của nó dưới dạng $u_n=u_1+\left(n-1\right)d$.
	\begin{listEX}[2]
		\item $u_n=3+5n$.
		\item $u_n=6n-4$.
		\item $u_1=2,u_n=u_{n-1}+n$.
		\item $u_1=2,u_n=u_{n-1}+3$.
	\end{listEX}
	\loigiai{
		\begin{listEX}
			\item $u_1=8;u_2=13;u_3=18;u_4=23;u_5=28$.\\
			Ta có $u_n-u_{n-1}=3+5n-\left[3+5\left(n-1\right)\right]=5; \forall x\ge 2$.\\
			Suy ra dãy số là cấp số cộng có $u_1=8$ và công sai $d=5$.\\
			Số hạng tổng quát $u_n=8+5\left(n-1\right)$.
			\item $u_1=2;u_2=8;u_3=14;u_4=20;u_5=26$.\\
			Ta có $u_n-u_{n-1}=6n-4-\left[6\left(n-1\right)-4\right]=6; \forall x\ge 2$.\\
			Suy ra dãy số là cấp số cộng có $u_1=2$ công sai $d=6$.\\
			Số hạng tổng quát $u_n=2+6\left(n-1\right)$.
			\item $u_1=2;u_2=4;u_3=7;u_4=11;u_5=16$.\\
			Ta có $u_2-u_1=2\ne u_3-u_2=3$.\\
			Suy ra đây không phải cấp số cộng.
			\item $u_1=2;u_2=5;u_3=8;u_4=11;u_5=14$.\\
			Ta có $u_n-u_{n-1}=3$.\\
			Suy ra đây là dãy cấp số cộng có $u_1=2$ và công sai $d=3$.\\
			Số hạng tổng quát $u_n=2+3\left(n-1\right)$.
		\end{listEX}
	}
\end{bt}

\begin{bt}%[1D2H2-4]
% nguon: Archive/SP_KhoLatex/ID6_old/1K2-6.tex #3496
[Bài $2.10$ - Sách giáo khoa]%[Trần Phú Hiếu]
	Một cấp số cộng có số hạng thứ $5$ bằng $18$ và số hạng thứ $12$ bằng $32$. Tìm số hạng thứ $50$ của cấp số cộng này.
	\loigiai{
		Gọi số hạng tổng quát của dãy là $u_n=u_1+\left(n-1\right)d$.\\
		Ta có $u_5=u_1+4d=18;u_{12}=u_1+11d=32$.
		Suy ra $u_1=10,d=2$.\\
		Số hạng tổng quát $u_n=10+2\left(n-1\right)$.\\
		Số hạng thứ $50$ là $u_{50}=108$.}
\end{bt}

\begin{bt}%[1D2H2-6]
% nguon: Archive/SP_KhoLatex/ID6_old/1K2-6.tex #3497
[Bài $2.11$ - Sách giáo khoa]%[Trần Phú Hiếu]
	Một cấp số cộng có số hạng đầu bằng $5$ và công sai bằng $2$. Hỏi phải lấy tổng của bao nhiêu số hạng đầu của cấp số cộng này đề có tổng bằng $2700$?
	\loigiai{
		Gọi $n$ là số các số hạng đầu cần lấy tổng.\\ Ta có
		$2700=S_n=\dfrac{n}{2}\left[2\cdot 5+\left(n-1\right)\cdot 2\right]=\dfrac{n}{2}\left(8+2n\right)$.\\
		Do đó $4n+n^2-2700=0$.\\ Giải phương trình bậc hai này ta được $n=-54$ (loại) hoặc $n=50$ (nhận).\\
		Vậy phải lấy $50$ số hạng đầu.}
\end{bt}

\begin{bt}%[1D2H2-7]
% nguon: Archive/SP_KhoLatex/ID6_old/1K2-6.tex #3498
[Bài $2.12$ - Sách giáo khoa]%[Trần Phú Hiếu]
	Giá của một chiếc xe ô tô lúc mới mua là 680 triệu đồng. Cứ sau mỗi năm sử dụng, giá của chiếc xe ô tô giảm $55$ triệu đồng. Tính giá còn lại của chiếc xe sau $5$ năm sử dụng.
	\loigiai{
		Giá của chiếc xe sau $n$ năm là $u_n=680-55\left(n-1\right)$.\\
		Vậy sau $5$ năm sử dụng giá của chiếc xe là $u_5=680-55\left(5-1\right)=460$ (triệu đồng).}
\end{bt}

\begin{bt}%[1D2H2-7]
% nguon: Archive/SP_KhoLatex/ID6_old/1K2-6.tex #3499
[Bài $2.13$ - Sách giáo khoa]%[Trần Phú Hiếu]
	Một kiến trúc sư thiết kế một hội trường với $15$ ghế ngồi ở hàng thứ nhất, $18$ ghế ngồi ở hàng thứ hai, $21$ ghế ngồi ở hàng thứ ba, và cứ như vậy (số ghế ở hàng sau nhiều hơn $3$ ghế so với số ghế ở hàng liền trước nó). Nếu muốn hội trường đó có sức chứa ít nhất 870 ghế ngồi thì kiến trúc sư đó phải thiết kế tối thiểu bao nhiêu hàng ghế?
	\loigiai{
		Số ghế ở mỗi hàng lập thành một cấp số cộng với số hạng đầu $u_1=15$ và công sai $d=3$.\\ Gọi $n$ là số các số hạng đầu của cấp số cộng cần lấy tổng.\\Ta có
		$870=S_n=\dfrac{n}{2}\left[2\cdot 15+\left(n-1\right)\cdot 3\right]=\dfrac{n}{2}\left(27+3n\right)$.\\
		Do đó $27n+3n^2-1740=0$, suy ra $n=20,n=-29$ (loại).\\
		Vậy cần phải thiết kế $20$ hàng ghế}
\end{bt}

\begin{bt}%[1D2H2-7]
% nguon: Archive/SP_KhoLatex/ID6_old/1K2-6.tex #3500
[Bài $2.14$ - Sách giáo khoa]%[Trần Phú Hiếu]
	Vào năm $2020$, dân số của một thành phố là khoảng $1{,}2$ triệu người. Giả sử mỗi năm, dân số của thành phố này tăng thêm khoảng $30$ nghìn người. Hãy ước tính dân số của thành phố này vào năm $2030$.
	\loigiai{
		Dân số mỗi năm của thành phố lập thành cấp số cộng có $u_1=1200$, công sai $d=30$.\\
		Dân số mỗi năm có dạng tổng quát là $u_n=1200+30\left(n-1\right)$.\\
		Dân số của năm $2030$ tức $n=11u_{11}=1200+30\left(11-1\right)=1500$ (nghìn người).}
\end{bt}

\begin{bt}%[1D2H2-7]
% nguon: Archive/SP_KhoLatex/ID6_old/1-GK1-KNTT-De17-NH23-24.tex #3690

Trong dịp nghỉ lễ 02/9 gia đình anh An cần thuê một xe Taxi để di chuyển từ thủ đô Hà Nội về thăm quê tại TP Bắc Giang. Biết giá của kilômét đầu tiên là $10.000$ đồng, kể từ kilômét thứ 2 giá của mỗi kilômét tăng thêm $500$ đồng so với giá của kilômét trước đó. Biết quãng đường Taxi di chuyển từ thủ đô Hà Nội về TP Bắc Giang là $80$ km. Hỏi gia đình anh An phải trả bao nhiêu tiền cho chuyến Taxi đó?
\loigiai{
Số tiền ở từng kilômet lập thành một cấp số cộng.\\
Số hạng đầu tiên là $u_1=10.000$, công sai $d=500$.\\
Cấp số cộng có $80$ số hạng.\\
Tổng số tiền là $S_{80}=80.10000+\dfrac{79.80 .500}{2}=2.380 .000$ đồng.
}
\end{bt}

\begin{bt}%[1D2H3-8]
% nguon: Archive/SP_KhoLatex/ID6_old/1-GK1-KNTT-De17-NH23-24.tex #3691

Biết một tế bào sau $1$ chu kì nguyên phân sẽ tạo ra thêm 1 tế bào mới. Tính tổng số lượng tế bào sau 10 chu kì nguyên phân từ 5 tế bào ban đầu?
\loigiai{
 Gọi tổng số tế bào sau 1 chu kì nguyên phân là $u_1=5+5=10$.\\
Do một tế bào sau $1$ chu kì nguyên phân sẽ tạo ra thêm 1 tế bào mới nên tổng số tế bào sau 1 chu kì tăng gấp 2 lần số lượng trước đó.\\
Suy ra tổng số lượng tế bào sau mỗi chu kì tạo thành 1 cấp số nhân có số hạng đầu tiên $u_1=10$ và công bội $q=2$.\\
Vậy số lượng tế bào sau 10 chu kì là: $u_{10}=u_1 \cdot q^9=10.2^9=5120$ tế bào.
}
\end{bt}
